% 本文件由 LaTeX 排版 Agent 根据有效 translation.json 自动生成。
% author=Cyril of Jerusalem; work_id=3101; title=教理讲授

\chapter{教理讲授}

% structure: p0001 source=h1 target=omitted reason=page-title-already-rendered-by-parent

\Emph{序言}\newline{}\Emph{第一讲}\newline{}\Emph{第二讲}\newline{}\Emph{第三讲}\newline{}\Emph{第四讲}\newline{}\Emph{第五讲}\newline{}\Emph{第六讲}\newline{}\Emph{第七讲}\newline{}\Emph{第八讲}\newline{}\Emph{第九讲}\newline{}\Emph{第十讲}\newline{}\Emph{第十一讲}\newline{}\Emph{第十二讲}\newline{}\Emph{第十三讲}\newline{}\Emph{第十四讲}\newline{}\Emph{第十五讲}\newline{}\Emph{第十六讲}\newline{}\Emph{第十七讲}\newline{}\Emph{第十八讲}\newline{}\Emph{第十九讲}\newline{}\Emph{第二十讲}\newline{}\Emph{第二十一讲}\newline{}\Emph{第二十二讲}\newline{}\Emph{第二十三讲}\par

\SourceRecord{Translated by Edwin Hamilton Gifford.；Nicene and Post-Nicene Fathers, Second Series；Vol. 7.；Edited by Philip Schaff and Henry Wace.；Buffalo, NY: Christian Literature Publishing Co.,；1894.；<http://www.newadvent.org/fathers/3101.htm>.}

% structure: p0001 source=h1 target=section reason=page-title

\section{\WorkTitle{预备教理（序言）}}

1. 你们这些即将受光照的人，已有福乐的芬芳临于你们身上；你们已采集属灵的花朵，编织天上的花冠；圣神的馨香已吹拂你们；你们已聚集在君王宫殿的门廊；愿你们也被君王引入！因为树上已现出花朵；愿果实也达至完美！迄今已有你们名字的登记、服役的召唤、婚宴队伍的灯火、对天上户籍的渴望、善良的意向以及随之而来的望德。因为那说\BibleQuote{\Emph{凡爱天主的人，事事为他们而合作，使之得益}}的，并不说谎。天主慷慨施恩，却等待每个人的真诚意愿；因此宗徒加上说：\BibleQuote{蒙召的，是按照祂的旨意}。意向的真诚使你们成为蒙召的；因为若你们的身体在此，而心意不在，便一无益处。\par

2. 甚至西满术士也曾来到圣洗池：他受了洗，却未受光照；他虽将身体浸入水中，却没有以圣神光照他的心；他的身体下去又上来，但他的灵魂并未与基督同埋，也未与祂同复活。我提及跌倒的事例，为使你们不致跌倒；\BibleQuote{\Emph{这些事作为鉴戒发生，并被记录下来，以警戒那至今仍前来的人}}。愿你们中无人试探祂的恩宠，以免\BibleQuote{有苦根长出而扰乱你们}。愿你们中无人进来说：\Quote{让我们看看信徒们在做什么；让我进去看看，好知道正在举行什么。}你以为你来看，却不被看吗？你以为当你探察所发生之事时，天主不在探查你的心吗？\par

3. 福音中曾有一个人窥探婚宴，穿着不合宜的衣服进来，坐下吃饮；新郎容许了。但当看见所有人都穿着白衣时，他本也该穿上同样的衣服；他虽吃同样的食物，却在样式和意向方面与他人不同。新郎虽然慷慨，却并非无辨识力；他巡视每位宾客并观察他们（他的关心不是他们吃得好，而是行为得体），看见一个陌生人\BibleQuote{\Emph{没有穿婚宴礼服}}，就对他说：\BibleQuote{朋友，你怎么到这里来了？}什么样的颜色！什么样的良心！纵然因为主人的慷慨，看门人没有拦阻你，又怎样？纵然你不知道该以何样式进入宴席，又怎样？你进来并看到了宾客闪耀的样式；难道没有被眼前所见教训吗？难道不该适时退出，好适时再进入吗？但你已不合时宜地进来，将不合时宜地被逐出。于是祂命令仆人：\BibleQuote{捆起他的脚，}那双胆敢擅入的脚；\BibleQuote{捆起他的手，}那不知如何穿上明亮衣服的手；\BibleQuote{把他丢在外面的黑暗中；}因他配不上婚宴的火炬。你们看到那人的遭遇：要确保你们自身的境况安全。\par

4. 因为我们，基督的仆役，接纳了每个人，并如同占据看门人之位，我们敞开了门：你可能灵魂沾满罪恶、意志不洁地进来了。你进来并被容许：你的名字被登记。告诉我，你见到教会这可敬的制度了吗？你见到她的秩序和纪律、圣经的诵读、圣职者的临在、教理的课程了吗？要因这地方而羞愧，并被所见所闻教导。现在适时地出去，明天最适时地进来。\par

若你灵魂的样式是贪婪，就换上另一样式再进来。脱去你从前的样式，不要掩饰。我恳求你，脱去淫乱和不洁，穿上最光明的贞洁礼服。我在灵魂的新郎耶稣进来察看他们样式之前，给你们这命令。你们被给予很长的通知；你们有四十天做补赎；你们有充分的机会脱去、洗涤、穿上并进入。但若你们坚持邪恶的意向，说话者是无过的，而你们不可指望恩宠；因为水会接纳，但圣神不会接纳你们。若有人意识到自己的创伤，让他敷上药膏；若有人跌倒了，让他起来。愿你们中没有西满，没有伪善，没有对这事的好奇。\par

5. 可能你也是因别的借口而来。可能有人想讨好一个女人，为此来到此处。同样，这话也适用于妇女。奴隶也许想讨好主人，朋友讨好朋友。我接受这钩上的饵，并欢迎你们，虽然你们怀着不良目的而来，却要因美好的希望而得救。也许你们不知自己来到何处，也不知落入了何种网罗。你们已进入教会的网中：要活捉，不要逃跑；因为耶稣在钓你们，不是为杀死，而是借杀死使你们活：因为你们必须死而复活。你们已听到宗徒说：\BibleQuote{死于罪恶，而活于正义}。死于你们的罪恶，活于正义，从今日起就活着。\par

6. 请看，我恳求你们，耶稣赐予你们何等大的尊荣。你们曾被称为慕道者，那时话语从外在环绕你们；你们听到希望，却不认识它；听到奥迹，却不理解它们；听到圣经，却不知道它们的深意。如今这回声不再在你们周围，而在你们内；因为从此以后，\Emph{内住的圣神}使你们的心灵成为天主的宫殿。当你们听到有关奥迹的记载时，你们就会明白先前所不知道的事。不要以为你们领受的是小事：虽然你们是卑微的人，却领受了一个天主的称号。听听圣保禄所说：\Quote{天主是忠信的}。听听另一段圣经所说：\Quote{天主是忠信而公义的}。圣咏作者预见这事，因为人要领受天主的称号，便以天主的口吻这样说了：\Quote{我曾说：\InnerQuote{你们是神，都是至高者的儿子}}。但要当心，免得你们拥有\Quote{忠信}的称号，却怀着不信者的心意。你们已进入一场竞赛，要奋力跑完赛程：你们不可能再有另一次机会。假如明天是你们的婚礼，你们岂不会放下一切，为婚宴做准备吗？如今在你们将要奉献灵魂给天上新郎的前夕，难道还不该停止肉体的事，好赢得属灵的事吗？\par

7. 我们不可两次或三次领受圣洗圣事；否则有人会说：\Quote{虽然我失败了一次，第二次我会改正}；但若你失败了一次，这事便无法补救；因为\Quote{只有一个主，一个信德，一个圣洗}；只有异端者才重洗，因为那第一次并不是圣洗。\par

8. 因为天主并不向我们要求别的，只要求善意。不要说：\Quote{我的罪如何被涂抹？}我告诉你：藉着愿意，藉着相信。还有什么比这更短的呢？但若你的口说愿意，而你的心却沉默，那审判你的主是知道人心的。从今天起，停止一切恶行。不要让你的口说出不当的言语，让你的眼睛远离罪恶，不去窥探无益之事。\par

9. 让你们的双脚赶快去参加要理讲授；要热切地接受驱魔礼：无论是被吹气还是被驱逐，这行为对你们都是救恩。假设你们有未经精炼的合金黄金，混杂着各种杂质，铜、锡、铁、铅；我们只想得到纯金，没有火，黄金能从杂质中分离出来吗？同样，没有驱魔礼，灵魂无法被净化；这些驱魔礼是神圣的，是从圣经中收集而来的。你们的脸被蒙住，为的是使你们的心思从此自由，免得眼睛游荡使心也游荡。但当你们的眼睛被蒙住时，你们的耳朵并未被阻碍接收救恩的工具。因为正如那些精通金匠工艺的人，通过某些精细的器具向火吹气，吹起隐藏于坩埚中的黄金，并搅动周围的火焰，从而找到他们所求之物；同样，当驱魔者藉着天主之神引起敬畏，并将灵魂置于身体这坩埚中，仿佛点燃它时，敌对魔鬼便逃跑了，留下救恩和永生的希望，灵魂从此被洗净罪恶而获得救恩。所以，弟兄们，让我们持守希望，并交付自己，盼望那万有的天主看见我们的心意，洗净我们的罪，赐予我们关于自身状态的美好希望，并赐给我们带来救恩的悔改。天主召叫了，这召叫是朝向你们的。\par

10. 要密切注意要理讲授，即使我们延长讲论，也不要让你们的内心感到疲倦。因为你们正在接受对抗敌对权势的盔甲，对抗异端、犹太人、撒玛黎雅人和外邦人的盔甲。你们有许多敌人，要取许多武器，因为你们要向许多敌人投射；你们需要学习如何击倒希腊人，如何与异端者、犹太人和撒玛黎雅人争战。武器已经预备好，最预备好的是\Quote{圣神的剑}；但你们也必须以良好的决心伸出右手，好能进行主的争战，克服敌对权势，并成为不可战胜的，对抗一切异端的企图。\par

11. 让我也给你们这嘱咐：研读我们的教导，并永远持守它。不要认为它们是普通的讲道；因为虽然它们也是好的、可靠的，但如果我们今天忽略它们，明天还可以研读。但如果关于重生的水池的教导是连续传授的，今天被忽略，何时才能纠正呢？假设这是栽种树木的季节：如果我们不挖土，且挖得深，那一次种得不好的树，何时才能种好呢？请设想，要理讲授好比一种建筑：如果我们不在建筑中用规则的连结把房子结合好，免得出现裂缝，建筑变得不稳固，那么我们先前的工作也白费了。石头必须按次序一块接一块，角石对正角石，并磨平不平之处，建筑才能这样均匀地上升。同样，我们是在把知识的石头带给你们。你们必须听论生活的天主，必须听论审判，必须听论基督，并听论复活。还有许多事要依次讨论，它们现在虽被一点一点地提出，但后来要和谐地连接起来。但除非你们把它们整合为一个整体，并记住什么是第一、什么是第二，建筑者可以建造，但你们会发现建筑是不稳固的。\par

12. 因此，当讲授进行时，若有慕道者问你们教师讲了什么，切不可告诉那外面的人。因为我们交给你们的是一个奥秘，是对来世生命的希望。请为那赐予赏报者保守这奥秘。绝不可有人对你说：\Quote{若我也知道，于你有何损害？} 同样，病人求酒，但若在不适当的时候给酒，会引发谵妄，就产生两种恶果：病人死亡，医生受责。慕道者从信徒口中听到什么也是如此：慕道者会变得神志不清（因为他不懂所听到的，反而抱怨这事，讥讽所说的话），而那信徒则被定罪为叛徒。但你们现在正站在边界上：当心，求你们不要向外泄露；这并非因为所说的话不值得传讲，而是因为他的耳朵不配领受。你们自己也曾是慕道者，我并未描述你们面前的事。当你们亲身经验到我们教导的事理何等崇高时，你们就会知道慕道者不配听这些事。\par

13. 你们这些已被登记的人，已成为同一位母亲的儿女。当你们在驱魔礼仪时辰之前进来时，你们每人都要说些有助于虔敬的事。若你们中间有人缺席，要寻找他。若你们被邀请赴宴，岂不等候同席的宾客？若你们有一位兄弟，岂不谋求兄弟的好处？\par

此后，不要忙于无益之事：既不要过问城里的事，也不要过问村庄的事，也不要过问君王、主教或司铎的事。要向上仰望，这才是你们此刻所需要的。\BibleQuote{你们要安静，要知道我是天主。} 若你们看见信徒在服务、毫无挂虑，他们享有平安，他们知道自己领受了什么，他们拥有恩宠。但你们此刻正站在天平上，将被接纳或不接纳：不要效法那些无忧无虑的人，而要怀有敬畏之心。\par

14. 驱魔结束后，直到其他被驱魔的人到来之前，男的要与男的在一起，女的要与女的在一起。因为我需用诺厄方舟的比喻：方舟内有诺厄和他的儿子们，以及他的妻子和儿媳们。因为虽然方舟是一个，门也关闭了，但一切都安排得妥善。同样，教会虽然关闭，你们都在内，仍要分开，男与男一起，女与女一起，免得救恩的借口变成毁灭的缘由。即使有靠近坐下的合宜理由，也要摒除情欲。此外，男性就座时，应有一本有益的书；一人读，另一人听。若无书，一人祈祷，另一人说些有益的话。同样，年轻女子们应一同就座，或歌唱，或安静阅读，使她们的嘴唇说话，但别人的耳朵听不见声音：\BibleQuote{我不许女人在教会中说话。} 已婚妇女也要效法这个榜样，祈祷，让她的嘴唇动，但声音不被听见，这样撒慕尔就会来到，你们不育的灵魂就会因天主垂听你们的祈祷而生出得救的果实；因为这就是撒慕尔名字的含义。\par

15. 我要观察每个人的热忱，每个女子的虔诚。愿你们的心思在虔敬中被火炼净；愿你们的灵魂如金属般被锻造；愿不信的顽硬被打掉；愿铁上多余的鳞片脱落，纯净的部分留下；愿铁锈被擦去，真金属存留。愿天主有朝一日向你们显示那夜晚，那如同白昼发光的黑暗，经上说：\BibleQuote{黑暗不能隐蔽你，黑夜必如同白昼发光。} 那时，愿天堂之门向你们中每个男女敞开。那时，愿你们享受带着芬芳的基督之水。那时，愿你们领受基督的名号，以及天主事物的权能。就在现在，我恳求你们，举起你们心灵的眼目；就在现在，想象天使的唱诗班，以及坐着的万有之主天主，和祂的独生子同祂坐在右边，圣神与他们同在；还有宝座与宰制者在服事，你们中每个男女都领受救恩。就在现在，愿你们的耳朵仿佛响起那荣耀的声音，那时天使要为你们的得救歌唱：\BibleQuote{罪恶蒙赦免、过犯得遮盖的人，是有福的。} 那时你们如同教会的星辰进入，身体光明，灵魂灿烂。\par

16. 摆在你们面前的圣洗圣事何其伟大：它是俘虏的赎价，罪过的赦免，罪恶的死亡，灵魂的重生，光明的衣服，神圣不可分解的印记，升天的战车，天堂的喜悦，进入王国的欢迎，义子之恩的赏赐！但是有一条蛇在路旁窥伺过路的人：当心，免得它用不信咬伤你们。它看见这么多人领受救恩，便\BibleQuote{寻找可吞食的人。} 你们正要来到众圣之神那里，却要经过那条蛇。那么你们如何通过它呢？要\BibleQuote{用和平福音的预备穿在脚上}，这样即使它咬，也不能伤害你们。要有内住的信德，坚定的望德，坚固的鞋子，好能经过仇敌，来到你们的主面前。预备你们自己的心，以领受道理，参与神圣的奥秘。要更频繁地祈祷，好使天主使你们配得上天上的、不朽的奥秘。不要日间或夜间停止祈祷：但当睡眠从你们眼中被驱散时，你们的心思就应自由地祈祷。若你们发现心中升起任何可耻的念头，要转向对审判的默想，以提醒你们得救的事。将你们的心思全部投入学习，好使它忘记卑贱的事。若你们发现有人对你们说：\Quote{那么你要进去，下到水里吗？难道城里现在没有澡堂吗？} 要认识到这是\Quote{海中的巨龙}在向你们设下这些圈套。不要听信那说话者的嘴唇，而要听信那在你们内工作的天主。保守你们的灵魂，使你们不被诱捕，好使你们在望德中存留，成为永生的继承人。\par

17. 我们作为人，这样命令并教导你们：但不要使我们的建筑成为\Emph{干草、禾秸}和糠秕，免得我们因\Emph{工程被焚}而\Emph{受亏损}，却要使我们的工程成为金子、\Emph{银子}和\Emph{宝石}！因为说话在于我，但用心思考在于你们，而使之成全在于天主。让我们振奋精神，鼓起灵魂，预备心灵。这场赛跑是为了我们的灵魂：我们的盼望是永恒的事物：天主知道你们的心，察看谁真诚、谁伪善，祂既能保护真诚的人，也能赐予伪善者信德；因为即使是不信的人，只要他肯交出他的心，天主也能赐予信德。愿祂\Emph{涂抹那反对你们的债卷}，并赐予你们先前过犯的赦免；愿祂将你们栽种在祂的教会中，征召你们为祂服役，给你们穿上\Emph{正义的武器}；愿祂以新盟约的天上事物充满你们，赐予你们圣神的印鉴，永世不朽，在基督耶稣我们的主内：愿光荣归于祂，直到永远！阿们。\par

(\Emph{致读者}。)\par

这些为即将受光照者准备的\WorkTitle{教理讲授}，你可以借给圣洗圣事的候选人以及已经受洗的信徒阅读，但绝不可交给慕道者或任何其他非基督徒，否则你将在主面前负责。如果你抄写一份，要在开头写下这话，如同在主眼前。\par

\SourceRecord{Translated by Edwin Hamilton Gifford.；Nicene and Post-Nicene Fathers, Second Series；Vol. 7.；Edited by Philip Schaff and Henry Wace.；Buffalo, NY: Christian Literature Publishing Co.,；1894.；<http://www.newadvent.org/fathers/310100.htm>.}

% structure: p0001 source=h1 target=section reason=page-title

\section{教理讲授第一讲}

\Emph{致那些将要受光照者，在耶路撒冷即席讲授，作为给前来领受圣洗者的入门讲授：}\par

读经：\BibleBook{依撒意亚}{先知书}（\BibleRef{依 1:16}）\par

\BibleQuote{你们要洗濯，自洁，从我眼前除掉你们灵魂的恶行}，以及其他（\BibleRef{依 1:16}）。\par

新约的门徒们，基督奥迹（\OriginalTerm{奥迹}{mysteries}）的分享者，目前只在蒙召上，但不久也将因恩宠而成：\BibleQuote{你们要造一个新心和新灵} \BibleRef{则 18:31}，好使天上的居民喜乐；因为\BibleQuote{如果有一个罪人悔改，为他就有喜乐}，照福音所说（\BibleRef{路 15:7}），那么许多灵魂得救岂不更使天上的居民欢欣。你们既已踏上这条美好而最荣耀的道路，就当怀着敬畏奔那虔敬的赛程。因为天主的独生子（\OriginalTerm{独生子}{Only-begotten Son}）就在这里，极愿救赎你们，说：\BibleQuote{凡劳苦和负重担的，你们都到我这里来，我要使你们安息。} \BibleRef{玛 11:28} 你们身着罪过的粗衣，被自己罪恶的绳索缠绕的，要听先知的话说：\BibleQuote{你们要洗濯，自洁，从我眼前除掉你们的恶行。} \BibleRef{依 1:16}：好使天使的歌咏团为你们歌唱：\BibleQuote{罪恶蒙赦免，过犯得遮盖的人，是有福的。} 你们刚刚点燃了信德之火的人，要小心不让它熄灭在你们手中；好使祂——祂昔日在这至圣的\Place{哥耳哥达}因那强盗的信德为他开启了乐园（\OriginalTerm{乐园}{Paradise}）——准许你们歌唱那婚礼之歌。\par

若这里有人是罪恶的奴隶，就让他藉着信德迅速准备自己，为了那重获自由的新生和义子地位（\OriginalTerm{新生}{new birth}；\OriginalTerm{义子地位}{adoption}）；脱去他罪恶的悲惨束缚（\OriginalTerm{束缚}{bondage}），穿上主的最有福的束缚，如此才配承受天国。藉着告解，脱去\BibleQuote{按照欺骗的情欲而败坏的旧人}，好能\BibleQuote{穿上新人，这新人是按照创造他的天主的肖像而更新的}。你们要藉着信德获得圣神的抵押（\OriginalTerm{抵押}{earnest}）（\BibleRef{格后 1:22}），好能被接进\BibleQuote{永远的帐幕}（\BibleRef{路 16:9}）。来领受奥迹的印记（\OriginalTerm{奥迹的印记}{mystical Seal}），好使你能轻易被主人认出；被列入基督的神圣而属灵的羊群中，站在祂右边，承受为你预备的生命。因为那些仍被罪恶的粗衣所缠的人，站在左边，因为他们没有藉着圣洗的新生来到那藉基督所赐的天主的恩宠面前：我所说的新生不是身体的新生，而是灵魂的属灵新生。因为我们的身体由可见的父母所生，但我们的灵魂藉着信德重生：\BibleQuote{风随意向哪里吹}（\BibleRef{若 3:8}）；然后，如果你被堪当，你可以听见：\BibleQuote{好！善良忠信的仆人}（\BibleRef{玛 25:21}），当你发现你的良心没有伪善的玷污时。\par

因为如果这里有任何人想试探天主的恩宠，他是自欺，不明白恩宠的能力。人啊，你要使你的灵魂远离伪善，因为祂\Emph{察看人心和肺腑}。正如那些准备征兵的人会检查应征者的年龄和身体，同样，主在招募灵魂时也审查他们的意图：如果有人存有隐秘的伪善，主就拒绝他，认为他不配为祂的真正服务；但如果祂发现一个堪当的人，祂便乐意赐予祂的恩宠。祂不将圣物给狗（\BibleRef{玛 7:6}）；但祂在哪里辨明良好的良心，就在哪里赐下救恩的印记（\OriginalTerm{救恩的印记}{Seal of salvation}），那奇妙的印记，魔鬼畏惧它，天使认识它；好使魔鬼被驱逐，天使如亲属般环绕它。因此，接受这属神而拯救的印记的人，也需要与它相称的意向。因为正如写字用的芦苇或标枪需要使用者，同样，恩宠也需要相信的心灵。\par

你们领受的不是一个可朽坏的盾牌，而是一个神性的盾牌（\OriginalTerm{神性的盾牌}{spiritual shield}）。从此你们被栽种在无形的乐园（\OriginalTerm{无形的乐园}{invisible Paradise}）中。你们领受一个新的名字，是你们从前没有的。从前你们是望教者，但现在将被称为信徒。你们从此被移植到属灵的橄榄树中，从野橄榄枝被接在好橄榄树上（\BibleRef{罗 11:24}），从罪恶到正义，从污秽到纯洁。你们成了神圣葡萄树的分享者。因此，如果你们住在这葡萄树上，你们就生长为结果的枝条；如果不住，你们将被火焚烧。让我们因而结出相称的果实。但愿在我们身上不发生那棵不结果的无花果树所遭遇的事（\BibleRef{玛 21:19}），免得耶稣现在就来到，因我们的不结果而诅咒我们。但愿人人都能运用另一句话：\Emph{我犹如天主殿中的茂盛橄榄树，永远信赖天主的慈爱}——一棵不是感官所感知，而是理智所见的、充满光明的橄榄树。因此，祂栽种和浇灌是祂的部分，你们结果实是你们的责任：天主赐予恩宠，但你们要接受并守护它。不要因恩宠是白白赐予而轻视它，而要虔诚地接受并珍藏它。\par

5. 现在是告解的时期：要告明你日夜在言语或行为上所做的事；\BibleQuote{在悦纳的时候告解，并在救恩之日} \BibleRef{格后 6:2}领受天上的宝藏。把你的时间投入驱魔礼：勤勉参加教理讲授，并记住所说的话，因为它们不仅是为你的耳朵而说，更是为使你藉信德将之密封在记忆中。从你心中抹去一切世俗的挂虑：因为你是为你的灵魂奔跑。你正在完全弃绝世事：你所弃绝的是小事，主所赐予的是大事。要弃绝现世之物，寄望于来世之物。你在世上徒劳忙碌地奔跑了这么多年的循环，难道不能为了你灵魂的缘故，拿出四十天来自由（祈祷）吗？\BibleQuote{你们要静止，并要知道我是天主}，圣经说。要避免说许多闲话：既不要毁谤，也不要乐意听毁谤者的话；反而要恳切祈祷。在苦修操练中显明你的心是坚强的。要洁净你的器皿，好能更丰富地领受恩宠。因为虽然罪过的赦免平等地赐给所有人，但圣神的相通是按各人的信德分量赐予的。如果你劳苦少，就领受少；如果你劳苦多，赏报就大。你是为自己奔跑，要留意你自己的益处。\par

6. 如果你对任何人有什么不满，要宽恕他：你来这里是要获得罪过的赦免，你也必须宽恕那得罪了你的人。否则，如果你自己没有宽恕你的同仆连那些小小的罪过，你还有什么脸面向主说：\Quote{赦免我众多的罪过}？要勤勉参加教会的聚会；不仅现在当神职人员要求你勤勉时如此，而且在你领受了恩宠之后也要如此。因为如果在领受之前，这操练是好的，难道在赐予之后不更好吗？如果在嫁接之前，浇水照料是稳妥的，难道在栽种之后不更好吗？要为你的灵魂奋斗，特别是在这样的日子里。用神圣的读经滋养你的灵魂；因为主为你预备了属灵的筵席；因此你也要像圣咏作者那样说：\BibleQuote{上主是我的牧者，我必一无所缺；在青草地上，祂使我安歇；在舒适的水边，祂牧养了我；祂使我的灵魂复新}：——好使天使们也分享你的喜乐，而基督自己，这位大司祭，既接纳了你的决心，将你们众人呈献给圣父，说：\BibleQuote{看，我和天主所赐给我的孩子们}。愿祂使你们众人蒙祂喜悦！愿光荣与权能归于祂，直到永远无穷之世。亚孟。\par

\SourceRecord{Translated by Edwin Hamilton Gifford.；Nicene and Post-Nicene Fathers, Second Series；Vol. 7.；Edited by Philip Schaff and Henry Wace.；Buffalo, NY: Christian Literature Publishing Co.,；1894.；<http://www.newadvent.org/fathers/310101.htm>.}

% structure: p0001 source=h1 target=section reason=page-title

\section{教理讲授第二讲}

\Emph{论悔改与罪过的赦免，并论及仇敌}\par

\BibleRef{则 18:20-23}\par

\BibleQuote{义人的义必归于他，恶人的恶也必归于他。但如果恶人转离他的一切罪过，等等。}\par

1. 罪是一件可怕的事，灵魂最严重的疾病是过犯，它暗中割断灵魂的筋，也成为永火的原因；这是人自选的那恶，意志的产物。因为我们出于自由意志犯罪，先知在某处清楚地说：\BibleQuote{我栽培你成为多产的葡萄树，全然真实；你何以变为苦涩，成了野葡萄？}\BibleRef{耶 2:21} 栽培是好的，来自意志的果实却是恶的；因此栽培者无可指摘，但葡萄树必被火烧，因为它被栽培为善，却因自己的意志结出恶果。\BibleQuote{因为天主}，按训道者所说，\BibleQuote{造了人正直，而人却自寻许多诡计。}\BibleRef{训 7:29} \BibleQuote{因为我们是祂的化工}，宗徒说，\BibleQuote{受造于善工之中，就是天主预先预备，使我们行在其中。}\BibleRef{弗 2:10} 所以造物主原是善的，为善工而造了人；但受造物却凭自己的自由意志转向了邪恶。因此，罪，如我们所说，虽是一件可怕的恶，却并非不可治愈；对紧握它的人可怕，但对藉悔改将它弃绝的人则易于医治。举例说，假如有人手中握着火；只要他紧握那火炭，就必然被烧，但若他扔掉那炭，他也将火焰一并抛弃了。但如果有人认为他犯罪时并未被烧，圣经就向他说：\BibleQuote{人若将火藏在怀中，他的衣服岂能不烧着？}\BibleRef{箴 6:27} 因为罪烧毁灵魂的筋，〔并折断心思的精神骨骼，且使心灵之光昏暗。〕\par

2. 但有人会说，罪是什么？是一个活物吗？是一位天使吗？是一个魔鬼吗？是什么在我们内工作？人哪，并非从外攻击你的仇敌，而是从你自身长出的恶枝。\BibleQuote{注视前方，用你的眼睛}\BibleRef{箴 4:25}，就没有贪欲。［保持你自己的，］不要夺取别人的东西，抢劫就停止了。记住审判，那么奸淫、淫乱、谋杀及任何违犯法律的事都不能胜过你。但每当你忘记天主，你就立刻开始图谋邪恶并行不义。\par

3. 然而你并非恶的唯一作者，还有另一个极邪恶的唆使者，就是魔鬼。他确实诱惑，但并不强制征服那些不同意的人。因此训道者说：\BibleQuote{若有掌权者的灵起来攻击你，不要离开你的本位。} 关上门，将他远远隔离，他就不能伤害你。但如果你漫不经心地接纳贪欲的念头，它便藉其暗示在你内扎根，并束缚你的心思，将你拖入邪恶的深渊。\par

但也许你说，我是信徒，贪欲不能胜过我，即使我常想它。难道你不知道，长久的坚持连石头也能穿破？不要接纳种子，否则它会撕裂你的信德：在恶开花之前连根拔出，免得起初漫不经心，日后不得不寻求斧子和火。当你的眼睛开始患病时，及时医治，免得你失明，然后不得不寻找医生。\par

4. 因此魔鬼是罪的第一作者，是恶者的父亲：这是主的话，不是我的，\BibleQuote{魔鬼从起初就犯罪}：没有谁在他之前犯罪。但他犯罪并非因本性必然有犯罪倾向，否则罪的缘由就要追溯到造他如此的那位了；但他被造是善的，却凭自己的自由意志成了魔鬼，并因他的行为得了这名。因为他原是总领天使，后来因毁谤被称为魔鬼；从良善的天主仆人，他恰当地被称为撒殚；因为\Quote{撒殚}意为\Emph{仇敌}。这不是我的教导，而是受默感的先知厄则克耳的教导：因为他为他作哀歌说：\BibleQuote{你是钦崇的封印，完美的冠冕；你生于天主的乐园中}；紧接着又说：\BibleQuote{在你的日子，从你受造之日起，你原是完美的，直到在你内发现了不义。} 他非常正确地说\BibleQuote{在你内发现了}，因为它们不是从外来的，而是你自己生了这恶。他随即也提到原因：\BibleQuote{你的心因你的美丽而高傲；因你罪恶繁多，你受了伤，我将你抛在地上。} 与此一致，主在福音中又说：\BibleQuote{我看见撒殚像闪电一样从天坠下。}\BibleRef{路 10:18} 你看到旧约与新约的和谐。他被逐出时，带走了许多跟他的人。就是他，将贪欲放入那些听从他的人心中：从他而来的是奸淫、淫乱和一切邪恶。藉着他，我们的祖先亚当因悖逆被逐出，将自发结出奇妙果实的乐园，换成了生出荆棘的土地。\par

5. 那么，有人会说，我们被诱骗而沉沦了。难道再没有救恩了吗？我们跌倒了：难道不能再站起来吗？我们瞎了眼：难道不能重见光明吗？我们瘸了腿：难道再也不能直立行走吗？简言之，我们死了：难道不能再复活吗？那唤醒了已死四天、甚至已经发臭的\Person{拉匝禄}的，岂不更轻易地使你活着的人复活？那为我们倾流宝血的主，将亲自解救我们脱离罪恶。兄弟们，我们不要对自己绝望，不要让自己陷入无望的境地。因为不相信悔改的希望是一件可怕的事。因为不看救恩的人，不会停止在恶上加恶；但希望得到医治的人，从此容易谨慎自守。不看赦免的强盗变得绝望；但若他期望宽恕，常常就会悔改。那么，蛇能蜕皮，难道我们不能脱去我们的罪吗？多荆棘的土地，若勤加耕种，也能变成肥沃；难道救恩对我们来说无法挽回吗？不，恰恰相反，我们的本性接受救恩，但还需要意志。\par

6. 天主是爱人的，而且爱得极大。因为不要说：\Quote{我犯了奸淫和通奸，我做了可怕的事，不仅一次，而且常常；祂会赦免吗？祂会宽恕吗？}听听圣咏作者说：\BibleQuote{上主，祢的恩慈何其广大！}你所累积的过犯超不过天主仁慈的广大；你的伤口超不过大医师的技艺。只管凭信德把自己交出来：告诉医师你的病痛：你也像\Person{达味}那样说：\BibleQuote{我决意向主承认我的罪过；}同样的事也要在你身上实现，正如他紧接着说的：\BibleQuote{祢就赦免了我心中的邪恶。}\par

7. 你最近才来听要理，你愿意看见天主的慈爱吗？你愿意看见天主的慈爱和祂宽忍的丰富吗？请听关于亚当的事。亚当，天主最初所造的人，犯了罪：祂难道不能立刻置他于死地吗？但请看主因对人大爱所行的：祂把他逐出乐园，因他因罪不堪在那里居住；但祂\Emph{使他住在乐园对面}，使他看到自己从何处坠落，从何等地位堕落到何等境地，以后可以因悔改而得救。加音，头生之人，成了他兄弟的凶手，邪恶的发明者，凶杀的首倡者，第一个嫉妒的人。然而，在杀害兄弟之后，他受了什么判决？\Emph{你在地上要哀叹和颤抖。}何等大的罪过，刑罚却是何等轻！\par

8. 即使这种情形，在天主确实是慈爱，但与随后相比，还算微小。请思考诺厄时代发生的事。巨人犯了罪，那时大地上邪恶遍布，因此洪水要临到他们；在五百年时，天主发出威胁；但在六百年时，祂使洪水降在地上。你看见天主的慈爱延展了百年之久吗？祂难道不能立刻行祂百年后所行的事吗？但祂特意延长时间，给予悔改的宽限。你看见天主的良善吗？如果那时的人悔改了，他们就不会错过天主的慈爱了。\par

9. 现在请跟随我来看另一类，那些因悔改而得救的人。但或许甚至在妇女中，有人会说，我犯了奸淫和通奸，我以各种放纵玷污了身体，还有救恩给我吗？妇人，请你转向辣哈布，你也寻求救恩；因为如果那公开作妓女的，因悔改而得救，那么，那在接受恩宠之前，曾一度犯过奸淫的人，岂不也能因悔改和斋戒而得救吗？请问她如何得救：她只说了这话：\Emph{\BibleQuote{你们的天主是上天和下地的主}}（\BibleRef{苏 2:11}）。\Emph{你们的天主}，因为她不敢说自己的天主，因她淫荡的生活。你若愿意从圣经中得到她得救的证明，在圣咏中记着说：\Emph{\BibleQuote{我要在认识我的人中提及辣哈布和巴比伦}}。哦，天主的慈爱何其伟大，甚至在圣经中提及妓女！不，不只是\Emph{我要提及辣哈布和巴比伦}，还加上了\Emph{在认识我的人中}。因此，无论男女，悔改都带来救恩。\par

10. 不仅如此，即使整个民族犯罪，这也不能超越天主的慈爱。百姓铸造了牛犊，但天主并未停止祂的慈爱。人否认了天主，但天主不能否认自己（\BibleRef{弟后 2:13}）。他们说：\Emph{\BibleQuote{以色列，这就是你的天主}}（\BibleRef{出 32:4}）；然而，按照祂的惯例，以色列的天主又成了他们的救主。不仅百姓犯罪，而且大司祭亚郎也犯罪。因为梅瑟说：\Emph{\BibleQuote{上主向亚郎发怒，我为他祈祷，天主就宽恕了他}}（\BibleRef{申 9:20}）。那么，梅瑟为犯罪的大司祭祈祷，尚且能感动天主，难道祂的独生子耶稣为我们祈祷，却不能感动天主吗？如果祂因亚郎的过犯，没有阻止他进入大司祭职，难道祂会阻止你，这从外邦人中出来的，进入救恩吗？人哪，只管同样悔改，恩宠必不向你禁止。从今以后，使你的生活无可指责，因为天主确实爱人，万古不足以述说祂的慈爱；不，即使所有人的口舌联合起来，也不能说出祂慈爱的任何重要部分。因为我们所讲的，只是圣经所载祂对人慈爱的一部分，但祂赦免天使多少，我们不得而知；因为祂也赦免他们，只有一位是无罪的，就是耶稣，祂洁净我们的罪。关于这些，我们已经说得够了。\par

11. 但是，如果关于我们人类，你面前也设置了别的例子，那么就来看蒙福的\Person{大卫}，以他为补赎的榜样。他虽伟大，却也跌倒了：他在睡后，傍晚在屋顶上行走，漫不经心的一瞥，感到了人性的冲动。他的罪完成了，但他对于认罪的那份坦诚并没有随之死亡。\Person{拿单}先知来了，是一位迅速的控告者，也是伤口的医治者。他说：\Emph{上主发怒}，\Emph{你犯了罪}。\EditorialNote{撒下 12}被控告者对在位君王如此说话。但\Person{大卫}王并不愤慨，因为他所看重的不是说话的人，而是那派遣他的天主。他并没有因周围站立的士兵阵列而自高：因为他心中看到了主的众天使，并且\Emph{如同看见那看不见的}（\BibleRef{希 11:27}）而战栗；他对使者，或者更确切地说，如同对那打发使者来的天主一样答复说：\Emph{我得罪了上主}（\BibleRef{撒下 12:13}）。你看到君王的谦卑了吗？你看到他的认罪了吗？岂有人定他的罪吗？有许多人知悉此事吗？这事很快就做成了，先知立即出现作为控告者，犯罪者承认了过错。因为他坦诚地承认，他得到了非常迅速的治愈。因为先知\Person{拿单}说了威胁的话，立刻说：\Emph{上主也赦免了你的罪}。你看到仁慈的天主迅速回心转意。然而，他说：\Emph{你大大惹起了上主的仇敌的} \Emph{上主}。虽然你因你的义德有许多仇敌，你的自制保护了你；但现在你交出了你最坚固的甲胄，你的仇敌起来，站好了阵势攻击你。\par

12. 先知就这样安慰了他，但是蒙福的\Person{大卫}，尽管听到了\Emph{上主赦免了你的罪}这句话，并没有停止补赎，尽管他是君王，他却穿上麻衣代替紫袍，代替金宝座，君王坐在灰堆中；不仅如此，他不仅坐在灰里，还以灰为食物，正如他自己说：\Emph{我吃灰烬如同吃饼}。他因流泪消耗了他好色的眼睛，说：\Emph{每夜我要洗我的床铺，用眼泪湿润我的床榻}。当他的官员恳请他吃面包时，他不听。他将斋戒延长了足足七天。如果一个君王这样认罪，你一个普通人难道不该认罪吗？又在\Person{阿贝沙隆}叛乱后，尽管有许多路可以逃脱，他选择从\Place{橄榄山}逃跑，仿佛在思想上呼求那位将要从此升入高天的救赎者。\BibleRef{撒下 16:10}当\Person{史米}辱骂他时，他说：\Emph{由他吧}，因为他知道\Quote{赦免人的，也必得赦免}。\par

13. 你看，认罪是好的。你看，悔改的人有救恩。\Person{撒罗满}也跌倒了，但他说什么？\Emph{后来我悔改了}。\Person{阿哈布}，\Place{撒玛黎雅}王，也成了最邪恶的偶像崇拜者，一个暴虐的人，先知们的谋杀者（\BibleRef{列上 18:4}），与虔诚无关，觊觎别人的田地和葡萄园。然而，当他因\Person{依则贝尔}的计谋杀害了\Person{纳波特}后，先知\Person{厄里亚}来了，只是威胁了他，他就撕裂衣服，穿上麻衣。仁慈的天主对\Person{厄里亚}说什么？\Emph{你看见\Person{阿哈布}在我面前如何痛心吗？}仿佛祂几乎要说服先知的火热热情怜悯这个悔改的人。因为祂说：\Emph{在他有生之日，我不降这灾祸。}虽然在这赦免之后，他肯定不会离开他的邪恶，然而宽恕的天主宽恕了他，不是因为祂不知道未来，而是按照他现在的悔改时期赐予相应的赦免。因为公义的审判者的职责是根据每一件发生的事作出判决。\par

14. 再者，\Person{雅洛贝罕}站在祭坛前向偶像献祭；他的手枯干了，因为他命令捉拿责备他的先知；但他通过经验认识了他面前这人的能力，便说：\Emph{请你恳求上主你的天主的慈颜}（\BibleRef{列上 13:6}）；因这句话，他的手复原了。如果先知治好了\Person{雅洛贝罕}，基督岂不能医治并救你脱离你的罪吗？\Person{默纳舍}也是极其邪恶的，他锯死了\Person{依撒意亚}，并沾染了各种偶像崇拜，\Emph{使\Place{耶路撒冷}充满无辜的血}（\BibleRef{编下 33:12}）；但被掳到\Place{巴比伦}后，他利用不幸的经历进行了治愈性的补赎；因为圣经说\Person{默纳舍}\Emph{在上主面前谦卑自己，祈祷，上主垂听了他，领他回到他的王国}。如果那锯死先知的人因悔改而得救，那么你既没有做如此大的恶事，岂不得救吗？\par

15. 当心，勿无端怀疑悔改的力量。你愿知道悔改有何等力量吗？你愿认识救恩的坚固武器，并学习告解的力量吗？希则克雅藉着告解，击溃了他十八万五千敌人。这诚然是大事，但较之将述说的事仍属微小：同一君王藉着悔改，挽回了已经发出的天主谕令。当他患病时，依撒意亚向他说：\Quote{处理你的家务，因为你必死，不能存活。}\BibleRef{列下 20:1} 还有什么指望，什么痊愈的希望，当先知说\Quote{你必死}？然而希则克雅并未停止悔改；他记着经上所写：\Quote{当你回转哀哭时，你必得救}\BibleRef{依 30:15}，于是转身面向墙壁，从病榻上举心向天（因为墙壁的厚薄不能阻挡虔诚的祈祷），他说：\Quote{上主，求祢记念我，因为祢记念我，便足以使我痊愈。祢不受时间限制，祢自己就是生命的立法者。因为我们的生命并非取决于出生或星宿的会合，如某些人妄言；而是取决于生命及其延续。祢自己按祢的旨意立法。}他因先知的预言而不敢指望存活，却得增寿十五年；作为征兆，太阳在轨道上后退。果然，为希则克雅的缘故，太阳回转；但为基督，太阳却遭蚀，不是后退，而是被遮蔽\BibleRef{依 38:8}，因此显出二者之间的差异——我是指希则克雅与耶稣之间。前者得以取消天主的命令，难道耶稣不能赐予罪的赦免吗？回转并哀哭吧，关上你的门，祈祷求赦，求祂从你身上除去燃烧的火焰。因为告解有能力熄灭烈火，有能力驯服狮子。\par

16. 但若你不信，请思量阿纳尼雅及其同伴所遇之事。他们倾倒了何等溪流？多少桶水能熄灭那高达四十九肘的火焰？不，当火焰稍稍升得过高时，信德如河涌流，他们念出对抗万恶的咒语：\Quote{上主，祢在我们身上所做的一切，都是公义的；因为我们犯罪，违背了祢的法律。}他们的悔改平息了火焰。若你不信悔改能熄灭地狱之火，请从阿纳尼雅的事迹学习。但有些敏锐的听者会说：那些人，天主在那种情况下公正地拯救了他们，因为他们拒绝敬拜偶像，天主赐予他们那力量。既有了这个念头，我接下来要举另一个悔改的例子。\par

17. 你认为拿步高怎样？你岂未从圣经中听说他嗜血、凶猛、性情如狮？你岂未听说他将诸王的骨骸从坟墓中取出暴露于光？你岂未听说他掳走百姓？你岂未听说他在国王目睹子女被杀之后，又剜了他的眼睛\BibleRef{列下 25:7}？你岂未听说他打碎了革鲁宾？我不是指那不可见的神体——人哪，摒弃这种想法——而是指雕刻的像和赎罪盖，天主曾在其间发声说话。圣所的帐幔被他践踏；香坛被他取走，带到偶像的庙中\BibleRef{编下 36:7}；所有祭品被他取走；圣殿被他从根基焚毁。他杀死君王、焚烧圣所、掳掠百姓、将圣器置于偶像庙中，该受何等重大的惩罚？他不该死一万次吗？\par

18. 你已看见他恶行的巨大；现在来观看天主的慈爱。他被变为野兽，住在旷野，受鞭打，为得救恩。他有狮爪，因他曾蹂躏圣所；他有狮鬃，因他是吞噬咆哮的狮子；他像牛一样吃草，因他是畜生，不认识那赐他王权者。他被露水浸湿，因他见过火焰被露水熄灭，仍不信。结果怎样？他说：\Quote{此后，我拿步高举目望天，我赞颂至高者，我赞美、光荣永生之主。}\BibleRef{达 4:34} 当他承认至高者，向天主发出这些感谢的话，并为自己所行的悔改，承认自己的软弱时，天主便将王国的尊荣归还于他。\par

19. 那么？当拿步高行过如此恶事后，告解了，天主便赐他赦免和王国；当你悔改时，难道祂不赐你罪的赦免和天国，倘若你度相称的生活？主是爱人、速于赦免、缓于惩罚的。因此，谁也不要绝望于自己的救恩。伯多禄，宗徒中最卓越、为首者，在一个使女面前三次否认了主；但他悔改了，痛哭流涕。哭泣显示内心的悔改；因此他不仅为自己的否认获得宽恕，而且保持了宗徒的尊位，未受剥夺。\par

20. 因此，弟兄们，你们既有许多犯罪后悔改而得救的榜样，也当诚心向主告解，好使你们既获得先前诸罪的赦免，又配得天赐恩典，并在耶稣基督内与众圣者一同继承天国；愿光荣归于祂，直到永远。阿们。\par

\SourceRecord{Translated by Edwin Hamilton Gifford.；Nicene and Post-Nicene Fathers, Second Series；Vol. 7.；Edited by Philip Schaff and Henry Wace.；Buffalo, NY: Christian Literature Publishing Co.,；1894.；<http://www.newadvent.org/fathers/310102.htm>.}

% structure: p0001 source=h1 target=section reason=page-title

\section{\WorkTitle{慕道训辞第三讲}}

\Emph{论圣洗圣事。}\par

\BibleRef{罗 6:3-4}\par

难道你们不知道：凡受洗归于基督耶稣的人，就是受洗归于他的死亡吗？我们藉着洗礼已归于死亡与他同葬，等等。\par

1. \Emph{\Quote{诸天，欢乐吧！大地，踊跃吧！}}因为那些将要被牛膝草洒净、并以属灵的牛膝草——就是那位在受难时被人用牛膝草和芦苇递上饮料之主的权能——所洁净的人。当天上的万军欢乐时，那些将要与属灵新郎结合的灵魂，也应预备自己。因为\Emph{有声音}在旷野中呼喊：\Quote{\InnerQuote{你们要预备上主的道路。}}\BibleRef{依 40:3} 因为这并非小事，也非按肉体的寻常、无区别的结合，而是按信德由洞察万有的圣神所拣选。因为世上的联姻和契约并不全然出于审慎：凡有财富或美貌之处，新郎便迅速认可；但在这里，并非容貌之美，而是灵魂的清白良心；并非被谴责的玛门，而是灵魂在虔敬上的财富。\par

2. 所以，正义的儿女们，请听若翰的劝诫，他说：\Quote{你们要修直主的道路。}除去一切障碍和绊脚石，使你们能笔直走向永生。预备好灵魂的器皿，以真诚的信德洁净，为迎接圣神。立即开始在忏悔中洗净你们的衣裳，使你们在被召入婚筵时，能被发现是洁净的。因为新郎毫无区别地邀请众人，因他的恩宠是丰盛的；高声呼喊的传报者召集他们所有人；但随后同一位新郎要分别那些前来参加比喻性婚筵的人。但愿那些名字已被登记的人中，没有一个听到：\Quote{朋友，你没有婚筵礼服，怎么到这里来了？} \BibleRef{玛 22:12} 但愿你众人都能听到：\Quote{好！善良忠信的仆人，你既在小事上忠信，我必委派你管理许多大事：进入你主人的福乐吧！} \BibleRef{玛 25:12}\par

因为现在你们暂时站在门外；惟愿天主恩赐你们众人都能说：\Quote{君王带我进入了他的内室。} \BibleRef{歌 1:4} \Emph{\Quote{愿我的灵魂因上主欢欣，因为他给我穿上救恩的衣裳，披上喜乐的外袍；他给我戴上冠冕，如同新郎，并以装饰点缀我，如同新娘。}}好使你们每个人的灵魂能被发现\Quote{没有瑕疵、没有绉纹，或其他类的缺陷} \BibleRef{弗 5:7}；我不是说在领受恩宠之前，因为那怎么可能呢？既然你们是因罪过的赦免而被召，而是说，当恩宠要被赐予时，你们的良心若不被定罪，便能与恩宠相合。\par

3. 弟兄们，这实在是一件严肃的事，你们必须以谨慎之心接近它。你们每个人将要被呈现于天主面前，在成千上万的天使军旅之前；圣神将要印封你们的灵魂；你们将被编入大君的军队。所以，要预备自己，装备自己——我所说的装备，不是指光鲜的衣服，而是指灵魂的虔敬与良好的良心。不要把洗礼池看作单纯的水，而要视同那与水一同赐予的属灵恩宠。正如那献给异教祭坛的供物，虽然本性简单，却因偶像的呼求而玷污；同样，这简单的水，一旦接受了圣神、基督和圣父的呼求，便获得一种新的圣化能力。\par

4. 因为人是由灵魂和身体两方面构成的，所以洁净也是双重的：一方面是为无形部分的无形洁净，另一方面是为身体的有形洁净；水洁净身体，圣神印封灵魂；这样我们就能亲近天主，\Quote{心里由圣神洒净，身体用清水洗净} \BibleRef{希 10:22}。因此，当你下到水中时，不要只想到那单纯的原质，而要因圣神的德能期待救恩；因为缺少两者中的任何一个，你们都不可能成为完美的。说这话的不是我，而是主耶稣基督，他在这件事上拥有权柄；因为他曾说：\Quote{人若不重生}（他又加上）\Quote{由水和圣神，就不能进天主的国。} \BibleRef{若 3:3} 那受水洗却未堪当领受圣神的人，并不能圆满地领受恩宠；同样，若有人在行为上圣善，却没有接受水的印封，也不能进入天国。这是一句大胆的话，但不是我的，因为这是耶稣所宣告的；以下是出自圣经的证明：科尔乃略是一个义人，曾蒙天使显现的荣耀，他的祈祷和施舍在天主面前已成为一份美好的纪念。伯多禄来了，圣神便倾注在那些信从的人身上，他们就说各种语言，并说预言；而且在圣神的恩宠之后，圣经记载伯多禄\Quote{吩咐他们因耶稣基督之名受洗} \BibleRef{宗 10:48}；为的是，灵魂既因信德而重生，身体也藉着水分享恩宠。\par

5. 但若有人想知道为何恩宠是借着水而不是其他元素赐予的，让他翻阅圣经，便会知晓。因为水是伟大的事物，是世上四种可见元素中最尊贵的。天是天使的居所，但诸天是由水而来；地是人的居所，但地是由水而来；在万物受造六天之前，\BibleQuote{圣神运行在水面上。}\BibleRef{创 1:2} 水是世界的开始，约旦河是福音佳音的开始：因为以色列从法郎手中得救是经过海，而世界从罪恶中得救是借着\BibleQuote{水洗与圣言}\BibleRef{弗 5:26} 的清除。凡与人订立盟约的地方，也有水。洪水之后，与诺厄订立盟约；以色列从西乃山所立的盟约，却是\BibleQuote{用水、朱红羊毛和牛膝草}\BibleRef{希 9:19}。厄里亚被接升天，也离不开水：因为他先渡过约旦河，然后乘战车升天。大司祭先洗净，然后献香；因为亚郎先洗净，然后成为大司祭：尚未被水洁净的人，怎能替别人祈祷？又作为圣洗的预象，会幕内设立了一个铜盆。\par

6. 圣洗是旧约的终结，新约的开始。因为它的创始者是若翰，\Quote{妇女所生的，没有一个大过他的}。他是先知们的终末：\BibleQuote{因为众先知和法律讲说预言，直到若翰为止}\BibleRef{玛 11:13}；但他是福音历史的初果。因为经上说：\BibleQuote{耶稣基督福音的开始，等等}：\BibleQuote{若翰在旷野施洗}。你可以提及提斯贝人厄里亚被接升天，但他并不大于若翰；哈诺客被提升天，但他并不大于若翰；梅瑟是伟大的立法者，所有先知都可敬，但都不大于若翰。我不敢拿先知与先知比较；但他们的主，也是我们的主，主耶稣亲自宣告：\BibleQuote{妇女所生的，没有兴起一个比若翰更大的}\BibleRef{玛 11:11}：祂不是说\InnerQuote{在童女所生的}，而是\Emph{在妇女所生的}。这是伟大的仆人与同侪的比较；但圣子的卓越和恩宠远非仆人可比。你看见天主选择了怎样伟大的人作为这恩宠的首席分施者吗？——他身无长物，喜爱旷野，却不憎恨人类；吃蝗虫，使他的灵魂飞向天堂；以蜂蜜为食，说的话比蜜更甘甜、更有益；身穿骆驼毛的衣服，在自身上显出苦修生活的榜样；他还在母腹中时，已被圣神祝圣。耶肋米亚在母腹中被祝圣，却没有预言\BibleRef{耶 1:5}；唯有若翰在母腹中跳跃欢喜\BibleRef{路 1:44}，虽未以肉眼看见，却以圣神认识他的主：因为圣洗的恩宠如此伟大，也需要它的创始者伟大。\par

7. 这人在约旦河施洗，\BibleQuote{全耶路撒冷都出去到他那里}\BibleRef{玛 3:5}，享受圣洗的初果：因为耶路撒冷有一切美事的特权。但是，耶路撒冷的居民啊，你们要学习那些出去的人是如何受他施洗的：经上说\BibleQuote{承认他们的罪}\BibleRef{玛 3:6}。他们先显出伤口，然后他施用药物，并赐给相信的人脱离永火的救赎。你们若想确信这一点，即若翰的圣洗是从火刑的威胁中得救赎，就要听他说：\Emph{毒蛇的种类！谁指示你们逃避那即将来临的忿怒？} \Emph{从此以后，不要再作毒蛇，但你们从前是毒蛇的孽种，他说，要脱去蛇皮} \Emph{你们从前罪恶的生活。因为每条蛇都爬进洞中脱去旧皮，擦掉旧皮后，身体又变得年轻。同样，你们也要从窄门进去}\BibleRef{玛 7:13}：借着禁食擦去你们的旧我，赶出那毁灭你们的东西。\BibleQuote{脱去旧人和他的行为}\BibleRef{哥 3:9}，并引用《雅歌》中的那句话：\Quote{我脱了外衣，怎能再穿上呢？}\par

但你们中也许有伪善者、取悦于人者、假装虔诚却心不相信者；具有西满术士的伪善；他来此不是为领受恩宠，而是窥探所赐的：让他也从若翰学习：\BibleQuote{现在斧子已经放在树根上，凡不结好果子的树，就砍下来丢在火里}\BibleRef{玛 3:10}。审判者是无情的；除去你们的伪善。\par

8. 那么你们该做什么？悔改的果实是什么？\BibleQuote{有两件衣裳的，就分给那没有的}\BibleRef{路 3:11}：这位老师是值得信赖的，因为他首先实行自己所教导的；他不耻于言说，因为良心没有束缚他的口舌：\BibleQuote{有食物的，也该照样做}。你愿意享受圣神的恩宠，却认为穷人不配得身体的食物吗？你寻求大恩赐，却不分施小恩赐吗？即使你是税吏或淫乱者，也有得救的希望：\BibleQuote{税吏和娼妓比你们先进天主的国}\BibleRef{玛 21:31}。保禄也作证说：\BibleQuote{淫乱者、奸夫等都不能承受天主的国。你们中有些人从前也是这样，但你们已经洗净了，已经圣化了}\BibleRef{格前 6:9}。他不是说\Quote{你们中有些人现在是这样的}，而是\Quote{从前是这样}。在无知中所犯的罪可得赦免，但坚持作恶必受谴责。\par

9. 你拥有圣洗的光荣，就是天主的独生子本身。因为何必再提及人呢？若翰是伟大的，但他与主相比算什么？他声音响亮，但与圣言相比算什么？那传报者十分高贵，但与君王相比算什么？那用水施洗的高贵，但与那\BibleQuote{\Emph{以圣神和火}}\BibleRef{玛 3:11}施洗的相比又算什么？救主以圣神和火为宗徒施洗，那时\BibleQuote{忽然从天上来了一阵响声，好像暴风刮来，充满了他们所在的全屋；又有散开的火舌显现给他们，停留在他们每人身上，他们都充满了圣神。}\BibleRef{宗 2:2}。\par

10. 任何人若不领受圣洗，就没有救恩；唯有殉道者除外，他们即使没有水也获得天国。因为当救主以十字架救赎世界时，肋旁被刺透，流出了血和水；这样，生活在和平时期的人可在水中受洗，而在迫害时期的人可在自己的血中受洗。因为救主也惯常称殉道为一种洗礼，说：\BibleQuote{你们能饮我所饮的杯吗？或者，你们能受我所受的洗吗？}\BibleRef{谷 10:38}。而殉道者们藉着\BibleQuote{成了一台戏，给世界、天使和世人观看}\BibleRef{格前 4:9}而宣认；你们也快要宣认——但现在还不是你们听这事的时候。\par

11. 耶稣藉着亲自受洗而圣化了圣洗。如果天主子受了洗，那么哪个虔敬的人会轻视圣洗呢？但祂受洗并不是为了领受罪的赦免，因为祂是无罪的；然而，无罪的祂受了洗，是为了赐给那些受洗者神圣而卓越的恩宠。因为\BibleQuote{儿女既然有同样的血肉，祂也照样取了同样的血肉}\BibleRef{希 2:14}，好使我们因分享了祂在肉身的临在，也能分享祂的神性恩宠；因此耶稣受了洗，好使我们藉着参与，再次获得救恩和尊荣。按约伯所说，水中曾有巨龙\BibleQuote{\Emph{吸约但河的水到口里}}\BibleRef{约 40:23}。因此，既然必须击碎巨龙的头，祂就下去并在水中捆住了强者，好使我们获得能力\BibleQuote{\Emph{践踏蛇和蝎子}}\BibleRef{路 10:19}。那兽巨大而可怕。\Emph{没有渔船能承载牠尾巴的一片鳞}；毁灭在牠面前奔跑，吞噬一切遇见者。生命与牠相遇，好使死亡的嘴从此被堵住，而我们所有得救的人都能说：\BibleQuote{死亡啊，你的胜利在哪里？死亡啊，你的刺在哪里？}\BibleRef{格前 15:55}。圣洗拔除了死亡的刺。\par

12. 因为你带着你的罪下到水里，但恩宠的呼求封印了你的灵魂，之后就不容你被可怕的巨龙吞没。你带着罪死于水中，却起来在正义中活了过来。因为如果你已\BibleQuote{与祂死亡的相似性结合}\BibleRef{罗 6:5}，你也将配得祂的复活。正如耶稣承担了世界的罪并死去，为藉着将罪置于死地而在正义中复活；同样，你下到水中，仿佛被埋葬在水里，正如祂被埋在岩石中，然后起来在新生活中行走。\par

13. 再者，当你堪当领受恩宠时，祂就赐给你力量与敌对权势搏斗。因为正如祂在受洗后四十天受试探（不是祂以前无法得胜，而是祂愿意按次序有秩序地完成一切），同样，你虽然在受洗前不敢与敌人搏斗，但在领受恩宠后，从此依靠\BibleQuote{正义的武器}\BibleRef{格后 6:7}而自信，就必须作战，若你愿意，还要宣讲福音。\par

14. 耶稣基督是天主子，但祂在受洗之前并未宣讲福音。如果主本身遵循了恰当的时间次序，祂的仆人岂可越序？那时\BibleQuote{\Emph{耶稣开始宣讲}}\BibleRef{玛 4:17}，当时\BibleQuote{\Emph{圣神藉着形体像鸽子一样降临在祂身上}}\BibleRef{路 3:22}；这并不是耶稣为了先看见祂，因为祂在形体降临之前就认识祂，而是为了使给祂施洗的若翰看见祂。因为他说：\BibleQuote{\Emph{我}也不认识祂，但那派遣我来以水施洗的，对我说：你看见圣神降下并停留在谁身上，那就是祂。}\BibleRef{若 1:33}。如果你也有真诚的虔敬，圣神也降临在你身上，父的声音从上面对你说——不是\Quote{这是我的爱子}，而是\Quote{这如今成了我的儿子}；因为\Quote{是}只属于祂，因为\Emph{在起初已有圣言，圣言与天主同在，圣言就是天主}。祂永远\Quote{是}，因为祂常是天主子；而你是\Quote{如今成了}，因为你没有儿子的本性，而是藉着义子地位获得。而你藉着进步获得恩宠。\par

15. 预备你灵魂的器皿，好使你成为天主的儿子，以及\Emph{天主承继者，与基督同为承继者}\BibleRef{罗 8:17}；如果你确实预备自己，好领受；如果你怀着信德走近，好得以成为信者；如果决心脱去旧人。因为凡你所行的一切都将获得赦免，无论是奸淫，或是通奸，或是其他任何形式的放纵。什么罪比钉死基督更大？但连这罪，圣洗也能净化。因为伯多禄对来到他那里的三千人，对那些钉死了主的人，当他们问他说：\Emph{诸位仁人弟兄，我们该作什么？}\BibleRef{宗 2:37}？因为这创伤巨大。伯多禄啊，你使我们想起我们的堕落，因为你说：\Emph{你们杀害了\OriginalTerm{生命的主宰}{Prince of Life}}。对如此大的创伤有什么药膏？对如此污秽有什么洁净？对如此丧亡有什么救恩？\Emph{悔改}，他说，\Emph{并各人以耶稣基督我们的主之名受洗，为赦免罪过，并要领受圣神的恩惠}。哦，天主不可言喻的慈爱！他们本无得救的希望，却竟被认为配得圣神。你看见圣洗的德能！如果你们中有人曾以亵渎的话钉死了基督；如果你们中有人因无知曾在人前否认了他；如果有人因恶行使教义受到亵渎；让他悔改并怀抱希望，因为同样的恩宠如今仍在。\par

16. \Emph{耶路撒冷，鼓起勇气；主将除去你的一切邪恶}。\Emph{主将以审判的神和焚烧的神，洗净他子女的污秽}\BibleRef{依 4:4}。\Emph{他将洒净水在你们身上，你们将洁净，脱离一切罪污}。\BibleRef{则 36:25}天使们将围绕你跳舞，并说：这位身穿白衣，倚靠着她所爱者上来的是谁？因为那从前是奴仆的灵魂，如今竟收纳她的主人自己为亲属；他接受真诚的心意，将回答说：\Emph{看，你是美丽的，我的爱人；看，你是美丽的：你的牙齿像刚剪毛的羊群}，（因为良心的宣认，并且）\Emph{它们个个都是双生的}；这是因为\OriginalTerm{双重恩宠}{twofold grace}，我是指那由水和圣神所完成的，或那由旧约和新约所宣告的恩宠。愿天主恩赐你们众人，在完成斋戒的课程后，能记住我所说的话，并在善工中结出果实，能无瑕地站在属灵新郎身边，并从天主获得罪过的赦免；愿光荣归于他，及子及圣神，至于永远。阿们。\par

\SourceRecord{Translated by Edwin Hamilton Gifford.；Nicene and Post-Nicene Fathers, Second Series；Vol. 7.；Edited by Philip Schaff and Henry Wace.；Buffalo, NY: Christian Literature Publishing Co.,；1894.；<http://www.newadvent.org/fathers/310103.htm>.}

% structure: p0001 source=h1 target=section reason=page-title

\section{教理讲授 四}

\Emph{论教义的十点}\par

\BibleRef{哥 2:8}\par

\BibleQuote{你们要谨慎，恐怕有人用哲学和虚妄的诡诈，按着人的传统，按着世俗的小学，把你们掳去，等等。}\par

恶习模仿德行，莠子竭力想被认作麦子，外表长得像麦子，但善于辨别的人从味道就能认出它们。\BibleQuote{魔鬼也装作光明的天使}\BibleRef{格后 11:14}，不是因为他能重新升到原来的地方，因为他使\Emph{他的心硬如铁砧}，从此不能悔改；而是为了把那些生活在天使般生活中的人包裹在盲目和不信的瘟疫状态中。\Emph{许多狼披着羊皮到处走动}，它们的披挂是羊的，但爪和牙不是：它们披着柔软的羊皮，用外表欺骗无辜的人，从牙尖上喷射出不虔敬的毁灭性毒液。因此我们需要天主的恩宠、清醒的心智和能看见的眼睛，免得因无知把莠子当作麦子吃而受害，免得把狼当作羊而成为它的猎物，免得把毁灭性的魔鬼当作仁慈的天使而被吞噬；因为正如圣经所说：\BibleQuote{它如同咆哮的狮子巡游，寻找可吞食的人。}\BibleRef{伯前 5:8}这就是教会训诫的原因，也是当前教导和所读课文的原因。\par

因为虔敬的方法包括这两件事：虔诚的教义和有德行的实践；教义没有善工不为天主所悦纳，同样，没有虔诚教义所成全的善工也不蒙天主悦纳。因为深知关于天主的教义却是一个卑鄙的奸淫者，有什么益处呢？同样，高尚地节制却不虔敬地亵渎，又有什么益处呢？因此，教义的知识是最宝贵的财富：同时也需要警醒的灵魂，因为有许多人通过哲学和虚妄的诡诈进行掠夺。\BibleRef{哥 2:8}希腊人一方面用圆滑的口舌引诱人，因为\BibleQuote{淫妇的嘴唇滴下蜂蜜}\BibleRef{箴 5:3}；而受割礼的人则用神圣的圣经欺骗前来的人，他们虽从幼年到老年研读圣经，却可怜地曲解它，在无知中变老。但异端者的子弟用甜言蜜语和圆滑的口舌欺骗无辜者的心，用基督的名字像蜂蜜一样包裹他们不虔敬教义的毒箭：关于所有这些，主说：\BibleQuote{你们要谨慎，免得有人迷惑你们}\BibleRef{玛 24:4}。这就是信经教导和阐释信经的原因。\par

但在将你们交付给信经之前，我认为现在最好先用一个简短的教义摘要；这样，要讲的众多内容以及整个神圣四旬期的漫长间隔，就不会使你们中较简单的人忘记；而是，现在以摘要方式播下一些种子，以后更广泛耕种时我们就不会忘记。但那些思想成熟、并且感官已经受过锻炼能辨别善恶\BibleRef{希 5:14}的人，请耐心聆听更适合儿童的内容，如同初级的奶食：这样，既可使需要教导的人受益，也可使有知识的人重新燃起对已知之事的记忆。\par

\Emph{一、论天主}\par

首先要放在你灵魂中作为基础的是关于天主的教义：天主是独一的、唯一非受生的、无始、无变易、无变化的；既非由另一个所生，也没有另一个在祂的生命中继承祂；祂既非在时间中开始生活，也永不终结；并且祂既是善的又是公义的；如果你曾听见异端者说，有一位公义的天主和另一位善良的天主，你就要立刻记住并认出差错的毒箭。因为有些人竟敢不虔敬地在他们的教导中分裂独一天主；有些说一位是灵魂的创造者和主宰，另一位是身体的创造者和主宰，这种教义既荒谬又不虔敬。因为当我们的主在福音中说：\BibleQuote{没有人能侍奉两个主人}，一个人怎能成为两个仆人的仆人呢？因此只有一个独一天主，灵魂和身体的创造者：独一的天地创造者，天使和总领天使的创造者：万物的创造者，但在万世以前只是一位父——独一的那位，祂的独生子，我们的主耶稣基督，借着祂创造了所有可见的与不可见的。\par

我们主耶稣基督的这位父不受任何地方的限制，也不小于天；但\BibleQuote{诸天是祂手指的化工，整个大地被祂握在手中}\BibleRef{依 40:12}：祂在万物之内，也在万物周围。不要以为太阳比祂更亮或与祂相等：因为最初形成太阳的那一位必然无可比拟地更大更亮。祂预知将要发生的事，比万物更有力，知道一切并随心所欲行事；不受任何必然事件顺序、出生、机遇或命运的支配；在一切事上完美，同样拥有一切德性的绝对形式，既不减少也不增加，但方式和状况恒常不变；祂为罪人准备了惩罚，为义人准备了冠冕。\par

6. 既然许多人以各种方式偏离了独一的天主，有的神化太阳，以至于当太阳西沉时，他们便在黑夜中无天主；有的神化月亮，以至于在白天没有天主；有的神化世界其他部分；有的神化技艺；有的神化各种食物；有的神化自己的快乐；有些疯狂追逐女人的人，竖起一尊裸女像，称之为\OriginalTerm{阿佛洛狄忒}{Aphrodite}，并在可见的形体中敬拜自己的情欲；还有的因黄金的光辉而眼花缭乱，将其及各种物质神化。——然而，若有人以天主的合一性为第一根基安放在心中，并信赖祂，他立刻根除了偶像崇拜的全部邪恶之果和异端者的谬误：因此，你要藉着信德将这信仰的第一教义作为根基安放在你的灵魂中。\par

\Emph{论基督。}\par

7. 也要相信天主子，独一的那一位，我们的主耶稣基督，祂是由天主所生的天主，由生命所生的生命，由光明所生的光明，在一切事上相似生祂的那一位，祂不是在时间中领受了存在，而是于万世之前、永恒地、不可理解地由圣父所生：那自存于位格中的天主的智慧、能力和正义：祂于万世之前坐在圣父的右边。\par

因为祂在右边的宝座并非像某些人所想的，因祂的坚忍而领受，仿佛在受难后天主给祂加冕；而是在祂整个存在中——那由永恒受生而有的存在——祂保持着王者的尊威，与圣父同坐，祂是天主、智慧和能力，如已所说；与圣父一同为王，为圣父创造万有，却丝毫不缺天主的尊威，并认识那生祂的，正如祂被那生祂的所认识；简而言之，你要记住福音中所记载的：\Quote{除了圣父，没有人认识圣子；除了圣子，也没有人认识圣父}。\par

8. 此外，你既不要将圣子与圣父分离，也不要因混淆而相信一种\OriginalTerm{圣子父性}{Son-Fatherhood}；但要相信，从独一的天主那里，有独一的独生子，祂是万世之前的天主圣言；不是那散发在空气中的出口之言，也不可比作非位格的话语；而是那位圣言、圣子，是一切有理性的受造物的创造者，那位聆听圣父并亲自说话的圣言。关于这些，若天主允许，我们将在适当的时候更详细地谈论；因为我们并没有忘记当前的目的，即对信德作一个概要性的介绍。\par

\Emph{论他的由圣童贞诞生。}\par

9. 那么，你要相信，天主的这独生子，为了我们的罪，从天降下到地上，取了与我们相同情欲的人性，由圣童贞和圣神所生，成了人，并非仅在外表和幻象中，而是真实地；也不是通过童贞如通过一个通道，而是从她真实地成了血肉，[并真实地用奶哺养]，并且真实地像我们一样吃，真实地像我们一样喝。因为如果降生成人是\OriginalTerm{幻影}{phantom}，那么救恩也是幻影。基督具有两性，在所见的是人，在未见的是天主；作为人，祂真实地像我们一样吃，因祂与我们有相同的肉身感觉；作为天主，祂用五个饼喂饱五千人；作为人，祂真实地死去；作为天主，祂复活那已死四天的人；作为人，祂在船上真实地睡觉；作为天主，祂在水面上行走。\par

\Emph{论十字架。}\par

10. 祂确实为了我们的罪被钉在十字架上。因为如果你否认这一点，那地方就明显反驳你，就是这有福的\Place{哥耳哥达}，我们现今为了那位在此被钉者而聚集在这里；并且自那以后，全世界都充满了十字架的木头碎片。但祂被钉并非因自己的罪，而是为了使我们从自己的罪中得释放。虽然作为人，祂那时\Quote{被人轻视}，并受笞，但祂被受造界承认为天主：因为太阳看见它的主受到侮辱，便暗淡颤抖，无法忍受这景象。\par

\Emph{论他的安葬。}\par

11. 祂作为人真实地被安放在岩石凿成的坟墓中；但岩石因祂而恐惧崩裂。祂下降到地下的区域，为的是也从那里救赎义人。因为请告诉我，你只愿活人——且他们大部分是不洁的——享受祂的恩宠，而不愿那些从亚当以来被囚禁已久的人现在获得自由吗？先知依撒意亚曾大声宣告了许多关于祂的事，难道你不愿那位君王下降去救赎祂的传报者吗？\Person{达味}在那里，\Person{撒慕尔}在那里，以及所有先知，连\Person{若翰}也在那里，他藉使者说：\BibleQuote{你就是要来的那一位，或者我们还要期待另一位？}\BibleRef{玛 11:3} 难道你不愿祂下降并救赎这样的一些人吗？\par

\Emph{论复活。}\par

12. 但祂降下到地下区域，又升上来了；而被埋葬的耶稣，第三天确实复活了。如果犹太人时常烦扰你，你就在此时问他们：\Person{约纳}第三天从鲸鱼腹中出来，基督第三天难道没有从地里复活吗？死人触摸\Person{厄里叟}的骨头就复活了，难道人类创造者因圣父的能力复活不是更容易得多吗？那么，祂确实复活了，并且复活后再次被门徒看见：十二位门徒是祂复活的见证人，他们作证不是用悦耳的言语，而是为复活的真理奋斗直至拷打和死亡。那么，\BibleQuote{凭两个或三个证人的口，每句话都得以成立}\BibleRef{申 19:15}，根据圣经，而且虽然有十二个人为基督的复活作证，你仍对祂的复活不信吗？\par

\Emph{论升天。}\par

13. 当耶稣完成了祂忍耐的历程，赎回了人类脱离他们的罪恶，祂再次升到天上，有云彩接祂上去；祂上升时，天使在祂身旁，宗徒们也在观看。若有人不信我所说的话，让他相信眼前所见事物的实际能力。所有君王死后，他们的权力随生命而消逝；但被钉十字架的基督却受全世界朝拜。我们宣讲被钉者，连魔鬼如今也战栗。许多人在不同时期被钉十字架；但有哪一个被钉者，呼求他的名曾驱走魔鬼？\par

14. 因此，我们不要以基督的十字架为耻；即使别人隐藏它，你却在额头上公开盖上它，使魔鬼看见这君王的标记就远远战栗逃跑。所以，在吃喝、坐卧、起身、说话、行走时，总之，在每一个行动中，都要作此记号。因为那位在此被钉十字架的，如今在天上。倘若祂被钉并埋葬后仍留在坟墓中，我们就有理由羞耻；但事实上，这位在哥耳哥达被钉十字架的，已经从东方的橄榄山升天去了。因为祂从此处下到阴府，又回到我们中间，再从我们这里升到天上，祂的父向祂说：\BibleQuote{你坐在我右边，直到我使你的仇敌作你的脚凳}。\par

论将来的审判。\par

15. 这位升了天的耶稣基督将再来临，不是从地上，而是从天上；我说\Quote{不从地上}，因为此时从地上将有许多假基督出现。因为正如你们所见，许多人已经开头说：\BibleQuote{我是基督}\BibleRef{玛 24:5}；并且那\Emph{招致荒凉的可憎之物}还将来到，妄自称基督。但你要期待真正的基督，天主的独生子，从今以后不再从地上，而是从天上降临，比任何闪电和光明的光辉更明亮地显现给众人，有天使护卫随行，为审判生者死者，并在一个没有终结的天上永恒国度中为王。因为关于这一点，我也请你确信，因为有许多人说基督的国度有终结。\par

论圣神。\par

16. 你也要相信圣神，并持守你对圣父和圣子所\Emph{领受的}同样意见，不要跟随那些对祂说亵渎话的人。但要学习：这一位圣神是唯一的，不可分的，具有多重能力；有许多行动，却并不分开；祂知道奥秘，\BibleQuote{洞察一切，就连天主的深奥也洞悉}\BibleRef{格前 2:10}；祂曾以鸽子的形状降临在主耶稣基督身上；祂曾在法律和先知中工作；祂如今也在洗礼的时刻印记你的灵魂；每个有理智的本性都需要祂的圣洁；若有人\BibleQuote{胆敢亵渎祂，必不得赦免，无论今世或来世}\BibleRef{玛 12:32}；\Quote{祂与圣父和圣子一同}享有天主性的光荣；\BibleQuote{上座者、宰制者、率领者、掌权者}也需要祂。\BibleRef{哥 1:16}因为只有一位天主，基督的父；只有一位主耶稣基督，唯一天主的独生子；只有一位圣神，万有的圣化者和神化者，祂曾在法律和先知中，在旧约和新约中说话。\par

17. 你要常将这印记放在心中，目前在我的讲论中以摘要的方式轻触，但若主允许，将尽我所能地以圣经的证明来陈述。因为关于信仰的神圣与神圣奥迹，连偶然的陈述也不可没有圣经而传达；我们也不可被单纯的动听和言辞的技巧所牵引。即使对向你们讲述这些事的我，也不要完全相信，除非你们从我宣布的事中得到了圣经的证明。因为我们所相信的这救恩，不倚靠巧妙的推理，而倚靠圣经的证明。\par

论灵魂。\par

18. 在认识了这可敬、光荣、至圣的信仰之后，还要进一步认识你自身是什么：作为人，你有两重本性，由灵魂和身体组成；正如刚才所说的，同一位天主是灵魂和身体的创造者。你也要知道，你有一个自主的灵魂，是天主最高贵的作品，按照其造物主的肖像受造；因天主赋予它不死性而是不死的；是一个有生命的、理性的、不朽坏的实体，因赐予这些恩赐的那一位；它有自由的能力去做它愿意做的事。因为你犯罪不是由于你的出生，也不是由于偶然的力量你犯奸淫，也不是像有些人胡说的那样，星辰的会合迫使你纵情放荡。你为什么不敢承认自己的恶行，却将罪责推给无辜的星辰？再不要听从占星家；因为圣经论到他们说：\BibleQuote{让那些观天象的站起来拯救你}，以及随后的话：\BibleQuote{看，他们都要像碎秸被火焚烧，不能救自己的灵魂脱离火焰}\BibleRef{依 47:13}。\par

19. 也要学习这一点：灵魂在进入此世之前未曾犯罪，而是无罪地进入，我们现今因自己的自由意志而犯罪。我恳求你，不要听信任何人曲解\BibleQuote{但我若做我不愿意的}\BibleRef{罗 7:16}；但要记住那说\BibleQuote{你们若愿意，且听从我，就必吃地上的美物；但若不愿听从，刀剑必吞灭你们等}\BibleRef{依 1:19}的那一位；并记住：\BibleQuote{你们从前怎样将肢体献给不洁和不法作奴仆，以致不法，现今也要照样将肢体献给义作奴仆，以致成圣}\BibleRef{罗 6:19}。也要记住圣经所说：\BibleQuote{他们既然不乐意存留天主在自己认识之中}\BibleRef{罗 1:28}，以及\BibleQuote{天主所显明的事，已显明在他们里面}\BibleRef{罗 1:19}；又：\BibleQuote{他们闭了自己的眼睛}\BibleRef{玛 13:15}。也要记住天主怎样再次指控他们，说：\BibleQuote{我栽种你如同全真的上等葡萄树，你怎么转向了我，成了野葡萄树呢？}\BibleRef{耶 2:21}\par

20. 灵魂是不死不灭的，所有灵魂无论男女都一样；因为只是身体的肢体有区别。并没有一类灵魂本性犯罪，另一类灵魂本性行义；两者都是出于选择，其灵魂的本体只有一种，且在所有灵魂中相同。然而我知道，我说了很多，时间已经很长了；但有什么比救恩更宝贵呢？你们不愿意费心在旅途中预备防备异端的干粮吗？你们不愿学习道路的岔径，以免因无知而坠入悬崖吗？如果你们的教师认为你们学习这些事是很大的收获，那么你们这些学习者不是应该欣然接受所告诉你们的众多事物吗？\par

21. 灵魂是自主的：虽然魔鬼能诱惑，但他没有能力违背意志去强迫。他向你呈现邪淫的念头：你若愿意，你就接受；你若不愿意，你就拒绝。因为如果你必然成为犯奸淫者，那么天主为何预备地狱？如果你本性行义而非出于意志，那么天主为何预备不可言喻的荣耀冠冕？绵羊是温顺的，但从未因其温顺而得冠冕：因为它的温顺并非出于选择，而是出于本性。\par

\Emph{论身体。}\par

22. 亲爱的，你们已学习了灵魂的本性，就目前时间而言已足够；现在当尽力也接受关于身体的道理。不要容忍那些说这个身体不是天主作品的人；因为他们相信身体独立于天主，灵魂住在其中如同在陌生的器皿里，便轻易滥用身体行邪淫。然而他们在这奇妙的身体上发现了什么缺点？它在俊美上缺少什么？在构造上哪一点不充满技巧？他们难道不该观察眼睛的光辉构造？耳朵如何倾斜放置，无阻碍地接收声音？嗅觉如何能分辨气味，感知气息？舌头如何服务于两个目的：味觉和言语的能力？肺如何隐蔽在看不见的地方，不停地呼吸空气？谁赋予了心脏不停的跳动？谁将血液分配成如此多的静脉和动脉？谁巧妙地将骨头与韧带编织在一起？谁将食物的一部分分配给我们身体，另一部分分离出来作适当的排泄，并将不美观的肢体隐藏在更美观的地方？当人类必然灭绝时，谁通过简单的交配使之永久延续？\par

23. 不要告诉我身体是罪的原因。因为如果身体是罪的原因，为什么死去的身体不犯罪？将剑放在刚死去的人右手上，不会发生谋杀。让各种美景经过一个刚死去的青年面前，不会产生不洁的欲望。为什么？因为身体本身不犯罪，而是灵魂通过身体犯罪。身体是灵魂的工具，如同衣服和袍子：若灵魂把它交给邪淫，它就受玷污；但若与圣洁的灵魂同住，它就成为圣神的宫殿。这不是我说的，而是宗徒保禄说的：\BibleQuote{难道你们不知道，你们的身体是住在你们内的圣神的宫殿吗？}\BibleRef{格前 6:19}。因此，要爱护你们的身体，如同圣神的宫殿。不要在邪淫中玷污你们的肉体：不要弄脏你们这件最美丽的袍子；如果你曾玷污过它，现在通过补赎洗净它：趁着还有时间，去洗濯吧。\par

24. 关于贞洁的道理，首先要注意的是独修者和贞女们，他们在世上保持天使般的生活；教会其余的人当跟随他们。因为你们，弟兄们，有伟大的冠冕存留：不要为一点小小的快乐失去极大的尊荣：听宗徒说话：\BibleQuote{免得有邪淫或亵圣的人，如厄撒乌，他为一碗食物出卖了自己长子的名分}\BibleRef{希 12:16}。你们既因宣认贞洁而被登记在天使册上，要小心不要因行邪淫而再次被涂抹。\par

25. 此外，另一方面，在持守贞洁时，不可对行走在较为卑微的婚姻道路上的人傲慢。因为正如宗徒所说：\BibleQuote{婚姻在众人中都当受尊重，床笫也不可玷污}\BibleRef{希 13:4}。你持守贞洁的人，难道不是由结婚的人所生的吗？因为你拥有金子，不可因此鄙视银子。但那些结婚而合法使用婚姻的人也该高兴；他们按照天主的安排结婚，而非出于放纵无度的欲望；他们承认禁欲的时期，\BibleQuote{以便专务祈祷}\BibleRef{格前 7:5}；他们在我们的集会中带着清洁的身体和清洁的衣服进入教会；他们结婚是为了生育子女，而非为了纵欲。\par

26. 但愿那些只结婚一次的人，也不谴责那些同意再婚的人：因为虽然节制是高贵而可赞美的事，但进入第二次婚姻也是允许的，以使软弱者不致陷入邪淫。因为宗徒说：\BibleQuote{对他们来说，如果他们能像我一样保持独身，那更好。但如果他们没有节制的恩赐，就让他们结婚吧：因为结婚比欲火中烧更好。}\BibleRef{格前 7:8} 但要让所有其他的行为都被远远抛弃：邪淫、奸淫以及各种放纵；要让身体为主保持纯洁，好使主也眷顾身体。要让身体以食物滋养，使之能存活并毫无阻碍地事奉；但不要让它沉溺于奢华。\par

\Emph{关于食物。}\par

27. 关于食物，你们当遵守这些规条，因为在饮食之事上也有许多人跌倒。因为有些人轻率地对待祭过偶像之物，而另一些人则克己自律，却指责那些吃的人；在饮食之事上，人的灵魂以各种方式被玷污，皆因不明了吃与不吃的有益理由。我们禁食，戒酒戒肉，并非因厌憎它们为可憎之物，而是因为我们期盼赏报；既轻视了感官之物，好能享受属神和理智的盛宴，并\Quote{现今含泪播种的，将在来世欢欣收割}。所以，不要轻视那些因身体软弱而进食的人；也不要指责那些\BibleQuote{因胃口和屡次患病而稍用酒}\BibleRef{弟前 5:23} 的人；既不要判定他们为罪人，也不要憎恶肉食为奇特之物；因为宗徒知道有这类人，他说：\BibleQuote{禁止嫁娶，戒绝食物——这些食物是天主创造，为使信者以感恩之心领受的。}\BibleRef{弟前 4:3} 因此，在戒绝这些事物时，不要像戒绝可憎之物那样，否则你们便没有赏报；而应视之为善物，为了摆在你面前的更卓越的属神事物而予以舍弃。\par

28. 你要稳妥地守护你的灵魂，免得你有时候吃祭过偶像之物：因为关于这类食物，不仅我现今如此，而且此前宗徒们和这教会的主教雅各伯也曾殷切关注。宗徒和长老们写信给所有外邦人，要他们首先\Quote{戒绝祭偶像之物}，然后也\Quote{戒绝血和勒死之物}。因为许多人本性野蛮，像狗一样生活，模仿最凶猛的野兽舔血，又贪婪地吞食勒死之物。但你，基督的仆人，进食时当心怀敬意。关于食物，就说到这里。\par

\Emph{论衣着。}\par

29. 但你的衣着应当朴素，不是为了装饰，而是为了必要的遮盖；不是助长你的虚荣，而是为了御寒和遮掩身体的不雅；免得在遮掩不雅的借口下，因奢侈的服装而陷入另一种不雅。\par

\Emph{论复活。}\par

30. 我恳求你，要温柔地对待这身体，并明白你将与这身体一同复活，以接受审判。但若有不相信的念头潜入你心，以为这事不可能，就凭你自身所发生的事来判断那未见之事吧。因为请告诉我：一百年前或更久以前，你自身何在？你是从何等微小卑微之物，长成如此高大的身材和如此美丽的尊荣？那么，那位使无中生有的，岂不能使那已存在而又腐朽的再次兴起吗？那位为我们的缘故使每年死去的种子复活的，岂会难以使我们复活——为了我们，那种子也被复活了？你岂不见树木一连数月既无果实也无叶子？但冬天一过，它们便如从死中复活，整体恢复生机。难道我们不会更容易地返回生命吗？梅瑟的棍杖因天主的意愿变成了不熟悉的蛇的本性；难道人陷入死亡后，就不能恢复自身吗？\par

31. 不要理会那些说这身体不会复活的人；它确实要复活。依撒意亚作证说：\BibleQuote{死人要复活，在坟墓中的要苏醒}\BibleRef{依 26:19}；达内尔说：\BibleQuote{许多睡在尘土中的人要醒来，有的得到永生，有的永远蒙羞。}\BibleRef{达 12:2} 然而，虽然复活是众人的共通之事，但复活并非对所有人都相同：因为我们所得的身体都是永恒的，但并非所有人都得相似的身体：义人领受它们，好能在永恒中加入天使的行列；罪人领受它们，好能永远忍受罪的痛苦。\par

\Emph{论洗礼池。}\par

32. 为此，主出于祂的慈爱，预先为我们安排了圣洗中的悔改，好使我们卸下罪恶的主要——甚至全部重担，并藉圣神领受印记，成为永生的承继者。既然前日我们已充分论及洗礼池，现在就回到入门教导的其余部分。\par

\Emph{论圣经。}\par

33. 如今，这些新旧约中由天主默感而写的圣经教导了我们。因为两约的天主是同一的，祂在旧约中预言了在新约中显现的基督；祂藉着法律和先知引领我们进入基督的学校。\Quote{因为在信德来临以前，我们被看守在法律之下}，并且\Quote{法律就成了我们的启蒙师，领我们归于基督}。倘若你听见任何异端者诋毁法律或先知，就用救主的声音回答说：耶稣\Quote{来不是为废除法律，而是为成全}。\BibleRef{玛 5:17} 你也要殷勤地从教会学习旧约和新约中的书卷。并且，请不要阅读任何伪经；因为既然你不认识那些大家都认可的书卷，何必徒劳地为那些有争议的书卷费心呢？请阅读神圣的圣经，就是旧约的二十二卷书，这些是由七十二贤士翻译的。\par

34. 因为在马其顿王亚历山大去世，他的王国分裂为四个主要部分——巴比伦、马其顿、亚洲和埃及之后，统治埃及的其中一位国王托勒密·斐拉德尔弗斯，是一位酷爱学问的君王。他在搜集各地书籍时，从他的图书馆馆长德米特里·法勒雷乌斯那里听说了关于法律和先知的神圣经书，便认为更崇高的做法不是用武力从拥有者那里夺书，而是通过礼物和友谊来争取他们；他知道强索之物常被篡改，因为是不情愿地给予，而自愿供给的却是真诚地自由赠送。于是他派遣使者到当时的大司祭厄肋阿匝尔那里，献上许多礼物给耶路撒冷的圣殿，并请他派遣以色列十二支派中各六位译者来进行翻译。然后，为了检验这些书是否神圣，他采取了预防措施，使被派来的人不能相互串通：他在名为法罗斯的岛上——该岛位于亚历山大港对面——为每位到来的译者分配了单独的房间，并将全部圣经交给每人翻译。当他们在七十二天内完成了任务后，他将所有译本收集起来——这些译本是他们在不同房间内完成的，彼此没有沟通——发现它们不仅意义相符，甚至连用词也一致。因为这一过程并非文字游戏，也不是人的设计；而是由圣神所说出的神圣经书的翻译，是由圣神所完成的。\par

35. 你要阅读这二十二卷书，但不要与伪经有任何关联。只专心研读那些我们在教会中公开诵读的书。宗徒们和古代的主教们，即传授这些书卷的教会首领，比你更明智、更虔诚。因此，作为教会的孩子，你不要触碰她的规条。关于旧约，正如我们所说，要研读二十二卷书，如果你渴望了解它们，请努力记住它们的名称，如同我列举它们一样。因为法律之书首先是梅瑟的五卷书：\WorkTitle{创世记}、\WorkTitle{出谷纪}、\WorkTitle{肋未纪}、\WorkTitle{户籍纪}、\WorkTitle{申命纪}。接着是农的儿子若苏厄，以及\WorkTitle{民长纪}（包括\WorkTitle{卢德传}），算作第七卷。在其他历史书中，\WorkTitle{列王纪}上、下在希伯来人中是一卷书；\WorkTitle{列王纪}三、四也是一卷。同样，\WorkTitle{编年纪}上、下在他们那里是一卷；\WorkTitle{厄斯德拉}上、下算为一卷。\WorkTitle{艾斯德尔传}是第十二卷；这些是历史书。而那些以诗歌形式写成的有五卷：\WorkTitle{约伯传}、\WorkTitle{圣咏集}、\WorkTitle{箴言篇}、\WorkTitle{训道篇}和\WorkTitle{雅歌}，这是第十七卷。此后是五卷先知书：小先知十二卷为一卷，\Person{依撒意亚}一卷，\Person{耶肋米亚}一卷（包括\WorkTitle{巴路克}、\WorkTitle{哀歌}和书信），然后是\WorkTitle{厄则克耳}和\WorkTitle{达尼尔书}，这是旧约的第二十二卷。\par

36. 至于新约，只有四部福音书，因为其余的都是标题虚假且有害的。摩尼教徒也写了一部按\Person{多默}的福音，它被福音标题的香气所浸染，却腐蚀了单纯之人的灵魂。也要接受十二宗徒的\WorkTitle{宗徒大事录}；此外，还有\Person{雅各伯}、\Person{伯多禄}、\Person{若望}和\Person{犹大}的七封公函；以及作为这一切的封印和门徒的最后著作，\Person{保禄}的十四封书信。但让其余的一切都放在次要位置。凡是在教会中不诵读的书卷，你连私下也不要读，正如你听我说过的那样。这些就是关于这些主题的内容。\par

37. 但要躲避每一个魔鬼的作为，不要相信那背道的蛇，它从善的本性中转变是由于它自己的自由选择；它能说服自愿的人，但不能强迫任何人。也不要注意星象的观测、占卜、预兆，或希腊人的虚构占卜。巫术、邪术和招魂术的恶行，连听都不要听。要远离每一种无节制，不要让自己沉溺于暴食或放荡，要超越一切贪婪和高利贷。不要冒险参加外邦人的公共集会，也不要在生病时使用护身符；还要躲避所有酒馆的粗俗行为。不要陷入撒玛黎雅人的派别或犹太主义；因为耶稣基督已经赎回了你。要远离一切安息日的遵守，以及把任何无足轻重的食物称为\Quote{不洁或不净}。特别要憎恶所有邪恶异端者的集会；并以各种方式确保你的灵魂安全：藉着禁食、祈祷、施舍和阅读天主的圣言；这样，在肉身中度过余生的清醒和虔诚的教训中，你就能享受来自圣洗的唯一救恩；如此，由天主父登记在天军中，你也可能被视为配得上天上的冠冕，在我们的主耶稣基督内，愿光荣归于祂，直到永远。阿们。\par

\SourceRecord{Translated by Edwin Hamilton Gifford.；Nicene and Post-Nicene Fathers, Second Series；Vol. 7.；Edited by Philip Schaff and Henry Wace.；Buffalo, NY: Christian Literature Publishing Co.,；1894.；<http://www.newadvent.org/fathers/310104.htm>.}

% structure: p0001 source=h1 target=section reason=page-title

\section{\WorkTitle{第五讲慕道释义}}

\Emph{论信德}\par

\BibleRef{希 11:1-2}\par

\BibleQuote{信德是所望之事的实质，是未见之事的确证。因这信德，先人们都得了称赞。}\par

一、主将你从慕道者的等级迁升到信者的等级，赐予你何等大的尊荣，宗徒\Person{保禄}表明说：\BibleQuote{天主是信实的，你们原是被祂所召，与祂的圣子耶稣基督共享团契}\BibleRef{格前 1:9}。既然天主被称为信实的，你领受这个称号时，也领受了极大的尊荣。因为正如天主被称为善、公义、全能、宇宙的创造者，祂也被称为信实的。因此，你要思量你将要上升到何等尊荣，因为你将分享天主的称号。\par

二、因此，这里还要求你们每一个人在良心上都要显出信实：\BibleQuote{因为忠信的人难以找到}\BibleRef{箴 20:6}。这不是要你向我显明你的良心，因为你\Emph{不该受人的判断}；而是要你向\Emph{那洞察肺腑心肠}、\Emph{知道人的思念}的天主显明你信德的真诚。忠信的人是一件伟大的事，是富人中最富有的。因为\Emph{全世界财富都属于忠信的人}，因为他藐视并践踏它。那些外表富有、拥有许多财产的人，灵魂却是贫穷的：因为他们聚敛得越多，就越渴望那尚欠缺的。但忠信的人，最奇妙的矛盾，在贫穷中却是富有的：因为他知道我们只要有\BibleQuote{食粮和衣服，就当知足}\BibleRef{弟前 6:8}，他已将财富踩在脚下。\par

三、并非只有我们这些背负基督之名的人，信德才有极大的尊荣；在世上所成就的一切事情，即使是那些远离教会的人所成就的，也是藉着信德。\par

藉着信德，婚姻的法律将素不相识的人结合为一；因着婚姻契约的信实，陌生人成为另一个人的伴侣，共享其人及财产。藉着信德，农业得以维持，因为若不信将获得收成，人便不能忍受劳苦。藉着信德，航海的人信赖最单薄的木板，将最坚实的土地换成波浪不息的海水，将自己托付于不确定的希望，且带着比任何铁锚更坚定的信德随行。因此，人类大部分事务皆因信德而维系；不仅在我们中间有这种信念，如前所述，在那些教外之人中也是如此。因为他们若不接受圣经，却提出自己的某些学说，连这些也是凭信德接受的。\par

四、今天所宣读的读经也邀请你们走向真信德，它向你们展示了你们必须怎样悦乐天主的途径：因为它肯定\BibleQuote{没有信德，不可能中悦天主}\BibleRef{希 11:6}。因为一个人何时才决心事奉天主，除非他相信\Emph{祂是赏报者}？何时少女选择贞女的生活，或青年持守节制，除非他们相信为贞洁有\BibleQuote{不朽的荣冠}\BibleRef{伯前 5:4}？信德是照亮每个良心的眼睛，并赋予理智；因为先知说：\Emph{你们若不相信，必不得了解}。\par

信德\BibleQuote{堵住了狮子的口}\BibleRef{希 11:34}，如同\Person{达尼尔}的情形：因为圣经论到他写道：\BibleQuote{达尼尔从坑里被拉上来，身上毫无损伤，因为他信靠了他的天主}\BibleRef{达 6:23}。还有什么比魔鬼更可怕的呢？然而，即使对抗它，我们也没有别的盾牌，只有信德——一面无形的盾牌，对抗无形的仇敌。因为它射出各种火箭，\Emph{在黑夜中射中}那些不警惕的人；但由于仇敌无形，我们有信德作为坚固的甲胄，正如宗徒所说：\BibleQuote{此外，拿起信德的盾牌，使你们能以此扑灭那恶者的一切火箭}\BibleRef{弗 6:16}。魔鬼常射出卑劣情欲的火箭；但信德提示审判的景象，使心思冷静，扑灭火箭。\par

五、关于信德有许多可说的，即使用一整天的时间也不足以完全描述。现在我们只满足于以\Person{亚巴郎}为例，取自旧约中的一个榜样，因为我们已藉着信德成了他的子女。他不仅因行为成义，也因信德成义：因为他虽然行了许多善事，但只有在相信时才被称为天主的朋友。此外，他的一切行为都是在信德中完成的。藉着信德，他离开了父母；藉着信德，他离开了故乡、地方和家园\BibleRef{希 11:8}。同样，你也要像他那样成义。他的身体对于生育已经如同死了，他的妻子\Person{撒辣}也已年老，没有希望再得子女。天主应许这老人一个儿子，\Person{亚巴郎}\BibleQuote{虽然考虑到自己的身体如同已死，信心却不软弱}\BibleRef{罗 4:19}，他不顾自己身体的软弱，而顾念那应许者的能力，因为\BibleQuote{他相信那应许者是信实的}\BibleRef{希 11:11}，因而超乎一切期望，从如同已死的身体得到了儿子。而且，当他得了儿子之后，奉命将他献上时，尽管他曾听到\Quote{藉以撒你的后裔才将传名}的话，他还是准备将他的儿子、他的独生子献给天主，相信\BibleQuote{天主有能力叫人从死者中复活}\BibleRef{希 11:19}。他捆了儿子，放在木柴上，在心意中确实献上了他，但因天主的仁慈，以一只羔羊代替了他的儿子，他得以领回活着的儿子。在这些事上，他是忠信的，遂被盖上了义德的印记，\BibleQuote{领受了割损作为他未受割损时因信德而获得的正义的印记}\BibleRef{罗 4:11}，并领受了\BibleQuote{他要成为多国之父}\BibleRef{创 17:5}的应许。\par

6. 那么，让我们看看亚巴郎如何成为多国之父。\BibleRef{罗 4:17} 他明认是犹太人的父亲，是按肉身传承的。但如果我们坚持按肉身传承，我们便不得不说那神谕是假的。因为按肉身，他不再是我们众人的父亲；但他的信德榜样使我们所有人都成为亚巴郎的子孙。如何？以何种方式？在人看来，死人复活是不可信的；同样，年迈如同已死的人生育后代也是不可信的。但当基督被宣讲为被钉在十字架上、死而复活时，我们相信了。因此，因我们信德的相似，我们被收纳为亚巴郎的义子。然后，紧随我们的信德，我们如同他一样领受了属神的印记，藉着圣洗被圣神施行割损，不是在肉身的包皮上，而是在心里，按照耶肋米亚说的话：\BibleQuote{你们要在你们心的包皮上为天主行割损}：并按照宗徒说的话：\BibleQuote{在基督的割损中，与祂同葬于圣洗}，以及其余经文。\BibleRef{哥 2:11}\par

7. 我们若保持这信德，就必免受定罪，且以各种美德为装饰。因为信德的力量如此之大，甚至能使人在海面上行走。伯多禄是像我们一样的人，由血肉组成，靠同样食物生活。但当耶稣说：\BibleQuote{来}\BibleRef{玛 14:29}，他就相信了，并在水面上行走，发现他的信德在水上比任何陆地更安全；他那沉重的身体被信德的浮力托住。但当他相信时，他在水上脚步稳妥；然而当他怀疑时，立刻开始下沉：因为随着他的信德逐渐松懈，他的身体也随之被拉下去。耶稣看见他的危难——祂医治我们灵魂的危难——就说：\BibleQuote{小信德的人哪，你为什么怀疑？}\BibleRef{谷 14:31} 然后，被那抓住他右手的主鼓起力量，他刚一恢复信德，就被主拉着手，重新开始在水面上行走：因为福音隐约提到这事，说：\BibleQuote{他们上了船}。因为经上没有说伯多禄游过去上了船，而是让我们明白，他走回了与去见耶稣相同的距离，然后上了船。\par

8. 是的，信德有如此大的力量，以致不仅是信者本人得救，有些人还因别人相信而得救。葛法翁的瘫子不是信者，但抬他来并从屋顶缒下他的人相信了\BibleRef{谷 2:4}：因为病人的灵魂与他的身体一同患病。不要以为我无故指责他：福音本身说，耶稣\BibleQuote{看见}——不是他的信德，而是\BibleQuote{他们的信德}——就对瘫子说：起来！抬的人相信了，瘫子就享受了痊愈的恩惠。\par

9. 你愿意更清楚地看见，有些人因别人的信德而得救吗？拉匝禄死了\BibleRef{若 11:14}：一天、两天、三天过去了：他的筋腱已经腐烂，朽坏已在侵蚀他的身体。一个死去四天的人怎能相信，并为自己恳求救赎者呢？但死人缺乏的，由他真正的姐妹们补足了。因为当主来到时，那姐妹俯伏在祂面前，当祂说：\BibleQuote{你们把他放在哪里？}她回答说：\BibleQuote{主，现在他必发臭，因为已经四天了}，主就说：\BibleQuote{你若相信，就必看见天主的荣耀}；这等于说：你补足死人的信德缺失。姐妹们的信德有如此大的力量，甚至使死人从地狱之门召回。那么，既然人能因着彼此相信而使死人复活，难道你，若真诚地为自己相信，不更得益处吗？不，即使你不信或小信，主是爱人的；祂因你的悔改而俯就你：只要你以诚实的心说：\BibleQuote{主，我信，求祢补助我的不信}\BibleRef{谷 9:24}。但如果你认为你确实有信德，却还没有圆满的信德，你也需要像宗徒们一样说：\BibleQuote{主，请增加我们的信德}\BibleRef{路 17:5}：因为有一部分是你自己的，但更大一部分是来自祂。\par

10. 因为\Quote{信德}一词，就言语形式而言是单一的，但有两个不同的含义。有一种信德是教义性的，涉及灵魂对某一特定观点的同意：这对灵魂是有益的，如主说：\BibleQuote{那听我的话，并相信派遣我者，就有永生，且不受审判}\BibleRef{若 5:24}：又说：\BibleQuote{那信从子的人，不受审判，已由死亡进入生命了}。哦，天主伟大的慈爱！义人多年来讨祂喜悦，但他们在多年善行中所获得的，如今耶稣在一小时内就赐给了你。因为如果你相信耶稣基督是主，且天主使祂从死者中复活，你必得救，并且被那曾领强盗进入乐园的祂带入乐园。不要怀疑这是否可能；因为在这神圣的哥耳哥达，祂在一小时的信德后拯救了强盗，同样在你相信时也必拯救你。\par

11. 但还有第二种信德，是由基督作为恩宠的恩赐所赐予的。\BibleQuote{因为这人藉圣神获得智慧的言语，另一人却藉同一圣神获得知识的言语；另一人信德，由同一圣神；另一人治病的奇恩}\BibleRef{格前 12:8}。这由圣神恩宠所赐的信德不仅是教义性的，也施行超乎人能力的事。因为凡有这信德的，\BibleQuote{必对这座山说：从这里移到那里，它就会移开}\BibleRef{谷 11:23}。因为无论谁凭信心说这话，\BibleQuote{相信必成就，心里不疑惑，就领受了那恩宠}。\par

关于这信德，经上说：\Emph{\BibleQuote{如果你们有信德像芥子那样大}}。\BibleRef{玛 17:20} 因为正如芥子虽小，但其作用如火，虽种在狭小之处，却能长出大枝，长成后甚至能遮蔽飞鸟\BibleRef{玛 13:32}；同样，信德在瞬息之间，在灵魂内产生最大的效果。因为灵魂被信德光照后，便看见天主，尽可能瞻仰天主，并遨游宇宙的边界，在现世终结之前，就已经看见审判和应许的赏报。所以，你要因自己而拥有对祂的信德，好使你也从祂那里获得那行超人之事的信德。\par

12. 但是在学习并宣认这信德时，务要获得并保持那唯一由教会现在传授给你们的、从全部圣经中坚固建立起来的信德。因为并非所有人都能阅读圣经，有些人因缺乏学识而受阻，另一些人因缺乏闲暇，为使灵魂不致因无知而丧亡，我们把信德的全部道理概括为几句话。这份摘要，当我在诵念时，我希望你们既牢记在心，又彼此勤加演练，不是写在纸上，而是凭记忆铭刻在你们的心版上，练习时注意不要让任何望教者偶然听到传授给你们的事。我还希望你们一生以此作为储备，此外不接受别的，即使我们自己改变并反对现在的教导，或是一位叛逆的天使，\Emph{\BibleQuote{伪装成光明的天使}}\BibleRef{格后 11:14}想要迷惑你们。\Emph{\BibleQuote{即使是我们，或是从天上来的天使，若向你们宣讲的福音与我们先前所宣讲的不同，当受诅咒。}}\BibleRef{迦 1:8} 所以现在，请听我简单地诵念\WorkTitle{信经}，并牢记在心；但在适当的时候，期待从圣经中获得每项内容的证实。因为信德的条文并非按人的意愿编写，而是从全部圣经中收集的最重要的点，构成信德的完整教导。正如芥子一粒小小的种子包含许多枝条，同样，这信德用几句话就包含了旧约和新约中一切虔诚的知识。所以，弟兄们，要谨慎，持守你们现在领受的传授，并将它们写在你心版上。\par

13. 你们要敬谨守护它们，免得仇敌突然劫掠任何懈怠的人，或免得异端者歪曲传授给你们的任何真理。因为信德就像把钱存入银行，正如我们现在所做的；但天主却要向你们索要存款的账目。\Emph{我在天主面前}，正如宗徒所说，\Emph{并在基督耶稣面前，祂在般雀比拉多面前作过美好的宣认，我嘱咐你}，要保守所托付给你的这信德，\Emph{毫无玷污，直到我们的主耶稣基督显现}。现在交托给你的是生命的宝藏，主在祂显现时要索回这笔存款，\Emph{到了时期，祂将显示那当受称颂的、唯一的全能者，万王之王，万主之主；唯有祂是不死的，居住在不可接近的光中，没有人见过，也不能见祂。愿尊威、光荣和权能归于祂}\BibleRef{弟前 6:15}，至于无穷之世。阿们。\par

\SourceRecord{Translated by Edwin Hamilton Gifford.；Nicene and Post-Nicene Fathers, Second Series；Vol. 7.；Edited by Philip Schaff and Henry Wace.；Buffalo, NY: Christian Literature Publishing Co.,；1894.；<http://www.newadvent.org/fathers/310105.htm>.}

% structure: p0001 source=h1 target=section reason=page-title

\section{教理讲授第六讲}

\Emph{关于天主的唯一性。论信经条文：\Quote{我信唯一的天主}。也论异端。}\par

\BibleRef{依 45:16-17}（七十贤士译本）\par

\BibleQuote{众海岛啊，你们要自洁归于我。以色列因上主获得拯救，获得永久的救恩；他们必不羞愧，也永不蒙羞。等等。}\par

1. \BibleQuote{愿我们的主耶稣基督的天主和父受赞美}。\BibleRef{格后 1:3} 愿他的独生子也受赞美。因为对于天主的思念，当立即与对父的思念相连，好使归与父和子的荣耀不可分割。因为父的荣耀并非一种，子的又是一种，而是同一的，因为他是父的独生子；父受赞美时，子也同享赞美，因为子的荣耀出自父的尊荣；同样，子受赞美时，如此大恩的父也大受尊崇。\par

2. 虽然思想极为迅速，但舌头仍需要言语和一系列的中介词句。因为眼睛瞬间即能囊括\Quote{星辰乐队}的众多星体；但若有人想逐一描述它们，哪个是晨星，哪个是昏星，哪个是其中的每一颗，他便需要许多言语。同样，心思在极短的时间内囊括大地、海洋和宇宙的所有界限；但瞬间所想到的，却需用许多言语来描述。虽然我所举的例子很有力，但毕竟还是软弱不足的。因为关于天主，我们所说的并非我们应说的一切（因为那只有他自己知道），而只是人性能力所领受的，以及我们的软弱所能承担的。因为我们并非解释天主是什么，而是坦诚承认我们没有关于他的确切知识。因为在天主的事上，承认无知便是最好的知识。因此，你们要与我一起赞美上主，让我们一同高举他的名——我们全体共同赞美，因为一人单独无力胜任；不，即使我们全都联合起来，也仍不能如应有的那样赞美。我所说的不仅是你们在场的这几位，即使全世界整个教会的所有儿女，无论是现在的还是将来的，都聚集在一起，也不能相称地歌颂他们的牧人。\par

3. 亚巴郎是一位伟大、可敬的人，但只是相对人而言伟大；当他来到天主面前时，他坦诚地说出了真理：\BibleQuote{我是尘土和灰烬}。\BibleRef{创 18:27} 他没有只说\Quote{尘土}就停止，免得他以这个巨大元素的名号自称；而是加上了\Quote{和灰烬}，以表明他易朽和脆弱的本性。他说，有什么比灰烬更小更轻呢？你们想，从灰烬到房屋，从房屋到城市，从城市到省份，从省份到罗马帝国，从罗马帝国到整个大地及其所有边界，再以整个大地到其所蕴含的苍天——大地与苍天的比例，如同轮心与整个轮缘的比例，因为大地与苍天相比不过如此；再思量：这第一重可见的苍天小于第二重，第二重又小于第三重，因为圣经就这样称呼它们，并非只有这么多，而是我们只需知道这么多。当你思想中巡视了所有苍天，即使苍天也不能相称地赞美天主，不，即使它们以比雷鸣更响的声音呼喊。如果这些伟大的穹苍不能相称地歌颂天主，那么\Quote{尘土和灰烬}，这存在中最小最微之物，何时才能向天主献上相称的赞美之歌，或相称地谈论\BibleQuote{那坐在大地之圈上，地上的居民好像蝗虫}的那位呢？\BibleRef{依 40:22}\par

4. 若有人试图谈论天主，让他先描述大地的边界。你住在大地上，却不知道你所居住的这大地的界限：你怎能对你的创造者形成相称的思想呢？你看见星辰，却看不见它们的创造者；数算这些可见的星辰，然后描述那不可见者，\Quote{他数点星辰的数目，一一称其名}。近日大雨倾盆，几乎毁了我们；你数算这城中雨滴的数目？不，我不说全城，你数算自家屋顶上在一个小时内的雨滴吧，如果你能的话，但你却不能。那么，请学到自己的软弱；从这件事上认识天主的全能：因为\Quote{他数过雨滴}，这雨滴降在整个大地上，不仅现在，而且永远。太阳是天主的工程，虽然巨大，但与整个苍天相比只是一点；你先定睛凝视太阳，然后好奇地审视太阳之主。\BibleQuote{不要探究过于深奥的事，也不要查考超越你能力的事：你要思想所命令的事}。\BibleRef{德 3:21}\par

5. 但有人会说：如果天主的本性是不可理解的，你为什么还谈论这些事呢？那么，既然我不能喝干整条河流，难道连适量取用对我有益的水也不行吗？既然我这双眼睛不能饱览整个太阳，难道连为了满足需要而看它一眼也不能吗？或者，因为我进入一个大花园，不能吃掉所有出产的水果，你就让我完全空着肚子离开吗？我赞美、光荣那造我们的主；因为有一条神圣命令说：\BibleQuote{一切有气息的，都要赞美上主}。我现在试图光荣上主，而不是描述他，同时知道我将不能相称地光荣他，但仍认为即使尝试也是一件虔诚工作。因为主耶稣鼓励我的软弱，说：\BibleQuote{从来没有人见过天主}。\par

6. 但有人会说：经上不是记载着\BibleQuote{小孩的天使在天上，常见我在天之父的面}(\BibleRef{玛 18:10})吗？是的，但天使看见天主，并非看见祂的本体，而是按他们各自所能承受的程度。因为耶稣亲自说：\BibleQuote{这不是说有人看见过父，只有那从天主来的，祂看见过父}(\BibleRef{若 6:46})。因此，天使看见的是他们所能承受的，总领天使看见的是他们所能承受的；上座者和宰制者看见的比前者多，但仍旧不及天主的尊荣；因为只有圣子同圣神能正当地看见祂：\BibleQuote{圣神洞察一切，连天主的深奥也洞悉}(\BibleRef{格前 2:10})；同样，独生子同圣神也完全认识父，因为祂说：\BibleQuote{除了子，没有人认识父；除了子和子所愿意启示的人，也没有人认识父}(\BibleRef{玛 11:27})。因为祂圆满地看见，并按各人所能承受的，藉着圣神启示天主；因为独生子同圣神共受父的天主性。那被生的认识那生祂的，那生祂的也认识那被生的。既然天使有所不知（因为如上所述，独生子藉圣神按各人能力启示他们），人也不该羞于承认自己的无知。我现在所说的，是人常说的；但我们如何说话，我们却说不清：那赐给我们言语的，我怎能说明祂呢？我有灵魂，却不能说出它的特性，我又怎能描述那赐予它的一位呢？\par

7. 为孝爱而言，我们只需知道我们有一位天主就够了；祂是独一的天主，生活的、永生的天主；常是一样；祂没有父，没有比祂更大的，没有后继者将祂赶出祂的王国；祂名称众多，能力无限，本体单一。因为虽然祂被称为善、义、全能者、\OriginalTerm{万军的上主}{Sabaoth}，祂并不因此多样多变；而是同一位、同一的，发出无数神性作为，并非在此超越、在彼欠缺，而是在一切事上都一样。祂不只在慈爱上伟大，在智慧上微小，而是智慧与慈爱同等能力；并非一部分看见，一部分无眼；而是全是眼，全是耳，全是理智；不像我们一样一部分知道、一部分不知道；因为这样的说法是亵渎的，与天主的本体不相称。祂预知未来的事；祂是圣的、全能的，在良善、尊威和智慧上超越一切；关于祂，我们既不能说开始，也不能说形象或形状。因为圣经说：\BibleQuote{你们从来没有听见祂的声音，也没有看见祂的形状}(\BibleRef{若 5:37})。因此梅瑟也对以色列人说：\BibleQuote{你们要谨慎，保全你们的灵魂，因为你们没有看见什么形状}(\BibleRef{申 4:15})。因为若完全不能想象祂的相似物，思想又如何能接近祂的本体呢？\par

8. 许多人有过许多想象，但全都失败了。有人以为天主是火；有人以为祂好像有翅膀的人，因为误解了一句真经文：\Quote{祢将把我藏在祢翼荫下}。他们忘记了我们的主耶稣基督——独生子——也以类似的方式论及自己，对耶路撒冷说：\BibleQuote{我多少次愿意聚集你的子女，如同母鸡把小鸡聚集在翅膀底下，但你却不愿意}(\BibleRef{玛 23:37})。因为天主的保护能力被比作翅膀，但他们不明白这点，就降到人的层次，以为那位不可揣测者具有人的形象。另有人竟敢说祂有七只眼，因为经上记着：\BibleQuote{主的七只眼遍察全地}(\BibleRef{匝 4:10})。若祂只有七只眼围绕祂的部分，那么祂的看见就是部分的，而非完全的；但这样说天主是亵渎的；因为我们应当相信天主在一切事上都是完美的，按我们救主的话说：\BibleQuote{你们在天上的父是完美的}(\BibleRef{玛 5:48})——在看见上完美，在能力上完美，在伟大上完美，在预知上完美，在良善上完美，在正义上完美，在慈爱上完美；不为任何空间所局限，却是所有空间的创造者，存在于万有之中，不被任何事物局限。\BibleQuote{上天是我的宝座}，但坐于其上的更高；\BibleQuote{大地是我的脚凳}(\BibleRef{依 66:1})，但祂的能力达到地下的事物。\par

9. 祂是独一的，无所不在，洞察一切，感知一切，藉着基督创造一切：\BibleQuote{万物是藉着祂造的；凡受造的，没有一样不是藉着祂造的}(\BibleRef{若 1:3})。祂是诸善之泉源，丰富而不竭，福泽之河流，永恒之光，恒久不灭之光辉，不可战胜之能力，却俯就我们的软弱；我们连听闻祂的名都不敢。约伯说：\Quote{你能寻着上主的踪迹吗？或能测透全能者所做的最小的事吗？}若祂的最小工作都不可理解，那创造这一切的祂，又怎能被理解呢？\BibleQuote{天主为爱祂的人所预备的，是眼所未见，耳所未闻，人心所未想到的}(\BibleRef{格前 2:9})。若天主所预备的，连我们的思想也不能理解，我们怎能以理智理解那预备这些的祂呢？\BibleQuote{啊，天主的富饶、上智和知识，是多么高深！祂的判断，是多么不可测量，祂的道路，是多么不可探察！}(\BibleRef{罗 11:33})宗徒这样说道。若祂的判断和道路是不可理解的，祂本身能被理解吗？\par

10. 天主既如此伟大，且更伟大，（因为即使我将整个身体变成舌头，也无法述说祂的卓越：不仅如此，即使所有天使聚集，也无法述说祂的尊贵），因此天主在美善与威严上如此伟大，人竟敢对刻过的石头说：\Quote{你是我的天主} \BibleRef{依 44:17}！何等可怕的盲目，竟从如此伟大的威严降得如此低下！那由天主种植、蒙雨水滋养、后被火焚烧成灰的树木，竟被称作天主，而真实的天主却被藐视。但偶像崇拜的邪恶变得更加放荡，猫、狗、狼被当作天主敬拜；食人狮也被当作天主敬拜，而这狮子本是人类最慈爱的朋友。蛇——那将我们逐出乐园者的仿冒品——被敬拜，而种植乐园的那位却被藐视。我羞于言说，却仍要说，甚至洋葱在某些地方也被敬拜。酒被赐予\Quote{悦人心}：而狄奥尼索斯（巴克斯）\OriginalTerm{狄奥尼索斯（巴克斯）}{Dionysus (Bacchus)} 被当作天主敬拜。天主藉说话造了五谷：\BibleQuote{让地生出青草，结种子的菜蔬，各从其类，各从其样} \BibleRef{创 1:11}，\Quote{使粮能养人心}：那么为何得墨忒耳（刻瑞斯）\OriginalTerm{得墨忒耳（刻瑞斯）}{Demeter (Ceres)} 被敬拜？火至今仍从击石产生：那么赫菲斯托斯（伏尔甘）\OriginalTerm{赫菲斯托斯（伏尔甘）}{Hephæstus (Vulcan)} 如何是火的创造者？\par

11. 希腊人的多神论谬误从何而来？天主没有身体：那么，在他们所谓的神当中，通奸之事从何而来？我不提宙斯\OriginalTerm{宙斯}{Zeus}变成天鹅的事：我羞于说他变成公牛的事：因为牛叫与神不相称。希腊人的神被发现是奸夫，他们却不以为羞：因为若他是奸夫，就不可被称为神。他们还讲述他们的神之死亡、坠落和雷击。你看见他们从多么高的高处跌落到多么低的地步吗？那么，天主子从天上降下是无缘无故的吗？或是为医治如此大的创伤？圣子降临是无缘无故的吗？或是为使圣父被认出？你已经知道独生子为什么从天主右手的宝座上降下。圣父被藐视，圣子必须纠正谬误：因为万物藉祂而造的那位，必须把万物作为祭品献给万有的主。创伤必须得医治：因为还有什么比这种疾病更糟糕：一块石头竟被当作天主敬拜？\par

\Emph{论异端}\par

12. 魔鬼不仅在外邦人中发动这些攻击；因为许多假称基督徒、被基督的甜美之名错误称呼的人，竟不敬地胆敢将天主从祂自己的创造中驱逐。我指的是异端之群，那些最不虔敬、声名狼藉的人，假装是基督的朋友却全然恨祂。因为凡亵渎基督之父者，就是圣子的仇敌。这些人胆敢讲论两个神性本体，一个善，一个恶！何等可怕的盲目！若是一个神性本体，则必定是善。但若不善，为何称为神性本体？因为若善是天主的属性；若慈爱、恩惠、全能属于天主，那么二者取一：要么在称祂为天主时，让名称与行动相符；要么他们若剥夺祂的行动，就不要给祂空洞的名称。\par

13. 异端者胆敢说有两个天主，一个善的源头和一个恶的源头，且这两个源头都是无始的\OriginalTerm{无始}{unbegotten}。若两者都是无始，则它们必定平等，且都有大能。那么光明如何消灭黑暗？或者它们并存，还是分开？它们不能并存；因为宗徒说：\BibleQuote{光明与黑暗有什么相通？}但如果它们彼此远离，则各自必有其位置；若它们各有分离的位置，那么我们必定是在一个天主的领域内，并且必定敬拜一个天主。因为即使我们同意他们的愚昧，也必须得出我们只敬拜一个天主的结论。让我们也审视他们对善天主所说的话。祂有能力还是无能？若祂有能力，邪恶如何违背祂的意愿产生？若祂不准，邪恶的本质如何侵入？因为若祂知道却不能阻止，他们指控祂无能；但若祂有能力却不阻止，他们又指控祂不忠。也注意他们缺乏理智。有时他们说恶者\OriginalTerm{恶者}{the Evil One}在世界的创造上与善天主毫无相通；但有时又说恶者只有四分之一。他们还宣称善天主是基督的父亲；但基督他们称为这个太阳。因此，照他们所说，世界由恶者创造，太阳在世界中，那么善天主之子如何成为恶者国度里不情愿的奴隶？我们谈论这些事是自陷泥沼，但这样做是怕在场的人因无知而落入异端的泥沼。我知道我玷污了自己的口和听者的耳：然而这是有益的。因为听到别人所犯的荒谬之事，远胜于因无知而陷于其中：你认识泥沼并憎恨它，远胜于不知不觉掉进去。因为异端的无神体系是一条多岔路，每当人偏离了唯一正路，就会一次次跌下悬崖。\par

14. 发明所有异端的是\Person{西满·玛古斯}：就是那位在\WorkTitle{宗徒大事录}中想用金钱购买圣神那不可卖的恩宠，并听见了这些话：\BibleQuote{\Emph{你在这事上既无份也无权}}（\BibleRef{宗 8:18}），以及其余的话。关于他也有记载：\BibleQuote{\Emph{他们从我们中间出去，但他们不属于我们；因为如果他们属于我们，就必会与我们一同留下}}（\BibleRef{若一 2:19}）。这个人被宗徒们驱逐后，来到\Place{罗马}，勾引了一个名叫\Person{赫肋纳}的妓女，是第一个敢用亵渎的口说：他自己就是曾在\Place{西乃山}上显现为圣父的那一位，后来又出现在犹太人中间，不是真肉身而是幻影，如同基督耶稣，后来又作为基督应许要派遣的护慰者圣神。他如此欺骗了\Place{罗马城}，以致\Person{克劳狄}为他立了雕像，并在下面用罗马文字写着\Quote{Simoni Deo Sancto}，翻译出来就是\Quote{致神圣的西满}。\par

15. 正当这谬误蔓延时，\Person{伯多禄}和\Person{保禄}，一对高尚伴侣，教会的主要首领，来到并纠正了错误；当这假神\Person{西满}想要炫耀自己时，他们立刻使他成为一具尸体。因为\Person{西满}许诺要升到天上，并乘坐一个邪魔的战车在空中行驶；但天主的仆人们双膝跪下，显示了耶稣所说的同心合意：\BibleQuote{\Emph{若你们中有两人在地上同心合意地祈求任何事，我在天之父必为他们成就}}（\BibleRef{玛 18:19}），他们发射同心合意的祈祷之箭对准\Person{玛古斯}，把他击落在地。这虽是奇事，却不足为奇。因为\Person{伯多禄}在那里，他带着天堂的钥匙；也不足为怪，因为\Person{保禄}在那里，他是\Emph{曾被提到第三层天上}，\Emph{被提到乐园里，听见了不可言传的话，那是人不许说的}。这两人把这个假神从天上传到地上，再从那里带下地下的区域。在这人身上首先出现了邪恶的毒蛇；但当一颗头被砍下后，邪恶的根源又显现出多头。\par

16. 因为\Person{则林托}曾\Emph{扰乱教会}，还有\Person{默南德尔}、\Person{卡尔波克拉特}、\Person{厄便尼派}以及\Person{玛尔基翁}，就是那个不敬神的代言人。因为那宣告不同神祇的人，一个为善神，另一个为公义之神，他与圣子的话相矛盾，因为圣子说：\BibleQuote{\Emph{公义的圣父}}（\BibleRef{若 17:25}）。而那又说圣父是一位，创造世界的是另一位的人，与圣子的话相悖，因为圣子说：\BibleQuote{\Emph{那么，如果今天还在的田间的草，明天就被丢在炉中的火里，天主尚且这样装饰它}}（\BibleRef{路 12:28}）；又说：\BibleQuote{\Emph{祂叫日头照好人，也照歹人；降雨给义人，也给不义的人}}（\BibleRef{玛 5:45}）。这里又出现了另一位更大的邪恶发明者，就是这位\Person{玛尔基翁}。因为他在新约中引用的旧约见证前被驳倒，他是第一个敢于删减那些见证的人，使信仰之言的宣讲失去见证，从而抹去了真天主；并试图破坏教会的信仰，好像教会中没有传报者一样。\par

17. 他之后又出现了另一位，\Person{巴西里得}，名声邪恶，性格危险，是不洁之事的宣讲者。邪恶之争也得到\Person{瓦伦提诺}的帮助，他是三十个神祇的宣讲者。希腊人只提到少数几个；而那个被称为基督徒——更确切地说，不是基督徒——的人将谬误扩展到整整三十个。他还说，\Person{比索斯}（即深渊）——因为作为邪恶的深渊，他从深渊开始他的教导是合适的——生出了沉默，从沉默生出了圣言。这个\Person{比索斯}比\Person{希腊人的宙斯}更坏，后者与他的姐妹结合：因为据说沉默是\Person{比索斯}的孩子。你看到这荒谬的假基督外表吗？稍等片刻，你会对他的不敬感到震惊；因为他声称从这个\Person{比索斯}生出了八个伊雍；从它们又生出十个；从它们又生出另外十二个，有男有女。但是这些事的证据从何而来？看看他们捏造的愚蠢之处。三十个伊雍的证据从何而来？他说，因为经上记载：\BibleQuote{\Emph{耶稣受洗}，\Emph{约有三十岁}}（\BibleRef{路 3:23}）。但即使祂在三十岁时受洗，这三十岁又能证明什么？难道因为祂用五个饼分给五千人，就有五位神吗？或者因为祂有十二位门徒，就必须有十二位神吗？\par

18. 而这与之后的种种不敬相比还算小的。因为最后一个神祇（他竟敢说）既是男性又是女性，他说，这就是智慧。何等不敬！因为\BibleQuote{\Emph{天主的智慧}}（\BibleRef{格前 1:24}）就是祂的独生子基督；而他用他的教义将天主的智慧贬低为女性元素，且是三十个之一，最后的捏造。他还说，智慧试图看见第一神，却因承受不住祂的光辉而从天上坠落，并被逐出她的第三十个位置。然后她叹息，从叹息中生出了魔鬼，又因哀哭自己被遗弃，用眼泪造了海。请注意这不敬。因为智慧怎能生出魔鬼？明智怎能生出邪恶？光明怎能生出黑暗？他还说，魔鬼生了其他东西，其中一些创造了世界；而基督降下来是为了使人类背叛世界的创造者。\par

19. 但是请听他们所说的耶稣基督是谁，好让你们更加厌恶他们。因为他们说，智慧被贬落后，为了使三十个数目完整，那二十九个伊雍各贡献了一小部分，形成了基督；他们还说祂也是男性兼女性。还有什么比这更不敬的吗？还有什么比这更可怜的吗？我向你们描述他们的谬误，是为了让你们更加憎恨他们。因此，要避免他们的不敬，\BibleQuote{\Emph{连问候也不要向这样的人}}（\BibleRef{弗 5:11}），免得你们\BibleQuote{与黑暗无益的作为有份}；不要好奇探问，也不要愿意与他们交谈。\par

20. 应憎恨一切异端者，但尤其要憎恨那位名副其实以\Quote{发狂}为名的人，他在普罗布斯皇帝在位时期兴起，距今不过七十余年，仍有见过他本人的在世者。但不要因他离现在不久而憎恨他，而要因他不敬神的教义憎恨这作恶的工人、一切污秽的容器，他收集了每种异端的污泥。因他渴望在恶人中出类拔萃，便采纳了所有人的教义，将它们合并成一个充满亵渎和不义的异端，蹂躏教会，更确切地说，是蹂躏教会外的人，像狮子一样四处游荡吞噬。不要听信他们动听的言辞，也不要相信他们所谓的谦卑：因为他们是蛇，\BibleQuote{\Emph{毒蛇的种}}。\BibleRef{玛 3:7} 犹大也曾\Quote{说：\InnerQuote{恭喜，拉比！}}，甚至就在他出卖祂的时候。不要听信他们的亲吻，但要提防他们的毒液。\par

21. 现在，为避免我似乎无端指控他，让我稍作离题，说说这个摩尼是何许人，并部分介绍他的教导；因为要完全描述他污秽的教导，全部时间也不够。但\BibleQuote{\Emph{为获得及时的助佑}} \BibleRef{希 4:16}，请将我对先前听众说过、现在也对在场者重复的话牢记在心，使不知者得以学习，已知者得以回想。摩尼并非源自基督徒，愿天主不许！他也不像西门那样被逐出教会，无论是他本人还是他之前的教师们。因为他窃取别人的恶行，使别人的恶行成为自己的；但如何、以何种方式，你们必须听听。\par

22. 在埃及曾有一个人名斯基提亚诺斯，是撒拉森人出身，与犹太教或基督教毫无共同之处。此人住在亚历山大里亚，效仿亚里士多德的生活方式，写了四本书：一本称为福音，却没有基督的作为，只有虚名；另一本称为\WorkTitle{篇章书}，第三本为\WorkTitle{奥秘书}，第四本是他们如今传播的\WorkTitle{宝藏书}。这人有一个门徒，名叫特雷宾托斯。但当斯基提亚诺斯打算进入犹太地，蹂躏那地方时，主用致命的疾病击打了他，制止了这场瘟疫。\par

23. 但特雷宾托斯，他这个邪恶错误的门徒，继承了金钱、书籍和异端，来到巴勒斯坦，在犹太地被人知晓并定罪后，他决定前往波斯；但为避免在那里也因名字被认出，他改名为布达斯。然而，他在那里也遇到了对手，即密特拉的祭司；在许多辩论和争辩中被驳倒，最后被紧逼时，他逃到某个寡妇那里。然后，他上了屋顶，召唤空中的恶魔——摩尼教徒至今在他们可憎的无花果仪式中仍召唤这些恶魔——他被天主击打，从屋顶摔下，断了气：这样第二只野兽被剪除了。\par

24. 然而，记载他邪恶的书籍仍在；寡妇继承了这些书和金钱。她既无亲属也无其他朋友，便决定用钱买一个名叫库布里科斯的男孩；收养并教育他，像儿子一样学习波斯的学问，如此磨砺了一件加害人类的邪恶武器。于是卑贱的奴隶库布里科斯在哲学家中长大，寡妇死后继承了书籍和金钱。然后，为避免奴隶之名成为耻辱，他不再叫库布里科斯，而自称摩尼，这在波斯语中意为\Quote{言论}。因他自认为是个善辩者，便自称摩尼，仿佛卓越的言论大师。但尽管他按波斯语为自己安排了一个尊贵的称号，天主的眷顾却使他即使不情愿也成为自我控诉者，使他在波斯自以为尊荣时，却在希腊人面前自称疯人（\OriginalTerm{疯人}{maniac}）。\par

25. 他还胆敢说自己是护慰者，尽管经上记载：\BibleQuote{\Emph{但凡亵渎圣神的，永远不得赦免}}。\BibleRef{谷 3:29} 因此他因说自己是圣神而犯了亵渎；与那些异端者为伍的人，看看自己与谁为伍。这个奴隶震撼了世界，因为\BibleQuote{\Emph{大地因三件事而震动，第四件事她不能承受——奴仆为王}}。\BibleRef{箴 30:21} 他进入公众视线后，开始许诺超乎人力之事。波斯国王的儿子病了，许多医生在旁；但摩尼许诺，仿佛他是个虔诚的人，要通过祈祷治愈他。医生离去后，孩子的生命也随之离去；此人的不敬神被揭露。于是这所谓的哲学家成了囚犯，被投入监狱，不是因为为真理而斥责国王，也不是因为他摧毁偶像，而是因为他许诺拯救却撒谎，更确切地说，若道出真相，是因为他犯了杀人罪。因为本来可以通过医疗得救的孩子，被这人赶走医生、因缺乏治疗而杀死了。\par

26. 既然我告诉你们关于他的许多邪恶之事，首先记住他的亵渎，其次他的奴隶身份（并不是说奴隶是耻辱，而是他本是奴隶却冒充自由人是邪恶的），第三他许诺的虚假，第四孩子的被杀，第五监禁的耻辱。而且不仅有监禁的耻辱，还有越狱。因为那自称护慰者和真理维护者的人逃跑了；他不是耶稣的继承者，耶稣欣然走向十字架，而此人正相反，是个逃跑者。此外，波斯国王下令处死狱卒：这样，摩尼因虚妄的夸口导致孩子死亡，因逃跑导致狱卒死亡。那么，这个与杀人同罪的人，应当受人敬拜吗？难道他不该效法耶稣，说：\BibleQuote{\Emph{你们若找我，就让这些人去吧}}？\BibleRef{若 18:8} 难道他不该像约纳那样说：\BibleQuote{\Emph{你们拿我，将我抛入海中：因为这风暴是因我的缘故}}？\BibleRef{纳 1:12}\par

27. 他逃出监狱，来到\Place{美索不达米亚}：但那里的主教\Person{阿凯劳斯}，正义的盾牌，与他相遇：并在哲学家面前控告他，以他们为审判，并召集了外邦人的听众，以免如果基督徒审判，审判官可能被认为偏袒——\Person{阿凯劳斯}对\Person{摩尼}说：告诉我们你所传讲的。而他，其\Emph{口如敞开的坟墓}，首先以对万有创造者的亵渎开始，说：旧约的天主是邪恶的作者，正如祂自称说：\BibleQuote{我是吞灭的火}\BibleRef{申 4:24}。但智慧的\Person{阿凯劳斯}通过说以下话瓦解了他的亵渎论点：\Quote{如果旧约的天主，如你所说，自称是火，那么那位说\BibleQuote{我来是把火投在地上}\BibleRef{路 12:49}的子是谁？如果你责备那说\BibleQuote{主使人死，也使人活}\BibleRef{撒上 2:6}的，你为什么尊敬\Person{伯多禄}，他复活了\Person{塔彼达}，却打击\Person{撒斐辣}致死？如果你又责备，因为祂预备了火，那你为何不责备那说\BibleQuote{离开我，进入永火里}\BibleRef{玛 25:41}的？如果你责备那说\BibleQuote{我是造和平、也造邪恶的上主}\BibleRef{依 45:7}的，请解释\Person{耶稣}如何说\BibleQuote{我来不是为送和平，而是刀剑}\BibleRef{玛 10:34}。既然两者说话相同，二者必居其一：要么两者都是善的，因为一致；要么如果\Person{耶稣}如此说话无可指责，你为什么责备那在旧约中说同样话的？}\par

28. 然后\Person{摩尼}回答他：\Quote{什么样的天主使人眼瞎？因为正是\Person{保禄}说：\BibleQuote{此世的天主蒙蔽了不信者的思想，免得基督福音的光照耀他们}。}但\Person{阿凯劳斯}作出了很好的反驳，说：请读前面一点：\BibleQuote{但如果我们的福音被蒙住，只是对丧亡的人蒙住}\BibleRef{格后 4:3}。你看见了吗？在丧亡的人中它是蒙住的。因为\BibleQuote{把圣物给狗}是不对的\BibleRef{玛 7:6}。再者，难道只有旧约的天主蒙蔽了不信者的思想吗？\Person{耶稣}自己不是说过：\BibleQuote{为此，我用比喻对他们讲，是因为他们看，却看不见}？是因恨他们而不愿他们看见吗？还是因为他们的不配，因为\BibleQuote{他们闭上了自己的眼睛}\BibleRef{玛 13:15}。因为哪里有故意的邪恶，哪里也有恩宠的扣留：\BibleQuote{因为凡有的，还要给他；凡没有的，连他自以为有的，也要从他夺去}。\par

29. \Quote{但如果有些人的解释是正确的，我们必须如下说（因为这不是不恰当的表达）——如果祂确实蒙蔽了不信者的思想，祂蒙蔽他们是为了一个良善的目的，好让他们以新的眼光看见善。因为祂没有说：祂蒙蔽了他们的灵魂，而是说\BibleQuote{不信者的思想}。其含义类似这样：\InnerQuote{蒙蔽淫秽之人的淫秽思想，那人便得救；蒙蔽强盗的贪婪和掠夺思想，那人便得救。}但你不会这样理解吗？那么还有另一种解释。太阳也蒙蔽那些视力模糊的人：眼睛有病的人被光线伤害并蒙蔽。不是因为太阳的本性是蒙蔽，而是因为眼睛的实质无法看见。同样，不信者因内心有病，无法看见天主性的光辉。祂也没有说：\BibleQuote{祂蒙蔽了他们的思想}，使他们不能听福音；而是说：\BibleQuote{免得基督福音荣耀的光照在他们身上}。因为听福音是允许所有人的，但福音的荣耀只为基督的真子女保留。因此主用比喻对那些不能听的人讲话\BibleRef{玛 13:13}，但对门徒，祂私下解释比喻\BibleRef{谷 4:34}：因为荣耀的光辉是为那些已受光照的人，蒙蔽是为那些不信的人。}这些奥秘，教会现在向你们这些即将离开\OriginalTerm{慕道者}{Catechumens}阶层的人解释，而不习惯向\OriginalTerm{外教人}{heathen}解释。因为对外教人，我们不解释关于\OriginalTerm{天主圣三}{Father, Son, and Holy Ghost}的奥秘，在慕道者面前也不明讲奥秘：但我们常常以隐晦的方式讲许多事，好让知道的人能明白，不知道的人不受伤害。\par

30. 藉由这些和许多其他论证，蛇被推翻：\Person{阿凯劳斯}就是这样与\Person{摩尼}角力并摔倒了他。再者，那个从监狱逃脱的人也从这地方逃走了：在逃离了他的对手后，他来到了一个非常贫穷的村庄，如同乐园中的蛇离开\Person{亚当}来到\Person{厄娃}。但善牧\Person{阿凯劳斯}为他的羊群着想，当他听到他的逃跑，立刻以全速赶去寻找这只狼。而当\Person{摩尼}突然看见他的对手时，他冲出去逃跑了：然而这是他的最后一次逃跑。因为波斯王的官员搜查各处，捕获了这逃亡者：而他本应在\Person{阿凯劳斯}面前接受的判决，由国王的官员下达。这位\Person{摩尼}，被他自己的门徒崇拜，被逮捕并带到国王面前。国王谴责了他的谎言和逃跑：对他奴隶的身份表示轻蔑，为他孩子的被杀报仇，也因狱卒的谋杀而定他的罪：他命令以波斯的方式剥他的皮。而他身体的其他部分被交给野兽作为食物，他的皮，他卑鄙思想的容器，像口袋一样挂在城门前。那位自称\OriginalTerm{护慰者}{Paraclete}并声称知晓未来的人，却不知道自己的逃跑和被擒。\par

31. 这人有三个门徒：\Person{多默}、\Person{巴达斯}和\Person{赫尔玛斯}。谁也不可读那按\Person{多默}的\WorkTitle{福音}：因为它不是十二宗徒中任何一位的作品，而是摩尼的三个邪恶门徒之一的作品。谁也不可与那使人丧亡的摩尼教徒来往，他们以谷壳的煎熬来伪装禁食的哀容，他们诽谤食物的创造者，却贪婪地吞食最精美的食物，他们教导人，说拔起这种或那种草药的人会变成这种植物。因为如果割取药草或任何蔬菜的人变成同样的东西，那么农夫和园丁之流会变成多少东西呢？我们看到，园丁用他的镰刀对付了那么多东西：那么他变成什么呢？实在，他们的教义是可笑的，充满自己的谴责和耻辱！同一个人，作为羊群的牧人，既宰杀绵羊又杀死豺狼：那么他变成什么呢？许多人既网鱼又粘鸟：那么他们变成什么呢？\par

32. 让那些懒惰的儿女，摩尼教徒，来回答吧；他们自己不劳动却吃劳动者的果实：他们笑脸迎接送食物给他们的人，却以咒骂代替祝福。因为当一个纯朴的人带给他们任何东西时，他说：\Quote{你在外面站一会儿，}他说，\Quote{我要祝福你。}然后他拿过面包（正如那些已经悔改并离开他们的人所承认的），摩尼教徒对面包说：\Quote{我没有造你。}并对至高者发出咒骂；他咒骂那造面包者，然后吃掉所造之物。如果你憎恨食物，为什么你笑脸相迎那带食物给你的人？如果你感谢那带来的人，为什么你向创造并制造它的天主发出亵渎？同样，他又说：\Quote{我没有播种你：愿那播种你的人被播种！我没有用镰刀收割你：愿那收割你的人被收割！我没有用火烘烤你：愿那烘烤你的人被烘烤！}真是对恩惠的美妙回报！\par

33. 这些是重大过失，但与其他相比还是很小的。他们的圣洗我不敢在男人女人面前描述。我不敢说他们分发给可怜的领受者的是什么。……实在，我们在说这些事时玷污了自己的口。异教徒比这些更可憎吗？撒玛黎雅人更悲惨吗？犹太人更不虔敬吗？行淫者更不洁吗？但摩尼教徒却把这些供品放在自己认为的祭台中间。而你，人啊，竟从这样的口中接受教导吗？遇见此人时，你还会以亲吻向他致意吗？且不说他其他的不虔敬，难道你不逃避这污秽，不逃避比放荡者更坏、比任何娼妓更可憎的人吗？\par

34. 对于这些事，教会告诫并教导你们，它触及污泥，为使你们不致沾污；它述说创伤，为使你们不致受伤。但你们只要知道就足够了：避免亲身体验。天主打雷，我们都战栗；他们却亵渎。天主闪电，我们都俯伏于地；他们却有亵渎的言论关乎诸天。这些事写在摩尼教徒的书中。这些事我们亲自读过，因为我们不能相信那些讲述的人；是的，为了你们的救恩，我们仔细查究了他们的丧亡。\par

35. 但愿主拯救我们脱离这样的迷惑：并且愿你们对蛇怀有憎恨，正如他们埋伏攻击脚跟，你们可以践踏他们的头。记住我说的话。我们的状况与他们之间能有什么协议呢？\BibleQuote{\Emph{光明与黑暗有什么相通？}}\BibleRef{格后 6:14}教会的庄严与摩尼教徒的可憎之物有什么相干？这里有秩序，这里有纪律，这里有庄严，这里有纯洁：这里甚至\BibleQuote{\Emph{注视妇女而有意贪恋她}} \BibleRef{玛 5:28}就是定罪。这里有神圣的婚姻，这里有坚定的节欲，这里有如同天使一样受尊重的童贞：这里有以感恩之心分享食物，这里有对世界创造者的感激。这里基督的圣父受到敬拜：这里教导对那降雨者敬畏和战栗：这里我们将荣耀归于那造雷和闪电者。\par

36. 要使你的羊群与羊同群：逃离豺狼：不要离开教会。也要憎恶那些曾在此类事上受怀疑的人：除非你及时看到他们悔改，不要轻率地置身于他们之中。天主唯一性的真理已传给你们：要学会辨别教义的牧场。作一个精明的钱庄主，\BibleQuote{\Emph{持守那善的，远离一切形式的恶}}。\BibleRef{得前 5:21}或者，如果你曾是他们这样的人，要认识并憎恶你的迷惑。因为有一条得救之路，如果你弃绝那所吐的，如果你从心里憎恶它，如果你离开他们，不单用嘴唇，也用灵魂：如果你敬拜基督的圣父，即法律和先知的天主，如果你承认善与公义是同一个天主。愿他保全你们众人，保护你们不致跌倒或失足，在信德中站稳，在基督耶稣我们的主内，愿光荣归于他，直到永远。阿们。\par

\SourceRecord{Translated by Edwin Hamilton Gifford.；Nicene and Post-Nicene Fathers, Second Series；Vol. 7.；Edited by Philip Schaff and Henry Wace.；Buffalo, NY: Christian Literature Publishing Co.,；1894.；<http://www.newadvent.org/fathers/310106.htm>.}

% structure: p0001 source=h1 target=section reason=page-title

\section{要理讲授第七篇}

\Emph{圣父}\par

\BibleRef{弗 3:14,15}\par

\BibleQuote{为此，我屈膝在圣父面前……天上和地上的一切父职都由此得名等等。}\par

1. 关于天主作为唯一的本原，我们昨天已经说得足够多了：我说\Quote{足够}，并非指与主题相称（因为受造的本性绝不可能达到那程度），而是指按我们软弱的程度所蒙赐予的。我也论及了不信神的异端者种种错误的歧途；但现在，让我们摆脱他们那污秽而毒害灵魂的教义，记念那些与他们相关的事（并非为了危害我们自己，而是为了更加憎恶他们），让我们回到自身，接受真正信仰的救恩教义，将父职的尊荣与唯一性相连，并相信唯一天主圣父。因为，我们不仅必须相信独一的天主，还应虔诚地领受：祂是独生子、我们的主耶稣基督之父。\par

2. 因为这样，我们的思想将超越犹太人，他们固然藉其教义承认唯一天主（尽管他们常常通过偶像崇拜否定这一点），但他们不承认祂也是我们主耶稣基督的父；他们与自己的先知心意相悖，因为先知在圣经中宣告：\BibleQuote{主对我说：\InnerQuote{你是我的儿子，我今日生了你。}}直到今日，他们还\BibleQuote{谋聚合谋攻击主和祂的受膏者}，以为可以离开对圣子的敬拜而与圣父为友，却不知道\BibleQuote{除非藉着圣子，没有人能到父那里去}\BibleRef{若 14:6}，祂说：\BibleQuote{我是门，我是道路。}因此，拒斥通往圣父之道的人，以及否认门的人，怎能被认为配得进入天主面前呢？他们也反对第八十八篇圣咏所写的：\BibleQuote{他要称呼我说：\InnerQuote{你是我的父，我的天主，我救恩的磐石。}我也要立他为长子，为地上众君王中最高的。}如果他们坚持认为这些话是指达味、撒罗满或他们的任何继承者，那么，请他们指出，按他们判断为预言中所描述的那位的\BibleQuote{宝座}，如何能\BibleQuote{如天上的日子，如太阳在天主面前，又如月亮永远立定}？他们又如何不因所写的\BibleQuote{从母胎中，在晨星以前，我已生了你}，以及\BibleQuote{他将与太阳共存，与月亮同久，从一代到另一代}而羞愧呢？\par

3. 然而，就任凭犹太人按其意愿，在他们原本的不信中继续如此吧，在这些和类似的说法上都是如此。但我们应采纳我们信仰的虔诚教义，敬拜那位独一的天主、基督的父（因为剥夺那赐予众人生育恩赐者的同样尊荣，实属不敬）；并相信唯一天主圣父，为的是在我们触及有关基督的教导之前，关于独生子的信德能先植入听者的灵魂，而不被关于圣父的中间教义所打断。\par

4. 因为圣父的名号，一提到这名，便提示了圣子的思想；同样，凡是称呼圣子的，也立刻想到圣父。因为如果是父，那祂必定是子的父；如果是子，那祂也必定是父的子。因此，当我们如此说：\Quote{唯一天主，全能圣父，天地创造者，有形无形万物的创造者}，然后又加上\Quote{唯一主耶稣基督}时，唯恐有人不敬地以为独生子在等级上仅次于天地——正是为此，我们在提及它们之前先提到了天主圣父，为的是在想到圣父时，同时也想到圣子；因为在圣子与圣父之间，没有任何受造物存在。\par

5. 因此，天主在非本义上是许多人的父，但按本性并在真理上只是一个的父，即独生子、我们的主耶稣基督的父；祂并非在时间过程中才成为父，而是始终是独生子的父。并非祂先前无子，后来因旨意改变而成为父；而是在一切实体、一切理智、一切时间与一切世代之前，天主就拥有父的尊荣，祂在这尊荣上比在其他尊荣上更显伟大；祂成为父，不是由于受生、结合、不知、流溢、减少或改变，\BibleQuote{因为一切美好的赠与、一切完美的恩赐，皆是自上而来，从光明之父降下，在祂那里没有变动，也没有转动的阴影}\BibleRef{雅 1:17}。完美的父生了完美的子，并将一切交给了所生的子（因为祂说：\BibleQuote{我父把一切都交给了我}\BibleRef{玛 11:27}）；并受独生子的尊崇：因为子说：\BibleQuote{我尊敬我的父}\BibleRef{若 8:49}；又说：\BibleQuote{正如我遵守了我父的命令，并常住在祂的爱内}\BibleRef{若 15:10}。因此，我们也如同宗徒一样说：\BibleQuote{愿我们的主耶稣基督的天主和父受赞美，祂是仁慈之父，和一切安慰的天主}\BibleRef{格后 1:3}；又说：\BibleQuote{我屈膝在圣父面前，天上和地上的一切父职都由此得名}\BibleRef{弗 3:14}；并与独生子一同光荣祂：因为否认父的，也否认子；又说承认子的，也有父；并知道\BibleQuote{耶稣基督为主，以光荣天主圣父}\BibleRef{斐 2:11}。\par

6. 因此，我们朝拜基督之父，天地的创造者，\newline{}\Quote{亚巴郎、依撒格和雅各伯的天主}（见：\BibleRef{出 3:6}）；从前那建在对面、与我们相对的圣殿，也是为光荣祂而建的。因为我们绝不容忍那些将旧约与新约割裂的异端者；相反，我们要相信基督所说关于圣殿的话：\Quote{难道你们不知道我必须在我父的家里吗？}（见：\BibleRef{路 2:49}）又说：\Quote{把这些东西拿出去，不要使我父的殿成为商场}（见：\BibleRef{若 2:16}），祂由此极其清楚地承认，耶路撒冷的前圣殿原是祂父的殿。若有人因不信而想获得更多证据，证明基督之父就是世界的创造者，让他再听祂说：\Quote{两只麻雀不是卖一个铜钱吗？但若没有你们在天之父的许可，一只也不会掉在地上}；还有：\Quote{你们看天空的飞鸟，它们不播种，也不收割，也不收聚在仓里，你们的天父尚且养活它们}（见：\BibleRef{玛 6:26}）；以及：\Quote{我父一直工作到现在，我也工作}（见：\BibleRef{若 5:17}）。\par

7. 但是，免得有人因单纯或乖僻的机巧，从祂所说的\Quote{我升到我的父和你们的父那里去}，而以为基督与义人仅享有同等的尊荣，最好事先作此区分：父的名号是同一的，但祂行动的权能却是多重的。基督自己知道这一点，曾确切无误地说：\Quote{我往我的父和你们的父那里去：}祂没有说\Quote{我们的父}，而是加以区分，先说属于祂自己的，即按本性是\Quote{我的父}，然后加上\Quote{和你们的父}，即按义子之名份。因为无论我们在祈祷中说\Quote{我们的天父}这恩宠何等崇高，这恩赐仍是出自慈爱。我们称祂为父，并非因我们按本性是由天父所生，而是因借着圣父的恩宠，通过圣子和圣神，我们从奴役转入儿子名分，才获准因不可言喻的慈爱如此称呼。\par

8. 若有人想知道我们如何称天主为\Quote{父}，让他听那卓越的师傅梅瑟说：\Quote{这不是你的父自己买了你，造了你，立了你吗？}（见：\BibleRef{申 32:6}）还有依撒意亚先知说：\Quote{如今，上主啊，你是我们的父；我们原是泥土，是你手中的工作}（见：\BibleRef{依 64:8}）。因为先知之恩极清楚地宣告，我们称祂为父，不是按本性，而是按天主的恩宠和义子名份。\par

9. 为使你们更准确地知道，在神圣的圣经中，并非只有生身的父才被称为父，请听保禄所说：\Quote{纵然你们在基督内有上万的导师，但你们没有许多父亲，因为在基督耶稣内，我借福音生了你们}（见：\BibleRef{格前 4:15}）。因为保禄是格林多人的父亲，不是因按肉身生了他们，而是因教导他们，并借圣神重生了他们。也听约伯说：\Quote{我成了穷人的父亲}（见：\BibleRef{约 29:16}）；他称自己为父亲，并非因生了他们众人，而是因照顾他们。天主的独生子自己，当祂的肉身被钉在十字架上时，看见按肉身祂自己的母亲玛利亚，以及祂最爱的门徒若望，就对若望说：\Quote{看，你的母亲！}又对玛利亚说：\Quote{看，你的儿子！}（见：\BibleRef{若 19:26}）以此教导她对他应有父母般的慈爱，并间接解释了路加福音所说：\Quote{他的父亲和母亲都惊奇于祂}（见：\BibleRef{路 2:33}）。这节经文被异端一派抓住，说祂是由一男一女所生。因为正如玛利亚被称为若望的母亲，是由于父母般的慈爱，并非因生了他，同样，若瑟也被称为基督的父亲，并非因生了祂（因为正如福音所说：\Quote{他不认识她，直到她生了她的长子}（见：\BibleRef{玛 1:25}）），而是由于对祂养育的照顾。\par

10. 目前暂且以此作为题外话，提醒你们。然而，我再加一个证言，证明天主被称为人的父，是在不严格的意义上。因为当依撒意亚书中如此向天主祈祷：\Quote{因为你是我们的父，虽然亚巴郎不认识我们}（见：\BibleRef{依 63:16}），\Quote{撒辣没有为我们受产痛}，对此我们还需要进一步查究吗？如果圣咏作者说：\Quote{愿他们在祂面前受惊扰，孤儿的父，寡妇的审判者}，难道不是众人皆知：当天主被称为刚失去生父的孤儿的父时，如此称呼祂，不是因祂生了他们，而是因祂照顾他们、保护他们。然而，正如我们所说，天主在不严格的意义上是人的父，唯独对基督，祂是按本性之父，而非因义子名份；祂是在时间内作人的父，但对基督，是在万世之前，正如祂说：\Quote{如今，父啊，请在你的面前荣耀我，赐我在世界未有之前，同你所有的荣耀}（见：\BibleRef{若 17:5}）。\par

11. 因此，我们信唯一的天主，圣父，是不可探究、不可言喻的，\Quote{没有人见过祂}（见：\BibleRef{弟前 2:16}），\Quote{只有独生子将祂彰显出来}（见：\BibleRef{若 1:18}）。\Quote{因为那出于天主的，祂见过天主}：诸天使在天上常瞻仰祂的面（见：\BibleRef{玛 18:10}），然而，每位天使是按自己品级的限度来瞻仰。但那毫无遮蔽地瞻仰圣父，是保留给圣子和圣神的纯粹特权。\par

12. 講到這裡，我記起剛才提到的經文，其中天主被稱為人的父親，我對人們的麻木不仁深感驚奇。因為天主以不可言說的慈愛，竟屈尊被稱為人的父親——祂在天上，他們在地上；祂是永恆的創造者，他們是時間中的受造物；祂是\Quote{以手掌托住大地}的那一位，他們在地上如蝗蟲。然而，人離棄了天上的父親，卻對木頭說：\Quote{你是我的父親}，對石頭說：\Quote{你生了我。}\BibleRef{耶 2:27} 為此，我想，聖詠作者對人說：\Quote{忘記你的民族和你父的家}，因為你選擇了那作為父親的，你為自己招致了滅亡。\par

13. 不僅是木頭和石頭，甚至連撒殫本身，靈魂的毀滅者，也有人曾選擇他為父親；主責備他們說：\Quote{你們行你們父親的事業}\BibleRef{若 8:41}，即魔鬼的事業，他成為人的父親不是出於本性，而是出於欺騙。正如\Person{保祿}因虔誠的教導而被稱為\Place{格林多}人的父親，同樣，魔鬼被稱為那些自願\Quote{隨從他}之人的父親。\par

因為我們不容忍那些曲解\Quote{從此我們認出誰是天主的子女，誰是魔鬼的子女}\BibleRef{若一 3:10}這句話的人，好像有些人天生就是得救的，有些人天生就是喪亡的。然而，我們成為如此神聖的義子，不是出於必然性，而是出於選擇；叛徒\Person{猶達斯}本不是魔鬼和喪亡之子，因為他絕不會以基督的名驅逐魔鬼：\Quote{撒殫不能驅逐撒殫}\BibleRef{谷 3:23}。另一方面，\Person{保祿}也不會從迫害者轉變為宣講者。但是，義子的身份在於我們自己的能力，正如若望所說：\Quote{凡接受祂的，祂賜給他們權能，成為天主的子女，就是那些信祂名的人}\BibleRef{若 1:12}。因為不是在他們相信之前，而是從他們相信之時起，他們才配得自願成為天主的子女。\par

14. 因此，我們既知道這事，就當按聖神行事，好能配得天主的義子名分。\Quote{因為凡被天主的聖神引導的，都是天主的子女}\BibleRef{羅 8:14}。因為徒有基督徒的稱號對我們毫無益處，除非行為也跟上；免得我們也聽到這話：\Quote{你們若是\Person{亞巴郎}的子女，就該作\Person{亞巴郎}所作的事}\BibleRef{若 8:39}。\Quote{因為我們稱那不偏待人、按各人行為審判的主為父，就當存敬畏的心，度我們寄居的時日}\BibleRef{伯前 1:17}，\Quote{不要愛世界，也不要愛世界上的事；誰若愛世界，天父的愛就不在他內}\BibleRef{若一 2:15}。因此，我親愛的孩子們，讓我們以行為光榮\Quote{我們天上的父，叫他們看見我們的好行為，便將榮耀歸給我們在天上的父}\BibleRef{瑪 5:16}。\Quote{讓我們將一切憂慮卸給祂，因為我們的父知道我們需要什麼}。\par

15. 但在尊敬我們天上的父親時，也當尊敬\Quote{我們肉身的父親}\BibleRef{希 12:9}，因為主自己明明在律法和先知中這樣規定說：\Quote{應孝敬你的父親和你的母親，好使你的日子在土地上得以長久}\BibleRef{申 5:16}。在座的各位若有父母，更當遵守這條誡命。\Quote{你們作子女的，要事事聽從父母，因為這是主所喜悅的}\BibleRef{哥 3:20}。因為主沒有說：\Quote{凡愛父親或母親的，不配作我的門徒}，免得你們因無知而曲解了正確寫下的話，而是加上\Quote{勝過愛我}\BibleRef{瑪 10:37}。因為當我們在地上的父親與我們在天上的父親意見相反時，我們就必須聽從基督的話。但若他們不阻礙我們的虔誠，而我們卻忘恩負義，忘記他們的恩惠，輕視他們，那麼這訓言就有效：\Quote{凡咒罵父親或母親的，應處以死刑。}\par

16. 基督徒虔敬的第一美德，就是尊敬父母，報答生養他們的辛勞，並盡全力為他們提供舒適（因為即使我們加倍報答，也永遠無法償還他們賜予生命的恩惠），好使他們也享受我們所提供的安慰，並為我們堅定那巧取之\Person{雅各伯}所奪得的祝福；這樣，我們在天上的父將接納我們的美善意願，並認為我們配得\Quote{在義人的國中，如同太陽一樣發光}\BibleRef{瑪 13:43}。願光榮歸於祂，並歸於獨生子我們的主\Person{耶穌基督}，以及聖神——生命之賦予者，從今世直到永遠，世世代代。阿們。\par

\SourceRecord{Translated by Edwin Hamilton Gifford.；Nicene and Post-Nicene Fathers, Second Series；Vol. 7.；Edited by Philip Schaff and Henry Wace.；Buffalo, NY: Christian Literature Publishing Co.,；1894.；<http://www.newadvent.org/fathers/310107.htm>.}

% structure: p0001 source=h1 target=section reason=page-title

\section{教理讲授第八讲}

\Emph{全能者。}\par

耶肋米亚 39:18, 19（七十贤士译本）。\par

\Emph{伟大、刚强的天主，大谋略的主，大能作为的天主，伟大的天主，全能的主和伟大的名。}\par

1. 藉着相信\Quote{唯一的天主}，我们斩断了对众多天主的错误信仰，以此作为抵挡希腊人的盾牌，并对抗一切异端者的反对势力；而加上\Quote{唯一的圣父天主}，我们则与那些否认天主独生圣子的受割损之人争辩。因为正如昨天所说，即使在解释有关我们主耶稣基督的真理之前，我们已藉着说\Quote{圣父}，立刻显明祂是一位圣子的父亲：好叫我们既知道天主存在，也能知道祂有一位圣子。但这些称号之外，我们还加上说祂也是\Quote{全能者}；这是为了同时对抗希腊人、犹太人以及所有的异端者。\par

2. 因为希腊人中，有的说天主是世界的灵魂；有的说祂的权能只达于天上，而不及于地上。也有些人共享他们的错误，并滥用那句经文说：\Quote{\Emph{祢的信实直到云霄}}，竟敢将天主的眷顾局限于云层和天界，并使地上的事物与天主隔绝；他们忘记了那篇圣咏说：\Emph{我若升到天上，祢在那里；我若下到阴间，祢也在那里。}因为若没有比天更高的，而阴间比地更深，那么掌管下界者也必及于地上。\par

3. 然而异端者，如我先前所说，不认识唯一全能的天主。因为\Quote{全能者}是统治万有、掌管一切的主。而那些说一位天主是灵魂的主宰，另一位是身体的主宰的人，使两者都不完美，因为彼此有所欠缺。因为若一位能掌管灵魂却不能掌管身体，他如何是全能者？若另一位能主宰身体却不能掌管灵魂，他如何是全能者？但主相反地责备这些人说：\Quote{\BibleQuote{宁可畏惧那能在地狱里毁灭灵魂和身体的那一位。}}\BibleRef{玛 10:28} 因为我们主耶稣基督的圣父若未拥有两者的权能，祂怎能将两者都置于惩罚之下？因为\Emph{除非他先捆绑那壮士，然后劫掠他的财物}，祂怎能夺取属于别人的身体而将它投入地狱呢？\par

4. 但神圣的圣经和真理的教义只认识一位天主，祂以祂的能力掌管万有，却出于祂的旨意容忍许多事。因为祂连拜偶像者也掌管，却因祂的宽容而容忍他们；祂连藐视祂的异端者也掌管，却因祂的忍耐而容忍他们；祂连魔鬼也掌管，却因祂的容忍而容忍它，并非出于能力的欠缺，好像失败一般。因为\Emph{它是主创造之始，被造来受戏弄}，并非由祂自己戏弄——这于祂不相称——而是由祂所造的\Emph{天使}戏弄。但祂容许它活着，有两个目的：使它在失败中更加蒙羞，并使人能戴上胜利的冠冕。哦，天主全智的眷顾！祂以邪恶的意图作为信徒救恩的基础。因为正如祂以若瑟弟兄们不友爱的意图作为祂自己安排的基础，并容忍他们因仇恨出卖自己的兄弟，借此使祂所愿的人成为君王；同样，祂也容许魔鬼角力，使得胜者得冠冕；而当胜利获得时，魔鬼因被较弱者征服而更蒙羞辱，人则因战胜了曾是总领天使的敌人而得大荣耀。\par

5. 因此，没有任何事物脱离了天主的权能；因为圣经论到祂说：\Emph{万物都是祢的仆役。}万物同样是祂的仆役，但其中只有一个，祂的独生子，以及一个，祂的圣神，是例外；而一切祂的仆役都藉着独生子并在圣神内服事主。因此天主统治万有，并因祂的容忍而甚至容忍杀人犯、强盗和奸淫者，为每个人指定了报应的时期，好让那些虽经长久警告仍存心不悔改的人，受更重的惩罚。他们是地上的君王，统治世上，但并非无权柄从上而来；这一点拿步高曾亲身体验过，他说：\Quote{\BibleQuote{祂的王国是永远的王国，祂的权能代代相传。}}\BibleRef{达 4:34}\par

6. 财富、金银，并不像有些人所想的那样是魔鬼的；因为\Emph{全部财富世界是为忠信之人，但不忠信的人连一文钱也没有。}没有什么比魔鬼更不忠信的了；天主藉着先知明确地说：\Emph{金子是我的，银子也是我的，我愿意给谁就给谁。}你只要善用它，金钱就无可指责；但每当你滥用了那美好的事物，你便不肯怪罪自己的管理，反而亵渎地归咎于造物主。人甚至可以因金钱而成义：\Quote{\BibleQuote{我饿了，你们给我吃的。}}\BibleRef{玛 25:35} 那当然来自金钱。\Quote{\BibleQuote{我赤身露体，你们给我穿的。}}那当然藉着金钱。而你想知道金钱可以成为天国的大门吗？祂说：\Quote{\BibleQuote{变卖你所有的，分施给穷人，你将有宝藏在天上。}}\par

7. 我说这些话是基于那些异端者，他们视财产、金钱和人的身体为可咒骂的。因为我既不愿你作金钱的奴隶，也不愿你憎恨天主赐给你使用的事物。因此绝不要说财富是魔鬼的；因为虽它说：\Emph{这一切我都要给你，因为它们是交付给我的}，我们甚至也可以拒绝它的断言；因为我们无需相信这个说谎者；然而或许它说了真话，因被祂临在的能力所迫；因为它没有说：\Emph{这一切我都要给你}，因为是我的，而是说：\Emph{因为它们是交付给我的。}它没有掌握对它们的统治权，而是承认曾被托付它们，并暂时管理它们。但释经者应在适当时候探究它的说法是虚假还是真实的。\par

8. 天主是独一的，圣父，全能者，异端者的后代竟敢亵渎祂。他们竟敢亵渎\Emph{坐在革鲁宾之上的}万军的上主；他们竟敢亵渎阿多乃上主；他们竟敢亵渎在先知中的全能天主。但你要崇拜独一的天主全能者，我们主耶稣基督的圣父。逃避多神的错误，也逃避一切异端，并像\Person{约伯}那样说：\Quote{\Emph{但我要呼求全能的上主，祂行伟大而不可测的事，荣耀而无数奇妙的事}}，以及\Quote{\Emph{这一切都从全能者得到尊荣}}。愿光荣归于祂，直到永远。阿们。\par

\SourceRecord{Translated by Edwin Hamilton Gifford.；Nicene and Post-Nicene Fathers, Second Series；Vol. 7.；Edited by Philip Schaff and Henry Wace.；Buffalo, NY: Christian Literature Publishing Co.,；1894.；<http://www.newadvent.org/fathers/310108.htm>.}

% structure: p0001 source=h1 target=section reason=page-title

\section{教理讲授第九讲}

\Emph{\Quote{论天地万物的创造者，以及一切有形与无形者。}}\par

\BibleRef{约 38:2-3}\par

\Emph{\Quote{谁用无知的言语，使我的计划模糊不明？}}\par

1. 以血肉之眼看见天主是不可能的，因为无形体者不能成为肉眼所见；天主的独生子亲自作证说：\BibleQuote{从来没有人见过天主。}\BibleRef{若 1:18} 如果有人根据厄则克尔书所载，认为厄则克尔看见了他，那么圣经说了什么？\BibleQuote{这是主荣耀的显现。}\BibleRef{则 1:28} 不是主本身，而是\Quote{他荣耀的显现}，并非荣耀本身的本相。当他仅看见\Quote{荣耀的显现}而非荣耀本身时，就因恐惧而俯伏在地。既然荣耀显现的幻象已令先知恐惧不安，那么任何试图看见天主本身的人，必定丧失生命，正如经上所说：\BibleQuote{没有人看见了我还能活着。}\BibleRef{出 33:20} 为此，天主出于极大的仁爱，铺张穹苍作为他神性本体的幔帐，免得我们灭亡。这不是我的话，而是先知的话：\Quote{你若撕裂诸天，山岭必在你面前战栗，熔化倾流。}为何你惊讶厄则克尔见了\Quote{荣耀的显现}而跌倒？当达尼尔见了加俾额尔——不过是天主的仆人——便立即惊惶，俯伏在地，作为先知竟不敢回答他，直到天使转变为人子的形像。如果加俾额尔的显现已使先知战栗，那么天主本身若按其本相显现，岂不令所有人灭亡？\par

2. 神性本体固然无法以肉眼看见，但从神圣的化工中，可以略略领悟他的大能，正如撒罗满所说：\Quote{由伟大的美丽和受造物的伟大，可以推想其创造者。}他没有说从受造物可以看见创造者，而是加了\Quote{推想}二字。因为人对受造物认识越广，天主在他眼中就显得越伟大；当他的心灵因这更广阔的认识而提升时，他对天主的认识也随之增长。\par

3. 你愿知道认识天主的本性是不可能的吗？在火窑中的三个青年歌颂天主时说：\Quote{你俯视深渊，坐在革鲁宾之上，是可赞美的。}请告诉我革鲁宾的本性是什么，然后看那坐在它们之上的天主。然而，厄则克尔先知曾尽可能描述它们，说\Quote{每个有四个面孔}：一个像人，一个像狮，一个像鹰，一个像牛犊；每个有六只翅膀，遍体都有眼睛；每个下面有四个轮子，轮子也是四个面。虽然先知作了说明，但我们阅读后仍不能明白。如果我们连他所描述的宝座都不能理解，又怎能理解那坐在其上、无形无像、不可言说的天主呢？因此，探究天主的本性是不可能的，但我们能因所见化工而颂扬他的荣耀。\par

4. 我说这些，是为了信经的上下文，因为我们宣认：\Quote{我们信唯一的天主，全能的圣父，天地万物的创造者，以及一切有形与无形者。}好使我们牢记：我们主耶稣基督的父，与创造天地的乃是同一位，并使自己免受无神异端者的歧途所害；他们竟敢诽谤这整个世界全智的工匠，这些人是肉眼看见，却蒙蔽了理智的眼睛。\par

5. 他们对天主如此浩大的创造有什么可指责的呢？——他们本应在看见穹苍的拱顶时惊叹，本应崇拜那位以穹苍为居所、从水的流动本性中形成坚固天空的天主。因为天主说：\BibleQuote{在水中间要有穹苍。}\BibleRef{创 1:6} 天主一次发言，穹苍便坚立不动，永不坠落。天是水，其中的天体——太阳、月亮、星辰——是火；这些火体如何在水中运行呢？如果有人因火与水本性相反而争论，让他记起梅瑟时代在埃及，火在冰雹中燃烧的情景，并观察天主全智的化工。既然需要水来耕种土地，他便在上方造了水的穹苍，以便大地需要雨水浇灌时，穹苍能因其本性而准备就绪。\par

6. 还有，观察太阳的结构岂不令人惊奇？它看似一个微小的形体，却蕴含着巨大的力量；从东方出现，将光芒照至西方；圣咏作者描述它清晨升起时说：\Quote{它像新郎一样走出洞房。}他描述太阳初见人类时的光辉与温良：因为当它在正午奔驰时，我们常躲避它的烈焰；但在升起时，它像新郎一样令众人喜爱。\par

也要注意祂的安排（更确切地说，不是祂的安排，而是那以命令规定其行程者的安排）：在夏季，祂升得更高，使白天更长，给人充足的时间工作；在冬季，祂缩短行程，使寒冷时期延长，黑夜变长，有助于人的休息，也有助于地上出产的结果。也要看白天如何按适当次序交替回应：夏季增长，冬季缩减，而在春秋两季则彼此给予相等的间隔。黑夜也完成类似的循环，以致圣咏作者论到它们说：\Quote{\Emph{\BibleQuote{日复一日吐露言语，夜复一夜传达知识}}。}因为对于那些没有耳朵的异端者，它们几乎在大声呼喊，并借其良好秩序说：除了创造他们的造物主之外，没有别的天主，是祂设定了他们的界限，并安排了宇宙的秩序。\par

7. 但不要让任何人容忍那些说光是一个造物主、黑暗是另一个造物主的人：因为要记住依撒意亚如何说：\Quote{\BibleQuote{我是造光、创造黑暗的天主}}。人啊，你为什么为此烦恼？你为什么为给你休息的时间而生气？若不是黑夜必然带来休息，仆人便无法从主人那里得到安歇。我们常常在白天劳累之后，夜间得到休息，昨天因辛劳而疲惫的人，因夜间休息而在早晨精力充沛。还有什么比黑夜更有助于智慧呢？因为此时我们常常将天主的事摆在心头；此时我们阅读并默想神圣的\WorkTitle{圣经}。我们的心何时最适合咏唱圣咏和祈祷？难道不是夜间吗？我们何时常想起自己的罪过？难道不是夜间吗？因此，我们不要接受那个恶念，即以为黑暗是另一个创造者：因为经验也表明，黑暗是善的、有益的。\par

8. 他们不仅应对太阳和月亮的安排感到惊奇和赞叹，也应对星辰有序的合唱、它们无阻的行程、以及各自适时地升起而惊叹：有些是夏天的标志，有些是冬天的标志；有些表示播种的季节，有些显示航海的开始。一个人坐在船里，航行在无边的波浪中，他仰望星辰来掌舵。关于这些事，\WorkTitle{圣经}说得很好：\Quote{\BibleQuote{让他们作为标记、时节和年月}}\BibleRef{创 1:14}，而不是作为占星术和命理的荒唐之说。但请注意，祂也怎样通过逐渐增加赐给我们白昼之光：我们不是一下就看到太阳升起，而是先有一点点光线跑在前面，使眼睛通过预先试炼能看见祂更强的光芒；也要看祂怎样借月光来减轻黑夜的黑暗。\par

9. \Quote{\BibleQuote{雨的父亲是谁？谁生了露珠？}}\BibleRef{约 38:28}是谁把空气凝结成云，命它们携带雨水，时而\Quote{从北方带来金黄色的云彩}，时而又将它们变成统一的外观，再变成多种圆圈和其他形状？\Quote{\BibleQuote{谁能用智慧数算云彩？}}\BibleRef{约 38:37}\BibleBook{约伯}{记}论到这些说：祂知道云彩的分离，并\Quote{使天垂向大地}；\Quote{祂以智慧数算云彩}；\Quote{云彩在祂以下不破裂}。因为如此大量的水压在云彩上，它们却不破裂，而是极有秩序地降在地上。谁\Quote{从府库中带出风来}？\Quote{\BibleQuote{谁，如我们先前所说，生了露珠？冰从谁的胎中出来？}}\BibleRef{约 38:28}因为它的实质像水，力量像石头。水一时变成\Quote{\BibleQuote{雪如羊毛}}，一时又服事那\Quote{\BibleQuote{散雾如灰}}者，一时又变成石质的物体，因为\Quote{\BibleQuote{祂随意支配水}}。水的性质是统一的，但其作用力是多方面的。水在葡萄树中变成\Quote{\BibleQuote{喜乐人心的酒}}；在橄榄中变成\Quote{\BibleQuote{滋润人面的油}}；也变成\Quote{\BibleQuote{强壮人心的大麦饼}}和祂所创造的各种果实。\par

10. 这些奇事应当产生什么效果？造物主应该被亵渎，还是更应该被崇拜？我还没有谈及祂智慧中不可见的工作。请注意，我请求你们，注意春天和各类花朵，它们彼此相似却又各不相同：玫瑰的深红和百合的纯白，这些出自同一雨水和同一土地，是谁使它们不同？谁塑造了它们？请观察一下这份精心的照料：同一棵树上有的是为遮阴，有的是为各种果实，而工匠是同一的。同一葡萄树，一部分用于焚烧，一部分用于嫩枝，一部分用于叶子，一部分用于卷须，一部分用于果串。\par

也要赞叹芦苇周围结节的粗大，正如工匠所造的。从同一片大地生出爬行动物、野兽、牲畜、树木、食物，以及金、银、铜、铁、石头。水的性质是唯一的，但从它产生鱼和鸟的实体；因此前者在水里游，后者在空中飞。\par

11. \Emph{这广大而宽阔的海洋，其中有不可胜数的爬行物}。谁能描述其中鱼类的美丽？谁能描述鲸鱼的巨大，以及两栖动物的本性，它们如何在陆地上和水中生活？谁能说出海的深度和广度，或它巨大波浪的力量？然而它停留在自己的界限内，因为那一位曾说过：\Emph{\BibleQuote{你只可到这里，不可越过，你狂傲的波浪要止息在你自身之内}}。\BibleRef{约 38:11} 这海也清楚地显示那加于它的命令的话语，因为当它涌上来之后，便在沙滩上留下一条可见的波浪线，仿佛向看见的人表明它没有越过指定的界限。\par

12. 谁能辨别空中飞鸟的本性？有的鸟发出悦耳的歌声，有的鸟翅膀上点缀着各种图案，有的鸟飞到半空中静止不动，如同鹰：因为出于天主的命令，\Emph{鹰展开翅膀，静止不动，向南观望}。若你连最无知的鸟的翱翔都无法辨别，又怎能理解万物的创造者呢？\par

13. 哪一个人知道所有野兽的名字？或者谁能准确辨别每一种野兽的生理构造？但若连野兽的名字都不知道，我们又怎能理解它们的创造者？天主的命令只有一个，它说：\Emph{\Quote{地要生出野兽、牲畜和爬行物，各从其类}}\BibleRef{创 1:24}，从同一块土地，因一个命令，产生了不同的本性：温顺的绵羊和食肉的狮子，以及无理性动物各种本能，与人类的不同性格相似；狐狸显明人里面的狡猾，蛇显明朋友中的毒害诡诈，嘶鸣的马显明年轻人的放荡\BibleRef{耶 5:8}，勤劳的蚂蚁则激励懒惰和迟钝的人：因为当人在青年时代虚度光阴时，就被无理性的动物所教导，受圣经的责备说：\Emph{\Quote{懒惰的人哪，你去察看蚂蚁，看它所行的，就可得智慧}}。因为当你看见它在好季节储存食物时，要效法它，为自己储存善工的果实以备来世。又如：\Emph{\Quote{你去察看蜜蜂，学习它何等勤劳}}：它如何环绕各种花朵，为你收集蜂蜜，使你也藉着遍览圣经，为自己获得救恩，并充满圣经，得以说：\Emph{\Quote{你的言语在我上颚何等甘甜，在我口中比蜜更甜}}。\par

14. 难道创造者不更应当受荣耀吗？为什么呢？如果你不知道万物的本性，所造之物就立刻变得无用了吗？你能知道所有草药的功效吗？或者你能知道从每种动物而来的所有益处吗？甚至直到如今，从毒蛇中也产生了为保全人的解毒剂。但你会对我说：\Quote{蛇是可怕的。}你要敬畏主，它就不能伤害你。\Quote{蝎子会蜇人。}你要敬畏主，它就不会蜇你。\Quote{狮子是嗜血的。}你要敬畏主，它就会卧在你旁边，如同在达尼尔身旁一样。但动物的行动也确实奇妙：有的，如蝎子，刺的锋利；有的以牙齿为力量；有的以爪子搏斗；而蛇怪的力量在于它的目光。所以，从这多样的工艺，要理解创造者的能力。\par

15. 但这些事你或许不知道：你与外面的受造物毫无共同之处。现在进入你自己，从你的本性思考它的创造者。在你身体的构造中有什么可挑剔的吗？你若管理好自己，你的任何肢体就不会产生罪恶。起初亚当和厄娃在乐园中赤身露体，但并非因他的肢体而该被驱逐。所以，肢体不是罪的起因，而是那些滥用自己肢体的人；它们的创造者是智慧的。谁预备了子宫的深处为生育？谁赋予其中无生命之物生命？\Emph{\Quote{谁用筋和骨编织了我们，用皮肤和肉包裹了我们}}\BibleRef{约 10:11}，并且孩子一出生，就从乳房涌出奶的溪流？婴儿如何长成孩童，孩童长成青年，然后成人，仍保持不变又进入老年，而没有人注意到每日精确的变化？食物的一部分如何变为血，另一部分分离为排泄物，另一部分变为肉？谁使心脏不间断地跳动？谁用眼睑的护栏智慧地保护了眼睛的娇嫩？因为关于眼睛的复杂而奇妙的构造，医生们厚厚的书卷几乎无法解释。谁将同一气息分给全身？人啊，你看到创造者，你看到智慧的造物主。\par

16. 我的论述现在对这些点进行了详尽的讨论，遗漏了许多，甚至上万的其他事物，特别是无形和不可见的事物，为叫你们憎恶那些亵渎智慧而良善的创造者的人，并从所听所读，以及你们自己能发现或构思的一切，\Emph{\Quote{从受造物的伟大和美丽，可以比例地看到它们的创造者}}\BibleRef{智 13:5}，并以虔敬的敬畏屈膝敬拜世界的创造者，我指感觉和思想的世界，有形无形的世界，你们要以感恩和圣洁的口舌，不停息的嘴唇和心，赞美天主说：\Emph{\Quote{上主，祢的作为何其奇妙，祢以智慧造了一切。}}因为尊荣、光荣和威严归于祢，从今世直到永远。阿们。\par

% structure: p0023 source=h2 target=subsection reason=relative-heading-rank; base=h2

\subsection{附录}

\EditorialNote{——在此讲道以\Quote{\WorkTitle{圣巴西略论天主之不可理解}}为题保存的手稿中，部分内容经过修改以适应该主题；但结尾尤其因大量增补和异文而突出，值得在此重录。因此，代替§15中\Quote{\OriginalTerm{有何可指摘之处？}{\GreekText{τ}\'\GreekText{ι}\ \GreekText{μεμπτ}\'\GreekText{ον}}}及后续文字，前述手稿作如下表述：}\par

\Quote{身体的构造有何可指摘之处？你出来，到人前说话。管束你的意志，你的任何肢体便不会生出邪恶。因为每一个肢体都是为我们的用途而必须被造的。你要以虔敬节制你的理性，顺从天主的诫命，那么这些肢体在它们被造所服务的用途上劳作服务时，就不会犯罪。你若不愿意，眼睛就不会妄视，耳朵就不会听不当听之言，手就不会为贪婪伸出，脚就不会走向不义，你便不会有异常的情爱，不犯奸淫，不贪图邻人的妻子。从你心中驱除邪恶的意念，如天主所造的那样，你反而要感谢你的造主。}\par

\Quote{起初亚当赤身露体，在乐园中过着安逸的生活；他领受了诫命，却没有遵守，并邪恶地伸出了手（并非因为手愿意如此，而是因为他的意志将手伸向被禁止之物），因他的悖逆，他也失去了所领受的美好事物。因此，肢体对那些使用它们的人来说，并非罪的缘由，而是邪恶的心思，如主所说：\BibleQuote{因为从心里发出来的是恶念、邪淫、奸淫、嫉妒等类}。你选择什么，你的肢体就为你服务；它们被造得极好，用以服侍灵魂：它们被提供给你，作为你理性的仆人。你要以虔敬的动向好好引导它们；以敬畏天主勒住它们；使它们服从于节制和禁欲的愿望，它们便永远不会起来反对你，欺凌你；反而要保护你，更大力地帮助你得胜魔鬼，同时期待那不朽而永恒的胜利冠冕。谁开启子宫的内室？谁，等等。}\par

在同一段的结尾，\Quote{智慧的造主}之后有如下文字：\Quote{你要因祂不可测度的化工光荣祂，而对于你无法认知的那一位，不要好奇探问祂的本质。你保持沉默，凭信德朝拜祂，依照天主的圣言，这比胆大妄为地追寻那些你既不能达到、圣经也未传授给你的事，更为有益。我的论述已对此详细阐述，好使你对那些亵渎智慧和美善的工匠的人心生憎恶，反而自己也这样说：\BibleQuote{上主，祢的化工何等奇妙；祢以智慧造成万物。}愿荣耀、权能及崇拜归于祢，偕同圣神，从今世直到永远。阿们。}\par

\SourceRecord{Translated by Edwin Hamilton Gifford.；Nicene and Post-Nicene Fathers, Second Series；Vol. 7.；Edited by Philip Schaff and Henry Wace.；Buffalo, NY: Christian Literature Publishing Co.,；1894.；<http://www.newadvent.org/fathers/310109.htm>.}

% structure: p0001 source=h1 target=section reason=page-title

\section{要理讲授第十讲}

\Emph{论\Quote{及在独一主耶稣基督}一句，并选读\BibleBook{格林多}{前书}}\par

\Emph{\Quote{因为虽有称为神的，或在天上，或在地上；然而我们只有一位天主，就是圣父，万物出于祂，而我们归于祂；也只有一位主耶稣基督，万物藉祂而有，我们也藉祂而有。}}\par

1. 凡受教而信\Quote{惟一天主全能圣父}者，亦当信祂的独生子。因为\BibleQuote{否认子的，也否认圣父}。\BibleRef{若一 2:23} \BibleQuote{我就是门}，耶稣说；\BibleQuote{除非藉着我，没有人能到圣父那里去}。因为若否认门，关于圣父的知识就被隔绝了。\BibleQuote{除了子，和子愿意启示的人，没有人认识圣父}。\BibleRef{玛 11:27} 因为若否认那启示者，你就停留在无知中。福音中有句话说：\BibleQuote{不信从子的，不会看见生命，天主的义怒却常在他身上。}\BibleRef{若 3:36} 因为当独生子被藐视时，圣父会发怒。因为若仅是一个士兵受侮辱，国王尚且恼怒；若是一位高贵的军官或朋友受侮辱，他的怒气更大；但若有人蔑视国王的独生子本人，有谁能平息圣父对其独生子的义怒呢？\par

2. 因此，若有人愿对天主表示虔诚，就当敬拜圣子，否则圣父不接受他的事奉。圣父从天上大声说：\BibleQuote{这是我的爱子，我所喜悦的。}\BibleRef{玛 3:17} 圣父已喜悦了，除非你也在祂内喜悦，否则你没有生命。不要被犹太人狡猾的话带走，他们说，只有一位天主；但要知道天主是唯一的同时，也要知道天主的独生子。我不是第一个说这话的，而是圣咏作者以子的身份说：\BibleQuote{主对我说：你是我的儿子}。因此，不要理会犹太人说什么，而要理会先知说什么。你奇怪吗？那些用石头打死并杀害先知的人，对先知的话置之不理。\par

3. 你要信\Quote{我等惟信独一主耶稣基督}，天主的独生子。因为当我们说\Quote{独一主耶稣基督}时，是为表明祂的子嗣是\Quote{独生}的；当我们说\Quote{独一}时，是为使你不以为别有另一位；当我们说\Quote{独一}时，是为使你不致亵渎地将祂行动的许多名号分散给许多子。因为祂被称为门；但不要按字面将这名号理解为木门，而是一道属灵的、活的门，分辨那些进入的人。祂被称为道路\BibleRef{若 14:6}，并非用脚踩踏的路，而是引向天上圣父的路；祂被称为羊，不是无理性的羊，而是那只藉其宝血洗净世界罪恶的羊，它在剪毛的人前被牵，并且知道何时缄默。这只羊又被称为牧人，祂说：\BibleQuote{我是善牧}\BibleRef{若 10:11}：称为羊是因为祂的人性，称为牧人是因为祂天主性的慈爱。你想知道有理性的羊吗？救主对宗徒们说：\BibleQuote{看，我派遣你们如同羊在狼群中}。此外，祂被称为狮子，不是吞吃人的狮子，而是藉这名号表明祂王者、坚定而自信的本性：祂被称为狮子，也是针对那咆哮着吞吃受骗者的仇敌狮子\BibleRef{伯前 5:8}。因为救主来临，并非改变自己本性的温柔，而是作为强壮的\BibleQuote{犹大支派的狮子}，拯救信祂的人，践踏仇敌。祂被称为石头，不是无生命的、人手所凿的石头，而是\BibleQuote{主要的角石}，凡信祂的，\BibleQuote{必不致蒙羞}。\par

4. 祂被称为基督，并非被人手所傅油，而是由圣父为人类的大司祭职而永远傅油。祂被称为死者，并非如所有在阴府中的死者那样住在死者中，而是作为\BibleQuote{在死者中独一自由的}。祂被称为人子，并非像我们每个人那样从地上出生，而是因为祂要驾云来临审判活人死人。祂被称为主，并非如人间那些被称为主的而不当，而是因为拥有本性的、永远的主位。祂被称为耶稣，此名正符合祂的医治救恩。祂被称为子，并非因被收养而提升，而是自然所生。我们救主的称号很多；因此，免得祂的众多名号使你以为有许多子，并因异端者的错误（他们说，基督是一位，耶稣是另一位，门又是另一位，等等），信德预先保护了你，正确地说：\Quote{在独一主耶稣基督}。因为虽然称号众多，但其主体只有一个。\par

5. 但救主以各种形态来到各人面前，为了每个人的益处。因为对那些需要喜乐的人，祂成为葡萄树；对那些想要进入的人，祂站在门前作为门；对那些需要献上祈祷的人，祂是中间的大司祭。同样，对于那些有罪的人，祂成为羊，好为他们牺牲。\BibleQuote{祂对一切人，成为一切}\BibleRef{格前 9:22}，在自己本性中保持祂所是。因为如此保持，并坚定地持守祂子嗣的尊严，祂适应我们的软弱，如同一位优秀的医生或慈悲的老师；尽管祂是真正的主，并非因晋升而得到主位，而是从本性拥有主位的尊严，并且祂被称为主并非如我们那样不当，而是真实如此，因为藉圣父的命令，祂是自己作品的主。因为我们的主权是对拥有平等权利和相似激情的人行使，而且常常是对我们的长辈，一个年轻的主人往往统治年老的仆人。但在我们的主耶稣基督身上，主权并非如此：祂先是创造者，后是主：祂首先藉圣父的旨意创造了万物，然后祂是祂所创造之物的主。\par

6. \Emph{\BibleQuote{主基督}}就是那\Emph{\BibleQuote{生於达味城}}的。\BibleRef{路 2:11} 你岂愿知道基督在降生成人之前就已与圣父同为天主，好使你不只以信德接受这句话，也能从旧约得到明证？到第一卷书 \BibleBook{}{创世纪} 看看：天主说：\Emph{\BibleQuote{让我们造人}}，不是\BibleQuote{按我的肖像}，而是\Emph{\BibleQuote{按我们的肖像}}。\BibleRef{创 1:26} 亚当受造后，神圣的作者说：\Emph{\BibleQuote{天主于是造了人；按天主的肖像造了他}}。因为他不将天主性的尊贵仅限于圣父，也包含了圣子：为表明人不仅是天主的工作，也是我们主耶稣基督的工程，祂自己也是真天主。这位与圣父一同工作的主，也在索多玛的事上与祂同工，正如圣经所言：\Emph{\BibleQuote{上主从天上的上主那里，把硫磺和火降在索多玛和哈摩辣}}。这位主后来也被梅瑟看见，尽他所能看到的。因为主是爱人至深的，常常迁就我们的软弱。\par

7. 此外，为使你确信这就是那被梅瑟看见的一位，请听保禄的证言，他说：\Emph{\BibleQuote{因为众人都饮过那随从他们的属神磐石；那磐石就是基督}}。\BibleRef{格前 10:4} 又说：\Emph{\BibleQuote{因着信德，梅瑟离开了埃及}} \BibleRef{希 11:27}，不久后他又说：\Emph{\BibleQuote{以基督的凌辱比埃及的宝藏更丰富}}。这位梅瑟向祂说：\Emph{\BibleQuote{求祢显示祢自己}}。你看见那些时代的先知也看见了基督，即每人按各自的能力所见的。\Emph{\BibleQuote{求祢显示祢自己，使我更清楚地认识祢}}。但祂说：\Emph{\BibleQuote{没有人看见我的面容而能生存}}。\BibleRef{出 33:20} 为此，因为没有人能看见天主性的面容而存活，祂就取了人性的面容，好使我们看见这面容而得以生存。然而，当祂愿意稍微显示那时，\Emph{\BibleQuote{祂的面容发光如太阳}} \BibleRef{玛 17:2}，门徒们吓得俯伏在地。如果祂的肉身容貌，所放射的光并非出于那做工者的全部权能，而是按门徒的容量，就已使他们惊惶，连这样都担当不起，那么谁能瞻仰天主性的威能呢？主说：\Quote{你所渴望的是一件大事，梅瑟；我赞许你不懈的渴望，\Emph{\BibleQuote{我要为你做这事}}，但按你所能承受的。\Emph{\BibleQuote{看，我要把你放在磐石的裂缝中}} \BibleRef{出 33:22}；因为你渺小，将居于窄处。}\par

8. 现在，为了犹太人，我希望你谨慎接受我所说的。因为我们的目的是证明主耶稣基督曾与圣父同在。于是主对梅瑟说：\Emph{\BibleQuote{我要在你面前经过，显示我的荣耀，并在你面前宣告上主的名号}}。祂自己就是主，却宣告哪一位主？你看祂是如何含蓄地教导关于圣父和圣子的虔敬教义。接着，记载的文字如下：\Emph{\BibleQuote{上主在云中降下，站在他那里，宣告上主的名号。上主从他面前经过，宣告说：\InnerQuote{上主，上主，慈悲宽仁的天主，缓于发怒，富于慈爱和真理，持守正义，向千万人施行仁慈，赦免过犯、愆尤和罪恶。}}}然后，接下来记载：\Emph{\BibleQuote{梅瑟俯首朝拜}} \BibleRef{出 34:8} 在宣告圣父的主面前，并且说：主，求你在我们中间前行。\par

9. 这是第一个证明：现在请接受第二个明证。\Emph{\BibleQuote{上主对我主说：你坐在我右边}}。上主这样说，不是对仆人，而是对万有之主，自己的儿子，他将万有置于他的脚下。\Emph{\BibleQuote{但当他说万物已置于他以下时，显然那将万物置于他以下的除外}}，以及后面的话：\Emph{\BibleQuote{为使天主成为万物之中的万有}}。\BibleRef{格前 15:27} 独生子是万有之主，但他是服从圣父的儿子，因为他没有强夺主权，而是因圣父自己的旨意由本性接受。因为圣子没有夺取它，圣父也不吝于赐予它。他就是那说：\Emph{\BibleQuote{一切都由我的父交给了我}}；\Quote{交给了我，并非好像我以前没有，我妥善保管，不抢夺赐予者的东西。}\par

10. 天主子就是主：祂是主，祂生于犹太的白冷，正如天使向牧人们所说的：\BibleQuote{\Emph{我给你们报告一个极大的喜讯，因为今天在达味城中，为你们诞生了一位主基督}}\BibleRef{路 2:10}；另一位宗徒也在别处说：\BibleQuote{\Emph{天主藉着耶稣基督——祂是万有的主——传和平的福音给以色列子民}}\BibleRef{宗 10:36}。但当他提到\Emph{万有}时，你不要从他的主权中排除任何事物：因为无论是天使、总领天使、宰制者、掌权者，或是任何宗徒所称述的受造物，都在圣子的主权之下。祂是天使的主，正如你在福音中所见：\BibleQuote{\Emph{于是魔鬼离开了祂，天使遂前来伺候祂}}\BibleRef{玛 4:11}；因为圣经没说他们帮助了祂，而是说他们\Emph{伺候了祂}，即如同仆人一样。当祂即将由一位童贞女诞生时，加俾额尔那时就是祂的仆人，领受祂的服事作为特殊的尊荣。当祂将要下到埃及，为推翻埃及人手造的神时，又\BibleQuote{\Emph{有天使在梦中显现给若瑟}}\BibleRef{玛 2:13}。当祂被钉十字架并复活后，一位天使报喜讯，并如同忠信的仆人对妇女们说：\BibleQuote{\Emph{你们去告诉祂的门徒，祂已复活，并在你们之前往加里肋亚去；看，我已经告诉了你们}}；几乎就好像他说：\Quote{我没有疏忽我的命令，我郑重地告诉你们；如果你们轻视它，罪责不在我，而在那些轻视的人身上。}这就是那唯一的主耶稣基督，正是刚才读经所论的：\BibleQuote{\Emph{因为虽然有被称为神的，不论在天上或在地上，}等等，\Emph{然而我们只有一个天主，就是圣父，万物出于祂，而我们也归于祂；也只有一个主，就是耶稣基督，万物藉着祂，我们也藉着祂}}\BibleRef{格前 8:5}。\par

11. 祂有两个名字：耶稣基督；耶稣，因为祂拯救；基督，因为祂是司祭。受默感的先知梅瑟知道这事，便将这两个称号赐给了两个出类拔萃的人：他将自己在治理上的继承人\OriginalTerm{欧瑟亚}{Auses}改名为耶稣；又将自己的兄弟亚郎称为基督，以便借着两位倍受认可的人，同时预表那将要来临的唯一耶稣基督的大司祭职和君王职。因为基督如同亚郎一样是大司祭，因为祂\BibleQuote{\Emph{没有自己争取成为大司祭，而是那向祂说话的那位：你永为司祭，按照默基瑟德的品位}}。农的儿子若苏厄在许多事上是祂的预像。因为当他开始治理百姓时，他从约但河开始\BibleRef{苏 3:1}，从此基督受洗后也开始宣讲福音。农的儿子指定十二个人来分产业；耶稣派遣十二位宗徒作为真理的宣讲者到全世界。预像中的耶稣拯救了相信的妓女辣哈布；真实的耶稣说：\BibleQuote{\Emph{看，税吏和娼妓要在你们之先进入天主的国}}\BibleRef{玛 21:31}。在预像的时代，仅凭呼喊耶里哥城墙就倒塌了；因为耶稣说：\BibleQuote{\Emph{这里将没有一块石头留在另一块石头上}}\BibleRef{玛 24:2}，那对面的犹太人的圣殿已经倒塌，其倒塌的原因不是诅咒，而是犯法者的罪。\par

12. 只有一位主耶稣基督，一个奇妙的圣名，由先知们间接预先宣告。因为依撒意亚先知说：\BibleQuote{\Emph{看，你的救主来临，带着祂的酬报}}。耶稣在希伯来语中的解释是\Quote{救主}。因为先知之恩预见犹太人对其主的杀害之心，便遮掩了祂的名字，免得他们提前清楚知道而轻易图谋害祂。但祂公开被称为耶稣，不是由人，而是由一位天使，这天使不是凭自己的权威来的，而是受天主德能的派遣，向若瑟说：\BibleQuote{\Emph{不要怕娶你的妻子玛利亚，因为那在她内受孕的是出于圣神。她要生一个儿子，你要给祂起名叫耶稣}}\BibleRef{玛 1:20}。他立刻解释了这名字的理由，说：\BibleQuote{\Emph{因为祂要把自己的民族从他们的罪恶中拯救出来}}。试想，祂尚未出生，怎能有一个\Emph{民族}，除非祂在出生以前就已经存在。先知也以祂的身份说：\BibleQuote{\Emph{从我母腹中祂就提了我的名字}}\BibleRef{依 49:1}；因为天使预言祂应称为耶稣。又关于黑落德的阴谋，他再次说：\BibleQuote{\Emph{并在祂手的荫庇下祂隐藏了我}}。\par

13. 因此，耶稣按希伯来语意为\Quote{救主}，而在希腊语中意为\Quote{医治者}；因为祂是灵魂和肉身的医师，精神的治愈者，治愈身体的瞎子，引导理智进入光明，医治明显瘸腿的，引导罪人的脚步走向悔改，对瘫痪者说：\BibleQuote{\Emph{不要再犯罪了}}，和\BibleQuote{\Emph{拿起你的床走吧}}。因为身体因灵魂的罪而瘫痪，祂先服事灵魂，以使痊愈延伸到身体。所以，如果有人灵魂因罪受苦，这里有医生为他；如果有人信德微小，让他对祂说：\BibleQuote{\Emph{帮助我的不信吧}}\BibleRef{谷 9:24}。如果还有人被身体疾病缠绕，不要不信，但要走近；因为耶稣也医治这样的疾病，让他知道耶稣就是基督。\par

14. 犹太人承认祂是耶稣，但不进一步承认祂是基督。因此宗徒说：\BibleQuote{\Emph{谁说谎呢？不是那否认耶稣是基督的吗？}}\BibleRef{若一 2:22}？但基督是大司祭，\BibleQuote{\Emph{祂的司祭职不传给别人}}\BibleRef{希 7:24}，祂的司祭职既非从时间开始，也无人继承祂的大司祭职；正如你们在主日听我们聚会中谈论\BibleQuote{\Emph{按照默基瑟德的品位}}这句话时所听见的。祂的大司祭职不是来自肉身的继承，也非由人配制的油所傅油，而是由圣父在万世之前所傅；祂远超过他人，正如\Quote{\Emph{带着誓言}}被立为司祭：\BibleQuote{\Emph{因为他们做司祭没有誓言，但祂却有誓言，因为那位向祂说：\InnerQuote{主发誓了，绝不后悔}的人。}}\BibleRef{希 7:21}圣父的意愿足以作为保证；但保证的方式是双重的，即意愿之后还有誓言，\BibleQuote{\Emph{好藉着两件不可更改的事，在这样的事上，天主不可能说谎，使我们这些接受基督耶稣为天主之子的人获得坚强的鼓励}}，好为我们的信德。\par

15. 这位基督来临后，犹太人否认了祂，但魔鬼却承认了祂。但祂的祖先达味并非不认识祂，他曾说：\Quote{\Emph{我为我的受傅者预备了一盏灯}}。这盏灯，有人解释为预言之光，也有人解释为祂从童贞女所取的血肉，按照宗徒的话：\Quote{\Emph{但我们有这宝贝放在瓦器中}}。先知并非不认识祂，他说：\Quote{\Emph{并向人类宣告祂的基督}}。梅瑟也认识祂，依撒意亚认识祂，耶肋米亚认识祂；没有一位先知不认识祂。甚至魔鬼也认出祂，因为祂斥责它们，经上说：\BibleQuote{\Emph{因为他们知道祂是基督。}}\BibleRef{路 4:41}大司祭们不认识祂，魔鬼却承认了祂；大司祭们不认识祂，一个撒玛黎雅妇人却宣讲了祂，说：\BibleQuote{\Emph{你们来看！有一个人说出了我所作的一切事。莫非这就是基督吗？}}\BibleRef{若 4:29}？\par

16. 这位耶稣基督来临，\BibleQuote{\Emph{作了未来美好事物的大司祭}}\BibleRef{希 9:11}；因祂天主性的慷慨，将祂自己的称号分施给我们众人。因为人间的君王有其尊号，他人不得共享；但耶稣基督是天主之子，赐给了我们被称为基督徒的尊荣。但有人会说，\Quote{基督徒}这个名称是新的，以前没有用过；而新造的词语常因其怪异而受人诟病。先知事先为此作了保证，说：\BibleQuote{\Emph{但我的仆人将被称为一个新名，这名字将在地上蒙受祝福}}。让我们质问犹太人：你们是主的仆人，还是不是？那么，请展示你们的新名字。因为你们在梅瑟时代、其他先知时代、从巴比伦回来之后直到如今，被称为犹太人和以色列人；那么，你们的新名字在哪里？但我们，既然是主的仆人，拥有这个新名字：确实是\Emph{新的}，但正是\BibleQuote{\Emph{将要在地上蒙受祝福的新名字}}。这个名字抓住了全世界：因为犹太人只限于某一地区，但基督徒却达于地极；因为所宣扬的正是天主独生子的名字。\par

17. 但你们想知道宗徒们是否认识并宣讲了基督的名字，或者说基督自己就在他们内吗？保禄向听众说：\BibleQuote{\Emph{或者你们寻求那在我内说话的基督的证据吗？}}\BibleRef{格后 13:3}？保禄宣讲基督，说：\BibleQuote{\Emph{因为我们不是宣传我们自己，而是宣传耶稣基督为主，并且我们因耶稣的缘故作你们的仆人}}。那么，这是谁？从前的迫害者。哦，伟大的奇迹！从前的迫害者自己宣讲基督。但为什么？他受了贿赂吗？不，没有人用这种方式说服他。或者是因为他看见祂在地上显现而羞愧了？祂已经被接到天上去了。他出去迫害，三天后，这个迫害者在大马士革成了宣讲者。靠着什么能力？别人请朋友为朋友作证，但我却向你们提出从前的敌人作证人：你们还怀疑吗？伯多禄和若望的见证虽然有力，但仍属可疑，因为他们是祂的朋友。但一个从前是祂的敌人，后来为祂而死的人，谁还能怀疑这真理呢？\par

18. 在我讲论的这一点上，我对圣神明智的安排充满惊奇：祂将其他人的书信限制在少数，却赐给从前的迫害者保禄写作十四封书信的特权。这并不是因为伯多禄或若望较次，祂才限制恩赐；绝对不可！而是为了使教义无可置疑，祂给从前的敌人和迫害者写作更多的特权，好使我们众人都得以成为信徒。因为众人\Emph{都惊奇}于保禄，并说：\BibleQuote{\Emph{这不是那从前迫害过……的人吗？}}\BibleRef{宗 9:21}？他不是到这里来，要把我们捆绑着带到耶路撒冷去吗？保禄说：不要惊奇，我知道\BibleQuote{\Emph{我难以反抗刺棒}}；我知道\BibleQuote{\Emph{我不配称为宗徒，因为我迫害了天主的教会}}\BibleRef{格前 15:9}；但我是\BibleQuote{\Emph{在无知中}}所作的\BibleRef{弟前 1:13}：因为我以为宣讲基督就是毁灭法律，却不认识祂亲自来是为成全法律，而不是毁灭它。\BibleRef{玛 5:17}\BibleQuote{\Emph{但天主的恩宠在我身上格外丰富}}\BibleRef{弟前 1:14}。\par

19. 我亲爱的诸位，关于基督的忠实见证为数众多。圣父从天上为祂的圣子作证；圣神以鸽子的形体降下，作证；总领天使加俾额尔向玛利亚报喜，作证；童贞天主之母作证；有福的马槽之地作证。埃及作证，在主尚年幼时接待了祂的身体；西默盎作证，他曾将祂抱在怀中，说：\BibleQuote{主啊！现在可照祢的话，让祢的仆人平安去了，因为我亲眼看见了祢的救恩，即祢在万民面前所准备的。}\BibleRef{路 2:29} 女先知亚纳，一位最虔诚的寡妇，生活刻苦，也为祂作证。洗者若翰作证，他是先知中最伟大者，新盟约的首领，在他身上新旧两盟约合而为一。约但河在众河间作证，提庇黎雅海在众海间作证；盲者与瘸子作证，死人复活作证，魔鬼说：\BibleQuote{我们与祢有什么相干？耶稣，我们知道祢是谁，祢是天主的圣者。}\BibleRef{谷 1:24} 风作证，因祂的命令而平静；五个饼变成五千人的食物为祂作证。十字架的圣木作证，至今仍在我们中间可见，由此处几乎充满整个世界，藉着那些以信心取得木块的人。山谷中的棕榈树作证，它曾为欢迎祂的孩童提供树枝。革责玛尼作证，至今仍使用心者几乎看见犹达斯。哥耳哥达，即我们眼前的圣山，为我们的眼目作证；圣墓作证，以及至今仍在那里的石头。现今照耀的太阳为祂作证，那时在祂救世苦难中被遮蔽；黑暗为祂作证，那时从第六时辰到第九时辰；光为祂作证，从第九时辰直到傍晚。橄榄山作证，那座圣山，祂从那里上升至圣父；载雨的云彩为祂作证，它们迎接了它们的主；是的，天国的门也为祂作证，关于这，圣咏作者说：\BibleQuote{诸位城门，你们请抬起头来！永久的门户，你们请被举起！光荣的君王要进来。}祂从前的敌人作证，其中有蒙福的保禄，他曾短暂为敌，却长久为仆；十二位宗徒为祂作证，不仅用言语，也用自己的苦难和死亡传扬真理；\Quote{伯多禄的影子}\BibleRef{宗 5:15} 作证，它因基督之名治愈病人。手帕与围裙作证，如同从前藉基督的能力，通过保禄行了治愈。波斯人和哥特人，以及所有归化的外邦人，为祂作证，为祂而死，他们从未用肉眼看见祂；至今仍被信徒驱逐的魔鬼也为祂作证。\par

20. 这么多、这么多种类，甚至比这些更多的见证属于祂；那么，被如此见证的基督还会被不信吗？反而，若有人从前不信，如今让他相信；若有人以前已信，让他获得更大的信德增长，因相信我们的主耶稣基督，并理解他背负着谁的名字。你被称为基督徒：要珍惜这名号；不要让我们的主耶稣基督、天主子，因你而受亵渎；反而\BibleQuote{你们的光当在人前照耀}\BibleRef{玛 5:16}，使他们看见，因我们的主耶稣基督，在天的圣父受光荣。愿光荣归于祂，从现在直到永远。阿们。\par

\SourceRecord{Translated by Edwin Hamilton Gifford.；Nicene and Post-Nicene Fathers, Second Series；Vol. 7.；Edited by Philip Schaff and Henry Wace.；Buffalo, NY: Christian Literature Publishing Co.,；1894.；<http://www.newadvent.org/fathers/310110.htm>.}

% structure: p0001 source=h1 target=section reason=page-title

\section{教理讲授第十一篇}

\Emph{论\Quote{天主的独生子，由圣父在万世之前所生，真天主，万物藉祂而造成}这些话。}\par

见：\BibleRef{希 1:1}\par

\BibleQuote{天主在古时，曾多次并以多种方式，藉着先知向我们的祖先说过话；在这末期内，祂藉着自己的儿子对我们说了话。}\par

昨天我们所讲的话，已按照我们的能力，充分说明了我们在耶稣基督内怀有希望。但我们不能只相信基督耶稣，不能把祂当作许多不恰当称为基督者之一来接受。因为那些是预像性的基督，而祂是真基督；祂并非由人间晋升而得到司祭职，而是从圣父那里永远拥有司祭职的尊严。为此，信德预先保护我们，免得我们以为祂是普通的基督之一，便在信德的宣认上加上这句：我们信\Quote{唯一的主耶稣基督，天主的独生子}。\par

再者，当听到\Quote{子}时，不要想到义子，而是本性之子，是独生子，没有兄弟。这就是为什么祂被称为\Quote{独生子}，因为在天主性的尊严中，以及祂由圣父所生的事上，祂没有兄弟。但我们称祂为天主之子，并非出于我们自己，而是因为圣父自称基督为祂的子：而真正的名字是父亲给予儿女的。\par

我们的主耶稣基督昔日成了人，但众人不认识祂。因此，祂愿意教导那未知之事，便召集门徒，问他们说：\BibleQuote{人们说人子是谁？}（见：\BibleRef{玛 13:16}）——并非出于虚荣，而是愿意向他们显示真理，恐怕他们与天主、天主的独生子同住，却轻看祂，以为祂仅是一个凡人。当他们回答说，有人说是厄里亚，有人说是耶肋米亚时，祂对他们说：他们不认识情有可原，但你们，我的宗徒们，你们因我的名洁净癞病人、驱逐魔鬼、复活死人，不该不认识那藉着你们行这些奇事者。当众人都默然不语时（因为此事高过人的智力），宗徒之长、教会首位宣讲者伯多禄，既非凭狡黠的发明，也非被人的推理说服，而是因父光照他的心灵，对祂说：\BibleQuote{祢是基督}，不仅如此，还是\BibleQuote{永生天主之子}。随后有祝福降于他的话（因为确实超乎人），作为对他所言盖印，表明是父启示了他。因为救主说：\BibleQuote{西满·巴尔约纳，你是有福的，因为不是血肉启示了你，而是我在天之父}（见：\BibleRef{玛 16:17}）。因此，凡承认我们的主耶稣基督为天主之子的人，就有分于此福；但否认天主之子的人，是可怜而又悲惨的人。\par

再者，我说，当听到\Quote{子}时，不要仅在不恰当的意义上理解，而要当作真实的子、本性的子、无始的子；不是从奴役中进入更高的义子地位，而是由一种不可测悟、不可理解的受生而永远受生的子。同样，当听到\Quote{首生子}时，不要以为这是按人的方式；因为人间的首生子还有别的兄弟。在某处记载着：\BibleQuote{以色列是我的儿子，我的首生子}（见：\BibleRef{出 4:22}）。但以色列，如勒乌本一样，是受排斥的首生子：因为勒乌本上了父亲的床榻；以色列把圣父的儿子赶出葡萄园，并钉死了祂。\par

对其他人，圣经也说：\BibleQuote{你们是上主你们天主的儿女}（见：\BibleRef{申 14:1}）；在另一处：\BibleQuote{我曾说：你们是神，你们都是至高者的儿子}。\Quote{我曾说}，而不是\Quote{我曾生}。他们，当天主如此\Quote{说}时，便领受了原先没有的义子地位；但祂并非受生成为与从前不同的存在；而是从起初就成为圣父的儿子，超越一切起始和万世，是圣父的儿子，在一切事上与生祂者相似，永恒者由永恒之父所生，生命由生命所生，光之光，真之真，智之智，君之君，天主的天主，能之能。\par

那么，如果你听到福音说：\BibleQuote{耶稣基督的族谱，达味之子，亚巴郎之子}（见：\BibleRef{玛 1:1}），要理解为\Quote{按肉身}。因为祂是达味之子\Quote{在时代的末期}（见：\BibleRef{希 9:26}），但祂是天主之子\Quote{在万世之前}，无始。前者是祂先前没有而领受的；后者是祂永远拥有的，作为由圣父所生。祂有两个父亲：按肉身，一个是达味；另一个是天主，祂按天主性之父。作为达味之子，祂受时间和触摸以及族谱的限制；但作为按天主性之子，祂既不受时间、也不受地点、也不受族谱的限制；因为\Quote{谁能述说祂的世代？} \BibleQuote{天主是神}（见：\BibleRef{若 4:24}）。那为神者，以属神的方式生，既无形体，其受生不可测悟、不可理解。圣子亲自论及圣父说：\Quote{上主对我说：祢是我的儿子，我今日生了祢}。这\Quote{今日}不是最近的，而是永恒的；是无时的今日，在万世之前。\Quote{自母胎，在晨星之前，我已生了祢}。\par

所以，你要相信耶稣基督，永生天主之子，而且是独生子，按照福音所说：\BibleQuote{天主如此爱了世界，甚至赐下祂的独生子，使凡信祂的人不至丧亡，反而获得永生}（见：\BibleRef{若 3:16}）。又说：\BibleQuote{信从子的，不受审判，已出死入生。不信从子的，不得见生命，天主的义怒常在他身上}。若翰为祂作证说：\BibleQuote{我们看见了祂的光荣，正如父独生子的光荣，满溢恩宠和真理}；魔鬼战栗着说：\BibleQuote{啊！祢与我们有什么相干？耶稣，永生天主之子}（见：\BibleRef{路 4:34}）。\par

7. 祂是天主本性所生的圣子，而非义子，由圣父所生。\Emph{\BibleQuote{凡爱生祂之主的，也必爱由祂所生的}} \BibleRef{若一 5:1}；但轻视那由祂所生的，就是把侮辱加于生祂者。每当你听到天主生了，不要沉溺于属肉体的思想，不要想到可朽坏的生产，免得你犯不虔之罪。\Emph{天主是神}，祂的生产是属神的：因为肉体生出肉体，肉体的生产需要时间介入；但在圣子由圣父所生这件事上，时间并不介入。在我们而言，所生的是不完美的：但圣子是完美而生的：因为祂现在是什么，从起初也是如此，无始而生。我们被生，是从婴孩的无知进入理性的状态：人啊，你的生产是不完美的，因为你的成长是渐进的。但不要以为祂也是如此，也不要把软弱归于生祂者。因为如果祂所生的不完美，并在时间中达到完美，你就把软弱归于生祂者；若果然如此，圣父并没有从起初就赐予那如你所说时间后来赐予的。\par

8. 所以不要以为这生子是属人的，也不要以为是如同亚巴郎生了依撒格。因为亚巴郎生依撒格时，所生的不是他想要的，而是别人所赐的。但在天主圣父的生子中，既没有无知，也没有中间商量。因为说祂不知道自己所生的是什么，这是最大的不虔；说祂经过时间考虑后才成为父亲，也是同样不虔。因为天主先前没有子，后来在时间中才成为父亲；而是永远有子，生了祂不是像人生人，而是唯有祂自己知道，祂在万世之前生了祂，真天主。\par

9. 圣父既是真天主，生了与自己相似的真天主圣子；不像教师生门徒，也不像保禄对某些人说的：\Emph{\BibleQuote{我在基督耶稣内借着福音生了你们}} \BibleRef{格前 4:15}。因为在这情况下，本来不是本性上的儿子，因作门徒而成了儿子；但在前一种情况下，祂是本性上的儿子，真儿子。不像你们这些即将受光照的人，现在成为天主的儿子：因为你们也成为儿子，但借着恩宠的义子身份，正如经上所载：\Emph{\BibleQuote{凡接受祂的，祂赐给他们权能，成为天主的子女，即那些信祂名字的人：他们不是由血气，也不是由肉欲，也不是由人意，而是由天主所生的}} \BibleRef{若 1:12}。我们确实是由水和圣神所生，但基督并非这样由圣父所生。因为在祂受洗时，天主向祂说：\Emph{\BibleQuote{这是我的爱子}} \BibleRef{玛 3:17}，并没有说：\Quote{现在成了我的儿子}，而是说：\Emph{这是我的爱子}，为表明甚至在洗礼实施之前，祂就是子。\par

10. 圣父生了圣子，不像人间的理智生出言语。因为理智在我们内是实体存在的；但说出的话语消散于空气而终止。但我们知道基督不是作为说出来的一句话而被生，而是作为实体存在且生活的圣言；不是嘴唇说出而消散，而是由圣父永远地、不可言喻地、在实体上所生。因为\Emph{\BibleQuote{在起初已有圣言，圣言与天主同在，圣言就是天主}} \BibleRef{若 1:1}，祂坐在天主的右边——这圣言明了圣父的旨意，并依照祂的命令创造万物：这圣言曾降下又上升，因为说出的话语被说出后不会降下，也不会上升；这圣言讲话并说：\Emph{\BibleQuote{我在我父那里所见的事，我就讲论}} \BibleRef{若 8:38}：这圣言拥有权能，统治万有：因为\Emph{父将一切交给了子}。\par

11. 圣父生了祂，不是任何人所能理解的，而是唯有祂自己知道。因为我们并不声称说明祂以何种方式生了祂，而是坚持不是以这种方式。不仅我们不知道圣子由圣父所生，就是一切受造的本性也不知道。\Emph{\BibleQuote{你问大地，大地必会教导你}} \BibleRef{约 12:8}：虽然你询问地上的一切，它们也不能告诉你。因为大地不能说它自己的陶工和塑造者的本体。不独大地无知，太阳也是如此：因为太阳是在第四天被造的，不知道它之前三天受造的是什么；那不知道之前三天所造之物者，不能说明造物主本身。天也不会宣告这事：因为因圣父的命令，\Emph{天也被圣子立定如烟云}。\Emph{\InnerQuote{诸天之上的天}}也不能宣告这事，\Emph{\InnerQuote{天上的水}}也不能。那么，人啊，你为何因不知道连天也不知道的事而沮丧呢？不但诸天不知道这个生产，就是一切天使的本性也不知道。因为如果一个人能够升到第一层天，看见那里天使的品级，前去问他们天主如何生了祂的圣子，他们大概会说：\Quote{我们之上有更伟大崇高的存在，问他们吧。}升到第二和第三层天吧；如果你能够，到达那些宝座、宰制、率领、掌权者；即使有人能够达到它们——这是不可能的——它们也会拒绝解释，因为它们不知道。\par

12. 就我而言，我常惊异于那些胆大之人的好奇心，他们因自己想象的虔诚而陷入不敬。因为他们虽然对基督的作品——\OriginalTerm{上座者}{Thrones}、\OriginalTerm{宰制者}{Dominions}、\OriginalTerm{率领者}{Principalities}和\OriginalTerm{掌权者}{Powers}——一无所知，却试图查究他们的创造者本身。最先告诉我，最胆大的人，上座者与宰制者有何不同，然后再查究关于基督的事。告诉我什么是率领者、什么是掌权者、什么是\OriginalTerm{异能者}{Virtue}、什么是\OriginalTerm{天使}{Angel}；然后探寻它们的创造者，\BibleQuote{万物是藉着祂造的}\BibleRef{若 1:3}。但你不愿，也不能问上座者或宰制者。除了那说话圣经的圣神，还有什么知道天主深奥的事呢？\BibleQuote{知道天主深奥的事}\BibleRef{格前 2:10}但连圣神自己也没有在圣经中论及圣子由圣父受生的事。那么，你为什么忙于连圣神都未在圣经中记载的事呢？你不知道所记载的事，却忙于未记载的事？圣经中有许多问题；所记载的我们都不理解，为什么我们还要忙于未记载的呢？我们只知道天主生了一位独生子，便足够了。\par

13. 不要羞于承认你的无知，因为你和\OriginalTerm{天使}{Angels}一样无知。只有生者知道受生者，而受生于祂的也知道生者。生者知道自己所生的；圣经也证实受生者是天主。因为\BibleQuote{就如父在自己内有生命，同样也赐给子在自己内有生命}\BibleRef{若 5:26}；并\BibleQuote{致使众人尊敬子如同尊敬父}；以及\BibleQuote{就如父随意赐予生命，子也随意赐予生命}。生者并未遭受任何损失，受生者也一无所缺（我知道这些事我说过多次，但为你们的安全，才反复说）：生者没有父，受生者没有兄弟。生者没有变为子，受生者也没有成为父。只有一位父，就有一位独生子：没有两个\OriginalTerm{非受生者}{Unbegotten}，也没有两个\OriginalTerm{独生子}{Only-begotten}；而是一位父，非受生（因为非受生者没有父）；和一位子，由父永恒而生；不是生于时间中，而是生于万世之前；不是因进步而增长，而是受生即如今所是。\par

14. 我们相信天主的独生子，由父所生，是真天主。因为真天主不会生出假神，正如我们所说的，祂并非先思量后生子，而是永恒地生子，其速度远超过我们的言语或思想：因为我们在时间中说话，消耗时间；但在天主的权能中，生子是超时间的。并且，正如我常说的，祂并未将子从无中带入存在，也未将不存在的纳入子嗣关系；但父是永恒的，永恒地且不可言说地生了一位独生子，祂没有兄弟。也没有两个第一本原；因为\BibleQuote{父是子的头}\BibleRef{格前 11:3}；本原只有一位。因为父生子是真天主，称为厄玛奴耳；\BibleQuote{厄玛奴耳，译出来就是\InnerQuote{天主与我们同在}}\BibleRef{玛 1:23}。\par

15. 你愿意知道那由父所生、后来成为人的那位是天主吗？听先知说：\BibleQuote{这是我们的天主，无人能与祂相比。祂发现了知识的一切途径，并将其赐给祂的仆人雅各伯和祂所爱的以色列。此后，祂在地上显现，与人来往。}你在此看见了天主成为人，是在梅瑟赐予法律之后吗？再听一个关于基督天主性的见证，就是刚才所读的：\BibleQuote{天主，你的宝座永永远远}\BibleRef{希 1:8}。恐怕因祂在此肉身中的显现，人们以为祂在此后才被提升至天主性，因此圣经清楚地说：\BibleQuote{所以天主，你的天主，用喜乐的油傅抹了你，胜过你的同伴}。你看见基督作为天主，被天主圣父傅油了吗？\par

16. 你愿意再听关于基督天主性的第三个见证吗？听依撒意亚说：\BibleQuote{埃及劳碌，厄提约匹亚的货物；}不久又说：\BibleQuote{在祢内他们祈求，因为天主在祢内，除祢以外没有神。因为祢是天主，我们却不知道，以色列的天主、救主。}你看见子是天主，祂内拥有天主父：祂所说的几乎与在福音中所说的一样：\BibleQuote{父在我内，我在父内}\BibleRef{若 14:11}。祂并没有说\Quote{我是父}，而是\Quote{父在我内，我在父内}。祂也没有说\Quote{我与父是一}，而是\Quote{我与父是同一}，好叫我们既不分离他们，也不将子—父混淆。他们是一，是由于属于天主性的尊严，因为天主生了天主。就他们的王国而言，他们是一；因为父并非统治这些，子统治那些，像阿贝沙隆一样反抗自己的父；而是父的王国同样也是子的王国。他们是一，因为他们之间没有不和与分裂；因为父所愿的，子也愿相同的事。他们是一，因为基督的创造工程无非是父的；因为万物的创造是同一的，父藉着子创造了一切：\BibleQuote{因为祂一发言，万物便造成；祂一命令，万物便创造}，圣咏作者说。因为发言者是对聆听者发言；命令者是对在祂面前的那一位发命令。\par

17. 圣子就是真实天主，祂将圣父包含在自己内，并非转变为圣父，因为圣父并未降生成人，而是圣子降生成人。让我们自由地宣讲真理：圣父并未为我们受苦，而是圣父派遣了那受苦者。我们既不要说\Quote{曾有一时圣子不存在}，也不要承认圣子就是圣父；而要走在君王的大道上，既不偏左也不偏右。既不要因想尊崇圣子而称祂为圣父，也不要因想尊崇圣父而以为圣子是受造物之一。当经由独生圣子敬拜独一圣父，敬拜不可分离。当宣告独一圣子，祂在万世之前坐在圣父右边，并非因苦难后时间上的提升而分享宝座，而是永恒地拥有。\par

18. \BibleQuote{谁看见圣子，就看见了圣父} \BibleRef{若 14:9}，因为圣子在一切事上都相似生祂的那一位：祂是出自生命的生命、出自光的光、出自能力的权能、出自天主的真实天主；天主性的特性在圣子内是不变的；那堪当在圣子内看见天主性的人，就达到了对圣父的享有。这不是我的话，而是独生子的話：\BibleQuote{斐理伯！这么长久的时间，我同你们在一起，你们还不认识我吗？谁看见了我，就看见了父。} \BibleRef{若 14:9} 简而言之，我们既不要将二者分开，也不要混淆：你永远不要说圣子与圣父无关，也不要接受那些说圣父有时是父、有时是子的人，因为这些都是奇怪且不虔敬的说法，不是教会的道理。圣父生了圣子，仍然保持为父，没有改变。祂生了智慧，并未失去智慧；祂生了权能，并未变得软弱；祂生了天主，并未失去自己的天主性；祂既没有因减少或改变而失去什么，生的那位也一无所缺。生祂的那位是完美的，被生的那位也是完美的：生祂的那位是天主，被生的那位也是天主；祂自己是万有的天主，却称自己的父为\Quote{我的天主}。因为祂不以说\BibleQuote{我升到我的父和你们的父那里去，到我的天主和你们的天主那里去。} \BibleRef{若 20:17} 为耻。\par

19. 但免得你以为圣父对圣子和对受造物是同一意义上的父，基督在接下来的话中作了区分。因为祂没有说\Quote{我升到我们的父那里}，以免受造物与独生子成为同伴；而是说\Emph{我的父和你们的父}：一方面是我的父，按本性；另一方面是你们的父，因义子名分。又说\Emph{我的天主和你们的天主}，一方面是我的天主，作为祂真实和独生的圣子；另一方面是你们的天主，作为祂的化工。因此，天主子是真实天主，在万世之前不可言喻地受生（我常向你们重复同样的话，为要铭刻在你们心中）。也当相信：天主有圣子；但关于方式，不要好奇，因为你寻找也找不到。不要抬高自己，免得跌倒；\BibleQuote{只应思念那些命令你的事。} \BibleRef{德 3:22} 先告诉我，生祂的那一位是什么，然后你才能认识祂所生的那一位；但若你无法理解那生者的本性，就不要再好奇地探求那被生者的方式。\par

20. 为虔诚而言，正如我们所说，你只需知道天主有一位独生子，一位按本性所生的子；祂的存在并非在生于白冷时才开始的，而是在万世之前。因为请听先知米该亚说：\Quote{你白冷啊，厄弗辣大！你在犹大诸城中虽是最小的，但将由你为我出生一位统治以色列的首领，祂的根源从太初，从永远就有了。} 因此不要想那位现在从白冷出来的，而要朝拜那位由父永恒所生的。不容许任何人谈论圣子有时间上的开始，而要承认父是永恒的开始。因为父是子的开始，无时间、不可理解、无始。正义之河的泉源，即独生子的泉源，就是父，祂以自己唯独知道的方式生了祂。你愿知道我们的主耶稣基督是永恒的王吗？请再听祂说：\BibleQuote{你们的父亲亚巴郎欢欣喜乐地期望看见我的日子，他看见了，且喜乐了。} \BibleRef{若 8:56} 然后，当犹太人对此难以接受时，祂说了对他们更为难的话：\Emph{在亚巴郎出现以前，我就是。} 祂又对父说：\Emph{现在，父啊！求祢以祢自己在我未有世界以前，我在祢面前所有的光荣，光荣我。} 祂清楚地说了\Quote{在世界未有以前，我在祢面前就有了光荣。} 又说：\BibleQuote{因为祢在我创世之前就爱了我。} \BibleRef{若 17:24} 祂清楚地宣告：\Quote{我在祢面前的光荣是永恒的。}\par

21. 我们信独一无二的主、耶稣基督、天主的独生子，在万世之前由父所生，出自天主的天主，借着祂万有都是受造的。因为\BibleQuote{无论是在天上的、掌权的、执政的、有能的，都是借着祂造成的} \BibleRef{哥 1:16}，受造之物没有一样不受祂权柄的管辖。愿每一引进不同创造者和造物者的异端闭口；愿那亵渎基督、天主子之舌缄默；愿那些说基督是太阳的人闭口，因为祂是太阳的创造者，而非我们所见的太阳。愿那些说世界是天使的作品、企图篡夺独生子尊严的人闭口。因为无论是可见的、不可见的、掌权的、执政的，还是其他任何名称，一切都是借着基督造成的。祂统治祂所造的一切，并非强夺别人的战利品，而是统治祂自己的化工，正如福音作者若望所说：\BibleQuote{一切都是借着祂造成的，凡受造的，没有一样不是借着祂造成的} \BibleRef{若 1:3} 一切都是借着祂造成的，天主圣父借着圣子工作。\par

22. 我也愿就我所说的给出一个比方，但我知道这比方是软弱的；因为在可见之物中，有什么能精确地比拟神圣能力呢？然而，既是软弱者，就向软弱者说软弱的话吧。正如一位国王，他的儿子也是国王，若他想建造一座城，可能会向他的儿子——王国的同享者——提出城的样式，儿子领受了蓝图后，便将设计付诸完成；同样，当圣父愿意创造万有时，圣子遵圣父之命创造了万有，这样的命令确保了圣父的绝对权威，而圣子也对自己亲手所造之物拥有权威，这样圣父未曾与对自己作品的统治权分离，圣子也并非治理由他人所造之物，而是由自己治理。因为，正如我所说，天使并未创造世界，而是独生子——如我所说，在万世之前受生——借着他万有得以被造，没有任何事物被排除在他的创造之外。迄今为止，愿我们靠基督的恩宠所说的话到此为止。\par

23. 现在让我们回到我们的信德宣认，如此暂时结束我们的讲论。基督创造了万有，无论你谈论的是天使、总领天使、宰制者或掌权者。并非圣父缺乏力量去创造工程本身，而是因为他愿意圣子统治自己亲手所造之物，是天主亲自将受造之物的蓝图赐予他。因为，为尊敬自己的圣父，独生子说：\BibleQuote{圣子不能由自己作什么，惟有看见父所作的，子才能作；父所作的事，子也照样作。}\BibleRef{若 5:19} 又说：\Quote{我父作工直到现在，我也作工}，其间毫无冲突。\Quote{凡是我的，都是你的；你的，也是我的}，主在福音中如此说。我们从旧约和新约中确知此事。因为他说：\BibleQuote{让我们照我们的肖像，按我们的模样造人}\BibleRef{创 1:26}，显然是在对在场的某一位说话。但最清楚的是圣咏作者的话：\Quote{他一发言，万有造成；他一出命，万物化生}，仿佛父命令发言，而圣子遵父命创造万有。约伯也曾奥秘地论及此事：\BibleQuote{惟有他独自铺张苍天，步行在海面上如履平地}\BibleRef{约 9:8}；向明悟之人指明，那曾在地上步行于海面上的，也正是那从太初创造诸天的一位。主又说：\Quote{你岂曾取过尘土，将泥塑成生灵？} 然后：\Quote{死亡的门岂为你洞开？地狱的守门者见了你岂不战栗？} 如此指明，那因慈爱降入地狱的，也曾在太初用泥土造人。\par

24. 基督是天主的独生子，是世界的创造者。因为\BibleQuote{他已在世界上，世界原是借着他造的；}\BibleQuote{他来到自己的领域}\BibleRef{若 1:10}，正如福音所教导我们的。不仅可见之物，连不可见之物，也都是基督遵圣父之命而创造的。因为依照宗徒所说：\BibleQuote{在天上和地上的一切，可见的与不可见的，或是上座者，或是宰制者，或是率领者，或是掌权者，都是在[他]内受造的；万有都是借着他，并为了他而受造的；他在万有之先，万有都赖他而存在}\BibleRef{哥 1:16}。即使你谈论世界，耶稣基督也是遵圣父之命而造世界的。因为\BibleQuote{在这末期内，天主借着他的儿子对我们说了话，并立他为万有的嗣子，借着他而创造了世界}\BibleRef{希 1:2}。愿光荣、尊威、权能归于他，从今世直到永远，世世无尽。阿们。\par

\SourceRecord{Translated by Edwin Hamilton Gifford.；Nicene and Post-Nicene Fathers, Second Series；Vol. 7.；Edited by Philip Schaff and Henry Wace.；Buffalo, NY: Christian Literature Publishing Co.,；1894.；<http://www.newadvent.org/fathers/310111.htm>.}

% structure: p0001 source=h1 target=section reason=page-title

\section{教理讲授第十二讲}

\Emph{论\OriginalTerm{降生成人}{Incarnate}与\OriginalTerm{成为人}{Made Man}。}\par

\BibleRef{依 7:10-14}\par

\Quote{\Emph{主又对阿哈次说：你向主要求一个征兆，等等}}，以及\Quote{\Emph{看！一位贞女将怀孕生子，给他起名叫厄玛努耳，等等}}。\par

纯洁的乳儿和节操的门徒，让我们以纯洁的口唇向由贞女所生的天主高唱赞歌。我们被视为配分享属神的羔羊的肉，就让我们连同头与脚一起领受，头领悟为天主性，脚领悟为人性。圣福音的听众，让我们聆听神学家若望。因为那说过\BibleQuote{在起初已有圣言，圣言与天主同在，圣言就是天主}(\BibleRef{若 1:1})的，接着又说\BibleQuote{圣言成了血肉}。因为仅仅崇拜这个人是不神圣的，说祂只是天主而没有人性也不算虔诚。因为如果基督是天主——祂的确是天主——却没有取人性于自身，我们便与救恩无分。那么，让我们以祂为天主来崇拜，但也相信祂成为了人。因为只称祂为人而不承认祂的天主性，毫无益处；拒绝承认祂的人性与天主性同在，也无救恩。让我们宣认那同为君王与医生的祂的临在。因为君王耶稣，当祂要成为我们的医生时，\Emph{以人类性的细麻布束腰}，治愈了那有病的。婴孩的完美导师(\BibleRef{罗 2:20})成为婴孩中的婴孩，为将智慧赐给愚蒙人。天上的食粮降到了地上，为养育饥饿的人。\par

2. 但是犹太人的子孙藐视那已经来的，期待那在邪恶中要求的那位，而拒绝了真正的默西亚，等候那欺骗者，自己受骗；在这里救主也被证实为真实的，祂曾说：\Quote{\Emph{我因我父的名而来，你们却不接纳我；如果另有人因自己的名而来，你们倒要接纳他}}。也宜向犹太人提出一个问题。那说厄玛努耳将由贞女诞生的先知依撒意亚，是真还是假(\BibleRef{依 7:14})？如果他们指控他说谎，并不奇怪：因为他们的习惯不仅是指控说谎，而且还要用石头砸死先知。但如果先知是真实的，请指出厄玛努耳，并且说：你们所期待的那位将来者，是由贞女诞生，还是不是呢？因为如果祂不是由贞女诞生的，你们就指控先知说谎；但如果你们在祂要来者身上期待这一点，为什么拒绝已经发生的事情呢？\par

3. 那么，就让犹太人被迷惑吧，既然他们愿意；但愿天主的教会得光荣。因为我们确确实实地领受了降生成人的天主圣言，不像异端者所说的，是由男女的意愿，而是按照福音，由童贞玛利亚和圣神所生，成为人，不是外表上，而是真实地。至于祂确实是由贞女所成的真人，请等待本讲中合适的教导时刻，你将得到证明：因为异端者的错误是多样的。有些人说祂根本就不是由贞女所生；另一些人说祂是由一个与丈夫同住的妻子所生，而不是由贞女所生。还有些人说基督不是降生成人的天主，而是成为天主的一个人。因为他们竟敢说，不是祂——那先存的圣言——成为了人，而是一个人以进步的方式被加冕。\par

4. 但你当记得昨天关于祂的天主性所说的话。要相信，祂——天主的独生子——祂自己再次由一位贞女所生。要相信福音作者若望所说的话：\BibleQuote{圣言成了血肉，住在我们中间}(\BibleRef{若 1:14})。因为圣言是永恒的，在万世之前由圣父所生；但祂最近为了我们而取了血肉。许多人反对这一点，说：\Quote{有什么如此重大的原因，需要天主降生成人？而且，天主与人类往来符合祂的本性吗？还有，一位贞女没有男子就能生育，这怎么可能？}既然争议如此之多，战斗形式多样，来吧，让我们藉着基督的恩宠和在座者的祈祷，逐一解决每个问题。\par

5. 首先让我们探究耶稣为何降来。现在不要在意我的论证，因为你可能被误导；除非你在每件事上接受先知的见证，不要相信我的话：除非你从圣经中得知关于贞女、地方、时间和方式的事，否则\BibleQuote{不要接受人的见证}(\BibleRef{若 5:34})。因为当前如此教导的人可能被怀疑；但有理智的人谁会怀疑一个在千年之前就已预言的人呢？如果你寻求基督来临的原因，请回到圣经的第一卷书。天主在六天内创造了世界：但世界是为了人。然而太阳虽闪耀着明亮的光芒，却是为给人带来光明；不仅如此，所有活物都是造来服侍我们的：草本和树木是为我们的享受而造的。所有受造之物都是好的，但没有一样是天主的肖像，唯独人除外。太阳仅凭一个命令就形成了，但人是凭天主亲手：\BibleQuote{让我们按我们的肖像，按我们的模样造人}(\BibleRef{创 1:26})。一位世上君王的木制肖像是受尊敬的，更何况天主的理性肖像呢？\par

但是，当这最伟大的受造之物在乐园中嬉戏时，魔鬼的嫉妒将他逐出。仇敌为他所嫉妒者的堕落而欢喜：你愿意仇敌继续欢喜吗？因不敢接近那有力量的男人，他就接近那较为软弱的女人，她仍是贞女：因为是在被逐出乐园之后，\Emph{亚当认识了他的妻子}。\par

6. 加音和亚伯尔出现在人类的第二世代；加音是第一个杀人者。后来，因着人们极大的邪恶，洪水泛滥在地上；因着索多玛人的过犯，火从天上降下。过了一段时间，天主拣选了以色列；但以色列也转离了，被拣选的民族受了伤。因为当梅瑟在山上立于天主面前时，百姓却在敬拜一只牛犊代替天主。在立法者梅瑟有生之年，他曾说：\BibleQuote{不可犯奸淫}，却有人胆敢进入娼妓之地而犯过。\BibleRef{户 25:6} 梅瑟之后，先知被派遣来医治以色列；但在他们的医治职分中，他们哀叹自己无法战胜疾病，以致他们中有一位说：\BibleQuote{祸哉，我！因为虔诚人从地上灭亡了，在世人中没有一个行正义的。}\BibleRef{米 7:2} 又说：\BibleQuote{他们都偏离了正道，一同变为无用的；没有行善的，连一个也没有。} 又说：\BibleQuote{咒骂、偷窃、奸淫和杀人弥漫在这地上。}\BibleRef{欧 4:2} \BibleQuote{他们将儿女祭献给魔鬼。} \BibleQuote{他们行占卜、邪术和法术。}\BibleRef{编下 33:6} \BibleQuote{再者，他们用绳索系住衣裳，在祭台旁挂起幔子。}\par

7. 人性的创伤极其深重；\BibleQuote{从脚掌到头顶，没有一处完好；}无人能敷上\BibleQuote{舒缓的油、油膏或绷带。}\BibleRef{依 1:6} 于是先知们哀叹并疲倦，说：\BibleQuote{谁将从熙雍赐下救恩？} 又说：\BibleQuote{愿祢的手临到祢右边的人，临到祢为自己所坚固的人子；这样我们就不离开祢。} 另一位先知恳求说：\BibleQuote{上主，求祢将天垂下，亲自降下。} 人性的创伤非我们所能医治。\BibleQuote{他们杀了祢的先知，拆毁了祢的祭坛。}\BibleRef{列上 19:10} 这恶事靠我们无法挽回，需要祢来挽回它。\par

8. 主垂听了先知的祈祷。圣父并未不顾我们种族的灭亡；祂派遣了祂的圣子，主从天上降下，作为医治者；一位先知说：\BibleQuote{你们所寻求的主，必然来临，且忽然来临。}\BibleRef{拉 3:1} 到哪里？\BibleQuote{主要来到祂自己的圣殿}，就是你们用石头打祂的地方。另一位先知听见这话，就对他说：你谈论天主的救恩，岂可低声细语？你宣讲天主来临施行救恩的喜讯，岂可隐秘？\BibleQuote{给熙雍报喜讯的，你要登上高山！要大声向犹大各城宣讲！} 我该讲什么？\BibleQuote{看哪！我们的天主！看哪！主带着权能来临！}\BibleRef{依 40:9} 主自己又说：\BibleQuote{看哪！我要来临，且要住在你们中间，上主说。许多民族要归向上主。}\BibleRef{匝 2:10} 以色列人拒绝了藉着我而来的救恩：\BibleQuote{我要聚集万民和万邦。}\BibleRef{依 66:18} 因为\BibleQuote{祂来到自己的地方，自己的人却没有接受祂。}\BibleRef{若 1:11} 祢来了，祢赐给万民什么？\BibleQuote{我来聚集万民，我要在他们身上留下一个记号。} 因为我从十字架上的争斗中，给每个士兵一个王室的印记，让他在额上佩戴。另一位先知说：\BibleQuote{祂使天垂下，亲自降临；祂的脚下是幽暗。} 因为祂从天上降下，不为人所知。\par

9. 后来，撒罗满听见他父亲达味说这些事，就建造了一座奇妙的殿，并预见了将要进入其中的那一位，惊讶地说：\BibleQuote{天主真的愿意与人同住在地上吗？} 是的，达味在题为\WorkTitle{为撒罗满}的圣咏中预先说：\BibleQuote{祂必像雨落在羊毛上}；\BibleQuote{雨}是因祂的天性，落在羊毛上则因祂的人性。因为雨落在羊毛上，是悄然无声的；因此，贤士们不知道诞生的奥秘，就说：\BibleQuote{那生下来作犹太人之王的在哪里？}\BibleRef{玛 2:2} 黑落德便不安，查问那诞生的孩子，说：\BibleQuote{基督将在何处诞生？}\par

10. 但降下来的这位是谁？祂在后续中说：\BibleQuote{祂与太阳并存，在月亮之前，直到万代。} 另一位先知说：\BibleQuote{熙雍女子，你应尽情欢乐！耶路撒冷女子，你应欢呼！看哪！你的君王来到你这里，正义的，带着救恩。}\BibleRef{匝 9:9} 君王很多，先知啊，你说的是哪一位？请给我们一个其他君王没有的标记。若你说是一位身穿紫袍的君王，这衣服的尊荣已被抢先用了。若你说是一位由持矛者护卫、坐在金车上的君王，这也被其他人抢先用了。请给我们一个你所宣布的君王独有的标记。先知回答说：\BibleQuote{看哪！你的君王来到你这里，正义的，带着救恩：祂是谦逊的，骑着一匹驴驹子}，不是战车。你们有了这君王独一无二的标记。耶稣是唯一骑在未经驯服的驴驹上、在欢呼中如君王般进入耶路撒冷的君王。当这位君王来了，祂做什么？\BibleQuote{因着盟约的血，求祢从无水之坑中释放你的囚徒。}\BibleRef{匝 9:11}\par

11. 但祂或许也能骑驴驹：请给我们一个标记，指明这位进来的君王将站立何处。请将这标记赐给城里不远之处，好使我们不至于不知道；请给我们一个眼前清楚的标记，好让我们即使在城里也能看见那地方。先知又回答说：\BibleQuote{在那一天，祂的脚必站在橄榄山上，这山在耶路撒冷东面。} 有谁站在城里而看不见那地方呢？\par

12. 我们有两个标记，还想知道第三个。请告诉我们，主来了之后做什么。另一位先知说：\Quote{看哪！我们的天主}，随后又说：\Quote{祂必来拯救我们。那时，盲人的眼睛要睁开，聋子的耳朵要听见；那时，瘸子要像鹿一样跳跃，口吃者的舌头要说话清晰。}\BibleRef{依 35:4} 但让我们再听另一个见证。先知啊，你说主来了，并行了前所未闻的奇迹：你还说了什么清楚标记呢？\Quote{主亲自与祂子民的长老和首领进入审判。}一个显著的标记！师傅被祂的仆人、长老们审判，并服从之。\par

13. 这些事犹太人读了，却不听从；因为他们塞住了耳朵，免得听见。但让我们相信耶稣基督，祂已降生成人，\Emph{成为人}，因为我们否则无法接受祂。既然我们不能按祂本来的样子注视或享见祂，祂就变成了我们所能及的，好让我们得以享见祂。因为我们若不能正眼观看在第四天受造的太阳，又怎能看见创造它的天主呢？主曾在西乃山在火中降临，百姓不能忍受，却对梅瑟说：\Quote{你同我们说话，我们肯听；不要天主同我们说话，免得我们死亡}\BibleRef{出 20:19}；又说：\Quote{凡有血肉的，谁曾听到永生天主从火中说话的声音，而能生存呢？}\BibleRef{申 5:26} 如果听到天主说话的声音就是死亡的原因，那么看见天主本身岂不带来死亡？这有什么奇怪呢？连梅瑟自己也说：\Quote{我甚是恐惧战兢}\BibleRef{希 12:21}。\par

14. 那么，你希望怎样呢？那为我们的救恩而来的，因人们不能忍受祂，反成了毁灭的工具吗？还是祂应当使祂的恩宠适应我们的尺度？达尼尔不能忍受一位天使的显现，你就能看见天使的主吗？加俾额尔显现了，达尼尔就倒下了：那显现的有着怎样的面貌和形态呢？他的面容像闪电\BibleRef{达 10:6}，不像太阳；\Quote{他的眼睛像火把}，不像火炉；\Quote{他说话的声音像群众的声音}，不像十二营天使的声音；然而，先知还是倒下了。天使来到他跟前说：\Quote{达尼尔，不要害怕，站直了！放心罢，你的话已蒙应允}\BibleRef{达 10:12}。达尼尔说：\Quote{我战战兢兢地站起}；即便如此，他还是没有回答，直到有人手的样式摸了他。当那显现者变成了人的模样，达尼尔才说话：他说什么呢？\Quote{我主，因为见到您的显现，我的内脏在我里面翻转，我毫无气力，也没有气息存留。}如果一位天使的显现夺去了先知的声音和力量，那么天主的显现岂能让他存活呢？直到\Quote{有一位像人的手摸了我}，圣经说，达尼尔才鼓起勇气。因此，在显出我们的软弱之后，主就取了人所需要的：因为人需要从一个面貌相同的人那里听到，救主就取了同样情感的本性，好使人更容易受教。\par

15. 再学习另一个原因。基督来了，为要受洗，并圣化圣洗圣事；祂来了，为要行奇迹，在水面上行走。既然在祂肉身显现之前，\Quote{海一见就逃，约但河也倒流}，主就取了祂的身体，好使海能忍受看见，约但河能不惧地接纳祂。这是一个原因；但还有第二个原因。死亡藉着仍是童贞的厄娃而来；生命必须藉着一位童贞女，更确切地说，从一位童贞女显现：如此，正如蛇欺骗了那一位，同样，加俾额尔向另一位报告喜讯。人离弃了天主，却塑造了人的雕像。既然人的雕像被虚假地当作天主崇拜，天主就真实地成了人，好使这虚假被除去。魔鬼曾用肉体作为攻击我们的工具；保禄知道这事，就说：\Quote{但我发现在我的肢体内，另有一条法律，与我理智的法律交战，并把我掳去}\BibleRef{罗 7:23}，等等。因此，藉着魔鬼用来击败我们的同样武器，我们获得了拯救。主从我们取了我们的相似，为要拯救人的本性；祂取了我们的相似，为要赐给那缺乏的更大的恩宠；使有罪的人性得以分享天主。\Quote{因为罪恶在哪里越多，恩宠在那里也越格外丰富。}主必须为我们受苦；但若魔鬼认识祂，就不敢接近祂了。\Quote{因为如果你们认识了，就不会将光荣的主钉在十字架上了}\BibleRef{格前 2:8}。因此，祂的身体成了死亡的诱饵，好使那希望吞噬它的毒龙，也吐出那些已被吞吃的人。因为\Quote{死亡得胜了，吞噬了}；但再次，\Quote{天主擦去了各人脸上的眼泪}。\par

16. 难道基督降生成人是没有原因的吗？我们的教导是巧妙的言辞和人的诡辩吗？难道圣经不是我们的救恩吗？难道先知的预言不是吗？那么，我祈求你们，保守这项恩赐不受搅扰，不要让人挪去你们：相信天主降生成人。但尽管已经证明祂降生成人是可能的，然而若犹太人仍然不信，我们要向他们提出这一点：我们说天主降生成人，有什么奇怪呢？因为你们自己说亚巴郎曾接待主为客。我们说这事有什么奇怪呢？当雅各伯说：\BibleQuote{我见了天主面对面，我的生命仍得保存}？那与亚巴郎一同进食的主，也与我们一起进食。那么，我们宣布的有什么奇怪呢？不仅如此，我们还提出两个证人，就是在西乃山上站在主面前的那两位：梅瑟曾在\BibleQuote{磐石穴中}\BibleRef{出 33:22}，厄里亚也曾一度在磐石穴中\BibleRef{列上 19:8}；他们在大博尔山上与主一同显圣容时，\Emph{谈到}门徒们\Emph{祂在耶路撒冷将要完成的结局}。但是，正如我先前说的，祂降生成人已被证明是可能的；其余的证明，可以留给勤勉的人去收集。\par

17. 然而，我的陈述应许也要宣告救主的时间和地点；我不能被判有谎言而离开，反而要使教会的初学之士确信无疑。因此，让我们探究我主来临的时间：因为祂的来临是近期的，且受争议；因为\BibleQuote{耶稣基督昨天、今天、直到永远，常是一样的}。那么，先知梅瑟说：\BibleQuote{上主你的天主要从你们兄弟中，给你兴起一位先知，像我一样}；但且把那\Quote{\Emph{像我一样}}暂时保留，在适当的地方查考。但这所期待的先知何时来临呢？他说，要回头看我写的：仔细查考雅各伯对犹大的预言：\BibleQuote{犹大！你的兄弟要赞美你}，之后，为免全部引用，\BibleQuote{权杖不离犹大，立法者不离他两膝之间，直到那应得者来到；祂不是犹太人的，而是外邦人的期待}\BibleRef{创 49:8}。因此，祂以犹太统治的终止作为基督来临的标记。如果他们现在不在罗马人统治下，基督还没有来；如果他们仍有犹大和达味后裔的王子，那期待的还没有来。因为我羞于讲述他们近来关于他们中如今称为族长的那些人所作的事，以及他们的出身和他们的母亲是谁；但我让知道的人去讲。但那作为\BibleQuote{外邦人的期待}而来的，还有什么进一步的标记呢？祂接着说：\BibleQuote{将他的驴驹拴在葡萄树上}\BibleRef{创 49:11}。你看到那驴驹，就是匝加利亚清楚宣告的。\par

18. 但你再次询问关于时间的另一个见证。\BibleQuote{上主对我说：你是我的儿子，我今日生了你}；稍后几句，\BibleQuote{你必用铁杖统治它们}。我先前说过，罗马帝国显然被称为铁杖；但关于这一点所缺少的，让我们进一步从达尼尔书中回忆。因为他在叙述并给拿步高解释那像的梦时，也告诉了他关于这像的全部异象：有一块未经人手所凿的石头（即不是由人的计谋竖立的）要征服整个世界；他非常清楚地说：\BibleQuote{在那列王在位的日子，天上的天主必要兴起一个永不灭亡的国家，祂的国度也不留给别的民族}\BibleRef{达 2:44}。\par

19. 但我们还要求更清楚地证明祂来临的时间。因为人难以说服，除非得到确切的年份来清楚计算，否则不信所说的。那么，是什么时节，时间的方式如何？就是当犹大后裔的君王失败，外邦人黑落德继位为王的时候。因此，与达尼尔交谈的天使说（你现在要留意这些话）：\BibleQuote{你应当知道，也应当明白：从发出命令回答和重建耶路撒冷，直到受傅者君王，是七个星期和六十二个星期}\BibleRef{达 9:25}。六十九个星期之年包含四百八十三年。因此，他说，在耶路撒冷重建之后，过了四百八十三年，统治者失败，那时另一个民族的一位君王来到，在祂的时候基督诞生。现在，玛待达理阿在自己统治的第六年，即希腊人的第六十六届奥林匹克运动会的第一年，建造了这城。奥林匹克运动会在希腊人中是每四年举行的比赛名称，因为每四年太阳运行的日数是由每年三个（额外的）小时累积而成的。黑落德王是在第一百八十六届奥林匹克运动会，即其第四年。从第六十六届到第一百八十六届奥林匹克运动会，中间有一百二十届奥林匹克运动会，并略多一些。因此，一百二十届奥林匹克运动会构成四百八十年：因为其余剩余的三年的间隔可能被第一年和第四年之间的间隔所占据。这样，你就有了根据圣经的证明，圣经说：\BibleQuote{从发出命令恢复和重建耶路撒冷，直到受傅者君王，是七个星期和六十二个星期}。所以，关于时间，目前你有这个证明，尽管对于达尼尔书中上述的星期之年还有其它不同的解释。\par

20. 但如今请听预许之地，如\Person{米该亚}所说：\BibleQuote{而你，\Place{白冷}，\Place{厄弗辣大}家，你在犹大千邑中何其渺小？因为将由你中为我出一位统御者，作以色列的领袖：他的根源来自亘古，来自永远之日。} 但就地方而言，你身为\Place{耶路撒冷}居民，也当预知圣咏第一百三十一篇所载：\BibleQuote{看！我们听说他在\Place{厄弗辣大}，我们在森林田野中找到他。} 因为几年前这地方还是林地。你又听见\Person{哈巴谷}对主说：\BibleQuote{当岁月临近，你将显明；时候来到，你将被显示。} 先知的标志为何，主来临的标记？他随即说：\BibleQuote{你将在两个生命之间被认识。} 他清楚地这样对主说：\Quote{你降生成人，活着也死去，并从死者中复活后，又活着。} 再者，他从\Place{耶路撒冷}周围何方向而来？东、西、南、北？请明确告知。他极其明白地回答：\BibleQuote{天主将从\Place{特曼}而来}（\Place{特曼}意为\Quote{南方}）\BibleQuote{圣者从\Place{帕兰}山而来，荫蔽多林}：正如圣咏作者以类似话语所说：\BibleQuote{我们在森林田野中找到他。}\par

21. 我们进一步问：祂来自谁，如何而来？\Person{依撒意亚}告诉我们：\BibleQuote{看！童贞女将在腹中怀孕，并生一子，他们要称祂的名字为\Person{厄玛奴耳}。} 犹太人反对这事，因他们自古惯于邪恶地抵触真理：他们说所写的不是\Quote{童贞女}，而是\Quote{少女}。但我同意他们所说，即便如此，我仍找到真理。因我们必须问他们：若童贞女被强迫，她在何时呼喊求救，是在凌辱之前还是之后？因此，若圣经别处说：\BibleQuote{那订了婚的少女呼喊，却无人救她} \BibleRef{申 22:27}，难道这不是指童贞女吗？\par

但为使你更清楚明白，在圣经中连童贞女也称为\Quote{少女}，请听\BibleBook{}{列王纪}中论叔能女子\Person{阿彼沙格}的话：\BibleQuote{那少女极其美丽}；因她作为童贞女被选，带到\Person{达味}面前，是公认的。\par

22. 但犹太人又说，这是对\Person{阿哈次}说的，指\Person{希则克雅}。那么，让我们读圣经：\BibleQuote{你向你的主天主求一个征兆，或求于深渊，或求于高天。} \BibleRef{依 7:11} 这征兆当然必是惊人的。因从磐石出水是征兆，海水分开，太阳回转等。但在我将提及的事中，有对犹太人更明显的反驳。（我知道我说得太长，听众已厌倦；但请容忍我详尽陈词，因这些问题是为基督的缘故而提出，且关系重大。）当\Person{依撒意亚}在\Person{阿哈次}王在位时说这话，\Person{阿哈次}只统治十六年，预言是在这期间对他说的，犹太人的异议被如下事实驳倒：继任国王\Person{希则克雅}是\Person{阿哈次}之子，他登基时二十五岁；因预言限于十六年内，他必在预言之前已由\Person{阿哈次}生育整整九年。那么，对于一个在父王\Person{阿哈次}登基之前就已生育的人，何必发出这预言？因他说的不是\Emph{已怀孕}，而是\Quote{\Emph{童贞女将怀孕}}，是预知而说。\par

23. 我们确知主将由童贞女诞生，但须指明童贞女出自哪一家族。\BibleQuote{主向\Person{达味}起了真誓，决不废止。我必使你身所生的，坐在你的宝座上。} \BibleQuote{并且，我要永远坚立你的后裔，你的宝座如天一样长久。} \BibleQuote{此后，我曾凭我的圣洁起誓，绝不向\Person{达味}说谎。他的后裔必永存，他的宝座在我面前如日头，又如月亮永远坚立。} 你看，这话是指基督，不是指\Person{撒罗满}。因\Person{撒罗满}的宝座不像日头一样永存。但若有人否认这点，因基督并未坐在\Person{达味}的木制宝座上，我们便提出那句话：\BibleQuote{经师和法利塞人坐在\Person{梅瑟}的座位上} \BibleRef{玛 23:2}；因为这并非指他的木座，而是指他教导的权威。同样，我愿你们寻求\Person{达味}的宝座，不是木座，而是王国本身。也请那些高声呼喊的儿童做我的证人：\BibleQuote{贺三纳归于\Person{达味}之子！} \BibleQuote{以色列的王当受赞美！} \BibleRef{若 12:13} 瞎眼者也说：\BibleQuote{\Person{达味}之子，可怜我们吧！} \BibleRef{玛 20:30} \Person{加俾额尔}也清楚向\Person{玛利亚}作证说：\BibleQuote{主天主将把祂祖先\Person{达味}的宝座赐给祂。} \BibleRef{路 1:32} \Person{保禄}也说：\BibleQuote{你要记得耶稣基督从死者中复活，出于\Person{达味}的后裔，照我的福音。} \BibleRef{弟后 2:8} 在致罗马人书开头，他说：\BibleQuote{祂按肉身是出于\Person{达味}的后裔。} \BibleRef{罗 1:3} 因此，你要接受那出自\Person{达味}者，相信那预言所说：\BibleQuote{到那日，\Person{叶瑟}的根必兴起，那要起来统治外邦人的；外邦人必寄望于祂。}\par

24. 但犹太人因这些事大为困扰。这也为\Person{依撒意亚}所预知，他说：\BibleQuote{他们必愿自己被火烧尽：因为有一婴孩为我们诞生了}（不是为他们），\BibleQuote{有一子赐给了我们。} \BibleRef{依 9:5} 你要注意，祂起初是天主子，然后赐给了我们。稍后他说：\BibleQuote{祂的平安没有穷尽。} 罗马人有边界；天主子的国没有边界。波斯人和玛待人有边界，但圣子没有边界。接着：\BibleQuote{祂必坐在\Person{达味}的宝座上，治理他的王国。} 因此，圣童贞女出自\Person{达味}。\par

25. 因为那至洁者、纯洁的导师，从纯洁的新娘之室出来，原是合宜的。因为若妥善履行耶稣司铎职务的人戒绝妻子，耶稣自己又怎能由男人和女人所生呢？因为他在圣咏中说：\Quote{祢是从母腹中把我取出来的}。请留意这话：\Quote{祢是从母腹中把我取出来的}，意指祂是无男人而生的，是从贞女的胎和肉中被取出来的。因为按婚姻常道所生的人，其方式是不同的。\par

26. 祂既不耻于从那些肢体中取得肉身，祂本就是那些肢体的创造者。但谁告诉我们这事？主对耶肋米亚说：\Quote{我尚未在母腹中形成你，我已认识了你；你尚未出母胎，我已圣化了你}。\BibleRef{耶 1:5} 那么，如果祂在塑造人时不嫌接触，难道祂在为自己塑造那神圣的肉体——祂天主性的面纱——时会嫌恶吗？如今仍是天主在母腹中创造儿女，正如约伯所记：\Quote{祢岂不是将我如奶倒出，使我凝结如酪？祢以皮和肉给我穿上，用骨和筋把我编织起来}。\BibleRef{约 10:10} 人的身体并无任何污秽，除非人用奸淫和邪淫玷污它。那形成亚当的，也形成了厄娃；男女都是天主的双手所造的。从起初形成的身体各肢体没有一件是污秽的。愿所有异端者的口被封闭，他们诽谤自己的身体，其实是诽谤那形成它们的天主。但让我们记住保禄的话：\Quote{你们不知道你们的身体是住在你们内的圣神的宫殿吗？}\BibleRef{格前 6:19} 先知也曾以耶稣的口吻预言说：\Quote{我的肉体是从他们而来}；另一处又写道：\Quote{因此，祂要将他们交出，直到那生产的时期}。\BibleRef{米 5:3} 那记号是什么？祂在后来告诉我们：\Quote{她要生产，他们兄弟中的残余必会归来}。那贞女——神圣新娘——的婚姻信物是什么？\Quote{我要以信实聘娶你归于我}。\BibleRef{欧 2:20} 依撒伯尔与玛利亚交谈时也这样说：\Quote{那相信的是有福的，因为上主对她所说的话，必要完成}。\BibleRef{路 1:45}\par

27. 但是希腊人和犹太人都困扰我们，说基督不可能由贞女所生。至于希腊人，我们将用他们自己的神话塞住他们的口。因为你们说扔出的石头变成了人，又怎能说贞女生子是可能的呢？你们虚构说从头脑中生出了一个女儿，又怎能说从贞女的子宫中生出一个儿子是不可能的呢？你们妄言狄俄倪索斯是从你们的宙斯的大腿所生，怎能藐视我们的真理？我知道我说的是与当前听众不相称的事，但为了让你们在适当的时候驳斥希腊人，我们提出这些，用他们自己的神话来回答他们。\par

28. 至于那些受割礼的人，你要用这个问题迎接他们：是年老、不育、过了生育期的妇人怀孕更难，还是年轻贞女怀孕更难？撒辣是不育的，虽然月经已停止，但违反自然地生了孩子。那么，如果不育妇人怀孕是反自然的，贞女怀孕也是反自然的，所以要么两者都拒绝，要么两者都接受。因为同一的天主行了前者，也安排了后者。你不敢说天主在前一事上是可能的，在后一事上是不可能的。再者，一个人的手在一小时内变成另外的样子又恢复原样，这怎么是自然的？梅瑟的手如何变得像雪一样白，又立刻复原？但你会说，是天主的意志造成了变化。在那件事上，天主的意志有能力，在这件事上就没有能力吗？那是一个关于埃及人的征兆，而这是一个给予全世界的征兆。犹太人哪，两者哪一个更困难？是贞女生子，还是杖变成活物？你承认，在梅瑟的事上，一根完全笔直的杖变成了蛇，使丢下它的人害怕；那先前紧握杖的人如躲蛇般逃开；因为它实在是条蛇。他逃跑并非因为害怕他所握的东西，而是因为畏惧那改变它的天主。杖有了牙齿和眼睛如同蛇：那么，看到眼睛从杖上长出，而如果天主愿意，孩子就不能从贞女的子宫生出吗？至于亚郎的杖，在一夜之间产生了其他树要几年才能结出的果实，这事我就不提了。因为谁不知道，一根剥了皮的杖永远不会发芽，即使种在河中也一样？但既然天主不依赖树木的本性，而是它们本性的创造者，那根不结果、干枯、无皮的杖就发芽、开花、并结出杏仁。那么，祂为了预像的大司祭，超自然地使杖结出果实，难道不会为了真正的大司祭，赐给贞女生子吗？\par

29. 这些叙述提供了出色的提示；但犹太人仍然反驳，不愿服从关于杖的叙述，除非他们被类似奇异超自然的出生所说服。因此，这样问他们：起初厄娃是谁生的？哪一个母亲怀了这无母之人？但圣经说她是出自亚当的肋旁。那么，厄娃是从男人的肋旁无母而生，而孩子就不能无父从贞女的子宫出生吗？这是一笔妇女亏欠男人的感恩债：因为厄娃出自亚当，并非由母亲怀胎，而仿佛是独自由男人所生。因此，玛利亚偿还了感恩的债，当她不是由男人，而是独自以无玷的方式，因圣神和天主的能力而受孕时。\par

30. 然而，让我们再思一件比这更奇妙的事。因为从身体孕育身体，即使奇妙，尚且可能；但尘土变成人，这更为奇妙。泥土被捏合后竟能形成眼睛的薄膜和光泽，这更为奇妙。从外观一致的尘土中产生出骨头的坚硬、肺的柔软以及各种不同的肢体，这真是奇妙。泥土竟能获得生命，自行游走世界，建造房屋，这真是奇妙。泥土竟能教导、说话、从事木工，甚至作王，这真是奇妙。那么，最愚昧的犹太人哪，亚当是如何造成的？难道天主不是从地上取土，塑造了这奇妙的形体吗？那么，泥土变成眼睛，难道处女就不能生子吗？那在人看来更不可能的事发生了，可能的事反倒永远不发生吗？\par

31. 兄弟们，让我们记念这些事，用这些武器来防卫自己。我们不要容忍那些教导基督以幻影降临的异端者。我们也要憎恶那些说救主是由夫妻所生的人，他们竟敢说祂是若瑟和玛利亚的儿子，因为经上记载：\\BibleQuote{\\Emph{若瑟便娶了他的妻子}}（\\BibleRef{玛 1:24}）。让我们记念雅各伯，他在接纳辣黑耳之前，对拉班说：\\BibleQuote{\\Emph{请给我我的妻子}}（\\BibleRef{创 29:21}）。因为正如她结婚前，仅仅因为有了婚约，就被称为雅各伯的妻子；同样，玛利亚因为已许配，也被称为若瑟的妻子。也请注意福音的准确性，它说：\\BibleQuote{\\Emph{到了第六个月，天使加俾额尔奉天主差遣，往加里肋亚一座名叫纳匝肋的城去，到一位童贞女那里，她已许配给一个名叫若瑟的男子}}（\\BibleRef{路 1:26}）等等。再者，当户口登记时，若瑟上去登记，圣经说：\\Quote{\\Emph{若瑟也从加里肋亚上去，与他已许配、怀有身孕的玛利亚一同登记}}。虽然她怀了孕，却没有说\\Quote{他的妻子}，而是说\\Quote{\\Emph{已许配给他的}}。因为保禄说：\\BibleQuote{\\Emph{天主派遣了祂的儿子}}，不是由男人和女人所生，而是\\BibleQuote{\\Emph{由女人所生}}（\\BibleRef{迦 4:4}），即由一位童贞女所生。因为童贞女也被称为女人，我们先前已指明。因为那使灵魂成为贞洁者，乃是由一位贞女所生。\par

32. 但你对此事件感到惊奇；甚至那生祂的童贞女自己也感到惊奇。因为她对加俾额尔说：\\BibleQuote{\\Emph{这事怎能成就？因为我不认识男人}}。但他回答说：\\BibleQuote{\\Emph{圣神要降临于你，至高者的能力要庇荫你，因此，那要诞生的圣者，将称为天主的儿子}}（\\BibleRef{路 1:34}）。祂的诞生是无玷无污的：因为圣神所嘘之处，一切污秽都被除去；独生子的肉身诞生来自童贞女，无玷污。如果异端者反驳真理，圣神将定他们的罪；至高者庇荫的能力将要发怒；加俾额尔在审判之日将当面斥责他们；那接纳了主的马槽之处将使他们羞愧。当时领受喜讯的牧人们将作证，以及那歌唱赞美、颂扬说\\BibleQuote{\\Emph{光荣颂赞于天主，在至高之处，在地上平安，善悦于人}}（\\BibleRef{路 2:14}）的天使大军；祂在第四十天被带上圣殿；为祂奉献的一对斑鸠；当时将祂抱在怀中的西默盎，以及在场的女先知亚纳。\par

33. 既然天主作证，圣神也一同作证，而基督说：\\Quote{\\Emph{为什么你们想要杀我，一个把真理告诉你们的人呢？}} 让异端者沉默吧，他们反对祂的人性，因为他们反对那说\\BibleQuote{\\Emph{你们摸摸我，看看！因为鬼神没有肉和骨头，你们看，我是有的}}（\\BibleRef{路 24:39}）的祂。愿由童贞女所生的主受朝拜，愿童贞女承认她们这身分的冠冕；也让隐修士群体承认贞洁的光荣，因为我们人类并非被剥夺了贞洁的尊荣。救主在童贞女腹中度过了九个月，但主作为人生活了三十三年。所以，如果童贞女因九个月而夸耀，我们更应因多年而夸耀。\par

34. 但是，让我们所有人都靠着天主的恩宠，奔跑贞洁的赛程，\\Quote{\\Emph{少年人、童贞女、老年人、孩童}}，不追求纵欲，而是赞美基督的名。我们不要忽略贞洁的光荣：因为它的冠冕是天使的，它的卓越超乎人类。我们要珍惜我们的身体，它们将来要像太阳一样发光；我们不要为短暂的快乐玷污如此伟大、如此高贵的身体：因为罪是短暂即过的，但耻辱却延至多年，直至永远。在地上行走的天使，就是那些实践贞洁的人；童贞女与童贞女玛利亚同得一份。要摒弃一切虚荣的装饰、一切有害的目光、一切放荡的步态、一切飘动的长袍、以及引向快乐的香膏。但在一切事上，要以祈祷作为馨香的香料，以善工和身体的圣化作为香料，好使由童贞女所生的主也能对我们——既包括守贞的男子，也包括戴冠的女子——说：\\BibleQuote{\\Emph{我要住在他们中间，在他们中间行走；我要作他们的天主，他们要作我的子民}}（\\BibleRef{格后 6:16}）。愿荣耀归于祂，直到永远。阿们。\par

\SourceRecord{Translated by Edwin Hamilton Gifford.；Nicene and Post-Nicene Fathers, Second Series；Vol. 7.；Edited by Philip Schaff and Henry Wace.；Buffalo, NY: Christian Literature Publishing Co.,；1894.；<http://www.newadvent.org/fathers/310112.htm>.}

% structure: p0001 source=h1 target=section reason=page-title

\section{要理讲授第十三篇}

论\Quote{被钉十字架}与\Quote{被埋葬}一语\par

\BibleRef{依 53:1,7}\par

\Quote{谁相信了我们的报道？上主的手臂又向谁启示了呢？……祂被牵去，像一只待宰的羔羊}等。\par

1. 基督的每件事迹都是天主教会夸耀的缘由，但她最夸耀的莫过于十字架；保禄深知此事，故说：\Quote{但至于我，我只以我们的主耶稣基督的十字架来夸耀}（\BibleRef{迦 6:14}）。诚然，一个生来就瞎眼的人在史罗亚池得见光明，是奇妙的事；但这与全世界瞎眼的人相比，又算什么呢？拉匝禄第四天复活，是伟大而超越本性的；但这恩惠只惠及他一人，与全世界死于罪恶中的人相比，又算什么呢？五饼供应五千人食物，是奇妙的；但这与全世界因无知而饥饿的人相比，又算什么呢？那被撒殚捆绑十八年的妇女得到释放，是奇妙的；但这与所有被罪恶锁链紧捆的我们相比，又算什么呢？然而，十字架的荣耀引领了那些因无知而瞎眼的人进入光明，释放了所有被罪恶束缚的人，并赎回了全人类世界。\par

2. 不要惊讶全世界被赎回；因为替它而死的，并非一个普通人，而是天主的独生子。此外，一个人的罪——即亚当的罪——尚且有能力给世界带来死亡；\Quote{如果因一人的过犯，死亡就因那一人作了王}，那么生命岂不更要因一人而作王吗（\BibleRef{罗 5:17}）？如果因那知善恶树的果子，他们当时被逐出乐园，那么因耶稣的十字架，信徒岂不更容易进入乐园？如果第一个由泥土形成的人带来了普遍的死亡，那么那由泥土形成他的人（即耶稣）本身是生命，岂不带来永生？如果丕乃哈斯因嫉恶如仇杀了那作恶的人，平息了天主的义怒，那么耶稣岂不更要舍己为人，作赎价（\BibleRef{弟前 2:6}），消除那针对人类的义怒？\par

3. 那么，让我们不要羞于我们救主的十字架，反倒要以此为荣。因为\Quote{十字架的道理，为犹太人是绊脚石，为外邦人是愚妄}，但为我们却是救恩；\Quote{为那些丧亡的人是愚妄，但为我们这些得救的人，却是天主的德能}。因为，正如我先前所说，为我们死的并非一个普通人，而是天主圣言成了血肉。再者，梅瑟的羔羊既能使毁灭者远离（\BibleRef{出 12:23}），那么\Quote{除免世罪的天主羔羊}（\BibleRef{若 1:29}）岂不更能解救我们脱离罪恶？一只愚笨羔羊的血能带来救恩，那么独生子的血岂不更能拯救？若有人不信被钉者的德能，让他去问魔鬼；若有人不信言语，让他相信所见之事。普世许多人被钉十字架，但魔鬼没有一个因他们而畏惧；然而，一见到那为我们被钉的基督的十字架标记，魔鬼便战栗。那些人是为了自己的罪而死，但基督是为了别人的罪而死，因为祂\Quote{没有犯过罪，口中也找不到诡诈}。说这话的不是伯多禄，因为我们可能会怀疑他偏袒老师；而是依撒意亚说的，他虽未在肉身中与主同在，却在圣神中预见了祂的肉身来临。但为什么只引先知为证呢？让比拉多自己作证吧，他判决了主，说：\Quote{我在这人身上查不出什么罪状}（\BibleRef{路 23:14}）；他交出主并洗了手，说：\Quote{对这义人的血，我是无罪的}（\BibleRef{玛 27:24}）。还有另一个证明耶稣无罪的人——那个强盗，第一个被接入乐园的人；他责备同伴说：\Quote{我们所受的，正配我们所行的；但是，这个人从未做过什么不正当的事}，因为我们二人，你和我，都在祂的审判中在场。\par

4. 所以耶稣确实为众人受苦；因为十字架不是幻象，否则我们的救赎也是幻象。祂的死不是一场虚构的表演，否则我们的救恩也是虚构的。如果祂的死只是表演，那么那些说\Quote{我们记得那个骗子活着的时候说过：三天以后我要复活}（\BibleRef{玛 27:63}）的人就对了。因此，祂的受难是真实的：祂确实被钉在十字架上，我们并不因此而羞耻；祂被钉十字架，我们不否认，反而我乐于谈及此事。因为即使我现在否认，有各各他在这里反驳我，我们正聚集在它附近；十字架的木头反驳我，它后来从这儿被分送到普世各地。我承认十字架，因为我认识复活；因为如果祂被钉后仍保持原样，我或许不会承认，因为我可能将祂和我的主都隐藏起来；但如今复活跟随十字架，我不羞于宣告它。\par

5. 既然祂像其他人一样有肉身，祂被钉十字架，但不是因同样的罪。因为祂不是因贪婪而被处死，因为祂是贫穷的老师；祂不是因贪欲而被定罪，因为祂自己明说：\Quote{凡注视妇女，有意贪恋她的，他已在心里奸淫了她}（\BibleRef{玛 5:28}）；不是因为打人或急躁打人，因为祂还把另一面脸颊转给打祂的人；不是因为藐视法律，因为祂是法律的成全者；不是因为辱骂先知，因为祂就是先知所预报的那一位；不是因为克扣工人工资，因为祂无偿而自由地服事；不是因为言语、行为或思想上的罪，因为祂\Quote{没有犯过罪，口中也找不到诡诈；祂被辱骂，却不还骂；祂受虐待，却不报复}（\BibleRef{伯前 2:22}）；\Quote{祂不是被迫，而是自愿地接受苦难；是的，若有人劝阻祂，甚至说：\InnerQuote{主，千万别这样}，祂仍会说：\InnerQuote{撒殚，退到我后面去！}}（\BibleRef{玛 16:22}）。\par

6. 你们可愿相信祂是甘愿受难的吗？别人因不预知而死，是出于无奈；而祂却提前论及受难说：\BibleQuote{看，人子将被出卖，被钉在十字架上。}但你们可知，这位人之友为何不逃避死亡？是为了避免全世界因罪而丧亡。\BibleQuote{看，我们上耶路撒冷去，人子将被出卖，被钉在十字架上。}\BibleQuote{再者，祂坚决面向耶路撒冷而行。}\BibleRef{路 9:5}你们可愿确知，十字架是耶稣的光荣？请听祂自己的话，而非我的。犹达斯对家主忘恩负义，即将出卖祂。他刚离开筵席，饮过祝福之杯，却以这救恩之饮为回报，图谋流无辜之血。\BibleQuote{吃过祂面包的人，却举脚踢祂。}他的手刚才还领受祝福的礼物，转眼间却为出卖的赏金而密谋置祂于死地。当他受责斥，听到那句话：\BibleQuote{你说的是}\BibleRef{玛 26:25}，他又出去了；于是耶稣说：\BibleQuote{时候到了，人子应受光荣。}\BibleRef{若 12:23}你们可看出祂如何视十字架为祂当有的光荣？那么，依撒意亚被锯成两半而不会羞耻，基督为世界而死却会羞耻吗？\BibleQuote{现在人子受了光荣。}\BibleRef{若 13:31}并非祂先前没有光荣：因为祂\BibleQuote{受光荣}于\BibleQuote{世界奠基之前}的光荣。祂作为天主，永远受光荣；但现在祂要藉佩戴祂忍耐的冠冕而受光荣。祂不是被迫交出生命，也不是被凶残暴力所杀，而是自愿的。请听祂说：\BibleQuote{我有权舍弃我的生命，也有权再取回它。}我自愿将它交给我的仇敌；因为除非我选择，这事不可能发生。因此祂是带着自己的明确目的来受难，为祂高尚的行为而喜乐，对这冠冕微笑，因人类的得救而振奋；不以十字架为耻，因为它是为拯救世界。因为受苦的不是一个普通人，而是成了人性的天主，为获得祂忍耐的奖赏而奋战。\par

7. 但是犹太人反驳这事，他们总是惯于吹毛求疵，不愿相信；因此刚刚读过的先知说：\BibleQuote{主，谁相信了我们的报道？}\BibleRef{依 52:15}波斯人相信，希伯来人却不信；\BibleQuote{那些未曾听闻祂的人将要看见，那些未曾听见的人将能明了。}而那些研究这些事的人，却鄙视他们所研究的。他们攻击我们说：\Quote{那么主会受苦吗？什么？人的手能控制祂的主权吗？}请读哀歌；因为在那哀歌中，耶肋米亚为你哀悼，写了值得哀悼的事。他看见了你的毁灭，目睹了你的倾覆，哀悼了那时的耶路撒冷；因为\BibleQuote{现在的}\BibleRef{迦 4:25}不会被哀悼；因为那时的耶路撒冷钉死了基督，而\BibleQuote{现在的}却朝拜祂。于是他在哀叹中说：\BibleQuote{我们气息的呼吸——主基督，在我们的腐败中被捕获。}\BibleRef{哀 4:20}我难道是在陈述自己的观点吗？看，他作证说主基督被人捉拿。由此将有什么结果？告诉我，先知。他说：\BibleQuote{我们曾说过：在他的荫下，我们将在列邦中生活。}因为他指明生命的恩宠不再住在以色列，而是住在外邦人中。\par

8. 但是，因着他们多有反对，来，让我藉你们的祈祷，（在时间短促所允许的限度内，）奉主的恩宠，陈述一些关乎苦难的见证。因为有关基督的事都已写入圣经，无一可疑，因为无一事无经文可据。一切都铭刻在先知们的记载上，清楚写就，不是写在石版上，而是藉圣神。既然你们听了福音论及犹达斯的话，难道不应接受这见证吗？你们听说祂被长矛刺透肋旁；难道不应查看这事是否也有记载？你们听说祂在园中被钉十字架；难道不应查看这事是否也有记载？你们听说祂被卖三十块银钱；难道不应学习是哪位先知说的这话？你们听说祂被给予醋喝；学习这事在哪里也有记载。你们听说祂的身体被安放在岩石中，并有一块石头盖在上面；难道不应也从先知接受这见证吗？你们听说祂与强盗同钉；难道不应查看这事是否也有记载？你们听说祂被埋葬；难道不应查看祂埋葬的细节是否在某处准确记载？你们听说祂复活了；难道不应查看我们教导这些事是否在嘲笑你们？因为\BibleQuote{我们的言语和宣讲，并不是凭人的智慧的劝导。}我们现在不再搅动诡辩的巧计；因为这些将被揭露；我们不以话语战胜话语，因为这些会过去；而是\BibleQuote{我们宣讲被钉在十字架上的基督}\BibleRef{格前 1:23}，这基督早已被先知们预讲过了。但我祈求你，接受这些见证，并印在你的心中。并且，由于它们很多，而我们剩下的时间已压缩为短暂时段，请趁时间允许，现在听一些更重要的；接受了这些开端后，要勤勉寻求其余的。不要让你的手只伸出接受，也要准备好工作。天主白白地赐予一切。\BibleQuote{因为你们中若有人缺少智慧，就应祈求那慷慨赐予的天主}\BibleRef{雅 1:5}，他必将获得。愿祂藉你们的祈祷，赐予我们说话者口才，也赐予你们听者信德。\par

9. 让我们寻求基督苦难的见证：因为我们聚在一起，现在不是要对圣经作推测性的讲解，而是要为我们已经相信的事得到确证。你们已经从我先前的讲论中领受了关于\Person{耶稣}来临的见证；以及关于他在海面上行走的见证，因为经上记述：\Emph{\Quote{你的道在海中}}。此外，关于各种治愈，你们也在另一场合领受了见证。所以现在我从苦难的开始讲起。\Person{犹达斯}是叛徒，他前来攻击他，站着说和平的话，却图谋战争。因此，圣咏作者说：\Emph{\Quote{我的朋友和我的邻人靠近我站着攻击我}}。又说：\Emph{\Quote{他们的话语比油还要柔和，其实却是刀剑}}。\Emph{\BibleQuote{师傅，你好！}}\BibleRef{玛 26:49}——然而他竟将他师傅出卖至死；主警告他时，他并不羞愧，因为主说：\Emph{\BibleQuote{犹达斯，你用亲嘴来负卖人子吗？}}\BibleRef{路 22:48}？因为主对他说的正是这话：记住你的名字；\OriginalTerm{犹达斯}{Judas}的意思是\Emph{\OriginalTerm{告解}{confession}}；你立了约，你收了钱，快些告解吧。\Emph{\Quote{天主，请勿沉默不语；因为邪恶和欺诈的口已张开攻击我；他们用诡诈的舌头对我说话，以仇恨的言语围绕我}}。至于一些大司祭也在场，他是在城门前被捆绑的，你们此前已经听过了，如果你们还记得那篇圣咏的讲解，它说明了时间和地点：\Emph{\Quote{他们傍晚返回，像狗一样饥饿，环绕着城}}。\par

10. 请听关于三十块银子的见证。\Emph{\Quote{如果你们认为好，就给我工价；不然，就作罢。}}\BibleRef{匝 11:12}以及后面的部分。你们欠我一个工价，因为我医治了瞎子和瘸子，但我却收到另一个工价；以感恩换羞辱，以敬拜换侮辱。你们看到圣经是如何预见了这些事吗？\Emph{\BibleQuote{他们称了我三十块银子作为我的工价。}}\BibleRef{匝 11:12}这预言何等精确！圣神的智慧何等伟大又无误！因为他说的不是十块，也不是二十块，而是三十块，恰如其数。先知啊，请告诉我们，这工价后来如何？那收下的人是否保留它？还是归还？归还后又变成什么？于是先知说：\BibleQuote{我就取了那三十块银子，扔进上主的殿中，进入熔炉。}将福音与预言比较：福音说：\BibleQuote{犹达斯后悔了，把那银子丢在圣殿里，出去离开了。}\par

11. 但现在我必须寻求这个看似矛盾之处的精确解答。因为那些轻视先知的人声称，先知说：\Emph{\BibleQuote{我就把它们扔进上主的殿中，进入熔炉。}}而福音却说：\Emph{\BibleQuote{他们就用那银子买了陶工的地。}}请听如何两者都是真实的。因为那些自以为义的犹太人，当时的大司祭们，看到\Person{犹达斯}后悔并说：\Emph{\BibleQuote{我犯罪了，因为我出卖了无辜的血。}}他们回答说：\Emph{\Quote{这与我们何干？你自己承担吧！}}那么，这血价与你们这些钉十字架的人无关吗？那收下并归还杀人价金的人要承担责任，而你们这些杀人者反而不承担责任吗？然后他们彼此商议说：\Emph{\Quote{把这血价放入圣库是不合法的。}}出自你们自己的口，你们定了自己的罪；如果那价金是污秽的，那行为也是污秽的；但如果你们钉死基督是履行正义，为什么不收下这价金呢？但邪恶的关键在于：如果福音说\Emph{\Quote{陶工的地}}，而先知说\Emph{\Quote{熔炉}}，怎能没有矛盾？不，不仅金匠或铜匠有熔炉，陶工也有熔炉处理他们的黏土。因为他们从砾石中筛出精细、肥沃、有用的泥土，并分离出废料，先用水调和黏土，以便容易捏成所需的形状。那么，当福音明确说\Emph{\Quote{陶工的地}}，而先知像谜语般说预言，因为圣经在许多地方是谜语式的，你们为何惊奇呢？\par

12. 他们绑了\Person{耶稣}，带他到大司祭的厅堂。你们愿意学习并知道这也是经上所记的吗？\Person{依撒意亚}说：\Emph{\BibleQuote{祸哉，他们的灵魂！因为他们对自己图谋了邪恶，说：让我们捆绑义人，因为他使我们厌烦。}}的确，\Emph{\BibleQuote{祸哉，他们的灵魂！}}我们来看\Person{依撒意亚}是如何被锯成两半的，但此后百姓得以恢复。\Person{耶肋米亚}被扔进水井的淤泥中，然而犹太人的创伤却得了痊愈；因为那是冒犯人的罪，是较轻的。但当犹太人犯罪，不是冒犯人，而是冒犯具有人性之天主时，\Emph{\BibleQuote{祸哉，他们的灵魂！}}——\Emph{\BibleQuote{让我们捆绑义人}}；有人会说，难道他不能使自己自由吗？他，在第四天从死亡的束缚中释放了\Person{拉匝禄}，并解开了\Person{伯多禄}在监狱中的铁链？天使们站在一旁，随时准备说：\Emph{\Quote{让我们挣断他们的锁链。}}但他们没有行动，因为他们的主愿意承受这一切。此外，他被带到长老们的审判台前；你们已经有此见证：\Emph{\BibleQuote{上主将亲自与他百姓中的长老和官长进行审判。}}\BibleRef{依 3:14}\par

13. 但 \OriginalTerm{大司祭}{High-priest} 审问了祂，听见了真理，就发怒；恶人的恶差役击打祂；那曾如太阳发光的容颜，竟忍受不法之手的击打。另有人前来，向那曾用唾沫治好生来瞎眼者的脸吐唾沫。\BibleQuote{你们就这样报答主吗？这人民愚昧无知。}\BibleRef{申 32:6} 先知大为惊讶，说：\BibleQuote{主，谁相信了我们的报告？}\BibleRef{依 53:1} 这事难以置信：天主，天主子，\OriginalTerm{上主的臂膀}{the Arm of the Lord}\BibleRef{依 53:1} 竟要受此苦难。但为使将要得救的人不致不信，圣神预先藉基督的口写下（因为那后来亲自经历这些事的，正是先前说这些话的那位）：\BibleQuote{我将我的背转给鞭打者}（因为比拉多 \Emph{鞭打了祂，就把祂交付去钉十字架}）；\BibleQuote{我的脸颊转给掌掴者；我并未转面回避唾污的羞辱；} 祂仿佛说：\Quote{我虽预知他们将击打我，但我连脸也未转开；因为我若自己畏惧这事，怎能激励我的门徒为真理视死如归？} 我说过：\BibleQuote{爱惜自己性命的，必要丧失性命。}\BibleRef{若 12:25} 我若爱惜自己的性命，怎能言行一致地教导人呢？所以，祂先以天主自身，忍受了来自人的这些苦难；为使我们这些日后为祂的缘故而忍受同样苦难的人，不至羞愧。你看见，先知们也已清楚预言了这些事。然而，如我先前所说，因时间有限，我略过了许多圣经见证；因为若有人仔细查考全部，关于基督的事，没有一件是没有见证的。\par

14. 祂被捆绑后，从盖法那里来到比拉多面前——这事也记载了吗？是的；\BibleQuote{他们绑了祂，将祂作为礼物送给 \OriginalTerm{雅里姆王}{king of Jarim}。} 但这里有些敏锐的听众会反驳：\Quote{比拉多不是王，}（暂且放下问题的主要部分）\Quote{那么他们如何绑了祂，将祂作为礼物送给王呢？} 但请读福音；\BibleQuote{比拉多听说祂是加里肋亚人，就把祂送到黑落德那里。}\BibleRef{路 23:6} 因为黑落德那时是王，正在耶路撒冷。现在请看先知预言的精确性：他说，祂被当作礼物送去；因为\BibleQuote{就在那一天，比拉多与黑落德成了朋友，以前他们是彼此为敌的。}因为那即将使天地和好的那一位，理应先使那些定祂罪的人彼此和好；因为主亲自在场，\BibleQuote{祂使地上君王的内心和好。}留心先知们的精确无误，以及他们真实的见证。\par

15. 当以敬畏之心看待受审判的主。祂任凭自己被士兵牵引和抬走。比拉多坐着审判，而那位坐在圣父右边的主，却站着受审。祂曾从埃及地、并多次从其他地区救赎的百姓，却向祂喊叫：\BibleQuote{除掉祂，除掉祂，钉祂在十字架上。}\BibleRef{苏 19:15} 为什么，你们犹太人？是因为祂治好了你们的瞎子吗？还是因为祂使你们的瘸子行走，并赐予其他恩惠？因此先知也惊讶地说：\BibleQuote{你们向谁张口，向谁伸舌？}\BibleRef{依 57:4} 主自己也在先知书中说：\BibleQuote{我的产业向我如同林中的狮子；它发声反对我，因此我憎恨它。}\BibleRef{耶 12:8} 我没有拒绝他们，他们却拒绝了我；因此我说：\BibleQuote{我撇弃了我的房屋。}\par

16. 当祂受审时，祂保持沉默；以致比拉多也为祂动容，说：\BibleQuote{你没有听见他们作证控告你什么吗？}\BibleRef{玛 27:13} 并非因为他认识受审者，而是因为他害怕别人向他报告的自己妻子的梦。耶稣仍缄默不语。圣咏作者说：\BibleQuote{我成了一个听不见的人，口中没有驳斥。}又说：\BibleQuote{我像一个聋子，听不见；像一个哑巴，不开口。}你以前曾听过关于这事的话，如果你还记得。\par

17. 但围拢的士兵戏弄祂，他们的主成了他们的笑柄，他们在自己的主人身上开玩笑。\BibleQuote{他们看着我，摇着头。}然而王权的形象显现；虽是戏弄，他们却屈膝下拜。士兵们在钉祂十字架以前，给祂穿上紫红袍，给祂戴上冠冕；即使那是荆棘的，又有什么关系？每个王都由士兵宣告；耶稣也必须以预象的方式由士兵加冕；因此圣经在雅歌中说：\BibleQuote{耶路撒冷的众女子啊，你们出去，观看所罗门王在他母亲给他戴上的冠冕。}\BibleRef{歌 3:11} 那冠冕本身是一个奥秘；因为它象征罪的赦免，诅咒的解除。\par

18. 亚当接受了判决：\BibleQuote{因你的劳作，地当受咒骂；它要给你生出荆棘和蒺藜。}为此耶稣接受了荆棘，为取消那判决；也为此祂被埋葬在地里，使那曾被诅咒的地，因祂而领受祝福而非诅咒。在犯罪之时，他们用无花果叶子蔽体；为此耶稣也使无花果树成为祂最后的神迹。因为当祂即将受难时，祂诅咒了那棵无花果树，不是所有的树，而是那一棵，为预象的缘故，说：\BibleQuote{今后永没有人吃你的果子。}\BibleRef{谷 11:1} 让那惩罚被取消。又因为他们从前用无花果叶子蔽体，祂来到之时正是无花果树上通常找不到食物的季节。谁不知道冬天无花果树不结果，只有叶子呢？耶稣难道不知道这众人皆知的事吗？不，祂虽知道，却仍像寻找什么而来；并非不知不会找到，而是要表明那象征性的诅咒只延及叶子。\par

19. 既然我们触及了与乐园相关的事，我确实对预像的真实性感到惊讶。在乐园中发生了堕落，而在一个园子里有了我们的救恩。从树而来的是罪，直到树罪才终止。\BibleQuote{傍晚，当主在园中行走时，他们便隐藏了。} \BibleRef{创 3:8} 而在傍晚，主将强盗带入了乐园。但有人会对我说：你在编造巧妙的说法；请从某位先知那里给我指出十字架的木头；除非你给我从先知来的证据，我不会被说服。那么请听耶肋米亚的话，并确信你自己：\BibleQuote{我好像一只温顺的羔羊被牵去宰杀；难道我不知道吗？}（因为应当像我已读过的那样，把它读作一个问句；因为那说\BibleQuote{你们知道两天后就是逾越节，人子将被出卖而被钉在十字架上} \BibleRef{玛 26:2} 的那一位，难道不知道吗？）\BibleQuote{我好像一只温顺的羔羊被牵去宰杀；难道我不知道吗？}（但那是怎样的一只羔羊？让若翰洗者来解释，当他说道：\BibleQuote{看，天主的羔羊，除去世界之罪的！} \BibleRef{若 1:29} ）\BibleQuote{他们图谋邪恶的计谋攻击我，说} \BibleRef{耶 11:19} ——（祂知道他们的计谋，难道不知道其结果吗？他们说了什么？）——\BibleQuote{来，让我们把一根木柴放在祂的饼上} ——（如果主认为你配得，日后你必将学到，根据福音，祂的身体带有饼的形象；）——\BibleQuote{那么，来，让我们把一根木柴放在祂的饼上，把祂从活人之地剪除；}——（生命不会被剪除，你们为什么徒劳？）——\BibleQuote{祂的名字将不再被记起。}你们的计谋是枉然的；因为\Emph{在太阳面前，祂的名字}常存于教会中。而挂在十字架上的是生命，梅瑟哭着说：\BibleQuote{你的生命必在你眼前悬着；你昼夜恐惧，不得深信你的生命。} \BibleRef{申 28:66} 同样地，刚才作为经文读的\BibleQuote{主，谁相信了我们的报告？}也是如此。\par

20. 这就是梅瑟通过把蛇固定在十字架上所完成的预像，使凡被活蛇所咬的，仰望铜蛇，便能因相信而得救。那么，铜蛇被钉在十字架上时能救人，难道降生成人的天主之子被钉在十字架上就不能救人吗？每一次生命都是借着木头而来的。因为在诺厄时代，生命的保存是借着木制的方舟。在梅瑟时代，海看见那象征性的棍杖，就对击打它的那人畏惧；那么，梅瑟的棍杖有力量，而救主的十字架就没有力量吗？但我略过大部分的预像，以保持适当篇幅。在梅瑟的情况下，木头使水变甜；而从耶稣的肋旁，水在木头上流出。\par

21. 在梅瑟之下，征兆的开端是血和水；而耶稣最后的一切征兆也是同样的。首先，梅瑟把河变成血；而耶稣在最后从祂的肋旁流出血和水。这或许是因为两种言论：审断祂的那人的言论，以及那些呼喊反对祂的人的言论；或者是因为信徒和不信者。因为比拉多说：\Quote{我是无罪的}，并用水洗手；那些呼喊反对祂的人说：\BibleQuote{祂的血归到我们身上} \BibleRef{玛 27:24}；因此这两样从祂的肋旁出来：水，或许是为那审断祂的人；而血，是为那些呼喊反对祂的人。再者，也可以从另一方式来理解：血是为犹太人，水是为基督徒：因为对于他们这些谋划者，从血而来的定罪；但对于你——现在相信的人——借着水而来的救恩。因为没有任何事情是没有意义的。我们的父辈——那些写过注释的——给了这件事另一个理由。由于在福音书中，救恩之圣洗的力量是双重的：一是借着水赐给受光照者的，二是借着他们自己的血赐给在迫害中的神圣殉道者的，从那救世的肋旁流出了血和水，以确认为基督所作的宣认的恩宠，无论是在圣洗还是在殉道之时。还有另一个提及肋旁的理由。那从肋旁被造的女人引入了罪；但耶稣——来将赦免的恩宠赐给男人和女人——在肋旁被刺，是为女人，为解除那罪。\par

22. 凡愿意探究的，还会发现其他理由；但所说的话已足够，因时间短促，且免使听众的注意力变得厌倦。然而，我们永远不会厌倦听有关我们主受冠冕的事，尤其是在这最神圣的哥耳哥达。因为别人只是听到，但我们既看见又触摸。不要让任何人厌倦；为了十字架本身的原因拿起你们的武装对抗仇敌；把十字架的信德作为战利品树立起来，以驳斥反对者。因为当你将要与不信者辩论基督的十字架时，首先用手作出基督十字架的记号，那反对者便会缄默不语。不要羞于承认十字架；因为天使以它为荣，说：\BibleQuote{我们知道你们寻找谁，那被钉在十字架上的耶稣。} \BibleRef{玛 28:5} 天使啊，难道你不可以说：\Quote{我知道你们寻找谁，我的主？} 但他大胆地说：\Quote{我知道那被钉十字架者。} 因为十字架是冠冕，而非耻辱。\par

23. 现在让我们回到我所说的来自先知的证据。主被钉在十字架上；你们已经接受了见证。你们看见这个哥耳哥达的地点！你们以欢呼的喊声回应，仿佛表示同意。务要在迫害之时不背弃你们所信。不要只在平安之时以十字架为喜乐，也应在迫害之时坚守同一信仰；不要在和平之时作耶稣的朋友，而在战争之时作祂的仇敌。你们现在领受罪过的赦免，和君王属灵恩赐的礼物；当战争来临时，要为主英勇作战。无罪的耶稣为你们被钉十字架；你们岂不愿为那为你们被钉十字架者而被钉十字架吗？你们不是在施恩，因为你们先已领受；你们是在报恩，偿还你们欠那在哥耳哥达为你们被钉十字架者的债。哥耳哥达被解释为\Quote{髑髅之地}。那么，是谁曾预言性地命名这个地方为哥耳哥达，而在此处基督——真正的头——忍受了十字架？如宗徒所说：\InnerQuote{祂是不可见的天主的肖像}；稍后又说：\InnerQuote{祂是身体——教会的头}。又说：\InnerQuote{男人的头是基督}（\BibleRef{格前 11:3}）；又说：\InnerQuote{祂是一切率领者和掌权者的头}（\BibleRef{哥 2:10}）。头在\Quote{髑髅之地}受了苦。哦，奇妙的预言性称谓！这名字本身也提醒你们，说：\InnerQuote{不要将受钉者想作一个凡人；祂是\Emph{一切率领者和掌权者的头}。那被钉的头是全能者的头，而祂的头是父；因为男人的头是基督，而基督的头是天主}（\BibleRef{格前 11:3}）。\par

24. 基督为了我们而被钉十字架，祂在夜里受审，当时天冷，所以生了一盆炭火（\BibleRef{若 18:18}）。祂在第三时辰被钉；从第六时辰直到第九时辰，遍地黑暗（\BibleRef{玛 27:45}）；但从第九时辰起又有光。这些事也有记载吗？让我们查考。先知匝加利亚说：\InnerQuote{到那日，必没有光，必有一日寒冷和霜冻；}（那寒冷就是伯多禄取暖的原因；）\InnerQuote{那日必被主知晓}（\BibleRef{匝 14:6}）；（什么，难道祂不知道别的日子吗？日子很多，但\InnerQuote{这是主忍耐的一天，是主所造的}；）——\InnerQuote{那日必被主知晓，不是白昼，也不是黑夜：}先知所说的这个隐晦的话是什么意思？那一日既不是白昼也不是黑夜？那么我们将如何称呼它？福音叙述事件来诠释它。那不是白昼；因为太阳从升起到落下并非始终照耀，但从第六时辰到第九时辰，正午有黑暗。因此黑暗被插入；但\InnerQuote{天主称黑暗为夜}（\BibleRef{创 1:5}）。所以那既不是白昼也不是黑夜：因为既不是全部光明，不可称为白昼；也不是全部黑暗，不可称为黑夜；但在第九时辰之后，太阳又照耀了。这位先知也预言了这一点；因为说了\InnerQuote{不是白昼，也不是黑夜}之后，他又加上\InnerQuote{到晚间必有光}。你们看见先知们的精确吗？你们看见预先写下的事情的真理吗？\par

25. 但你是否要问，太阳究竟在什么时辰亏缺？是第五时辰、第八时辰，还是第十时辰？先知啊，请告诉不愿听的犹太人确切的时间：太阳何时落下？亚毛斯先知回答说：\InnerQuote{到那日，上主天主说，太阳要在正午落下}（因为从第六时辰起有黑暗）；\InnerQuote{白昼，光明要在世上昏暗。}先知啊，这是什么样的季节，什么样的日子？\InnerQuote{我要使你们的庆节变为悲哀}；因为这事发生在无酵节和逾越节期间；随后他又说：\InnerQuote{我要使他像哀悼独生子一样，使那些与他同在的人痛苦如一日}（\BibleRef{亚 8:10}）；因为在无酵节和庆节，他们的妇女哀号哭泣，宗徒们躲藏起来，心中痛苦。这预言真是奇妙！\par

26. 但也许有人说：请给我另一个征兆；对于所发生的事，还有什么其它确切征兆？耶稣被钉十字架；祂只穿了一件上衣和一件外衣：士兵们将祂的外衣撕成四份分了；但祂的上衣没有撕开，因为撕开就不再有用；于是士兵们为此拈阄；因此他们分了一件，对另一件拈阄。这也有记载吗？那些勤勉的教会歌咏者知道，他们效法天使军旅，不断赞美天主；他们配得在这个哥耳哥达歌唱圣咏，说：\InnerQuote{他们分了我的衣服，为我的外衣拈阄。}\InnerQuote{阄}是士兵所拈的。\par

27. 再者，当祂在比拉多面前受审后，祂被穿上红衣；因为他们在那里给祂披上紫红袍。这也有记载吗？依撒意亚说：\InnerQuote{那从厄东而来的是谁？祂的衣服因波责辣而染红}（\BibleRef{依 63:1}）；（这位在屈辱中穿紫红袍的是谁？因为波责辣在希伯来文中有类似含义。）\InnerQuote{为何你的衣服是红的，你的衣裳像从榨酒池中出来的？}但祂回答说：\InnerQuote{我终日向悖逆和对抗的百姓伸开我的手。}\par

28. 他在十字架上伸开双手，为要怀抱地极；因为此哥耳哥达正是大地的中心。这不是我的话，而是一位先知曾说的：\Quote{你在大地中间行了救恩}。祂伸出了人类的手，那曾以属神之手建立诸天的一位；它们被钉子穿透，使祂的人性，那背负了人的罪，被钉在树上，且死去，罪恶也随之而死，我们得以在义中复活。\Quote{因为死是藉着一人而来，生命也是藉着\Emph{一人}而来}；藉着那一位，救主，自愿死去：要记得祂曾说：\BibleQuote{我有权舍去我的生命，也有权再取回它} \BibleRef{若 10:18}。\par

29. 然而，祂虽为众人的救恩而来，忍受了这一切，但百姓却以恶报回应祂。耶稣说：\Quote{我渴}——祂曾为他们在坚岩中引出水来；祂向祂所栽种的葡萄树求果实。但这葡萄树做了什么呢？这葡萄树，按本性出自圣祖，按心意却属索多玛——（\BibleQuote{他们的葡萄树是索多玛的，他们的葡萄枝是哈摩辣的} \BibleRef{申 32:32}；）——这葡萄树，当主渴时，竟用海绵蘸了醋，放在芦杆上，递给祂。\Quote{他们拿苦胆给我当食物，我渴时，他们拿醋给我喝}。你看先知预言的清晰描述。他们往我口中放的是哪种苦胆呢？经上说：\BibleQuote{他们拿没药调和的酒给祂} \BibleRef{谷 15:23}。没药的味道如同苦胆，非常苦。这些就是你们报答主的事吗？葡萄树啊，这就是你献给主人的供品吗？先知依撒意亚曾恰如其分地哀叹你说：\BibleQuote{我所爱的在肥沃的山丘上有一个葡萄园；}（不必引全）他说：\Quote{我期望它结葡萄}；我渴望它产出酒；\Quote{它却结出荆棘} \BibleRef{依 5:1}；因为你们看见我所戴的冠冕。那么我现在要颁布什么判决呢？我要命令云彩不再降雨在其上。因为那曾被称为先知的云彩已从他们那里移去，今后将归于教会；如保禄所说：\BibleQuote{先知可以作两次或三次发言，其他的要辨别} \BibleRef{格前 14:29}；又说：\BibleQuote{天主在教会中设立的，有宗徒，有先知} \BibleRef{弗 4:11}。阿加波，他曾捆缚自己的手脚，正是一位先知。\par

30. 论及与祂一同被钉的强盗，经上记载：\BibleQuote{祂被列于罪犯之中} \BibleRef{依 53:12}。这两人先前都是罪犯，但其中一个不再如此。因为那一个始终是罪犯，顽固抵抗救恩；他虽然双手被钉，却用舌头发出亵渎。当路过的犹太人摇头嘲笑被钉者，应验了经上所载：\Quote{他们望着我，摇头}，他也同他们一起辱骂。但另一个斥责了那辱骂者；这于他是生命的终结与复兴的开始；他交出灵魂，作为救恩中的最初一份。在斥责了另一个之后，他说：\Quote{主，求祢记念我}；因为我的账在祢那里。不要理会这人，因为他悟性的眼目被蒙蔽；但求祢记念我。我不是说，记念我的行为，因为我害怕这些。每个人对自己的旅伴都有感情；我正与祢一同走向死亡；求祢记念我，祢的同行者。我不是说现在记念我，而是\Quote{当祢来临时，在祢的王国里}。\par

31. 强盗啊，是什么力量引你到光明？谁教导你去朝拜那位被轻视的人，你十字架上的同伴？永光啊，祢光照在黑暗中的人！因此他也恰当地听到了这话：\Quote{放心吧}；不是因你的行为值得放心，而是因君王在此，广施恩惠。这请求指向遥远的时日，但恩宠却极迅速。\BibleQuote{我实在告诉你，今天你就要同我在乐园里}；因为\Quote{今天}你\Quote{听了我的声音}，没有\Quote{硬着心}。我极速地判定了亚当，也极速地宽恕了你。对他曾说的：\BibleQuote{你吃的那一天，必定要死} \BibleRef{创 2:17}；但你今天服从了信德，今天就是你的救恩。亚当因树而堕落；你因树而被领入乐园。不要怕那蛇；它不能赶出你，因为它已\BibleQuote{从天上坠落} \BibleRef{路 10:18}。我不对你说，今天你要离去，而是\BibleQuote{今天你要同我在一起}。鼓起勇气：你不会被赶出。不要怕那旋转的剑；它在主人面前退缩。哦，伟大而不可言喻的恩宠！忠信者亚巴郎尚未进入，但强盗进入了！梅瑟和先知们尚未进入，而这强盗，虽是一个违犯法律者，却进入了。保禄也在此前惊叹说：\BibleQuote{罪恶在哪里增多，恩宠在那里也格外丰富} \BibleRef{罗 5:20}。那些承受了白日炎热的人尚未进入，而第十一时辰的人进入了。不要让任何人埋怨家主，因为他说：\Quote{朋友，我没有亏待你；难道我不能拿我的财物做我所愿意的吗？} 强盗有意行义，但死亡阻止了他；我不单单等待行为，我也接纳信德。我来了，那\BibleQuote{在百合花中牧放我的羊群} \BibleRef{歌 6:3}，我来在花园中牧放它们。我\Quote{找着了}那\Quote{迷失的羊} \BibleRef{路 15:5}，但把它放在我肩上；因为他相信，既然他自己曾说：\Quote{我如同亡羊走迷了路}；\Quote{主，当祢进入祢的王国时，求祢记念我。}\par

32. 关于这个园子，我从前在\WorkTitle{雅歌}中向我的净配歌唱，并这样对她说：\BibleQuote{\Emph{我进入了我的园子，我的姊妹，我的净配}}\BibleRef{歌 5:1}；（\BibleQuote{\Emph{在祂被钉的地方有一个园子}}\BibleRef{若 19:41}；）然后你从那里带走什么呢？\BibleQuote{\Emph{我采集了我的没药}}；祂喝了掺有没药的酒和醋，在喝了以后，祂说：\Quote{完成了}。因为奥秘已应验；所记载的已成就；罪过已被赦免。因为\BibleQuote{\Emph{基督已来到，作未来美好事物的大司祭，经过了更大更全备的帐幕，不是人手所造的，就是说，不属于这个受造界；也不是藉着公羊和牛犊的血，而是藉着自己的血，只一次进入至圣所，获得了永远的救赎；因为如果公牛和公山羊的血，以及小母牛的灰烬，洒在不洁的人身上，能圣化他们，使身体洁净，那么基督的血更会怎样}}\BibleRef{希 9:11}？又如：\Quote{\Emph{所以，弟兄们，我们藉着耶稣的血，得以坦然进入至圣所，是藉着祂给我们开创的一条又新又活的道路，穿过帐幕，即祂的肉身}}。由于祂的肉身，这帷幔，受了侮辱，因此圣殿的预像帷幔也裂为两半，正如经上所记：\BibleQuote{\Emph{看，圣殿的帷幔从上到下裂为两半}}\BibleRef{玛 27:51}；没有留下一丝半缕；因为主人说过：\Quote{看，你们的房屋将留给你们荒芜}之后，整座房屋就都粉碎了。\par

33. 救主忍受了这一切，\BibleQuote{\Emph{藉着祂十字架的血，成就了和平，包括天上的和地上的}}\BibleRef{哥 1:20}。因为我们因罪成了天主的仇敌，而天主已定下罪人必死。因此必然要发生两件事中的一件：要么天主按祂的真理毁灭所有人，要么按祂的慈爱取消判决。但请看天主的智慧：祂既保全了判决的真理，也施行了慈爱。基督\BibleQuote{\Emph{亲身在十字架上承担了我们的罪，使我们藉着祂的死向罪而死，并向义而活}}\BibleRef{伯前 2:24}。为我们而死的那一位并非无足轻重；祂不是一只具体的羊；祂不仅仅是一个人；祂超过天使；祂是降生成人的天主。罪人的过犯并不及那为他们而死者的义；我们所犯的罪并不及那为我们舍命者所行的义——祂随己意舍命，也随己意再取回来。你想知道祂并非被迫舍命，也非违背己意交出灵魂吗？祂向父呼喊说：\BibleQuote{\Emph{父啊，我把我的灵魂交托在祢手中}}\BibleRef{路 23:46}；我交托它，是为了再取回来。说完这些话，\BibleQuote{\Emph{祂便断了气}}\BibleRef{玛 27:50}；但时间不长，因为祂很快就从死者中复活了。\par

34. 太阳昏暗了，因为\BibleQuote{\Emph{正义的太阳}}\BibleRef{拉 4:2}。岩石崩裂，因为那属神的磐石。坟墓开了，死人复活了，因为那一位是\Quote{\Emph{在死人中自由的}}；\BibleQuote{\Emph{祂从无水之坑中释放了祂的囚徒}}\BibleRef{匝 9:11}。所以不要以被钉者为耻，反倒要勇敢地说：\BibleQuote{\Emph{祂承担了我们的罪，为我们忍受痛苦，因祂的创伤我们获得了痊愈}}\BibleRef{依 53:4}。我们不要对恩人忘恩负义。又说：\Quote{\Emph{因我百姓的过犯，祂被引到死亡；我要使恶人作祂的坟墓，使富人作祂的死亡}}。因此保禄明确地说：\BibleQuote{\Emph{基督照经上所记载的，为我们的罪死了，被埋葬了，并且照经上所记载的，第三天复活了}}\BibleRef{格前 15:3}。\par

35. 但我们想要清楚知道祂被埋葬在哪里。祂的坟墓是手所造的吗？像君王的坟墓一样高出地面吗？坟墓是由石头砌成的吗？上面覆盖着什么？先知们，请告诉我们关于祂坟墓的精确真相，祂被放在哪里，我们将在哪里寻找祂？他们回答说：\BibleQuote{\Emph{看看你们凿出的磐石}}\BibleRef{依 51:1}。\Emph{看进去}并观看。在福音中记载：\Quote{\Emph{在一个由磐石中凿出的石墓里}}。接下来发生什么？坟墓有什么样的门？又一位先知说：\Quote{\Emph{他们在牢狱中砍断了我的性命}}，并\Quote{\Emph{把一块石头扔在我身上}}。我，那\BibleQuote{\Emph{精选的、宝贵的房角石}}\BibleRef{伯前 2:6}，暂时躺在石头内——我是犹太人的绊脚石，信徒的救恩之石。因此\Quote{\Emph{生命之树}}被栽种在地上，使那受诅咒的地得以领受祝福，并使死人得释放。\par

36. 所以我们不要羞于承认被钉者。愿十字架成为我们的印记，我们敢于用指头划在额上，并划在一切事物上：在我们所吃的饼上、所饮的杯上；在我们出入之时；在睡前、躺下和起床时；在路途中、静默时。这护卫何等伟大；它是无价的，为了穷人的缘故；无须劳苦，为了病人的缘故；因为它的恩宠来自天主。它是信者的记号，也是魔鬼的恐惧：因为祂\BibleQuote{\Emph{在十字架上战胜了它们，公开羞辱了它们}}\BibleRef{哥 2:15}；因为它们看见十字架，便想起被钉者；它们害怕那一位，\Quote{\Emph{击碎了巨龙的头}}。不要因礼物是白送的而鄙视这印记；反而因此更要尊敬你的恩人。\par

37. 若你一旦陷入争辩，又没有确凿的证据，仍要让信德在你内坚定；或者，你更应成为博学之人，然后以先知之言使犹太人噤声，以他们自己的神话使希腊人无言。他们自己崇拜那些被雷击的人；但雷电从天而降，并非偶然。如果他们不以崇拜那些被雷击、被天主憎恶的人为耻，你岂会羞于崇拜那为你被钉在十字架上的天主爱子？我羞于讲述他们所谓的神祇的故事，因时间所限，姑且不提；让那些知道的人去说吧。也让所有的异端者闭口。若有人说十字架是幻象，你当远离他。憎恶那些说基督只是我们想象中被钉死的人；因为若是如此，且救恩来自十字架，那么救恩也是幻想了。若十字架是幻想，复活也是幻想；但\BibleQuote{如果基督没有复活，我们仍在我们的罪恶中}\BibleRef{格前 15:17}。若十字架是幻想，升天也是幻想；若升天是幻想，那么第二次来临也是幻想，从此一切皆虚无。\par

38. 因此，首先以十字架作为不可动摇的根基，在其上建造信德的其他条款。不要否认被钉者；因为你若否认祂，有许多人要指控你。叛徒犹达斯将首先指控你；那出卖祂的人知道祂被司祭长和长老判处死刑。那三十块银钱是见证；革责玛尼是见证，那里发生了背叛；我暂且不提橄榄山，那夜他们与祂在那里祈祷。黑夜中的月亮是见证；白昼和那被遮蔽的太阳是见证；因为它不忍目睹谋害者的罪行。火是见证，伯多禄曾站在它旁边取暖；你若否认十字架，永恒之火正等待你。我说严厉的话，好使你不受严厉的苦楚。要记住那在革责玛尼中朝着祂而来的刀剑，好使你免受永罚的剑。盖法的房子是见证，它今日的荒凉显明了那曾在那里受审者的权能。是的，盖法本身在审判之日要起来反对你，那用手掌击打耶稣的仆人也要起来反对你；还有那些捆绑祂、押送祂的人。甚至黑落德也要起来反对你；还有比拉多，仿佛在说：你为何否认那被犹太人在我们面前诬告、而我们明知他无罪的那一位？因为我比拉多那时洗了手。假证人要起来反对你，还有那些给祂穿上紫红袍、给祂戴上荆棘冠冕、在哥耳哥达钉死祂、为祂的衣服抽签的兵士。基勒乃人西满要向你呼喊，他曾在耶稣之后背负十字架。\par

39. 星辰中，那被遮蔽的太阳要向你呼喊；地上的事物中，那掺了没药的酒要呼喊；芦苇中，那芦苇要呼喊；草本中，那牛膝草要呼喊；海中之物中，那海绵要呼喊；树木中，那十字架的木头要呼喊；还有那些兵士，如我所说，他们钉了祂，为祂的衣服抽签；那用长矛刺透祂肋旁的兵士；当时在场的妇女们；那时被撕裂的圣殿帐幔；比拉多的厅堂，如今因那被钉者的权能而荒废；这圣哥耳哥达，高耸于我们之上，至今仍显现，还显示着当时因基督而岩石崩裂的情景；附近的坟墓，祂曾安葬其中；以及那放在墓门上的石头，至今仍躺在墓旁；当时在场的天使们；复活后朝拜祂的妇女们；奔往坟墓的伯多禄和若望；还有多默，他把手探入祂的肋旁，把手指探入钉痕。因为他为我们而如此仔细地触摸了祂；而你，当时不在场，本应寻求的，他在场时，因天主的眷顾，却已寻求了。\par

40. 你有十二宗徒，是十字架的见证；还有整个大地和相信那挂在十字架上者的世人。你在此处的临在本身，就当劝服你相信被钉者的权能。因为如今是谁带你来到这个集会？什么兵士？你被何种锁链约束？什么判决使你现在停留在此？不，是救恩的凯旋纪念、耶稣的十字架，把你们全都聚集在一起。正是它奴役了波斯人，驯服了斯基泰人；正是它赐给了埃及人，为了猫、狗和他们的众多谬误，对天主的认识；正是它，至今仍治愈疾病，至今仍驱逐魔鬼，推翻药物和符咒的骗术。\par

41. 这十字架将随耶稣从天再度显现；因为那纪念物要先行于君王之前：好使犹太人看见\BibleQuote{他们所刺透的那一位}\BibleRef{匝 12:10}，并因十字架而认识那受辱者，从而忏悔哀恸；（但他们\Quote{要一支派一支派地哀恸}，因为他们要忏悔，那时却再无忏悔之时；）好使我们得荣耀，因十字架而欢欣，朝拜那被派遣、为我们被钉的主，也朝拜那派遣祂的天主圣父，偕同圣神：愿光荣归于祂，至于无穷之世。阿们。\par

\SourceRecord{Translated by Edwin Hamilton Gifford.；Nicene and Post-Nicene Fathers, Second Series；Vol. 7.；Edited by Philip Schaff and Henry Wace.；Buffalo, NY: Christian Literature Publishing Co.,；1894.；<http://www.newadvent.org/fathers/310113.htm>.}

% structure: p0001 source=h1 target=section reason=page-title

\section{要理讲授第十四讲}

\Emph{论这句话：\Quote{第三日从死者中复活，升了天，坐在圣父之右。}}\par

\BibleRef{格前15:1-4}\par

\BibleQuote{弟兄们，我如今把先前所传与你们的福音通知你们。……祂照圣经所说，第三日复活了，等等。}\par

\Quote{耶路撒冷啊，欢乐吧！凡爱慕耶稣的，都踊跃庆祝吧！}因为祂复活了。\Quote{凡先前因听见犹太人的狂妄邪恶而悲伤的，都该欢乐。}因为那曾在此受他们恶待的，已复活了。既然关于十字架的讲论是令人悲伤的，那么就让复活的喜讯给听众带来喜乐吧。愿悲哀转为喜乐，哀叹化为欢欣；愿我们的口充满喜乐和欢唱，因为那在复活后说了\Quote{喜乐吧}的那一位。我知道在过去的日子里，基督的朋友们的忧伤；因为我们的讲论讲到死亡和埋葬便停止了，没有传出复活的喜讯，你们心中悬盼着想要听的事。\newline{}（见：\BibleRef{依66:10}）\par

如今，那死者复活了，祂是在死者中自主的，也是死者的拯救者。那曾忍受凌辱、戴上荆棘冠冕的，祂已复活，并以战胜死亡的冠冕加冕了祂自己。\par

既然我们已陈述了关于祂十字架的见证，那么现在让我们也验证祂复活的证据：因为宗徒在我们面前确认说：\Quote{祂被埋葬了，并于第三日按圣经复活了。}因此，既然宗徒已引导我们回到圣经的见证，我们就应当充分认识我们救恩的希望；并且首先学习圣经是否告诉我们祂复活的时节，是在夏天、秋天还是冬天之后；以及救主是从怎样的地方复活的，在可钦佩的先知书中，复活的地方被称为什么名称；还有那些寻找祂却找不到的妇女，后来是否因找到祂而欢欣；这样，当福音被诵读时，这些圣经的叙述就不会被认为是神话或狂想曲了。\par

关于救主被埋葬，你们在前面的讲论中已清楚听到，如依撒意亚说：\Quote{他的埋葬必在平安中。}因为藉着祂的埋葬，祂使天地和好，将罪人带向天主；以及\Quote{义人被从不义的道路上取去}；以及\Quote{他的埋葬必在平安中}；以及\Quote{我要将恶人作为他的埋葬。}还有雅各伯的预言，在圣经中说：\Quote{他卧下如雄狮，如母狮之子；谁能唤醒他？}（见：\BibleRef{创49:9}）以及民长纪中类似的一段：\Quote{他卧下如雄狮，如母狮之子。}（见：\BibleRef{户24:9}）还有你们常听到的圣咏说：\Quote{你把我带进了死亡的尘埃。}此外我们还留意了地点，当我们引用\Quote{看这岩石，你们所凿出的}这句话时。但如今，让关于祂复活的见证伴随我们的旅程。\par

首先，在第十一圣咏中祂说：\Quote{因贫苦者的委屈和穷困人的叹息，我现在要起来，主说。}但这段经文对某些人仍存疑问：因为祂也常发怒起来，向祂的敌人复仇。\par

那么来到第十五圣咏，它清楚地说：\Quote{主，求祢保护我，因为我投靠了祢。}之后：\Quote{我不参与流血的集会，我的嘴唇也不提他们的名字；}因为他们拒绝了我，而选择了凯撒作他们的王。接着又说：\Quote{我常将主置于我眼前，因祂在我右边，我决不动摇。}随后：\Quote{甚至夜间我的肺腑也责教我。}此后祂最清楚地说：\Quote{因为祢不会将我的灵魂撇在\OriginalTerm{地狱}{hell}，也不让祢的圣者见到\OriginalTerm{腐朽}{corruption}。}祂没有说\Quote{也不让祢的圣者见到死亡}，因为那样祂就不会死了；但祂说\Quote{我见不到腐朽，也不会留在死亡中。}\Quote{祢已将生命的道路指示给我。}看，这里清楚宣告了死后的生命。也来到第二十九圣咏：\Quote{主，我要称扬祢，因为祢提拔了我，没有让我的仇敌向我夸耀。}发生了什么事？祢从仇敌手中被拯救，还是即将被打时得到释放？祂自己最清楚地说：\Quote{主，祢从地狱把我的灵魂提上来。}那里祂预言说\Quote{祢不会遗弃}，而这里祂把将要发生的事当作已经发生来说：\Quote{祢已提上来。祢拯救了我免于下到深渊的人们。}这事件将在何时发生？\Quote{一宿虽然有哭泣，早晨便必欢呼。}因为晚上是门徒的悲伤，早晨是复活的喜乐。\par

但是你们也想知道地点吗？祂又在雅歌中说：\Quote{我下到核桃园}（见：\BibleRef{歌4:11}）；因为祂被钉的地方是一个园子。虽然现在它已用皇家礼物装饰得极其华丽，但从前它是一个园子，而且这些迹象和残留仍存在。\Quote{封闭的园子，密封的泉源}（见：\BibleRef{歌4:12}），这是针对那些犹太人的，他们曾说：\Quote{我们记得那个骗子活着的时候说过：三天后我要复活。所以请下令把坟墓把守妥当。}接着又说：\Quote{他们就去，将坟墓把守妥当，用石头封住，并派了卫兵。}（见：\BibleRef{玛27:63}）这些经文都说得准确。有一位说：\Quote{在安息中祢要审判他们。}但谁是那封闭的泉源，或谁被解释为\Quote{活水的泉源}（见：\BibleRef{歌4:15}）？就是救主自己，关于祂经上记载：\Quote{因为生命的泉源在祢那里。}\par

6. 但索福尼亚先知以基督的身份向门徒说了什么？\Emph{\Quote{你当预备，在黎明时兴起：他们所有的遗穗都毁了}}：这遗穗，即犹太人的遗穗，他们那里没有一串葡萄，连救恩的遗穗也没有留下，因为他们的葡萄树已被砍断。请看祂如何对门徒说：\Emph{\Quote{预备你自己，在黎明时兴起}}：应在黎明期待复活。\par

在同一上下文中的圣经里，祂又说：\Emph{\Quote{所以你当等候我，直到我复活得证的日子，}主说。}你看，先知也预见了复活的地点，这地点要被称为\OriginalTerm{证}{Testimony}。为什么哥耳哥达和复活的地点不像其他教会那样被称为教会，而被称为证？或许是因为先知曾说：\Emph{\Quote{直到我复活得证的日子。}}\par

7. 那么，这位是谁？复活的记号是什么？在同一上下文中的先知话里，祂明确说：\Emph{\Quote{那时我要转变给万民一种言语}}：因为复活后，当圣神被派遣时，赐下了舌音的恩赐，\Emph{\Quote{使他们都在同一轭下事奉主}}。同一位先知还提出什么别的标记，说他们要在\Emph{\Quote{同一轭下}}事奉主？\Emph{\Quote{从埃塞俄比亚的河流之外，他们要给我带来祭品}}。你知道宗徒大事录中记载，当埃塞俄比亚的太监从埃塞俄比亚的河流之外来时。\BibleRef{宗 8:27}因此，既然圣经既说明了时间与地点的特性，又说明了复活后所追随的记号，你从此当坚定信德于复活，不可让任何人动摇你宣认\Emph{\Quote{基督从死者中复活}}\BibleRef{弟后 2:8}。\par

8. 现在再取第八十七篇圣咏中的另一个证言，基督在先知们中说话（因为那时说话的那位后来来到我们中间）：\Emph{\Quote{上主，我救恩的天主，我昼夜在祢面前呼号}}，再稍后，\Emph{\Quote{我好像成了一个无人援助的人，在死者中自由}}。祂不是说，我成了无人援助的人；而是说，\Emph{\Quote{好像一个无人援助的人}}。因为祂被钉十字架不是由于软弱，而是出于自愿，祂的死不是出于不自愿的软弱。\Emph{\Quote{我被列于下到坑中的人中}}。那记号是什么？\Emph{\Quote{祢使我的相识远离我}}（因为门徒都逃走了）。\Emph{\Quote{祢要向死者显示奇迹吗？}}然后稍后：\Emph{\Quote{上主，我曾向祢呼号；我的祈祷要在清晨到达祢面前}}。你看见他们如何指出时辰、苦难和复活的精确时刻吗？\par

9. 救主从何处复活？祂在\BibleBook{}{雅歌}中说：\Emph{\Quote{起来，前来，我的佳偶}}：接着又说，\Emph{\Quote{在岩石的洞穴中}}！祂称那曾是在救主坟墓门前的洞穴为岩石的洞穴，是从岩石本身凿出来的，如同此地坟墓前常见的那样。因为如今已看不见了，当时外面的洞穴为了现今的装饰而被切除了。因为在君王慷慨装饰坟墓之前，岩石前面有一个洞穴。但那藏有洞穴的岩石在哪里？是在城中心附近，还是在城墙和郊外附近？是在古墙之内，还是在后来建的外墙之内？于是祂在\BibleBook{}{雅歌}中说：\Emph{\Quote{在岩石的洞穴中，靠近外墙}}。\par

10. 救主在什么季节复活？是夏季，还是别的？在\BibleBook{}{雅歌}中，紧接上述引文之前，祂说：冬天已过，雨水已止息；地上百花开放，修理葡萄树的时期已来到。那么，现在地上不正是百花开放，葡萄树正在修剪吗？你看祂也说冬天已过。因为当这个\OriginalTerm{克桑提科月}{Xanthicus}到来时，已经是春天了。这正是希伯来人的第一个月，逾越节庆节就在这个月发生，以前是预像，现在却是真实。这是世界受造的季节：因为那时天主说：\Quote{让地上生出青草，按种类和相似结种子的蔬菜。}现在，如你所见，每种植物都已结出种子。正如当时天主造了太阳和月亮，给了它们昼（和夜）相等的周期，同样，几天前正是春分的季节。\par

那时天主说：\Quote{让我们按我们的肖像和相似造人。}人获得了肖像，但因他的不顺从而模糊了相似。那么，就在他失去这一切的那个季节，恢复也发生了。就在受造的人因不顺从被逐出乐园的那个季节，相信的人因顺从被带回。我们的救恩就发生在与堕落相同的季节：当百花开放，修剪时期到来之时。\par

11. 一所园子就是祂被埋葬的地方，那里栽种着一棵葡萄树；祂曾说过：\BibleQuote{我是葡萄树}！祂被栽种在地里，为要拔出因亚当而来的咒诅。大地曾受罚长出荆棘和蒺藜；真葡萄树却从地里长出，使这句话得到应验：\BibleQuote{真理从地里生出，正义从天上俯视}。那被埋在园中的祂将说什么？\BibleQuote{我收集了我的没药和我的香料}；又说：\BibleQuote{没药与沉香，以及一切上等的香料}。这些是埋葬的象征；在福音中记载：\BibleQuote{妇女们来到坟墓，带着她们所预备的香料}（\BibleRef{路 24:1}）；\BibleQuote{尼苛德摩也带着没药与沉香调和的混合物}（\BibleRef{若 19:39}）。再往后又写道：\BibleQuote{我以我的蜂蜜吃了我的饼}：受难前是苦的，复活后是甜的。然后祂复活后，穿过了紧闭的门；但他们不相信是祂，因为\BibleQuote{他们以为看见了一个鬼魂}（\BibleRef{路 24:37}）。但祂说：\BibleQuote{你们摸摸我，看看}。把你们的手指放进钉痕里，正如多默所要求的。\BibleQuote{他们因欢喜还不敢相信，且惊奇；祂便对他们说：你们这里有吃的吗？他们给了祂一片烤鱼和一块蜜脾}（\BibleRef{路 24:41}）。你看见这话如何应验：\BibleQuote{我以我的蜂蜜吃了我的饼}？\par

12. 但在祂穿过紧闭的门之前，灵魂的新郎与追求者被那些高贵而勇敢的妇女所寻求。她们，那些有福的人，来到坟墓，寻找那已复活的主，眼泪仍从她们眼中滴落，而她们本应为复活的主欢欣舞蹈。玛利亚前来寻祂，按照福音，却没有找到；立刻她从天使那里听见，然后看见了基督。那么这些事也记载了吗？祂在\WorkTitle{雅歌}中说：\BibleQuote{我在床上寻见我灵魂所爱的那一位}。在什么时候？\BibleQuote{夜间我在床上寻见我灵魂所爱的那一位：玛利亚}，经上说，\BibleQuote{她是在天还黑的时候来的。我在夜间在床上寻祂，我寻祂，却寻不见}。而在福音中玛利亚说：\BibleQuote{他们把我主搬走了，我不知道他们把祂放在哪里}（\BibleRef{若 20:13}）。但当时在场的天使弥补了她们的无知；因为他们说：\BibleQuote{你们为什么在死人中找活人呢}（\BibleRef{路 24:5}）？祂不仅复活了，而且复活时还带着死者（\BibleRef{玛 27:52}）。但她不知道，并且在她身上，\WorkTitle{雅歌}对天使说：\BibleQuote{你们看见了我灵魂所爱的那一位吗？}我不过稍微离开他们（即两位天使），\BibleQuote{直到我寻见了我灵魂所爱的那一位。我拉住祂，不容祂走}（\BibleRef{歌 3:3}）。\par

13. 因为在看见天使之后，耶稣亲自作为自己的先驱而来；福音说：\BibleQuote{看，耶稣迎上她们说：愿你们平安！她们就上前抱住祂的脚}（\BibleRef{玛 28:9}）。她们抱住祂，为使这话应验：\BibleQuote{我要拉住祂，不容祂走}。虽然那女人身体软弱，但她的心灵是刚毅的。\BibleQuote{爱情，大水不能熄灭，江河不能淹没}（\BibleRef{歌 8:7}）；她们所寻的祂是死的，但对复活的希望并未熄灭。天使又对她们说：\BibleQuote{你们不要害怕}；我不是对兵士说\Quote{不要害怕}，而是对你们；至于他们，让他们害怕，好让他们受教后作见证说：\BibleQuote{这人真是天主子}（\BibleRef{玛 27:54}）；但你们不应害怕，\BibleQuote{因为完美的爱驱除恐惧}（\BibleRef{若一 4:18}）。\BibleQuote{你们去，告诉祂的门徒说祂复活了}（\BibleRef{玛 28:7}）；还有其他话。她们就带着快乐又充满恐惧离去；这事也记载了吗？是的，那论及基督受难的第二篇圣咏说：\BibleQuote{你们当以敬畏事奉主，以战栗而欢欣}——\Quote{欢欣}，因为复活的主；但\Quote{战栗}，因为地震和那如闪电般显现的天使。\par

14. 因此，虽然司祭长和法利塞人藉比拉多之手封了坟墓，但妇女们却看见了那位复活的主。依撒意亚深知司祭长的软弱和妇女信德的力量，于是说：\Quote{\Emph{你们这些从观看出来的妇人，到这里来}; \Emph{因为这百姓没有悟性}}；——司祭长缺乏悟性，而妇女却成了目击证人。当士兵们进城到他们那里，将所发生的事都告诉了他们，他们就对他们说：\BibleQuote{\Emph{你们就说：祂的门徒夜间来了，趁我们睡觉的时候把祂偷去了。}}\BibleRef{玛 28:13} 因此依撒意亚也预言了这事，好像是代表他们说的：\BibleQuote{\Emph{告诉我们，向我们述说另一个诡计。}}\BibleRef{依 30:10} 那位复活的主已经起来，他们用钱贿赂了士兵；但他们却无法说服我们时代的君王。那时，士兵们为银子出卖了真理；但今日的君王却以虔诚建造了这座我们救主天主复活的神圣圣殿，用银镶嵌，用金雕琢，我们正聚集在其中；并用金银和宝石的珍宝装饰了它。\BibleQuote{\Emph{倘若这事传到总督耳中}}，他们说，\BibleQuote{\Emph{我们会说服他。}}\BibleRef{玛 38:14} 是的，你们可以说服士兵，却无法说服世界；因为，当伯多禄逃出监狱时，他的守卫被定罪，而看守耶稣基督的人为何未被定罪呢？前者被黑落德宣判，因为他们无知，无话可说；而后者看见了真理，却为钱隐瞒了它，被司祭长包庇。然而，虽然当时只有少数犹太人被说服，但世界却顺从了。那些隐瞒真理的人自己反被隐藏；而那些接受真理的人却因救主的能力而显扬出来，祂不仅自己从死者中复活，还使死者与祂一同复活。欧瑟亚先知代表这些人明确地说：\BibleQuote{\Emph{两天后祂必复兴我们，第三天我们必复活，在祂面前生活。}}\BibleRef{欧 6:2}\par

15. 既然悖逆的犹太人不肯被圣经说服，反而忘记了一切记载，反对耶稣的复活，那么答复他们如下是合宜的：你们说厄里叟和厄里亚复活了死人，却反对我们救主的复活，这凭的是什么理由？难道是因为我们现在没有那一代的活见证人证明我们所说的事吗？那么，你们也提得出那一时代历史的见证人。但那事有记载——这事也有记载：你们为何接受一个，而拒绝另一个呢？写那历史的是希伯来人；所有的宗徒也都是希伯来人：你们为何不信这些犹太人呢？写福音的玛窦是用希伯来语写的；宣讲者保禄是希伯来人中的希伯来人；十二位宗徒全是希伯来血统；随后从希伯来人中相继任命了十五位耶路撒冷主教。那么，你们有什么理由容许你们自己的记述，却拒绝我们的，尽管这些也是由你们当中的希伯来人所写？\par

16. 但有人会说：死人复活是不可能的；然而厄里叟两次复活了死人——一次在他活着时，一次在他死后。那么，我们相信：当厄里叟死后，一个死人被扔在他身上，碰触了他，就起来，而基督却没有复活吗？但在那个事例中，碰触厄里叟的死人起来了，而使他复活的人仍然死了；但在这个事例中，我们所说的这位死者自己复活了，而且许多死人甚至没有碰触祂就复活了。因为\BibleQuote{\Emph{许多长眠的圣徒的身体复活了，在祂复活后，他们从坟墓里出来，进了圣城}}\BibleRef{玛 27:52}（显然是指我们现在所在的这座城），\BibleQuote{\Emph{并向许多人显现。}}厄里叟曾复活一个死人，但他没有战胜世界；厄里亚曾复活一个死人，但魔鬼并不因厄里亚的名而被驱逐。我们并非在说先知们的坏话，而是更崇高地颂扬他们的主；因为我们不是贬低他们的奇迹来抬高自己的奇迹；他们的奇迹也是我们的；但我们借着他们当中发生的事，来赢得人们对我们自己的奇迹的信任。\par

17. 但他们又说了：\Quote{一个刚死的尸体被活人复活了；但请向我们证明，一个死了三天的人可能复活，一个人被埋葬后三天能复活。} 如果我们为此类事实寻求圣经的证据，主耶稣基督亲自在福音中提供了，祂说：\BibleQuote{\Emph{因为约纳在鲸腹中三天三夜，人子也要在地心里三天三夜。}}\BibleRef{玛 12:40} 当我们审视约纳的故事时，相似之处极其有力。耶稣被派去宣讲悔改；约纳也被派去；但前者逃避，不知道将要发生的事；后者却甘愿前来，为赐予悔改以致救恩。约纳在船上睡觉，在汹涌的海浪中打鼾；而耶稣也睡了，海按照天主的眷顾开始翻腾，为要随后显示那睡着者的权能。对前者他们说：\BibleQuote{\Emph{你为什么打鼾？起来，呼求你的天主，或许天主会拯救我们。}}\BibleRef{纳 1:6} 而在后者的事上，他们对主说：\BibleQuote{\Emph{主啊，拯救我们。}}\BibleRef{玛 8:25} 前者说：\BibleQuote{\Emph{呼求你的天主；}}而这里他们说：\BibleQuote{\Emph{拯救我们。}}但前一人说：\BibleQuote{\Emph{把我拿起来，投进海里，海就会为你平静。}}\BibleRef{纳 1:12} 后一位自己\BibleQuote{\Emph{斥责了风和海，便风平浪静了。}}\BibleRef{玛 8:25} 前者被扔进鲸腹；但后者自愿下到那无形的死亡之鲸所在的地方。祂是自愿下去的，为叫死亡吐出它所吞没的人，正如经上记载：\BibleQuote{\Emph{我要救赎他们脱离阴间的权势，我要赎他们脱离死亡之手。}}\BibleRef{欧 13:14}\par

18. 在这一点上，让我们思考哪一件事更困难：一个人被埋葬后从地里复活，还是一个人在鲸鱼腹中，进入活物的炽热之中，竟能免于朽坏？因为哪一个人不知道，腹中的热量如此之大，连被吞下的骨头都会朽烂呢？那么，约纳在鲸鱼腹中三天三夜，如何能免于朽坏？而且，既然所有人的本性都是：我们无法像在空气中那样不呼吸而生存，他如何能在三天之内不吸入一点空气而活着？但犹太人回答说：天主的权能同约纳一起下降，当他在地狱中被抛来抛去时。那么，上主派遣祂的权能与祂的仆人同在，赐予他生命，难道不能同样赐予自己生命吗？如果那可信，这也可信；如果这不可信，那也不可信。因为对我来说，两者同样值得相信。我相信约纳得蒙保全，因为\BibleQuote{在天主，一切都是可能的}\BibleRef{玛 19:26}；我也相信基督从死者中复活了；因为我对此有许多见证，既来自神圣的圣经，也来自那位复活者至今仍在运作的德能——祂独自下到地狱，却带着一大群人从那里升上；因为祂下到死亡中，\BibleQuote{许多睡眠的圣人的躯体也复活了}\BibleRef{玛 27:52}，藉着祂。\par

19. 死亡看见新的访客下到阴间，不被那地方的锁链束缚，就惊惶失措。阴间的门卫啊，你们为何看见祂就害怕？是什么不寻常的恐惧占据了你们？死亡逃跑了，它的逃跑暴露了它的怯懦。圣先知们奔向祂，还有立法者梅瑟、亚巴郎、依撒格、雅各伯；达味、撒慕尔、依撒意亚，以及洗者若翰，他曾询问道：\BibleQuote{祢就是那位要来的，还是我们该期待别人呢？}\BibleRef{玛 11:3}所有的义人都被赎回了，他们曾被死亡吞没；因为被他们所宣告的君王，理应成为祂高贵先驱的救赎者。于是每个义人都说：\BibleQuote{死亡，你的胜利在哪里？坟墓，你的毒刺在哪里？}因为那征服者救赎了我们。\par

20. 论到我们的救主，先知约纳构成了预像，当他在鲸鱼腹中祈祷时说：\BibleQuote{我在患难中呼求}等等；\BibleQuote{从地狱的腹中}\BibleRef{纳 2:2}，然而他仍在鲸鱼中；虽然他在鲸鱼中，他却说自己在阴间；因为他是基督的预像，基督将要下到阴间。稍后，他以基督的身份极其清楚地预言说：\BibleQuote{我的头下到了山根}；然而他还在鲸鱼腹中。那么，什么山环绕着你呢？我知道，他说，我是祂的预像，祂将被安放在从岩石凿成的坟墓中。虽然他还在海中，约纳却说：\BibleQuote{我下到地里}，因为他是基督的预像，基督下到了地心。他预见犹太人说服士兵说谎，告诉他们：\Quote{你们就说，他们把祂偷去了}，他说：\BibleQuote{他们因看重浮虚的假神，就抛弃了自己的仁慈}\BibleRef{纳 2:8}。因为那怜悯他们的主来了，被钉十字架，又复活了，将自己的宝血为犹太人和外邦人倾流；但他们却说：\Quote{你们就说，他们把祂偷去了}，看重\Quote{浮虚的假神}。但关于祂的复活，依撒意亚也说：\BibleQuote{那从地上领回羊群的大牧者}；他加上了\Emph{大}字，免得祂被认为与以前的牧者同等。\par

21. 既然我们有了预言，就让信德常与我们同在。让那些因不信而跌倒的人跌倒吧，既然他们愿意；但你们已站立在复活信德的磐石上。决不要让任何异端者说服你们毁谤复活。因为直到今日，摩尼教徒还说，救主的复活是幻影式的，并非真实的，他们不理会保禄所说：\BibleQuote{天主在肉身上由达味的后裔所生}；又说：\BibleQuote{因耶稣基督我们的主从死者中的复活}。他又针对他们说：\BibleQuote{你不要心里说：谁要升到天上去？或者谁要下到深渊里去？——那就是把基督从死者中领上来}\BibleRef{罗 10:6}；同样，他在别处警告说：\BibleQuote{你要记得耶稣基督从死者中复活了}\BibleRef{弟后 2:8}；又说：\BibleQuote{如果基督没有复活，那么我们的宣讲便是空的，你们的信德也是空的。而且，我们也被发现是天主的假证人，因为我们曾为天主作证说，祂复活了基督，其实祂并没有复活祂}\BibleRef{格前 15:14}。但在后面他说：\BibleQuote{但是，基督确实从死者中复活了，是那些睡眠者初熟的果子}\BibleRef{格前 15:20}；\BibleQuote{——祂显现给刻法，然后显现给那十二位}；（因为你们若不相信一个证人，你们有十二个证人；）\BibleQuote{然后祂一次显现给五百多位弟兄}\BibleRef{格前 15:5}；（他们若不信十二位，就让他们接受那五百位；）\BibleQuote{此后，祂显现给雅各伯}，祂的兄弟，也是此教区的第一位主教。既然这样一位主教最初就看见了复活的基督耶稣，你，他的弟子，不要不信他。但你们说，祂的兄弟雅各伯是一个偏颇的证人；\BibleQuote{最后也显现给我}\BibleRef{格前 15:8}保禄，祂的仇敌；当一个仇敌宣告时，什么见证还会被怀疑呢？\Quote{\BibleQuote{我从前是迫害者}\BibleRef{弟前 1:13}，如今宣讲复活的喜讯。}\par

22. 救主复活的证人众多：黑夜和满月的光（因为那一夜是十六日）；接纳祂的坟墓岩石；那石头也要起来反对犹太人的脸，因为它看见了主；当时被滚开的石头，至今仍躺在那里，它本身也为复活作证。在场的天主天使为独生子的复活作证：\Person{伯多禄}、\Person{若望}和\Person{多默}和其余众宗徒。有的跑到坟墓，看见埋葬祂时包裹的殓布在复活后仍留在那里；有的摸祂的手脚，看见钉痕；大家同享救主的气息，并蒙恩因圣神的德能赦罪。妇女们也是证人，她们曾抱住祂的脚，看见大地震和站在旁边的天使的光辉；包裹祂的殓布，祂复活时留下的；士兵们和给他们的金钱；那地方本身至今仍可见——还有这神圣教会的殿宇，是可纪念的君士坦丁皇帝出于对基督的爱而建造并装饰得如你们所见。\par

23. 耶稣复活的证人还有\Person{塔彼达}，她因祂的名字从死者中复活（\BibleRef{宗 9:41}）。我们怎能不信基督复活了，既然连祂的名字都能使死人复活？大海也作证，正如你们先前听到的。一网鱼的神迹也作证，那里的炭火和上面放的鱼。\Person{伯多禄}也作证，他先前三次否认主，后又三次承认主，并被命令牧放主的灵羊。\Place{橄榄山}至今屹立，信德之眼几乎仍能看到祂驾云上升，以及祂升天的天门。因为祂从天上降临\Place{白冷}，却从\Place{橄榄山}升天：在开始人间争战之地，在争战后受冠冕之所。你们因此有许多证人：有这复活之地，也有东方的升天之地；还有在那里作证的天使，以及祂上升时的云彩和从那里下来的门徒。\par

24. 信仰的讲授课程本应也讲升天，但天主的恩宠如此安排，使你们昨日在主日，按我们软弱的限度，已最详尽地听到了；因为由于天主恩宠的眷顾，教会的读经安排包括了救主升天的记述。当时的话主要是为了众人和聚集的信友团体，但特别为了你们。问题是，你们是否留意了所讲的话？因为你们知道，信经接下来的话教训你们信\Quote{第三日复活，升天，坐在圣父之右}。我猜想你们当然记得那解释，但我现在再简要提醒你们当时所说的话。要记住《圣咏》清楚记载的：\Emph{天主上升，有欢呼声}；记住天军彼此说：\Emph{众城门哪，你们要抬起头来}，等等；还要记住《圣咏》说：\Emph{祂上升高处，掳掠了俘虏}；记住先知说：\Emph{祂在天空建造楼阁}；以及昨日因为犹太人的反驳而讲的其他细节。\par

25. 因为他们反对救主的升天，认为不可能，便应记住\Person{哈巴谷}被提的故事：若\Person{哈巴谷}被天使抓住头发提去，更何况先知和天使的主，岂不能凭自己的权能，从\Place{橄榄山}驾云升天呢？你们可追忆此类奇事，但要将超越性归给行奇事的主：因为其他人是被提举，祂却提举万有。记住\Person{哈诺客}被提升（\BibleRef{希 11:5}），但耶稣是升天；记住昨日关于\Person{厄里亚}的话：\Person{厄里亚}被火马车接去（\BibleRef{列下 2:11}），但\Emph{基督的战车成千累万}；\Person{厄里亚}是在\Place{约旦河}以东被接去，基督却在\Place{克德龙溪}以东升天；\Person{厄里亚}是\Quote{好像}升天，耶稣则真正升天；\Person{厄里亚}说双倍的圣神要赐予他的圣徒，基督却赐予祂的门徒如此丰富的圣神恩宠，不仅自己拥有，还能藉覆手使信者分享圣神的共融。\par

26. 当你们这样与犹太人争辩，并以平行的例子胜过了他们之后，再进一步谈到救主荣耀的优越之处：即他们是仆人，而他却是天主子。这样，你们将因想起基督的一位仆人曾被提到三重天，而记起他的优越地位。因为如果厄里亚只达到了第一重天，而保禄却达到了第三重天，那么后者便获得了更尊贵的地位。你们不要为你们的宗徒感到羞愧；他们并不逊于梅瑟，也不次于先知；他们是在高贵者中的高贵者，甚至更为高贵。因为厄里亚确实被接到天上；但伯多禄却拥有天国的钥匙，他接受了这话：\BibleQuote{\Emph{凡你在地上所释放的，在天上也要被释放。}}\BibleRef{玛 16:19} 厄里亚只被接到天上；但保禄却被接到\Emph{天上}和\Emph{乐园}（因为耶稣的门徒理当领受更丰盛的恩宠），并且\Emph{听到了不可言传的话，是人不能说的。}但保禄又从上面下来，并不是因为他配不上留在三重天，而是为了在享用了超越人所能达到的事物，并光荣地降下，传扬基督，并为他殉道之后，也接受殉道者的荣冠。但我略过这个论证的其他部分，这些是我昨天在主的日子的聚会上所讲的；因为对于有悟性的听者，仅需提醒就足以教导了。\par

27. 但也请记得我常说的关于圣子坐在圣父右边的事，因为信经的下一句说：\Quote{并升了天，坐在圣父的右边。}我们不要好奇地探究宝座的确切含义，因为它是不可理解的；但也不要容忍那些虚假的说法，即认为圣子是在十字架、复活、升天之后才开始坐在圣父的右边。因为圣子不是因提升而获得他的宝座；而是（他的存在是由永恒的生）他本来就和圣父一同坐着。先知依撒意亚在救主降生成人之前看见了这宝座，说：\BibleQuote{\Emph{我看见主坐在崇高的宝座上}}\BibleRef{依 6:1}，等等。因为\BibleQuote{\Emph{从来没有人见过天主}}\BibleRef{若 1:18}，而那时显现给先知的乃是圣子。圣咏作者也说：\Emph{你的宝座由太初就已预备，你从永远就存在。}虽然这方面的见证很多，但由于时间已晚，我们就以这些为满足。\par

28. 但现在我必须提醒你们关于圣子坐在圣父右边所说的许多事情中的几件。因为第109篇圣咏清楚地说：\BibleQuote{\Emph{主对我主说：你坐在我右边，直到我使你的仇敌作你的脚凳。}}救主在福音中证实这话，说达味不是凭自己说的，而是因圣神的感动，说：\BibleQuote{\Emph{那么，达味怎么在圣神内称他为主，说：主对我主说：你坐在我右边}}\BibleRef{玛 22:43}？等等。而且在\BibleBook{宗徒}{大事录}中，伯多禄在五旬节日与十一位宗徒站在一起\BibleRef{宗 2:34}，向以色列人讲论时，引用了来自第109篇圣咏的这段证言。\par

29. 我也必须提醒你们另外几处类似的证言，关于圣子坐在圣父右边。因为在\BibleBook{玛窦}{福音}中记载：\BibleQuote{\Emph{我却告诉你们：从今以后，你们要看见人子坐在大能者的右边}}\BibleRef{玛 26:64}，等等。与此一致，宗徒伯多禄也写道：\BibleQuote{\Emph{借着耶稣基督的复活，他已进入天堂，坐在天主的右边}}\BibleRef{伯前 3:22}。宗徒保禄致罗马人书说：\BibleQuote{\Emph{基督死了，但更复活了，现今坐在天主的右边}}\BibleRef{罗 8:34}。他又劝勉厄弗所人这样说：\BibleQuote{\Emph{按照他大能的力量，这力量在基督身上所运行的，当他从死人中复活他，叫他在自己右边坐下}}\BibleRef{弗 1:19}；等等。他教训哥罗森人这样说：\BibleQuote{\Emph{你们既然与基督一同复活了，就该追求天上的事，在那里基督坐在天主的右边}}\BibleRef{哥 3:1}。在\BibleBook{希伯来}{书}中他说：\BibleQuote{\Emph{当他涤除了我们的罪恶之后，便坐在高天至尊者的右边}}\BibleRef{希 1:3}。又说：\BibleQuote{\Emph{他曾对哪一位天使说过：你坐在我右边，等我使你的仇敌作你的脚凳？}}又说：\BibleQuote{\Emph{但他只奉献了一次祭品，为所有人永远坐在天主的右边，从此等待他的仇敌成为他的脚凳。}}又说：\BibleQuote{\Emph{仰望耶稣，他是我们信德的创始者和完成者，他为了那摆在他面前的喜乐，轻视了羞辱，忍受了十字架，如今坐在天主宝座的右边。}}\par

30. 虽然关于独生子坐在天主右边还有许多其他经文，但目前这些已足供我们使用；我再重复说，祂并不是在降生成人之后才获得这一座位的尊荣；不，因为早在万世之前，天主独生子、我们的主耶稣基督，就永远拥有圣父右边的宝座。但愿万有的天主，就是基督的圣父，以及我们的主耶稣基督——祂曾降下，又升上去，并与圣父同坐——护守你们的灵魂；使你们对复活之主的望德坚定不移、毫无改变；使你们与祂一同从死的罪恶中复活，达于祂天上的恩赐；使你们堪当被\BibleQuote{提到云中，到空中迎接主}\BibleRef{得前 4:17}，在祂适当的时机；并且，直到祂荣耀第二次来临的时刻到来，将你们所有人的名字写在生命册上，写了之后永不涂抹（因为许多跌倒的人的名字会被涂抹）；并恩赐你们众人相信那复活的主，并仰望那升天且将要再来的那一位——祂来，但不是从地上来；人啊，你要提防，因为将有欺骗者来到——祂坐在至高之处，却又与我们同在，\BibleQuote{察看各人的秩序和信德的坚定}\BibleRef{哥 2:5}。因为你们不要以为祂如今在肉身中缺席，就在圣神中也缺席。祂此刻就在我们中间，听人们论及祂的话，察看你们内心的思念，并\Emph{试探肺腑和心肠}——祂也如今准备在圣神内将那些前来领受圣洗的人，以及你们众人，呈献给圣父，并说：\Quote{看，我和天主所赐给我的孩子们}。——愿光荣归于祂，直到永远。阿们。\par

\SourceRecord{Translated by Edwin Hamilton Gifford.；Nicene and Post-Nicene Fathers, Second Series；Vol. 7.；Edited by Philip Schaff and Henry Wace.；Buffalo, NY: Christian Literature Publishing Co.,；1894.；<http://www.newadvent.org/fathers/310114.htm>.}

% structure: p0001 source=h1 target=section reason=page-title

\section{慕道讲授 第十五讲}

\Emph{关于\Quote{祂将光荣地降来，审判生者死者；祂的王国万世无疆}这一句。}\par

\BibleRef{达 7:9-14}\par

\Emph{\Quote{我观望，直到安置了宝座，有一位万古常存者坐上去}，然后，\Quote{我在夜间的神视中观望，看见一位相似人子者，乘着天上的云彩而来，等等。}}\par

1. 我们宣扬基督的来临不仅一次，而是第二次，且远比第一次更为光荣。因为前者显明了祂的忍耐，后者却带来了神圣王国的冠冕。因为在我们主耶稣基督内，大部分事物都是双重的：双重生——一次是在万世之前由天主所生，一次是在时期之末由童贞女所生；双重降临——一次是隐秘的，\Quote{像落在羊毛上的雨}；第二次则是公开的来临，尚待实现。在祂初次来临时，祂被裹在襁褓中，躺在马槽里；在第二次来临时，祂\Quote{以光为衣披戴}。在祂第一次来临时，\Quote{祂忍受了十字架，轻看羞辱}（\BibleRef{希 12:2}）；在第二次来临时，祂由天使的军旅伴随，接受荣耀。因此，我们不仅安于祂的第一次来临，也期待祂的第二次来临。正如在祂第一次来临时我们曾说：\Quote{奉主名而来的，当受赞美}，同样，在祂第二次来临时我们也要重复这话；好使我们与天使一同迎接我们的主，朝拜祂并说：\Quote{奉主名而来的，当受赞美}。救主来临，不是为了再受审判，而是要审判那些审判祂的人；祂在受审判时曾缄默不语，如今祂要提醒那些在十字架上做出狂妄之事的悖逆者，并要说：\Quote{你做了这些事，而我缄默不语}。当时祂因天主的救恩计划而来，用劝诫教导世人；但这一次，即使他们不愿意，也必然要认祂为君王。\par

2. 关于这两次来临，玛拉基亚先知说：\BibleQuote{你们所寻求的主，必忽然进入祂的殿}（\BibleRef{拉 3:1}）；这是一次来临。关于第二次来临，他又说：\BibleQuote{你们所喜悦的盟约使者，看，祂来了——万军的上主说。但祂来临之日，有谁能支持得住？祂显现时，有谁能站立得住？因为祂像炼金者之火，又像漂布者之碱。祂将坐下，如同炼净者和精炼者}。接着，救主亲自说：\BibleQuote{我必临近你们，施行审判；我要迅速作证，攻击那些行邪术的、犯奸淫的、起假誓的}（\BibleRef{拉 3:5}），以及其他等等。为此，保禄预先警告我们说：\BibleQuote{若有人在根基上建造金银、宝石、木草、禾秸，各人的工程必显明出来；因为那日子要将其揭示，它要在火中被揭露}（\BibleRef{格前 3:12}）。保禄也知道这两次来临，当他写信给弟铎时说：\BibleQuote{天主的恩宠已显现，给众人带来救恩，教导我们弃绝不虔敬和世俗的情欲，在今世度着清醒、正义和虔敬的生活，期待那幸福的希望，以及我们伟大的天主和救主耶稣基督光荣的显现}。你看，他提到了第一次来临，为此他感恩；也提到了第二次来临，那是我们所期盼的。因此，我们所宣认的信经刚才就这样传授了：我们相信那位升了天、坐在圣父之右的，将来要带着光荣降临，审判生者死者；祂的王国万世无疆。\par

3. 因此，我们的主耶稣基督将从天降来；在世界的终结、最后的日子，祂带着荣耀降临。因为这个世界将有终结，这受造的世界将被重新塑造。既然\BibleQuote{奸淫、盗窃、凶杀}和一切罪恶\BibleQuote{充满了大地，流血事件接踵而至}（\BibleRef{欧 4:2}），遍布世界，为使这奇妙的居所不至于永远充满不义，这世界将要过去，好让更美好的世界显现出来。你愿意从圣经的话语中接受这个证据吗？请听依撒意亚说：\BibleQuote{诸天要像书卷一样卷起；众星要像葡萄叶、无花果叶落下}（\BibleRef{依 34:4}）。福音也说：\BibleQuote{太阳将要昏暗，月亮不再发光，星辰要从天上坠落}（\BibleRef{玛 24:29}）。我们不要悲伤，好像只有我们会死；星辰也要死，但或许会复活。主卷起诸天，不是为了毁灭它们，而是为了使它们更美丽地复活。请听达味先知说：\BibleQuote{主，您起初奠定了大地的根基，诸天是您手的化工；它们将会毁灭，而您永存}。但有人会说：看，他明确说\Quote{它们将会毁灭}。请听他是在什么意义上说\Quote{它们将会毁灭}；从下文可知：\BibleQuote{它们都要像衣服一样陈旧；您要将它们卷起，如同披肩，它们就被更换}。因为一个人被称为\Quote{毁灭}，正如经上所载：\BibleQuote{看，义人怎样毁灭了，没有人放在心上}（\BibleRef{依 57:1}），虽然期待复活；同样，我们也期待诸天好像复活一样。\BibleQuote{太阳要变为黑暗，月亮要变成血红}（\BibleRef{岳 2:31}）。让那些从摩尼教归正的人在这里受到教导，不要再把那些光体当作他们的神，也不要不虔敬地认为这将要昏暗的太阳就是基督。再次听主说：\BibleQuote{天地要过去，我的话却绝不会过去}（\BibleRef{玛 24:35}）；因为受造之物不如师傅的话珍贵。\par

4. 那看得见的事物将要逝去，那所期待的更美好的事物将要来临；但至于时间，谁也不可好奇。因为祂说：\Quote{那日子和时刻，不是你们应该知道的，因为圣父已将它们归于自己的权下。} \BibleRef{宗 1:7} 你不可冒然宣告这些事何时发生，也不可因此懈怠沉睡。因为祂说：\Quote{你们要警醒，因为在你们料想不到的时辰，人子就来了。} 然而，由于我们有必要知道末世的征兆，并且我们正在期待基督，因此，为了我们不致受骗而死，被那假基督引入歧途，宗徒们因神圣的旨意感动，凭天主的安排向真导师发问说：\Quote{请告诉我们：什么时候要发生这些事？你又来临和世界终结的预兆是什么？} 我们期待你再次来临，但\Quote{撒但常把自己伪装成光明的天使}；求你因此提醒我们，使我们不致敬拜别人而非你。于是祂开启祂神圣而蒙福的口说：\Quote{你们要谨慎，免得有人迷惑你们。} 我的听众们，你们也当如同现在以心眼看见祂一样，听祂对你们说同样的话：\Quote{你们要谨慎，免得有人迷惑你们。} 这句话劝勉你们所有人留意所说的话；因为这并非过去的历史，而是对将来必然发生之事的预言。并非我们在说预言——我们是不配的——而是要将所写的事摆在你们面前，将征兆宣告出来。你们要观察，哪些已经应验，哪些尚待发生，从而保证自己的安全。\par

5. \Quote{你们要谨慎，免得有人迷惑你们：因为将来有许多人冒我的名来说：我是基督，并且要迷惑许多人。} 这事已部分应验：因为西蒙·马古斯已经这样说了，墨南德以及其他一些不虔敬的异端首领也说过；在我们这个时代或我们之后，还有人会这样说。\par

6. 第二个征兆。\Quote{你们要听见战争和战争的风声。} 现在难道不是波斯人和罗马人为了美索不达米亚而交战吗？或者不是？民族不是要起来攻击民族，国家攻击国家吗？\Quote{并且到处有饥荒、瘟疫和地震。} 这些事已经发生；此外，\Quote{从天上出现可怕的异兆和巨大的风暴。} 因此祂说：\Quote{所以你们要警醒，因为你们不知道你们的主什么时候来到。} \BibleRef{玛 24:42}\par

7. 但我们寻求祂来临的真正征兆；我们教会人士寻求教会独有的征兆。救主说：\Quote{那时，许多人要跌倒，彼此出卖，互相仇恨。} \BibleRef{玛 24:10} 如果你们听说主教反对主教，神职人员反对神职人员，平信徒反对平信徒，甚至流血，不要惊慌；因为这早已预先写下了。不要在意现在发生的事，而要在意所写的事；即使我——教导你们的这个人——丧亡了，你们也不至于与我一同丧亡；相反，听众甚至可能比导师更好，后来的可能成为最先的，因为连在第十一时辰来的人，主人也接纳。如果在宗徒中出了叛徒，你们岂能奇怪在主教中也有仇恨弟兄者呢？但这征兆不仅涉及统治者，也涉及百姓；因为祂说：\Quote{由于罪恶的增多，许多人的爱德就冷淡了。} \BibleRef{玛 24:12} 那么，在座的人中，有谁能夸耀自己对邻人怀有真诚的友谊呢？难道不是嘴唇常常亲吻，面容常常微笑，眼睛似乎闪烁，而心里却策划诡计，口中说着和平的话而心里图谋恶事吗？\par

8. 你们也有这个征兆：\Quote{这天国的福音必先在全世界宣讲，给万民作证，然后结局才来到。} \BibleRef{玛 24:14} 正如我们所看到的，如今几乎整个世界都充满了基督的教导。\par

9. 此后发生什么事？祂接着说：\BibleQuote{所以，几时你们见到达尼尔先知所说的\InnerQuote{招致荒凉的可憎之物}立在圣地，让诵读的领悟罢。} 又说：\BibleQuote{那时，若有人对你们说：看，基督在这里，或说：在那里！你们不要相信。} 兄弟间的憎恨接着为假基督腾出空间；因为魔鬼事先预备了百姓间的分裂，使那将临者被他们接纳。但愿基督的仆人，无论是此处或别处的，没有一个投奔仇敌！关于此事，宗徒保禄给出了一个明显的标记，说：\BibleQuote{因为那日子不会来临，除非背教先来到，并出现了罪恶的人、丧亡之子，他相反并高举自己超过一切称为神或受崇拜的，以致他坐在天主的殿中，显示自己就是天主。你们不记得我还与你们在一起时，给你们说过这些事吗？现在你们也知道那阻止他的，为的是使他在自己的时候被显露。因为那不法的奥秘已经运行，只等那现在阻止的从中间被除去。然后那无法的人将被显露，主耶稣要用自己口中的气息将他杀死，并以自己来临的显现消灭他。那无法无天的人来到，是依靠撒殚的工作，具有各种异能、征兆和欺骗的奇事，并以各种不义的诡计，陷害那些丧亡的人。} \BibleRef{得后 2:3} 保禄如此写道，如今正是背教。因为人们已从正确的信仰中背教了；有些宣讲子与父是同一的，另一些人竟敢说基督是从无中被造的。从前异端者是显明的；但现在教会充满了伪装的异端者。因为人们已背离真理，并且\BibleQuote{耳朵发痒}。\BibleRef{弟后 4:3} 是动听的言论吗？所有人都欣然听从。是矫正的话吗？所有人都回避。大多数人已偏离正言，宁愿选择恶，而不渴慕善。因此这就是\Emph{背教}，仇敌即将出现：同时他已开始部分派遣自己的先驱，以便他随后预备好时突袭猎物。因此，人哪，你要看顾自己，使你的灵魂稳妥。教会如今在永生天主前叮嘱你；她在假基督来临之前将关于假基督的事宣告给你。这些事是否发生在你的时代，我们不知道；或是发生在你之后，我们也不知道；但知道了这些事，你预先使自己稳妥，这是好的。\par

10. 真基督，天主的独生子，不再从地而来。若有人假作显现于旷野，不要出去；若他们说：\BibleQuote{看，基督在这里，或说：在那里！你们不要相信。} \BibleRef{玛 24:23} 不再向下观看大地；因为主从天降下；不像从前独自一人，而是与许多天使同行，由数以万计的天使簇拥；也不像羊毛上的露水那样隐秘，而是像闪电那样公开照耀。因为祂亲自说过：\BibleQuote{犹如闪电从东方发出，直照到西方，人子的来临也要这样。} \BibleRef{玛 24:27} 又说：\BibleQuote{他们要看见人子带着威能和大光荣乘云降来，祂要派遣祂的天使，吹起大号角。}等等。\par

11. 但是，正如从前当祂要摄取人性，人们期待天主由一位童贞玛利亚出生时，魔鬼对此制造了偏见，狡猾地在崇拜偶像者中预备了假神的传说，说他们生子并被女子所生，以便虚假先出现，真理（如他所认为）便不被相信；同样，如今既然真基督要第二次来临，仇敌利用头脑简单者的期望，特别是受割礼之人的期望，带来一个精通巫术和迷惑人的咒语的魔法师；他将为自己夺取罗马帝国的权力，并假称自己为基督；以此基督之名欺骗那些期待受傅者的犹太人，并以他的魔法幻象迷惑外邦人。\par

12. 但前述的假基督要在罗马帝国的时期届满、世界的终结临近时来临。将有十个罗马王一同兴起，可能在不同地区统治，但大致同时；之后有第十一个，即假基督，他凭自己的魔法技艺夺得罗马权力；在他之前统治的诸王中，\Emph{他要打倒三个}，其余七个他使臣服于自己。起初他确实会装出温和的样子（仿佛他是一位博学审慎的人），以及节制与仁慈；并以他魔法欺骗的虚假征兆和奇事迷惑犹太人，好像他就是所期待的基督，此后他将以各种不人道和无法无天的罪行显其特征，以致超过所有先前的不义和不敬虔之人；对所有人，尤其对我们基督徒，表现出一种嗜杀、极其残忍、无情和狡猾的灵。他只做了这些事三年零六个月，便要被天主的独生子、我们的主、救主耶稣，即真基督，从天上光荣的第二次来临所摧毁；祂要以\Emph{自己口中的气息}杀死假基督，并把他投入地狱之火。\par

13. 这些事并非我们自己杜撰，而是从教会使用的圣经中学来的，尤其从刚才宣读的达尼尔先知预言中得知的；正如总领天使加俾额尔所解释的，他说：\Quote{第四只兽是在地上将要兴起的第四个王国，它要超越所有王国。}这王国就是罗马帝国，这是教会解释者们的传统。因为第一个著名的王国是亚述，第二个是玛待和波斯，第三个是马其顿，那么第四个就是现在的罗马帝国。接着加俾额尔继续解释说：\Quote{它的十个角是指将要兴起的十个君王；在他们之后，将兴起另一个君王，他要比他以前的所有的君王更邪恶}\BibleRef{达 7:24}；他说，不仅比那十个，也比在他以前的所有君王；\Quote{他要制服三个君王}；显然是从那十个先前的君王中制服三个；但很明显，他制服了这十个中的三个后，将成为第八个君王；\Quote{他要向至高者说傲慢的话。}这人是亵渎者、无法无天者，不是从祖先承受了王国，而是藉巫术篡夺了权柄。\par

14. 这人是谁？又凭借何种能力？保禄，请为我们解释。他说：\Quote{他的来临，是依照撒殚的行动，具有各种德能、各种奇迹和欺骗的标记}；意思是，撒殚把他用作工具，亲自藉他行事；因为撒殚知道自己的审判不再有缓期，就不再像往常那样藉仆役作战，而是从此更公开地亲自作战。所谓\Quote{各种奇迹和欺骗的标记}，因为谎言之父将展示虚假的奇迹，使众人以为看见死人复活（其实并未复活）、跛子行走、盲人看见（其实未得治愈）。\par

15. 又说：\Quote{他反对且高举自己于一切称为神或受崇拜者之上}（\InnerQuote{反对一切神}；假基督自然憎恶偶像），\Quote{以致坐在天主的殿里。}\BibleRef{得后 2:4} 是哪座殿？他指的是已经毁灭的犹太人的圣殿。因为我们所在的殿决不是它！我们为什么说这话？免得有人认为我们偏袒自己。因为如果他以基督的身份来到犹太人中间，并希望被犹太人崇拜，他就会特别重视圣殿，以便更彻底地欺骗他们；使人以为他是达味后裔中那位将要重建所罗门所建圣殿的人。假基督将在犹太人的圣殿没有一块石头留在另一块石头上时来临，正如救主所预言的；因为当时间的腐蚀、或假借新建筑之名的拆毁、或任何其他原因，使所有石头都被摧毁时——我不仅指外围，也指内殿，即革鲁宾所在之处——那时他将带着\Quote{各种奇迹和欺骗的标记}来临，高举自己于一切偶像之上；起初他假装仁慈，但后来露出其无情的性情，尤其针对天主的圣人。他说：\Quote{我观看，见那角同圣徒作战}；又在别处说：\Quote{那时必有大灾难，自从有国以来直到此时，没有这样的。}这兽是可怕的，是一条巨大的龙，人无法战胜，它预备吞食；关于它，我们本可从圣经中更多讲述，但为了控制篇幅，目前满足于此。\par

16. 为此，主知道那仇敌的厉害，便准许虔诚人得宽容，说：\Quote{那时，在犹太的，应当逃到山上。}\BibleRef{玛 24:16} 但若有人自觉非常勇敢，能与撒殚对抗，就让他站立（因为我不对教会的神经绝望），并且说：\Quote{谁能使我们与基督的爱隔绝呢等等}\BibleRef{罗 8:35}。但我们中胆小的人要为自己的安全打算；而勇敢的人要站稳：\Quote{因为那时必有大灾难，是从世界开始直到如今从未有过的，将来也决不会再有。}\BibleRef{玛 24:21} 但感谢天主，他限制了那大灾难的时日；因为他说：\Quote{但为了那些被选的人，那些日子要缩短；}假基督只统治三年半。我们不是根据伪经说的，而是根据达尼尔；因他说：\Quote{他们将被交在他手里，直到一载、二载、半载。}\Quote{一载}是指他来临后暂时扩张的一年；\Quote{二载}是指剩下的两年不义，加起来共三年；\Quote{半载}是六个月。达尼尔在另一处又说同样的话：\Quote{他指着永生的者起誓说，直到一载、二载、半载。}\BibleRef{达 12:7} 有人或许也将下面的话引申于此，即\Quote{一千二百九十天；}以及\Quote{那些忍耐到底直到一千三百三十五天的，是有福的。}为此我们必须隐藏并逃跑；因为\Quote{你们还没有走完以色列的城邑，人子就来了。}\BibleRef{玛 10:23}\par

17. 那么，谁是有福的人，将在那时虔诚地为基督作证呢？因为我说，那时的殉道者超乎所有殉道者。迄今的殉道者只与人搏斗，但在假基督时期，他们将亲自与撒殚作战。先前的迫害君王只处死人，他们不假装复活死人，也不制造虚假的征兆和奇迹；但在他的时期，将有恐惧和欺骗的邪恶诱惑，\Emph{以致如果可能，连被选者也会被迷惑}。\BibleRef{玛 24:24} 那时活着的人心中绝不可起意问：\Quote{基督做了什么更多的事？此人凭何能力行这些事？若非天主的旨意，祂不会允许它们。}宗徒警告你，并预先说：\Emph{为此，天主给他们派遣一种错误的效力}（\Emph{派遣}，即\Emph{允许发生}），不是要他们推脱，而是\Emph{叫他们被定罪}。为何？他说：\Emph{因为他们不信真理}，即真基督，\Emph{反而喜爱不义}，即假基督。正如在时常发生的迫害中，那时天主也将允许这些事，并非因祂无力阻止，而是按祂的惯例，祂要藉忍耐加冕自己的勇士，如同祂加冕先知和宗徒一样；为使他们劳苦片刻后，继承永恒的天国，正如达尼尔所说：\Quote{那时你的百姓必得拯救，凡写在书上的（显然是生命册）都必得救；睡在尘埃中的许多人必醒起，有些得永生，有些受羞辱、永被憎恶；智慧者必发光如穹苍的光辉，众多义人必如星辰，直到永远。}\par

18. 所以，人啊，你要谨慎你自己；你有假基督的记号，不仅自己要记住，也要慷慨地传授给众人。如果你有肉身所生的孩子，现在就要劝诫他；如果你藉教理讲授生了孩子，也要提醒他，以免他误认假者为真者。因为\Emph{罪恶的奥秘已经在运作}。\BibleRef{得后 2:7} 我惧怕这些民族的战争，惧怕教会的分裂，惧怕弟兄间的彼此仇恨。但关于此事已足够；只愿天主不允许它在我们这个时代应验；然而，让我们保持警惕。关于假基督，就说这些。\par

19. 但让我们等候并期待主从天上云彩中的来临。那时将响起天使的号角；\Emph{在基督内死了的人要先复活} \BibleRef{得前 4:16}——活着虔诚的人将被提到云中，领受超越人性的荣誉作为他们劳苦的报酬，因为他们进行了超越人性的战斗；正如宗徒保禄所写：\Quote{因为主必亲自从天降下，有呼喊的声音、天使长的声音和天主的号角；在基督内死了的人要先复活。然后我们这些活着还存留的人，必和他们一同被提到云中，在空中与主相遇；这样我们就永远与主同在。}\par

20. 主的来临和世界的终结，训道者早已知道；他说：\Quote{青年人，你在青春时欢乐罢}等等；\Quote{因此，你要从心中除去烦恼，从肉体上驱除邪恶}……\Quote{要记着你的造主}……\Quote{在坏日子来到以前}……\Quote{在太阳、光明、月亮、星辰变暗以前}……\Quote{在窗户向外看的人昏暗以前}（意指视觉官能）；\Quote{在银链折断以前}（意指星辰的集合，因为它们的形状像银子）；\Quote{金花破碎以前}（以此暗示金色的太阳；因为甘菊是众所周知的，有许多放射状的叶子环绕它）；\Quote{他们因麻雀的声音而起立，并从高处观望，路上有恐怖}。他们将看见什么？\Quote{那时，他们要看见人子带着威能和大光荣乘云降来；各族都要哀号。}当主来临时将发生什么？\Quote{杏仁树要开花，蚱蜢要沉重，续随子要散开。}\BibleRef{训 12:5} 正如解释者所说，杏仁开花象征冬季的结束；那时我们的身体在冬季后将绽放天上的花朵。\Quote{蚱蜢要增大}（意指有翅膀的灵魂穿上身体），\Quote{续随子要散开}（即像荆棘一样的悖逆者要被驱散）。\par

21. 你看他们如何都预言了主的来临。你看他们如何知道\Emph{麻雀的声音}。让我们知道这是什么样的声音。\BibleQuote{因为主亲自要从天降下，有发令的声音，有总领天使的声音，并有天主的号声}。\BibleRef{得前 2:16} 总领天使要宣告并向众人说：\Quote{起来迎接主}。我们主人的降临将是可怕的。\Person{达味}说：\Emph{天主必显现，我们的天主必来临，决不缄默；在他面前有火燃烧，四周有强烈的暴风}等等。人子要来到圣父面前，正如刚才宣读的圣经所说，\Emph{乘着天上的云彩}，被\Emph{一条火河}所牵引，这火要试炼人。那时，若谁的工程是金的，他要变得更光亮；若谁的生活像碎秸、不坚实，就要被火烧尽。圣父\BibleQuote{要坐着}，\BibleQuote{他的衣服洁白如雪，他的头发如纯净的羊毛}。\BibleRef{达 7:9} 但这是按人的方式说的；为什么？因为他是那些没有被罪玷污之人的君王；\Emph{因为}，他说：\BibleQuote{我要使你的罪白如雪，白如羊毛}。\BibleRef{依 1:18} 这是罪恶赦免或无罪本身的象征。但那乘着云彩从天降下的主，正是乘着云彩升天的那一位；因为他亲自说过：\BibleQuote{他们要看见人子带着大威能和光荣，乘着天上的云彩来到}。\BibleRef{玛 24:30}\par

22. 但他来临的记号是什么呢？免得敌对的力量敢于假冒它。\BibleQuote{那时，}他说，\BibleQuote{人子的记号要出现在天上}。\BibleRef{玛 24:30} 基督自己的真记号是十字架；一个发光的十字架记号要走在君王前面，清楚宣告那曾被钉在十字架上的那一位：为使那些先前\Emph{刺透了他}并谋害他的犹太人，当他们看见时，可能\BibleQuote{各族都要哀悼}，\BibleRef{匝 12:12} 说：\Quote{这就是那被掌掴的，这就是那被吐唾沫在脸上的，这就是那被绳索捆绑的，这就是那昔日被钉在十字架上并受藐视的。他们将说：我们从你震怒的面容前逃到哪里去呢？}但天使的军旅要包围他们，使他们无处可逃。十字架的记号要对他的仇敌是恐怖，但对他的朋友——那些相信他、宣讲他或为他受苦的人——是喜乐。那么，那有福的人是谁，谁在那时要被发现是基督的朋友呢？那伟大而荣耀的君王，由天使卫队侍奉，与圣父共享宝座，不会轻视自己的仆人。为使他的选民不与他的仇敌混淆，\BibleQuote{他要派遣他的天使，用大号角，从四方把他的选民聚集起来}。\BibleRef{玛 24:31} 他没有轻视独一的\Person{罗特}；那么他怎么会轻视许多义人呢？\Quote{来吧，我父所祝福的！}他将对那些那时乘坐云彩战车、由天使聚集的人这样说。\par

23. 但有人会说：\Quote{我是穷人，}或者又说：\Quote{那时我或许卧病在床；}或者：\Quote{我只是一个女人，我要在磨坊被捉住：我们会被轻视吗？}鼓起勇气，人啊；审判者不偏待人；\Emph{他不凭外貌审判，也不凭言辞训斥}。他不看重学者过于朴实的人，也不看重富人过于贫穷的人。即使你在田野里，天使也会带走你；不要以为他会带走地主，而留下你农夫。即使你是奴隶，即使你贫穷，也不要丝毫忧愁；那\BibleQuote{取了奴仆的形体}的，\BibleRef{斐 2:7} 并不轻视仆人。即使你卧病在床，然而经上记载：\BibleQuote{那时，两个人同在一张床上；一个被带走，另一个被留下}。\BibleRef{路 17:34} 即使你被迫推磨，无论是男是女；即使你戴着脚镣，坐在磨旁，那\BibleQuote{以他的大能释放被囚者}的，\BibleRef{出 11:5} 必不会忽略你。那曾将\Person{若瑟}从为奴和监牢中领出，使他登基称王的，也要将你从苦难中赎回，进入天国。只管鼓起勇气，只管工作，只管努力奋斗；因为没有什么会失落。你的每一次祈祷，你咏唱的每一篇圣咏都被记录；每一次施舍，每一次禁食都被记录；每一次合宜的婚姻都被记录；为天主所持守的节欲都被记录；但在记录中首要的冠冕是童贞和纯洁的冠冕；你要如同天使般闪耀。但正如你乐意聆听了善事，也要同样毫不畏缩地聆听相反的事。你的每一件贪念都被记录；你的每一件邪淫都被记录，你的每一件虚誓都被记录，每一件亵渎、巫术、偷窃、凶杀都被记录。这些事从此都要被记录，如果你领洗后仍做同样的事；因为你先前的行为已被涂抹。\par

24. \Emph{当人子}，祂说，\Emph{在自己的光荣中，与众天使一同降来}。\BibleRef{玛 25:31} 请看，人啊，你将在多么庞大的群众面前受审判。那时所有人类的种族都要在场。因此，请计算罗马民族有多少人；计算现今活着的蛮族部落有多少，以及过去一百年内死去了多少人；计算在过去一千年中埋葬了多少民族；计算从亚当直到今日的所有人。数量确实庞大；但还算小，因为天使更多。他们是\Emph{那九十九只羊}，而人类只是那单独的一只。因为按照宇宙空间的范围，我们必须估计其居民的数目。整个地球在诸天之中不过是一个点，却已容纳如此众多的群众；那么环绕地球的天庭该容纳何等众多的群众呢？何况诸天之上的天，岂不更容纳不可想象的数目吗？经上记载：\BibleQuote{有千千万万事奉祂，亿亿众侍立在祂面前}\BibleRef{达 7:10}；这并不是说群众只有这么多，而是因为先知无法表达更多。所以在审判之日，将有所有人的天主圣父在场，耶稣基督与祂同坐，圣神也与他们同在；天使的号角将召唤我们所有的人，带着我们的事迹前来。我们难道不该从此刻起就感到极大的不安吗？人哪，不要以为只在不受到惩罚的情况下，在如此众多的人面前被定罪是轻微的处罚。我们难道不该宁愿死许多次，也不愿被朋友定罪吗？\par

25. 因此，弟兄们，让我们害怕，以免天主定我们的罪；祂定罪不需要审查或证据。不要说：夜间我犯了奸淫，或行了巫术，或做了任何其他事，却无人看见。你将被自己的良心审判，你的\BibleQuote{思想互相控告或辩护，在天主审判人隐秘事的那一天}\BibleRef{罗 2:15}。审判者那可怕的面容将迫使你说出真相；或者，即使你不开口，祂也定你的罪。因为你要披戴自己的罪恶或自己的义行复活。审判者自己已宣告了这事——因为审判者是基督——\BibleQuote{因为父不审判任何人，而是把审判权完全交给了子}\BibleRef{若 5:22}，并非放弃自己的权能，而是通过子审判；因此，子是按父的旨意审判；因为父和子的旨意没有不同，而是同一的。那么，审判者关于你是否要背负你的行为，说了什么呢？\BibleQuote{万民都要聚集在祂面前}\BibleRef{玛 25:32}（因为在基督面前，\BibleQuote{一切在天上的、地上的和地下的，无不屈膝叩拜}\BibleRef{斐 2:10}）；\Emph{祂要把他们彼此分开，如同牧人分开绵羊和山羊一样}。牧人如何分开？他是否查阅书册来分辨哪只是绵羊，哪只是山羊？还是根据明显的标记来区分？难道不是羊毛显示绵羊，多毛粗糙的皮肤显示山羊吗？同样，如果你刚被洗净了罪恶，你的行为从此就应如纯洁的羊毛；你的外衣应保持无污，并且你要常说：\Emph{我脱了外衣，怎能再穿上呢？} 你的外衣将表明你是绵羊。但如果你被发现是多毛的，如同厄撒乌，他浑身是毛，心术不正，为了一餐食物失去了长子权，出卖了特权，你将是左边的人。但愿在场的任何人都不被恩宠所弃，或因恶行而被列入左边罪人的行列！\par

26. 审判实在是可怕，所宣告的事也可怕。天国摆在我们面前，永火已经预备好了。那么，有人会说，我们如何逃避永火？又如何进入天国？祂说：\Emph{我饿了}，\Emph{你们给了我吃的}。由此可知途径；这里不需要比喻，而是要实行所说的话。\BibleQuote{我饿了，你们给了我吃的；我渴了，你们给了我喝的；我作客，你们收留了我；我赤身露体，你们给了我穿的；我患病，你们看顾了我；我在监里，你们来探望了我。}\BibleRef{玛 25:35} 你们若做了这些事，就要与祂一同为王；若不做，就要被定罪。因此立刻开始做这些事，并要坚持信德，以免像那些糊涂的童女，去买油时被关在门外。不要因你仅仅拿着灯就自信，而要常使它燃烧。\BibleQuote{你们的光也当在人前照耀}\BibleRef{玛 5:16}，不要使基督因你们而被亵渎。你要穿上不朽的外衣，在善行中发光；凡天主托付你管理的事，要妥善管理。你是否受托财富？要好好分配。你是否受托教导之言？要作好管家。你能感动听者的灵魂吗？要勤勉去做。善尽管家的门路很多。只愿我们中无人被定罪、被抛弃；使我们能坦然无惧地迎接永远为王的基督，祂统治直到永远。因为祂永远为王，祂要审判活人死人，因为祂为活人和死人死了。正如保禄所说：\BibleQuote{为此基督死了，又复活了，为作活人和死人的主。}\BibleRef{罗 14:9}\par

27. 若你曾听人说基督的王国将有终结，当憎恶此异端；这是龙的另一头，晚近在……出现。有一人竟敢断言，世界末日之后基督不再为王；他还敢说，自父出来的圣言将再被吸收归于父，不再存在；他说这样亵渎的话自取灭亡。因为他没有听从主说：\Quote{子永远存在。}（\BibleRef{若 8:25}）他也没有听从加俾额尔说：\Quote{祂要为王统治雅各伯家，直到永远，祂的王国没有终结。}（\BibleRef{路 1:33}）请思量这段经文。今日的异端者讲论基督时加以贬抑，而总领天使加俾额尔却教导救主永远常存；那么你愿意相信谁呢？难道你不更相信加俾额尔吗？请听达尼尔在经文中作证：\BibleQuote{我夜间在异象中观看，看见有一位相似人子者，乘着天上的云彩而来，来到万古常存者前……于是给祂赐予尊荣、统治权与王国：一切人民、种族、语言都要事奉祂；祂的统治权是永远的统治权，永不消逝，祂的王国不会灭亡。}（\BibleRef{达 7:13}）你要紧握这些事，相信这些事，并摒弃异端的言论；因为你已十分清楚地听说了基督无终的王国。\par

28. 对那\Emph{\BibleQuote{从山中凿出的非人手所凿的石头}}的解释中，你也能找到同样的道理，这石头是\Emph{按肉身说的基督}；并且\Emph{\BibleQuote{祂的王国不会归于另一个民族}}。达味也在\Emph{一处说：\BibleQuote{天主，你的宝座永远常存}}；又在另一处说：\Emph{\BibleQuote{主，你在起初奠定了大地根基，等等，它们都要消灭，你却常存，等等；你却永远同一，你的岁月没有穷尽}}：保禄将这些话解释为关于子的（\BibleRef{希 1:10}）。\par

29. 你想知道那些教导相反道理的人是如何陷入如此的疯狂吗？他们错误地读宗徒这句好话：\Emph{\BibleQuote{因为祂必须为王，直到把一切仇敌放在祂脚下}}（\BibleRef{格前 15:25}）；然后他们说，当祂的仇敌被放在祂脚下后，祂就不再为王，这是错误而愚蠢的断言。因为那位在制伏仇敌之前就已经是君王的主，在获得了对仇敌的统治之后，岂不更是君王吗？\par

30. 他们还竟敢说，经文\Emph{\BibleQuote{当万物都顺服于祂时，子自己也要顺服于那使万物顺服于祂的那位}}——这表明子也要被吸收归于父。那么，你这最不虔敬的人，你这基督的受造物，难道要继续存在，而基督——你和万物都是藉着祂造的——却要灭亡吗？这样的话是亵渎的。再者，万物如何服于祂？是通过消灭，还是通过常存？那么，其他事物在顺服于子时得以常存，而子在顺服于父时却不能常存吗？因为祂要顺服，并非因祂那时才开始行父的旨意（因为从永远祂就总是\Emph{\BibleQuote{常做祂所喜悦的事}}（\BibleRef{若 8:29}）），而是因为祂那时如同以前一样，服从父，不是勉强的服从，而是自愿的同意；因为祂不是仆人，不会被迫顺服，而是子，以自由的选择和本性的爱来顺从。\par

31. 但让我们审视他们；\Quote{直到}或\Quote{在……期间}是什么意思？因为我要用这词组本身来反驳他们，试图推翻他们的错误。既然他们敢说\Emph{\BibleQuote{直到祂把仇敌放在祂脚下}}这句话表明祂自己将有终结，并胆敢给基督永恒的王国设限，并以言语终止祂无尽的统治，那么，让我们读一读宗徒中类似的表述：\Emph{\BibleQuote{然而，从亚当直到梅瑟，死亡为王了。}}难道那些人只死到那时，而在梅瑟之后，或法律之后，人就不再死亡了吗？那么你看到，\Quote{直到}这个词并不是限制时间；而是保禄更愿意表明这一点——\Quote{然而，虽然梅瑟是一个正义而奇妙的人，但针对亚当所宣告的死亡判决，还是临到了他和他的后代；而且，尽管他们没有犯像亚当因吃树上的果子而违背命令那样的罪。}\par

32. 再举另一个类似的经文。\Emph{\BibleQuote{因为直到今日，每逢诵读梅瑟时，还有帕子盖在他们的心上}}（\BibleRef{格后 3:14}）。\Emph{直到今日}仅指\Quote{直到保禄}吗？它岂不是指\Emph{直到今日}的现在，甚至到终结？如果保禄对格林多人说：\Emph{因为我们宣讲基督的福音，一直来到你们那里，希望当你们的信德增长时，将福音传到你们以外的地区}，你清楚看到，\Emph{直到}并不表示终点，而是有后续。那么，你应如何记得那句经文\Emph{直到祂把一切仇敌放在祂脚下}（\BibleRef{格前 15:25}）？正如保禄在另一处说：\Emph{\BibleQuote{趁还有今日，你们要天天互相劝勉}}（\BibleRef{希 3:13}），意思是\Quote{继续不断}。因为正如我们不可谈论基督的\Quote{日子的开始}，同样你也不容许任何人谈论祂王国的终结。因为经上记载：\Emph{\BibleQuote{祂的王国是永远的王国。}}\par

33. 虽然我还有许多出自圣经的见证，论及基督的王国永无穷尽，但我暂且满足于上述引文，因为天色已晚。至于你，听众啊，要唯独敬拜祂为你的君王，并逃避一切异端邪说。如果天主的恩宠允许我们，信仰的其他信条也会适时向你宣告。愿全世界的天主保全你们众人，铭记末日的征兆，不被假基督征服。你们已领受了那将到来的欺骗者的标记；你们已领受了那真基督的明证，祂将公开从天降临。因此，要逃避那一个——那假者；并期待那另一个——那真者。你们已学会了如何在审判中被列入右边的方法；要保守\BibleQuote{\Emph{你所受托的}}\BibleRef{弟前 6:20}关于基督的事，并显扬于善工，使你能够满怀信心地站在审判者面前，继承天国。藉着祂，并与祂一起，愿光荣归于天主及圣神，至于无穷之世。阿们。\newline{}\par

\SourceRecord{Translated by Edwin Hamilton Gifford.；Nicene and Post-Nicene Fathers, Second Series；Vol. 7.；Edited by Philip Schaff and Henry Wace.；Buffalo, NY: Christian Literature Publishing Co.,；1894.；<http://www.newadvent.org/fathers/310115.htm>.}

% structure: p0001 source=h1 target=section reason=page-title

\section{\WorkTitle{教理讲授第十六讲}}

\Emph{论信经条文：\InnerQuote{我信唯一圣神，护慰者，曾借先知们发言。}}\par

\BibleRef{格前 12:1, 4}\par

\BibleQuote{弟兄们，论到神恩的事，我不愿意你们不明白。…… 神恩虽有区别，却是同一圣神赐予的。}\par

1. 我们实在需要属神的恩宠，才能谈论圣神；不是要说得恰到好处——这是不可能的——而是通过宣讲圣经的话，使我们安全地前行。因为在福音中记载了一件真正可怖的事，基督明明地说：\Quote{凡说话干犯圣神的，在今世及来世，总不得赦免。} \BibleRef{玛 12:32} 常有这样的恐惧，唯恐有人因无知或自以为是，说了不当说的话，而受此惩罚。审判活人死人的耶稣基督已宣告他不得赦免；因此，若有人冒犯，他还有什么希望呢？\par

2. 因此，必须属于耶稣基督的恩宠本身，既赐予我们无缺失地说话，也赐予你们有辨别地聆听；因为辨别力不仅为说话者需要，也为听者需要，免得他们听到一件事，心里却误解另一件。那么，让我们关于圣神只讲圣经所记的事；凡圣经未记的，我们不要自寻烦恼。圣神亲自说了圣经，关于祂自己，祂说了祂所愿意的，或我们所能够接受的。因此，让我们说祂所说的话；因为凡祂没有说的，我们不敢说。\par

3. 圣神只有一位，就是护慰者；正如只有一位天主圣父，没有第二位父；——也正如只有一位独生子、天主的圣言，祂没有兄弟——同样，只有一位圣神，没有第二位与祂同尊的神。圣神是一种至大能的力量，一个神圣而不可测度的存有；因为祂是活生生的、有理智的，是万物的圣化本原，万物都是天主通过基督造的。正是祂光照义人的灵魂；祂曾在先知们内，祂也在新约的宗徒们内。愿那些胆敢分离圣神行动的人受咒骂！有一位天主，父，新旧约的主；有一位主耶稣基督，旧约预言祂，新约中祂来了；有一位圣神，祂曾借先知们宣讲基督，当基督来到后，祂降临并显明了祂。\par

4. 因此，不要有人将旧约与新约分开；不要有人说，在旧约中的圣神是一位，在新约中的是另一位；因为这样他就是冒犯圣神本身，祂与圣父和圣子同受尊荣，并在圣洗时与他们一起被纳入圣三之中。因为天主的独生子明明对宗徒们说：\Quote{你们要去使万民成为门徒，因父、及子、及圣神之名给他们授洗。} 我们的希望在于父、子、圣神。我们不是宣讲三位神祇——愿马尔基翁派闭口——而是借着子同圣神，我们宣讲一位天主。信德是不可分的；敬拜是不可分的。我们不像某些人那样分离圣三，也不像撒伯里乌那样造成混淆。但我们虔诚地知道一位父，祂派遣了祂的子作我们的救主；我们知道一位子，祂应许要从父那里派遣护慰者；我们知道圣神，祂曾在先知们内发言，并在五旬节那天以火舌的形状降临在宗徒们身上——就在这里，在耶路撒冷，在宗徒们所在的\Place{楼厅}；因为在一切事上，最特选的恩惠与我们同在。在这里基督从天上降下；在这里圣神从天上降下。确实，最相宜的是：正如我们在此地的哥耳哥达谈论基督和哥耳哥达，同样我们也应在楼厅谈论圣神；然而，既然那降临那里的这位，与那在此被钉的这位共享荣耀，我们就在这里也谈论那降临那里的那位：因为对他们的敬拜是不可分的。\par

5. 我们现在要讲一点关于圣神的事；不是要精确地陈述祂的本体——因为这是不可能的——而是要提及一些人对祂的各种错误看法，免得我们因无知而陷入其中；并堵住错谬的道路，使我们能行走在君王的唯一大道上。如果我们现在为了谨慎起见，重复异端者的任何说法，愿它回到他们自己头上，但愿我们——说话者和听者——都无罪。\par

6. 因为极其亵渎的异端者也在一切事上磨尖了他们的舌头反对圣神，竟敢说出不敬之言；正如\Person{依勒内}这解释者在反对异端的著作中所写的。因为他们中有些人胆敢说他们自己就是圣神——其中第一个是\Person{西满}，那位在\WorkTitle{宗徒大事录}中提到过的术士；因为当他被驱逐出去后，他竟敢教导这样的学说。那些被称为诺斯替派的亵渎之人，说了别的反对圣神的话；邪恶的瓦伦提尼安派又说了别样的话；亵渎的\Person{摩尼}竟敢称自己为基督所派遣的\OriginalTerm{护慰者}{Paraclete}。另有些人教导说，在先知们中和在新约中的圣神是不同的。是的，他们的错误——毋宁说是他们的亵渎——是极大的。因此，你们要憎恶并躲避那些亵渎圣神的人，他们不得赦免。因为既然你们现今将要受洗归于圣神，你们与这些绝望之人有什么相通呢？如果一个人与窃贼为伍并与他同谋，就要受惩罚，那么冒犯圣神的人还能有什么希望呢？\par

7. 愿\OriginalTerm{马尔西翁派人}{Marcionists}也遭憎恶，他们从新约中撕去旧约的言语。因为\OriginalTerm{马尔西翁}{Marcion}首先——那最不虔敬的人，他首先主张三位神祇——知道新约中包含着先知们关于基督的证言，便删除了取自旧约的证言，使君王没有见证。憎恶上述的\OriginalTerm{诺斯替派人}{Gnostics}，名为有知识，却充满无知；他们竟敢论及圣神说些我不敢重复的话。\par

8. 愿\OriginalTerm{弗吕家派人}{Cataphrygians}也成你们的憎恶，以及他们的恶魁\OriginalTerm{蒙塔努斯}{Montanus}，和他的两个所谓女先知\OriginalTerm{马克西米拉}{Maximilla}与\OriginalTerm{普里斯奇拉}{Priscilla}。因为这位\OriginalTerm{蒙塔努斯}{Montanus}，神志不清、真正疯狂（若他未疯，便不会说这样的话），竟敢宣称他自己就是圣神——他，可怜的人，充满一切不洁和淫荡；仅暗示这一点就够了，出于对在场妇女的尊重。他占据了\Place{弗吕家}的一个很小村落\Place{佩普扎}，假称其为\Place{耶路撒冷}；并割断可怜幼童的喉咙，剁成不洁的食物，用于他们所谓的奥秘——（正因如此，直到不久前在迫害时期，我们曾被怀疑做这事，因为\OriginalTerm{蒙塔努斯派人}{Montanists}虽被假称为基督徒，却共用基督徒之名）——然而他竟敢称自己为圣神，充满一切不虔和残忍，已被不可撤销的判决定罪。\par

9. 并且如前所述，那最不虔敬的\OriginalTerm{摩尼}{Manes}也支持他，他汇合了每个异端中的邪恶；他是毁灭的最深渊，收集所有异端者的教义，炮制并传授更新奇的错误，竟敢宣称他自己就是基督所应许派来的护慰者。但救主在应许祂时，对宗徒们说：\BibleQuote{但你们要留在耶路撒冷城里，直到你们领受自高天而来的能力。}\BibleRef{路 24:49} 那么怎样？宗徒们已死了二百年，岂非要等待摩尼，\BibleQuote{直到他们领受能力}？有人敢说他们未立即充满圣神吗？而且经上记载：\BibleQuote{那时他们就按手在他们身上，他们就领受了圣神。}\BibleRef{宗 8:17} 这事岂非在摩尼之前，甚至许多年之前，当圣神在五旬节降临时吗？\par

10. 那么术士\OriginalTerm{西满}{Simon}为何被定罪？岂非因他来到宗徒们面前说：\BibleQuote{你们也把这权柄给我，叫我手按谁，谁就可以领受圣神。}？因为他没有说：\Quote{你们也给我圣神的共融，}而是说：\Quote{给我权柄，}好叫他能出售那不可出售、他自己也未拥有的东西。他还向一无所有的人献上银钱；尽管他看见人们把变卖之物价银带来，放在宗徒们脚前。他却不思想：那些践踏为供养穷人而带来的财富的人，是不大可能为贿赂而赐予圣神之权的。但他们对西满说了什么？\BibleQuote{你的银子和你一同丧亡！因为你以为天主的恩赐可用钱买得。}因为你成了又一个\OriginalTerm{犹达斯}{Judas}，竟期望用钱购买圣神的恩宠。那么，如果西满因想用代价获得此权而\Emph{丧亡}，摩尼自称是圣神，其不虔有多大？让我们憎恶那堪受憎恶的；让我们避开天主所避开的；让我们自己也满怀着胆量，对天主说，论及一切异端者：\BibleQuote{上主，憎恨祢的人，我怎不憎恨他们？对祢的敌人，我怎不厌恶？}因为有一种仇恨是正当的，正如经上所载：\BibleQuote{我要把仇恨放在你和她的后裔之间。}\BibleRef{创 3:15} 因为与蛇为友，成了与天主为敌，并导致死亡。\par

11. 关于那些被弃绝者，就此说这些已足够；现在让我们回到神圣的圣经，让我们\BibleQuote{喝我们自己水池的水}[即圣教父们]，并\BibleQuote{饮自我们自己的活井}。\BibleRef{箴 5:15} 让我们饮用\BibleQuote{活水，涌到永生}；但救主说这话是指那信祂的人将要领受的圣神。因为请留意祂说：\BibleQuote{信我的人}（不单如此，而是）\BibleQuote{就如经上所说}（这样祂把你引回旧约），\BibleQuote{从他的腹中要流出活水的江河}，不是感官所觉的江河，仅仅浇灌大地及其荆棘树木，而是将灵魂引向光明。在另一处祂说：\BibleQuote{我所赐的水，要在他里头成为泉源，直涌到永生}。——一种新的水，活的、涌流的，涌向那堪当的人。\par

12. 祂为何称圣神的恩宠为水？因为万物都靠水而存；因为水带来青草和生物；因为雨水从天而降；因为水以同一种形式降下，却以多种形式运作。同一个泉源浇灌整个乐园，同一场雨降在全世界，却在百合花中变成白色，在玫瑰中变成红色，在紫罗兰和风信子中变成紫色，在每一种植物中不同而多样：同样，在棕榈树中是一种，在葡萄树中是另一种，在万物中是一切；而自然界中却是同一的，并不自相差异；因为雨不改变自身，不以一种形式降下，然后以另一种形式，而是适应接受它的每一事物的构造，对每一事物变成适宜的。同样，圣神是唯一的，同一本性，不可分割，却按祂的意愿分赐各人的恩宠 \BibleRef{格前 12:11}；正如干枯的树获得水后便发出嫩芽，同样，灵魂在罪恶中，藉着悔改成为堪当领受圣神时，就结出正义的果实。祂在本性上虽是唯一的，却因天主的旨意和基督的名而成就许多德行。因为祂使用一个人的舌头为智慧；另一个人的灵魂祂以预言启迪；给另一个人驱魔的能力；再给另一个人解释圣经。祂加强一个人的自制；教导另一个人施舍的方式；另一个祂教导禁食和克己；另一个祂教导轻视肉身的事物；另一个祂训练为殉道：在众人中各不相同，却与自身无异，正如经上所载：\BibleQuote{圣神显示在每人身上虽不同，但全部是为人的好处。这人从圣神领受了智慧的言语，那人也由同一圣神领受了知识的言语；有人在同一圣神内领受了信德，另有人在同一圣神内领受了治病的奇恩；有的能行奇迹，有的能说预言，有的能辨别神恩，有的能说各种语言，有的能解释语言：可是，这一切都是这唯一而同一的圣神所行的，随祂的心愿，个别分配给人。}\par

13. 但是，由于论到\Quote{灵}的许多不同事物在圣经中被记载，恐怕有人因无知而陷入混乱，不知道经文所指的\Quote{灵}是哪一种；所以，现在最好使你们明白，圣经所宣告的圣神是哪一种。因为亚郎被称为基督，达味和撒乌耳及其他人都被称为基督（受傅者），但只有一位真基督；同样，既然\Quote{灵}这个名称被用于不同事物，就应该分辨什么是那独特地被称为圣神的。因为许多事物都称为灵。例如，被称为灵的有：天使、我们的灵魂、正在吹的风；伟大的德行也被称为灵；不洁的行为被称为灵；仇敌魔鬼被称为灵。因此，当你们听到这些时，要当心，免得因它们共享一个名字而混淆。关于我们的灵魂，圣经说：\BibleQuote{他的灵将离去，他归于尘土}；又论到同一灵魂说：\BibleQuote{祂铸成人灵在其中} \BibleRef{匝 12:1}。关于天使，圣咏说：\BibleQuote{祂使祂的天使成为灵，使祂的仆役成为火焰。}关于风，它说：\BibleQuote{你要以暴灵打破塔史士的船只}；又\BibleQuote{如同树林中的树被灵摇动} \BibleRef{依 7:2}；以及\BibleQuote{火、冰雹、雪、冰、风暴之灵}。关于良善的教导，主亲自说：\BibleQuote{我对你们所说的话，就是灵} \BibleRef{若 6:63}，\BibleQuote{就是生命}；这相当于说\Quote{是属灵的}。但圣神不是由舌头说出的；祂是一个活生生的灵，赐予言语的智慧，祂亲自说话和讲论。\par

14. 你们愿知道祂讲话和发言吗？斐理伯藉着天使的启示，下到往迦萨的路，那时太监正在来；圣神对斐理伯说：\BibleQuote{你上前去，靠近这辆车} \BibleRef{宗 8:29}。你们看见圣神对那听见祂的人说话吗？厄则克耳也这样说：\BibleQuote{上主的神临到我，对我说：\InnerQuote{上主这样说}} \BibleRef{则 11:5}。再者，\BibleQuote{圣神向安提约基雅的宗徒们说：\InnerQuote{你们为我选拔巴尔纳伯和扫禄，去作我召叫他们作的工}} \BibleRef{宗 13:2}。你看见这位圣神是活的、选拔、召唤、并且有权柄派遣吗？保禄也说：\BibleQuote{除了圣神在各城里向我指证，说锁链和患难等待我。}因为这位教会善良的圣化者、施助者和教师，圣神，护慰者，救主论到祂曾说：\BibleQuote{祂要教训你们一切}（祂不只说\InnerQuote{祂要教训}，还说\InnerQuote{祂要使你们想起我对你们所说的一切} \BibleRef{若 14:26}；因为基督的教训和圣神的教训不是不同的，而是同一的）——我说，祂预先向保禄见证将要临到他身上的事，好让他因预知而更勇敢。我现在向你们说了这些话，是因为经文\BibleQuote{我对你们所说的话，就是灵}；好使你们明白这并非指嘴唇的言语，而是指此处的良善教导。\par

15. 但罪也被称为神，正如我已说过；只是另一种相反的意义，如经上所说：\Emph{\OriginalTerm{邪淫之神}{spirit of whoredom}使他们误入歧途}。\BibleRef{欧 4:12}\Quote{神}这一名称也用来指\Emph{\OriginalTerm{不洁之神}{unclean spirit}}，即魔鬼；但加上了\Quote{不洁}；因为每个事物都附有其区别性名称，以标示其本性。如果圣经论及人的灵魂，就说\Emph{神}加上\Emph{人}；如果指风，就说\Emph{\OriginalTerm{暴风之神}{spirit of storm}}；如果指罪，就说\Emph{邪淫之神}；如果指魔鬼，就说\Emph{不洁之神}：好让我们知道所说的是哪一个，而你不至于以为它是指圣神；决非如此！因为神的名称是许多事物共用的；凡没有坚实身体的东西，一般都称为神。既然魔鬼没有这样的身体，它们便被称为神：但有很大的区别；因为那不洁的魔鬼，当他进入人的灵魂（愿主从一切听我言者及不在场者的灵魂中解救他们），他就像狼进入羊群，贪嗜鲜血，随时准备吞噬。他的来临极其凶猛；感觉极其压抑；心灵变得黑暗；他的攻击也是一种不义，侵占他人所有。因为他强行使用他人的身体和他人的工具，仿佛是他自己的；他推倒站立的人（因为他与那\Quote{从天上坠落的}相近\BibleRef{路 10:18}）；他扭曲舌头，歪斜嘴唇；泡沫代替言语；那人充满黑暗；眼睛睁着，但灵魂却看不见；那可怜的人在死亡边缘抽搐。魔鬼实在是人的仇敌，粗暴而无情地虐待他们。\par

16. 圣神却不是这样；决非如此！因为祂的行动倾向于相反的方面，即良善和有益的。首先，祂的来临是温柔的；对祂的感知是芬芳的；祂的担子极其轻省；光明和知识的光线在祂来临之前闪耀。祂以真正守护者的心肠而来：因为祂来是为了拯救、医治、教导、劝诫、加强、鼓励、启迪那接受祂者的心灵，首先是他自己，然后通过他也启迪他人。正如一个人先前处于黑暗之中，突然看见太阳，他身体的眼睛就被照亮，清楚看见以前看不见的东西；同样，那被赐予圣神的人，他的灵魂被照亮，看见超越人视界的事物，即他所不知道的；他的身体在地上，而他的灵魂却映照出天堂。他像依撒意亚一样看见\BibleRef{依 6:1} \BibleQuote{主坐在高高举起的宝座上}；他像厄则克耳一样看见\BibleRef{则 10:1} \BibleQuote{那在革鲁宾之上的那位}；他像达尼尔一样看见\BibleRef{达 7:10} \BibleQuote{千万的侍从，亿万的服务者}；而这如此渺小的人，却看见了世界的起源，知道世界的终结，以及中间的时代和君王的继承——那些他从未学过的事：因为真正的启迪者与他同在。那人虽在房屋的墙内，但他知识的能力却广及远方，甚至看见别人在做什么。\par

17. 伯多禄与阿纳尼雅和撒斐辣出售产业时并不在一起，但他藉圣神临在；\Emph{为什么}，他说，\Emph{撒殚充满了你的心，使你欺骗圣神} \BibleRef{宗 5:3}？没有控告者；没有证人；他怎知道所发生的事？\Emph{产业未变卖时，不是你的吗？既卖了，不是由你支配吗？你为什么心里起这意念？}那\Emph{不学无术}的伯多禄，藉圣神的恩宠，知道了连希腊智者都不知道的事。在厄里叟的例子中也有类似的事。因为他白白治愈了纳阿曼的癞病，而革哈齐却接受了报酬，即他人成就的报酬；他从纳阿曼那里拿了钱，藏在暗处。但\Emph{黑暗}瞒不过圣人。当他来时，厄里叟问他；就如伯多禄所说：\Emph{告诉我，你们卖田地得了这些钱吗？} \BibleRef{宗 5:8}；他也询问：\Emph{革哈齐，你从哪里来？} \BibleRef{列下 5:25} 我问你\Emph{从哪里来？}，不是出于无知，而是出于悲伤。你从黑暗中来，也要往黑暗里去；你出卖了癞病人的痊愈，癞病便成了你的遗产。他说，我遵行了那对我说\Emph{你们白白得来的，也要白白施舍} \BibleRef{玛 10:8} 的那位的命令；但你出卖了这恩宠，现在就接受这出卖的代价吧。但厄里叟对他说什么？\Emph{我的心岂没有与你同去吗？}我虽被身体关在这里，但天主赐给我的神却看见了遥远的事，清楚向我展示了别处所行的事。你看见圣神不仅消除无知，而且赋予知识吗？你看见祂如何启迪人的灵魂吗？\par

18. 依撒意亚活在近一千年前；他看见熙雍\Emph{如同茅屋}。那城仍然矗立，装饰着公共场所，披着威严；但他却说：\BibleQuote{熙雍必将像田地一样被耕种}，预言了在我们今日应验的事。注意预言的精确性；因为他曾说：\BibleQuote{熙雍的女子必像一个葡萄园中的茅屋，像黄瓜园中的草棚}。如今那地方到处是黄瓜园。你看见圣神如何启迪圣人吗？所以不要被共同术语的力量牵引到其他事物上去，而要持守确切的意义。\par

19. 當你坐在這裡時，若心中曾有關於貞潔或童貞的思想，那便是祂的教導。豈不是時常有已到婚嫁門檻的少女逃離了，因祂教導她童貞的教義？豈不是時常在宮廷中顯赫的人，蔑視財富與地位，因聖神的教導？豈不是時常有少年人，見美色而閉目，逃離眼目，避開玷污？你問這如何成就？聖神教導了那少年人的靈魂。世上有許多貪婪的方式，但基督徒拒絕財產：為什麼？因聖神的教導。那聖神確實是善而可敬的；我們合理地受洗歸於父、子及聖神。一個仍穿著肉身的人，與許多最兇猛的魔鬼搏鬥；魔鬼常被許多人用鐵鍊也無法制伏，卻被那人自己以祈禱的言語，藉著他內裡聖神的能力制伏了；驅魔者的氣息對那看不見的仇敵成了火。因此，我們從天主那裡得到了一位大能的盟友與保護者，一位教會的大導師，一位為我們作戰的大力士。讓我們不要怕魔鬼，也不要怕魔鬼；因為為我們作戰的那一位更強大。只要我們向祂敞開我們的門；\Emph{因為祂周遊尋覓配得的人}，尋找祂可賜予恩賜的人。\par

20. 祂被稱為護慰者，因為祂安慰並鼓勵我們，\Emph{且扶助我們的軟弱；因為我們不知道該如何祈禱才好，但聖神親自以無可言喻的歎息，為我們轉求}\BibleRef{罗 8:26}，即是向天主轉求。人常因基督的緣故，無理地受辱與被輕慢；殉道已在眼前；酷刑四圍，火、劍、猛獸與深坑。但聖神輕聲向他說：\Quote{\Emph{你要等候主}，人啊；現在臨到你的是小事，賞報是大的。稍忍耐片刻，你將永遠與天使同在。\Emph{現今的苦楚，與將來在我們身上要顯示的光榮，是不能較量的}\BibleRef{罗 8:18}。}祂為那人描繪天國；祂讓他瞥見樂園的福樂；殉道者們，面容不得已轉向審判官，但心神已在樂園中，便輕看這些可見的痛苦。\par

21. 你是否確信殉道者是藉著聖神的能力作見證？救主對祂的門徒說：\Emph{當人把你們帶到會堂、官長和有權柄者面前時，不要憂慮如何辯護，或說什麼話；因為在那時刻，聖神必要教給你們該說的話。}\BibleRef{路 12:11} 因為除非人藉著聖神，不可能為基督作殉道的見證；如果\Emph{除非藉著聖神，沒有人能說耶穌基督是主}，那麼，人怎能為耶穌的緣故捨棄自己的生命，除非藉著聖神？\par

22. 聖神實在偉大，在恩賜上全能，且奇妙。請想想，你們現在有多少人坐在這裡，我們有多少靈魂在場。祂合宜地為每人工作，並在我們中間，洞察各人的心情，也洞察各人的理性與良心，以及我們所說、所想、所信的。我剛才所說的確是大事，但仍是小事。因為我求你，用祂所啟發的心思想想，在這整個教區中有多少基督徒，在整個巴勒斯坦省中有多少，再從這省擴展到整個羅馬帝國；之後，思想整個世界：波斯人的種族，印度人的民族，加爾布人與薩爾馬提亞人，高盧人與西班牙人，摩爾人、利比亞人、埃塞俄比亞人，以及其他我們不知其名的人；因為許多民族，連名字也未傳到我們這裡。我求你想想，每個民族的主教、司鐸、執事、獨修士、貞女，以及平信徒；然後再注視他們的大保護者，恩賜的分施者——祂如何在全世界賜予一人貞潔，另一人永恆的童貞，另一人施捨，另一人自願貧窮，另一人驅逐敵對魔鬼的能力。正如光以其光輝一觸間照亮萬物，聖神也照樣啟示那些有眼目的人；因為若有人因盲目而不獲賜祂的恩寵，不要責怪聖神，而要責怪自己的不信。\par

23. 你已看見祂的權能，遍及全世界；不要再滯留於地上，而要上升高處。我說的上升，是指用想像升至第一層天，那裡有無數萬萬的天使。再用你的心思，若能，更往上升；想想總領天使，也想想諸靈；想想德能，想想宰制，想想異能，想想寶座，想想主權——這一切都是護慰者從天主而來的統治者、導師和聖化者。在人當中，厄里亞、厄里叟和依撒意亞需要祂；在天使當中，彌額爾和加俾額爾需要祂。沒有任何受造物與祂同等尊榮；因為天使的家族及其所有的軍旅集合起來，也與聖神不相等。護慰者至卓越的權能籠罩這一切。他們\Emph{奉派為服役的}\BibleRef{希 1:14}，但祂卻洞察天主的深奧事，正如宗徒所說：\Quote{聖神洞察一切，連天主的深奧事也洞悉。因為除了人內裡的心神，誰知道人的事？同樣，除了天主聖神，也沒有人知道天主的事}\BibleRef{格前 2:10}。\par

24. 祂藉先知们讲论基督；祂在宗徒们身上工作；祂至今仍在圣洗中印刻灵魂。的确，圣父给予圣子；圣子与圣神分享。因为正是耶稣自己，而非我，说：\BibleQuote{万物皆由我父交付于我}\BibleRef{玛 11:27}；关于圣神，祂说：\BibleQuote{当那真理之神来到时}，以及其余……\BibleQuote{祂要光荣我，因为祂要从我领受，并将启示给你们}\BibleRef{若 16:13}。圣父藉圣子，同圣神，是一切恩宠的施予者；圣父的恩赐无异于圣子的恩赐，亦无异于圣神的恩赐；因为只有一个救恩、一个德能、一个信德；一个天主圣父；一个主，祂的独生子；一个圣神，施慰者。我们知晓这些就足够了；但不要好奇探究祂的本性或实体：因为假若经上记载了，我们必会论及；凡未经记载的，我们切勿贸然妄言；我们得救所需的知悉，便是知道有圣父、圣子、圣神。\par

25. 这圣神在梅瑟时代降于七十位长老身上。（亲爱的，愿我们讲论的篇幅不要使你们疲倦；但愿我们讲论的主题本身赐予众人力量，不论是对我们宣讲者，还是对你们听者！）这圣神，如我所说，在梅瑟时代降于七十位长老身上；我对你们说这话，是为了证明祂知晓一切，并\BibleQuote{随祂的意愿行事}\BibleRef{格前 12:11}。七十位长老被选出来；\Emph{主在云中降下，取了梅瑟身上的圣神，将祂分给那七十位长老}；并非圣神被分割，而是祂的恩宠按器皿和接受者的容量分配。当时在场的六十八人说了预言；但埃尔达和摩达不在场：因此，为表明不是梅瑟赐予恩赐，而是圣神工作，被召唤却尚未到场的埃尔达和摩达也说了预言。\par

26. 农的儿子若苏厄惊讶，来到他跟前说：\Quote{你听说埃尔达和摩达在说预言吗？他们被召唤，却未前来；\BibleQuote{我的主梅瑟，请禁止他们}\BibleRef{户 11:28}}。\Quote{我不能禁止他们，}他说，\Quote{因为这恩宠来自天上；不，我非但不禁止他们，反而为此感恩。然而，我不认为你说这话是出于嫉妒；\Emph{你可是为我而嫉妒}？因为他们说预言，而你还未说预言？等待适合的时候吧；\BibleQuote{但愿主的全体百姓都是先知，当主将祂的圣神赐予他们时}\BibleRef{户 11:29}！}——难道亚巴郎、依撒格、雅各伯和若瑟没有得到吗？他们古人不也得到了吗？不，但这话：\Quote{\Emph{当主将赐予时}}显然意思是\Quote{赐予众人；如今恩宠只是局部的，那时将慷慨地赐予。}他暗暗预示了在我们中间五旬节那一天所发生的事；因为祂自己降临在我们中间。然而祂从前也降临在许多人身。因为经上记载：\BibleQuote{农的儿子若苏厄充满了智慧的圣神，因为梅瑟曾按手在他身上}\BibleRef{申 34:9}。\par

27. 祂也降临于所有义人和先知身上：即厄诺士、哈诺客、诺厄等等；亚巴郎、依撒格和雅各伯；至于若瑟，连法郎也发现他\BibleQuote{内有天主的圣神}\BibleRef{创 41:38}。论到梅瑟，以及圣神在他时代所行的奇事，你们已多次听闻：那至坚韧的约伯也拥有这圣神，以及所有圣人，尽管我们不一一列举他们的名字。在建造会幕时，祂也被遣发，充满智慧于那些与贝匝肋耳同处的智者心中。\par

28. 凭这圣神的德能，正如我们在\WorkTitle{民长传}中所见，敖特尼耳判事了\BibleRef{民 3:10}；基德红强盛了；依弗大获胜了；德波辣，一位女子，打了仗；三松，只要他行事正直，不使圣神担忧，便行了超乎人力的事。至于撒慕尔和达味，我们在\WorkTitle{列王纪}中清楚读到，他们怎样藉圣神自己说预言，并作先知的领袖——撒慕尔被称为\Emph{先见者}\BibleRef{撒上 9:9}；达味明确说：\BibleQuote{上主的圣神藉我说话}\BibleRef{撒下 23:2}，并在圣咏中说：\BibleQuote{不要从我夺走你的圣神}，又说：\BibleQuote{你的善神要在正直之地引导我}。我们在\WorkTitle{编年纪}中读到，阿匝黎雅\BibleRef{编下 15:1}在阿撒王时代，以及雅哈齐耳在约沙法特王时代，分受了圣神；还有另一位阿匝黎雅，那位被石头砸死的。厄斯德拉说：\BibleQuote{你也赐下你的善神来教训他们}。至于被接升天的厄里亚和厄里叟，这些受灵感而显奇迹的人，无需我们多说，显然他们充满了圣神。\par

29. 再者，若有人翻阅先知书全部，包括十二小先知及其他先知，必能找到许多关于圣神的证言。如米该亚以天主的口吻说：\Quote{我必藉着上主之神，满全能力}\BibleRef{米 3:8}；岳厄尔呼喊：\Quote{此后——上主说——我要将我的神倾注在一切有血肉的人身上}\BibleRef{岳 2:28}，以及其余的话；哈盖说：\Quote{因为是我与你们同在——万军的上主说}\BibleRef{盖 2:4}；\Quote{我的神仍留在你们中间}；同样，匝加利亚也说：\Quote{但你们应接纳我的言语和我藉我的神，对我的仆人先知们所命令的法令}\BibleRef{匝 1:6}；以及其他段落。\par

30. 依撒意亚也用他庄严的声音说：\Quote{上主的神将住在他身上，即智慧与聪敏的神，超见与刚毅的神，知识与孝爱的神；并以敬畏上主之神充满他}\BibleRef{依 11:2}；这说明圣神是惟一面不可分，但祂的化工却是多样的。同样又说：\Quote{我的仆人雅各伯……我已将我的神放在他身上。}又说：\Quote{我要将我的神倾注在你的苗裔上。}又说：\Quote{如今上主天主和祂的神派遣了我。}又说：\Quote{这是我与他们所立的盟约——上主说——我的神在你身上}\BibleRef{依 59:21}；又说：\Quote{上主的神临于我身上，因为祂给我傅了油}\BibleRef{依 61:1}，以及其余的话；又在祂对犹太人的责斥中说：\Quote{但他们叛逆，使祂的圣神忧郁。}以及\Quote{那将他们置于祂圣神中的在哪里呢？}此外，在厄则克耳书中（若你现在还未听倦）有已引用的话：\Quote{神便降到我身上，向我说：\InnerQuote{你说话：上主这样说。}}\BibleRef{则 11:5} 但\Quote{降到我身上}这话我们应理解为好的意思，即\Quote{慈爱地}；正如雅各伯找到若瑟时，伏在他颈上；同样，福音中那位慈爱的父亲看见从流浪中回来的儿子时，\Quote{动了怜悯的心，跑上前去，扑到他的脖子上，热情地亲吻他。}再者，在厄则克耳书中：\Quote{在神视中，天主打发我往加色丁地，到被掳的人那里。}\BibleRef{则 11:24} 其他经文你们也听过了，关于圣洗时所讲的：\Quote{那时我要在你们身上洒清水}等等；\Quote{我还要赐给你们一颗新心，在你们五内注入一种新精神}；然后紧接着：\Quote{我要把我的神置于你们五内。}又说：\Quote{上主的手临于我，上主的神引我出去}\BibleRef{则 37:1}。\par

31. 祂以智慧充满达尼尔的心，使他虽年轻却能审判长老们。贞洁的苏撒纳被诬为淫妇；无人为她辩护，谁能从权贵手中解救她？她被押赴刑场，已在刽子手手中。但她的救助者就在眼前，那就是护慰者、圣化每个理性本性的圣神。\Quote{到这里来，}祂对达尼尔说，\Quote{你虽年轻，却要指证那些染有青年罪愆的长老们。}因为经上记载：\Quote{天主兴起一位少年人，使圣神降在他身上。}尽管如此（我们迅速略过），因达尼尔的判决，那贞洁的妇人得救了。我们提出这事作为证言，因为现在不是讲解的时候。拿步高也知道达尼尔内有圣神，因为他对他说：\Quote{贝尔特沙匝，术士之长，我知道你身上有圣神。}\BibleRef{达 4:9} 他说对了一句，也错了一句。他说达尼尔有圣神是对的，但说他是\Quote{术士之长}却不对，因为他不是术士，而是藉圣神获得智慧。此前，他还为他解释了那梦像的异象，而梦者本人却不知道；因为他说：\Quote{告诉我这异象，我不知道。}你们看圣神的能力：那些看见异象的人不知道，而没看见的人却知道并加以解释。\par

32. 事实上，从旧约中收集极多经文并就圣神作更详尽的论述并不难。但时间有限，我们须顾及训诲辞的适当长度。因此，我暂且满足于旧约中的这些段落；若天主愿意，我们将在下一篇训诲辞中继续讨论新约中的其余经文。愿和平的天主，藉我们的主耶稣基督，并藉圣神的爱，使你们众人堪当祂属神和属天的恩赐：愿光荣与权能归于祂，至于无穷之世。阿们。\par

\SourceRecord{Translated by Edwin Hamilton Gifford.；Nicene and Post-Nicene Fathers, Second Series；Vol. 7.；Edited by Philip Schaff and Henry Wace.；Buffalo, NY: Christian Literature Publishing Co.,；1894.；<http://www.newadvent.org/fathers/310116.htm>.}

% structure: p0001 source=h1 target=section reason=page-title

\section{要理讲授第十七篇}

\Emph{论圣神的延续}\par

\BibleRef{格前 12:8}\par

\BibleQuote{这人蒙圣神赐他智慧之言}等。\par

1. 在前一篇要理讲授中，我们按自己的能力，为亲爱的听众们陈述了关于圣神的少许见证；如今，若蒙天主喜悦，我们将尽可能继续讨论新约中其余的部分。当时我们为不使你们注意力过度疲劳而克制了自己的热忱（因为谈论圣神永无餍足），如今我们同样只能从遗留的内容中讲说一小部分。因为如今和当初一样，我们坦诚承认，我们的软弱被所写之事的众多所压倒。今天我们也不使用人的诡辩，因为那无益处；而仅仅记起来自圣经的言语，因为这才是最稳妥的道路，正如蒙福的宗徒保禄所说：\BibleQuote{并且我们讲论这些事，不是用人的智慧所教的言语，乃是用圣神所教的言语，以属神的事解释属神的事。} \BibleRef{格前 2:13} 我们如同行路或航海的人，虽有一共同目标向极远的旅程，虽急切前行，却因人性软弱，惯于沿途在各城或各港口停靠。\par

2. 因此，虽然我们关于圣神的言论是分开的，但祂自己却是不可分的，同为一位。因为正如论及圣父时，我们时而教导祂是唯一的根源，时而教导祂被称为父或全能者，时而教导祂是万有的创造者；然而讲授的分割并未分割信仰，因为敬拜的对象祂曾是且仍是一。同样，论及天主的独生子时，我们时而教导祂的天主性，时而教导祂的人性，将关于吾主耶稣基督的道理分为许多讲论，却仍宣讲对祂不可分割的信仰。如今亦然，虽然关于圣神的讲授是分开的，但我们仍宣讲对祂不可分割的信仰。因为同一位圣神\BibleQuote{随祂的意愿，个别地分赐恩赐} \BibleRef{格前 12:11}，祂自己却保持不可分。因为护慰者并非不同于圣神，乃是同一位，有许多不同名称；祂生活、存在、说话、工作；祂是凡经由基督所造的一切理性受造物——天使与人——的圣化者。\par

3. 但恐怕有人因学识浅薄，从圣神的众多称号以为有多种不同的神，而非唯一同一位，因此天主教会预先保护你们，在信仰宣言中传授给你们，要你们\Quote{信唯一圣神，护慰者，曾藉先知发言}；好使你们知道，虽然祂的名称众多，圣神却只有一位。这些名称中，我们如今将列举众多中的少数几个。\par

4. 祂被称为\Quote{神}，如刚才读的圣经：\BibleQuote{这人蒙圣神赐他智慧之言。} \BibleRef{格前 12:8} 祂被称为\Quote{真理之神}，如救主所说：\BibleQuote{当那真理之神来到时。} \BibleRef{若 16:13} 祂又被称为\Quote{护慰者}，如祂说：\Quote{因为我若不去，护慰者就不会到你们这里来。} 但祂虽是同一位，却有不同称号，由以下经文清楚显明。因为圣神与护慰者是同一位，在以下话中声明：\BibleQuote{但那护慰者，就是圣神} \BibleRef{若 14:26}；而护慰者与真理之神是同一位，在以下话中声明：\Quote{我要另赐给你们一位护慰者，使祂永远与你们同在，就是真理之神}；又：\Quote{当护慰者，就是我从父那里要差给你们的，真理之神，来到时。} 祂被称为\Quote{天主之神}，如经上所记：\BibleQuote{我看见圣神仿佛鸽子从天降下} \BibleRef{若 1:32}；又：\BibleQuote{因为凡被天主的圣神引导的，都是天主的子女。} \BibleRef{罗 8:14} 祂又被称为\Quote{父之神}，如救主说：\BibleQuote{因为说话的不是你们，乃是你们父的圣神在你们内说话。} \BibleRef{玛 10:20} 保禄又说：\Quote{因此，我在父面前屈膝}，等等：\BibleQuote{求祂按祂荣耀的丰富，藉着祂的圣神叫你们得着能力。} \BibleRef{弗 3:14} 祂又被称为\Quote{主之神}，如伯多禄所说：\BibleQuote{你们为什么同心试探主的圣神呢？} \BibleRef{宗 5:9} 祂又被称为\Quote{天主和基督之神}，如保禄所写：\BibleQuote{你们不属于血肉，乃属于圣神，只要天主的圣神住在你们内。人若没有基督的圣神，就不属于基督。} \BibleRef{罗 8:9} 祂又被称为\Quote{天主子之神}，如经上所说：\BibleQuote{因为你们是儿子，天主就派遣了祂儿子的圣神。} \BibleRef{迦 4:6} 祂又被称为\Quote{基督之神}，如经上所写：\Quote{考证在他们内的基督之圣神，是指明什么时候或怎样的时候}；又：\BibleQuote{藉着你们的祈祷，和耶稣基督的圣神的供应。} \BibleRef{斐 1:19}\par

5. 你还会发现圣神的许多其他称号。例如，祂被称为圣德之神，正如经上所载：\BibleQuote{按圣德之神。}\BibleRef{罗 1:4} 祂也被称为义子之神，如保禄说：\BibleQuote{因为你们没有领受奴役的灵，再度恐惧，但你们领受了义子的灵，使我们呼号：阿爸，父啊！}祂也被称为启示之神，正如经上所载：\BibleQuote{赐给你们智慧和启示的神，使你们认识祂。}\BibleRef{弗 1:17} 祂也被称为恩许之神，如同保禄所说：\BibleQuote{在祂内，你们也因相信，受了所恩许的圣神的印证。}祂也被称为恩宠之神，如同他再次说：\BibleQuote{并侮辱了恩宠之神。}\BibleRef{希 10:29} 还有许多类似的称号。在先前的讲道中，你们清楚听见，在\BibleBook{圣}{咏}中祂有时被称为\Quote{善神}，有时被称为\Quote{尊威之神}；在依撒意亚书中，祂被称为\BibleQuote{智慧、聪敏、超见、刚毅、明达、孝爱和敬畏天主之神}\BibleRef{依 11:2}。所有这些经文，无论先前所引还是现在所引，都证实：尽管圣神的称号不同，但祂是同一且不变的；祂是生活的、自存的，常与圣父和圣子同在；并非出自圣父或圣子的口唇和气息，也不是消散于空气中，而是有实在的本体，亲自说话、运作、分施和圣化；正如那来自圣父、圣子和圣神、赐予我们的救恩工程，是不可分割、和谐一致的，如同我们先前所说。因为我希望你们牢记最近所说的话，并清楚知道：法律和先知中的圣神，与福音和宗徒中的圣神，并非两个不同的神；而是同一、自同的圣神，在旧约和新约中都发言了圣经。\par

6. 这位圣神降临于童贞玛利亚身上；因为所怀的是独生子基督，\BibleQuote{至高者的能力要荫庇她}，并且\BibleQuote{圣神要降到她身上}\BibleRef{路 1:35}，圣化了她，使她能够接纳祂，\BibleQuote{万物是由祂而造成的}\BibleRef{若 1:3}。但我无需多言来教导你们，那次怀孕是无玷无污的，因为你们已学过。是加俾额尔对她说：\Quote{我是所成之事的传报者，但这工程我一无所份。我虽为总领天使，但也知道自己的身份；虽然我欢欣地向你请安，但你的生产方式并非出于我的任何恩惠。}\BibleQuote{圣神要临到你身上，至高者的能力要荫庇你，因此那要诞生的圣者，将称为天主子。}\BibleRef{路 1:35}\par

7. 这位圣神在依撒伯尔身上工作；因为祂不单眷顾贞女，也眷顾已婚妇女，只要她们的婚姻是合法的。\BibleQuote{依撒伯尔充满了圣神}\BibleRef{路 1:41}，并且说预言；那高贵的婢女论及自己的主说：\Quote{吾主的母亲驾临我这里，这是从哪里来的呢？}因为依撒伯尔自认有福。若翰的父亲匝加利亚也充满了这圣神而说预言，讲述独生子将带来多少恩惠，以及若翰要籍着洗礼作祂的先驱。也是因这圣神，正义的西默盎得到启示：\BibleQuote{他未见主基督以前，决见不到死亡}\BibleRef{路 2:26}；他亲手接抱祂，并在圣殿中为祂作了明确的见证。\par

8. 若翰从母腹中就充满了圣神，因此被圣化，得以给主施洗；他不是自己赐予圣神，而是宣讲那赐予圣神者的喜讯。因为他说：\BibleQuote{我固然以水洗你们，为叫你们悔改；那在我以后来的，等等；祂要以圣神及火洗你们。}\BibleRef{玛 3:11} 为什么用火？因为圣神的降临是在火舌之中；关于此事，主喜悦地说：\BibleQuote{我来是为把火投在地上，我是多么切望它已经燃烧起来！}\BibleRef{路 12:49}\par

9. 当主受洗时，这位圣神降下，为使受洗者的尊贵不致隐藏；正如若翰所说：\BibleQuote{那派遣我来以水施洗的，曾对我说：\InnerQuote{你看见圣神降下，住在谁身上，谁就是那要以圣神施洗的。}}\BibleRef{若 1:33} 但请看福音所说：\Quote{天开了}；天开了，因为降临者的尊贵；因为他说：\BibleQuote{看，天开了，他看见天主的圣神，有如鸽子降下，来到祂上面}\BibleRef{玛 3:16}；意思是祂下降是出于自愿的运动。因为有人解释，圣神应许给受洗者的首位和初果，应当赋予救主的人性，祂是这恩宠的赐予者。也许祂以鸽子的形状降下，如有些人所说，是为了显示那鸽子纯洁、无玷、无瑕的形象，并帮助她所生儿女的祈祷和罪过的赦免；正如以象征方式预示了基督在祂眼睛的外表上要这样显现；因为在\BibleBook{雅}{歌}中，新娘论及新郎说：\BibleQuote{你的两眼像河边的鸽子。}\par

10. 关于这只鸽子，按某些人的说法，\Person{诺厄}的鸽子在某种程度上是一个预像。因为正如在\Person{诺厄}时代，藉着木头和水，他们得到了救恩，并开始了新的世代，鸽子在傍晚衔着一根橄榄枝回到他那里；同样，他们说，圣神也降临在真正的\Person{诺厄}身上，即第二次诞生的创始者，祂将万民的心意聚拢为一，这预表在方舟中各种动物的性情：——在祂来临时，属神的豺狼与羔羊同食，在祂的教会中，牛犊、狮子和牛在同一牧场上吃草，正如我们今天看见世界的统治者由教会人士引导和教导。因此，按某些人的解释，这属神的鸽子在祂受洗的时刻降下，为要表明：正是祂藉着十字架的木头拯救相信的人，祂要在傍晚时分藉着祂的死亡赐予救恩。\par

11. 这些事或许应以别的方式解释；但现在我们仍要听救主自己关于圣神的话。因为祂说：\BibleQuote{人除非由水和圣神而生，不能进天主的国。}\BibleRef{若3:5} 而这恩宠来自圣父，祂如此说：\BibleQuote{何况你们的天父，岂不更将圣神赐与求祂的人吗？}\BibleRef{路11:13} 我们应当以圣神崇拜天主，祂这样表明：\BibleQuote{然而时候到来，现在就是，那些真正朝拜的人，将以心神和真理朝拜父；因为父也寻找这样朝拜祂的人。天主是神；朝拜祂的人应当以心神和真理朝拜祂。}\BibleRef{若4:23} 又说：\BibleQuote{如果我仗赖天主的圣神驱魔，}\BibleRef{玛12:28} 随后立即说：\BibleQuote{因此我告诉你们：一切罪过和亵渎，人都可得赦免；但是亵渎圣神的罪，必不得赦免。凡说话干犯人子的，可得赦免；但说话干犯圣神的，在今世和来世都不得赦免。}\BibleRef{玛12:31-32} 祂又说：\BibleQuote{我将祈求父，他必会赐给你们另一位护慰者，使他永远与你们同在；他是真理之神，世界不能领受他，因为世界看不见他，也不认识他；你们却认识他，因为他与你们同在，并在你们内。}\BibleRef{若14:16} 祂又说：\BibleQuote{我还与你们同在的时候，给你们讲了这些事。但那护慰者，就是圣神，父因我的名将要派遣的，他必要教训你们一切，也要使你们想起我所对你们说的一切。}\BibleRef{若14:26} 又说：\BibleQuote{当护慰者来到时，即我从父那里要派遣给你们的，那发自父的真理之神，他必要为我作证。}\BibleRef{若15:26} 救主又说：\BibleQuote{我若不去，护慰者便不会到你们这里来。……他来到时，就要指证世界关于罪、关于义、关于审判的事。}\BibleRef{若16:7-8} 随后又说：\BibleQuote{我还有许多事要告诉你们，但你们现在不能负担。当那真理之神来到时，他要把你们引入一切真理；因为他不凭自己讲论，只把他所听到的讲出来，并把未来的事传告给你们。他必要光荣我，因为他要从我那里领受而传告给你们。凡父所有的一切，都是我的；为此我说：他要从我那里领受而传告给你们。}\BibleRef{若16:12-15} 我现在已向你们宣读了独生子自己的话，使你们不要听从人的言语。\par

12. 祂将这圣神的共融赐给了宗徒们；因为经上记载：\BibleQuote{说了这话，就向他们嘘了一口气，说：你们领受圣神吧！你们赦免谁的罪，谁的罪就赦免了；你们留存谁的罪，谁的罪就存留了。}\BibleRef{若20:22} 这是祂第二次向人嘘气（第一次的嘘气因故意犯罪而被窒息了）；为应验经上的话：\Quote{祂上升，吹拂你的面，将你从苦难中解救出来。} 但祂从何处上升？从阴间；因为福音如此叙述，那时祂复活后向他们嘘气。然而，祂虽在那时赐下了恩宠，却还要更丰盛地倾注；祂对他们说：我现在就愿意赐下，但器皿尚不能容纳；你们暂且领受你们能承受的恩宠，并期待更多；\BibleQuote{但你们要留在耶路撒冷城中，直到你们领受来自高天的能力。}\BibleRef{路24:39} 现在领受一部分；然后，你们将穿戴上圆满的恩赐。因为领受的人，往往只拥有部分的恩赐；但穿戴上的人，则被袍服完全包裹。祂说：\Quote{不要害怕魔鬼的武器和飞镖；因为你们将携带圣神的能力。}但要记住刚才的话：圣神不是分开的，只有由祂赐予的恩宠是分开的。\par

13. 耶稣于是升了天，并履行了祂的许诺。因为祂对他们说：\BibleQuote{我将祈求父，他必会赐给你们另一位护慰者。}\BibleRef{若14:16} 所以他们坐着，等待圣神的降临；\BibleQuote{五旬节日一到，}\BibleRef{宗2:1} 在这耶路撒冷城中——（因为这份尊荣也属于我们；我们所说的不是发生在别人身上的好事，而是我们自己所蒙受的恩惠）——在五旬节日，他们坐着，而护慰者从天上降下，他是教会的守护者和圣化者，灵魂的治理者，风浪中的舵手，带领迷失者走向光明，主持战斗，并为胜利者加冕。\par

14. 但祂降临，为赋予宗徒们权能，并为他们施洗；因为主说：\BibleQuote{你们不多几日以后，要受圣神的洗}\BibleRef{宗 1:5}。这份恩宠不是局部的，而是祂的权能圆满全备；因为正如那跳入水中受洗的人，被水从四面环绕，同样，他们也被圣神完全浸洗。然而，水只从外面流经，而圣神也内在地浸润灵魂，且是完全的。你为何感到惊奇？取一个来自物质界的例子，虽卑微而常见，却有益于较简单的人。如果火穿过铁块，使整块铁成为火，以致冷的变为灼热，黑的变为明亮——如果这有形的火能如此无碍地渗透并作用在有形的铁上，为何惊奇圣神能进入灵魂的最深处呢？\par

15. 为免人们对于降临到他们身上的伟大恩赐之浩大一无所知，遂有如天上号角响起，因为\BibleQuote{忽然从天上来了一阵响声，好像暴风刮来}\BibleRef{宗 2:2}，预示那位将赐予人能力用暴力夺得天主之国者的临在；好使他们的眼睛看见火舌，耳朵听见响声。\BibleQuote{它就充满了他们坐着的整座房屋}；因为房屋成了属神之水的容器；当门徒坐在里面时，整座房屋都被充满。这样，他们按照应许完全受洗，并以一种神圣的救恩之衣裹覆了灵魂和身体。\BibleQuote{又有舌头如火焰显现，分开落在他们每人身上，他们就都充满了圣神}。他们所领受的火，不是焚烧的火，而是拯救的火；是烧尽罪恶荆棘、却使灵魂发光的火。这火如今也要降临在你们身上，为要除去并烧尽你们如荆棘般的罪恶，并在你们灵魂的宝贵产业上增添光彩，且赐给你们恩宠；因为祂当时就赐给了宗徒们。祂以火舌的形态落在他们身上，好使他们能以头上的火舌为自己加冕，戴上新的属神冠冕。从前，一把火剑把守了乐园的门；如今，一把带来救恩的火舌重新恢复了这份恩赐。\par

16. \BibleQuote{他们就都充满了圣神，照圣神赐给他们的话，说起外方话来}。加里肋亚人伯多禄或安德肋说波斯语或玛待人。若望和其他宗徒对外邦人后裔说各种语言；因为并非在我们时代，才有众多外方人初次从各地聚集于此，而是自那时起便如此。有哪位老师能如此伟大，一下子教人们从未学过的东西？他们花费多年，借助语法和其他技艺，才只学好希腊语；而且并非所有人都说得同样好：修辞学家或许说得不错，语法学家有时却说得不好，而高明的语法学家却对哲学题目一无所知。但圣神一下子教会他们许多语言，这些语言是他们一生从未学过的。这实在是无上的智慧，这是神圣的大能。他们过去长久的无知，与如今突然、完全、多样且不习惯地运用这些语言之间，有多么大的对比啊！\par

17. 听众多感困惑——这是第二次混乱，取代了巴比伦首次邪恶的混乱。因为在言语混乱中，目的分裂，因为他们的意念与天主为敌；但在这里，心灵得以恢复并合一，因为关注的对象是虔敬的。堕落的手段也是恢复的手段。因此他们惊奇地说：\BibleRef{宗 2:8}，\BibleQuote{我们怎么听见他们说话呢？}你若不知，不足为奇；因为连尼苛德摩也对圣神的来临无知，他曾被告知：\BibleQuote{风随意向哪里吹，你听到风的响声，却不知道它从哪里来，往哪里去}；但若我听到祂的声音，却不知祂从哪里来，我又怎能解释祂本身本质是什么呢？\par

18. \BibleQuote{但另有些人却讥笑说：他们喝醉了新酒}\BibleRef{宗 2:13}，他们虽在讥笑，却说了真话。因为那酒确实是新的，即新约的恩宠；但这新酒出自属神的葡萄树，这葡萄树以前屡次在先知身上结果，并在新约中萌芽。正如在可感的事物中，葡萄树永远相同，但在各季节结出新果；同样，那永远不变的同一圣神，如同祂常在先知身上工作，如今显出了新而奇妙的工作。因为虽然祂的恩宠以前也临于祖先们，但这里却丰沛地涌流；因为以前人们只分沾圣神，如今却完全受洗。\par

19. 然而，拥有圣神且知道自己所拥有的\Person{伯多禄}说：\Emph{以色列人哪}，你们宣讲岳厄尔，却不知道所记载的事：\Emph{这些人并非如你们所想的那样喝醉了}。他们确是醉了，却不像你们所想的那样，而是正如经上所记：\Emph{他们必因你殿中的肥甘而醉；你必使他们畅饮你喜乐的河川}。他们醉了，是一种清醒的醉，杀死罪恶，赋予心灵生命，与肉身的醉相反；后者使人忘掉已知之事，前者却使人知道未知之事。他们醉了，因为他们饮了属灵葡萄树的酒，这葡萄树说：\Emph{我是葡萄树，你们是枝条}。\BibleRef{若 15:5}但若你们不信我的话，从一天的时刻就可以明白，因为\Emph{现在是白天第三时辰}。因为按\Person{马尔谷}所载，祂在第三时辰被钉十字架，如今在第三时辰倾注了祂的恩宠。因为祂的恩宠不是别的，正是圣神的恩宠；而那位当时被钉十字架、也是赐予许诺者，祂实现了所许诺的。若你们还愿意接受一个证词，\Emph{请听}，他说：\Quote{ \Emph{但这正是先知岳厄尔所说的：此后，天主说：我要倾注我的圣神} \BibleRef{岳 2:28} ——（这个词语\Emph{我要倾注}暗示丰富的恩赐；\Emph{因为天主赐圣神没有限量，因为父爱子，已把一切都交在祂手中} \BibleRef{若 3:34}；并且祂也将权能赐给祂，使祂能随意将至圣圣神的恩宠赐予任何人）—— \Emph{我要将我的圣神倾注在一切血肉上，你们的儿子和女儿要说预言}；此后，\Emph{甚至在那些日子，我也要将我的圣神倾注在我的仆人和婢女身上，他们要说预言} \BibleRef{岳 2:29}。}圣神不看人的情面；因为祂不寻求地位，只求灵魂的虔敬。但愿富人不自傲，穷人不沮丧，只要每人预备自己以接受天上的恩赐。\par

20. 今天我们说了许多，或许你们已厌倦听讲；但尚有许多待讲。的确，为了关于圣神的道理，需要第三次演讲，甚至更多次。但两件事我们都须请你们包容。因为圣洁的逾越节已近，我们今日延长了讲论，却仍未能将应当引用的新约证词全部呈于你们面前。因为尚有\WorkTitle{宗徒大事录}中许多段落，圣神的恩宠在\Person{伯多禄}和众宗徒身上大能运作；又有\WorkTitle{公函}及\Person{保禄}的十四封书信中的许多段落；从这一切中，我们现在如同从大草原上采撷几朵鲜花一般，略作回忆而已。\par

21. 因为藉圣神的大能，并因圣父与圣子的旨意，\Person{伯多禄}与十一宗徒一同站起，扬声（正如经上所记：\Emph{你这给耶路撒冷报喜讯的，要大力扬声} \BibleRef{依 40:9}）用他话语的灵网捕获了\Emph{约三千人}。临于众宗徒的恩宠如此之大，以至在犹太人——那些钉死基督的人——中，有这么一大群人相信了，并因基督之名受洗，且\Emph{恒心遵守宗徒的训诲和祈祷}。\BibleRef{宗 2:42}又藉同一圣神的大能，\Emph{\Person{伯多禄}和\Person{若望}在第九时辰祈祷时辰上圣殿}，并以耶稣之名治愈了那位生来瘸腿四十年的坐在美门的人；应验了经上所记：\Emph{那时瘸子将跳跃如鹿}。\BibleRef{依 35:6}就这样，正如他们藉教义灵网一次捕获五千信徒，他们也驳斥了迷途的民长和司祭长，并非靠自己的智慧，因为\Emph{他们是无学识的小民} \BibleRef{宗 4:13}，而是靠圣神的大能；因为经上记载：\Emph{那时\Person{伯多禄}充满圣神，对他们说}。圣神的恩宠藉十二宗徒在信者中所行也是如此浩大，以致\Emph{他们一心一意}，财物共享，拥有者虔敬地变卖产业，他们中间没有一人缺乏什么；而\Person{阿纳尼雅}和\Person{撒斐辣}因企图欺骗圣神，遭受了应有的惩罚。\par

22. \Emph{藉宗徒的手，在民间行了许多征兆和奇事}。\BibleRef{宗 5:12}环绕宗徒的属灵恩宠如此浩大，以至他们虽温和，却令人畏惧；因为\Emph{其他的人没有一个敢贴近他们，但百姓却尊敬他们；信主的人，男女都越来越多}；街道上满有病人在床榻上，\Emph{为的是当\Person{伯多禄}经过时，至少他的影子能遮住他们中的一些人}。\Emph{周围城市的人}也都来到这圣城耶路撒冷，\Emph{带着病人和受邪魔折磨的人，他们全都得了痊愈}，藉此圣神的大能。\par

23. 再者，十二宗徒因宣讲基督被司祭长下监，夜间奇迹般地被天使救出，并从圣殿被带到公议会前时，他们无所畏惧地以关于基督的言论斥责了他们，并补充说：\Emph{天主也将圣神赐给了服从祂的人}。当他们受了鞭打后，欢喜地离去，不停地\Emph{教导并宣讲耶稣是基督}。\BibleRef{宗 5:42}\par

24. 圣神的恩宠不仅运行在十二位宗徒身上，也运行在这位一度不生育的教会初生子身上，我是指那七位执事；因为他们正如经上所载，被选为\Quote{充满圣神和智慧的人}。其中司提反，名副其实的殉道者初果，一位\Quote{充满信德和圣神的人，在民间行了伟大的奇事和征兆}，并战胜了与他争辩的人；\Quote{因为他们不能抵抗他说话时所有的智慧和圣神}。但当他被恶意控告并带到审判庭时，他放射出天使般的光辉；因为\Quote{凡坐在公议会的人，都定睛看他，见他的面貌好像天使的面貌}。他用明智的辩护驳倒了犹太人，那些\Quote{硬着颈项、心与耳未受割礼、时常反抗圣神的人}，他看见了\Quote{天开了}，并看见\Quote{人子站在天主的右边}。他看见祂，并非靠自己的力量，而是如神圣的圣经所说：\Quote{充满圣神，举目望天，看见天主的荣耀，又看见耶稣站在天主的右边}。\par

25. 凭借圣神同样的能力，斐理伯也曾在撒玛黎雅城，因基督之名赶走污灵，它们\Quote{大声喊叫}；并医治了瘫痪的与瘸腿的，带领大批信众归向基督。伯多禄和若望下到他们那里，以祈祷和覆手，分施了圣神的共融，只有西满魔术师被公正地宣布为局外人。另一次，斐理伯为了一位极其虔诚的厄提约丕雅太监，被主的天使召唤在路上，并清楚地听到圣神亲自说：\Quote{你上前去，贴近这辆马车。}他教导了太监，为他施洗，从而将基督的使者派往厄提约丕雅，正如经上所载：\Quote{厄提约丕雅将急忙向天主伸手}，然后他被天使提去，在沿途各城宣讲福音。\par

26. 保禄被我们的主耶稣基督召唤后，也充满了这位圣神。让虔诚的阿纳尼雅来为我们的话作证，他在大马士革对保禄说：\Quote{那在来路上显现给你的主耶稣，打发我来，叫你看见，并充满圣神。}\BibleRef{宗 9:17} 圣神的大能立刻将保禄眼睛的瞎转变为重见光明；并在他的灵魂上赐予印记，使他成为\Quote{蒙选的器皿}去\Quote{在外邦人和君王并以色列子民面前}传扬那位向他显现的主的\Quote{名}，使从前的迫害者成为大使和忠仆——\Quote{从耶路撒冷直转到依里黎苛，处处传了基督的福音}，\BibleRef{罗 15:19} 甚至教导了帝国罗马，将宣讲的热忱带到西班牙，经历了无数的战斗，行了征兆与奇事。关于他，目前说这些就够了。\par

27. 凭借同一圣神的能力，伯多禄——宗徒之长、天国钥匙的持有者——在吕达（即现在的迪奥斯波里斯）因基督之名医治了瘫子艾乃阿，在约培使乐善好施的塔彼达从死者中复活。他在房顶上神魂超拔，看见天开了，有一块像大布的物体降下，里面装满各种四足兽，他清楚地学到不该称任何人为凡俗或不洁，即便他是外邦人。\BibleRef{宗 10:11} 当科尔乃略派人来请他时，他清楚地听到圣神亲自说：\Quote{看，有人找你；你起来，下去，跟他们同去，不要疑惑，因为是我派他们来的。}为了清楚显示信主的外邦人也得以分享圣神的恩宠，当伯多禄来到凯撒勒雅，讲论基督的事时，圣经论到科尔乃略和与他同在的人说：\Quote{伯多禄还在讲这些话的时候，圣神降在所有听道的人身上；那些与伯多禄一同来的受过割礼的信徒都惊讶，因为他们明白，圣神的恩赐也倾注在外邦人身上了。}\par

28. 在叙利亚最著名的城市安提约基雅，当基督的宣讲奏效时，巴尔纳伯被派往那里协助善工，他是一位\Quote{好人，充满圣神和信德}；他看到信基督的庄稼很多，就从塔尔索带来了保禄，作为他的战友。当群众受他们教导并在教会聚集时，\Quote{门徒称为\InnerQuote{基督徒}是从安提约基雅开始的}；我想，是圣神将主先前应许的这个新名赐给了信徒。天主要丰厚地在安提约基雅倾注圣神的恩宠，那里有先知和教师，阿加波是其中之一。\BibleRef{宗 11:28} 当他们\Quote{敬礼主并禁食的时候，圣神说：\InnerQuote{你们为我选出巴尔纳伯和扫禄，去作我召叫他们所作的工。}}在他们被覆手后，\Quote{他们就由圣神派遣出去了}。显然，那说话并派遣的圣神，是活生生的、自立体、运作的，正如我们所说。\par

29. 这同一的圣神，与圣父圣子同心，在公教会内建立了新约，将我们从难负的律法重担中解放出来——我指那些关于洁净与不洁之物、食物、安息日、月朔、割损、洒洗和祭献的规条；这些原是暂时赐下的，\BibleQuote{是将来美事的影子} \BibleRef{希 10:1}，但真理既已来临，便正当地被撤去。因为当\Person{保禄}和\Person{巴尔纳伯}被派往宗徒那里，因\Place{安提约基雅}有人争议说必须受割损并遵守\Person{梅瑟}的规条时，在\Place{耶路撒冷}的宗徒们以书面谕令使普世从一切法律和预像的规条中得自由；然而在这等大事上，他们并未将全权归于自己，而是发出书面命令，并如此承认：\Quote{因为圣神和我们决定，不再将更大的重担加给你们，除了这几件必要的事：就是禁食祭偶像之物、血、勒死的牲畜和奸淫。}他们以此明示，虽然书写是出于人手，但这道命令却来自圣神；\Person{保禄}和\Person{巴尔纳伯}便将这道命令带到并确认给普世。\par

30. 如今，在演讲进行到这一步时，我恳请你们的爱德谅解，更确切地说，是恳请那住在\Person{保禄}内的圣神谅解，倘若我因自身的软弱和听众的疲惫，不能一一详述。因为何时我才能以相称的言辞，宣告圣神因基督之名所行的一切奇事？在\Place{塞浦路斯}对术士\Person{厄吕玛}所行的，在\Place{吕斯特辣}治好瘸子，在\Place{基里基雅}、\Place{夫黎基雅}、\Place{迦拉达}、\Place{米息雅}、\Place{马其顿}所行的？或在\Place{斐理伯}所行的（我是指宣讲，以及因基督之名驱逐占卜之神；以及地震后夜间狱警全家借圣洗而得救）；或在\Place{得撒洛尼}所发生的事；以及在\Place{阿勒约帕哥}对雅典人的演讲；或在\Place{哥林多}及全\Place{阿哈雅}的训导？我如何能相称地述说圣神藉\Person{保禄}在\Place{厄弗所}所行的异能 \BibleRef{宗 19:1}？那城的人起初不认识祂，但藉\Person{保禄}的教导认识了祂；\Person{保禄}给他们覆手后，圣神便降临在他们身上，\BibleQuote{他们就说起各种语言，也说预言}。而且\Person{保禄}身上有如此大的神恩，不仅他的触摸能治病，甚至\BibleQuote{从他身上取来的手巾和围裙}，也能治愈疾病，驱走邪魔；最后\BibleQuote{那些行巫术的人，也将他们的书籍拿来，当众焚烧了}。\par

31. 我略过\Place{特洛阿}关于\Person{厄乌提苛}的事：他因\BibleQuote{沉睡而坠下三层楼，扶起时已死了}，却藉\Person{保禄}而复活。我也略过\Person{保禄}在\Place{米肋托}向\Place{厄弗所}的长老们所说的预言，他召他们来，公开说：\BibleQuote{圣神在各城向我指明，说}——等等；因为\Person{保禄}说\BibleQuote{在各城}，表明他在各城所行的奇事，都是由于圣神的运行之力，出于天主的旨意，并藉着在他内说话的基督之名。因这圣神之能，\Person{保禄}自己正赶往这圣城\Place{耶路撒冷}，尽管\Person{阿加波}藉圣神预言了将临到他身上的事；但他仍满怀信心地向民众说话，宣讲有关基督的事。当他被解到\Place{凯撒勒雅}，置身法庭之前，时而面对\Person{斐理斯}，时而面对总督\Person{斐斯托}和\Person{阿格黎帕}王时，\Person{保禄}从圣神得了如此大的恩宠，智慧超群，以致犹太王\Person{阿格黎帕}自己说：\BibleQuote{你差一点就劝服我作了基督徒}。这圣神也使\Person{保禄}在\Place{默利特}岛上，被毒蛇咬时安然无恙，并为患病的人行了各种医治。这圣神引导他——昔日迫害教会的人，成为基督的宣讲者——一直到了帝国的\Place{罗马}；在那里他说服了许多犹太人相信基督，对反驳的人清楚地说：\BibleQuote{圣神藉依撒意亚先知向你们的祖先说得真好，等等}。\par

32. \Person{保禄}充满了圣神，他所有的同工宗徒们也一样，以及后来信从父、子、圣神的人，请听他在书信中清楚写下的：他说：\Quote{我的言词和宣讲，不是用巧辩的智慧之言，而是用圣神和能力的明证}。\BibleRef{格前 2:4} 又说：\Quote{那为此事给我们盖上印记，并赐给我们圣神作为抵押的，就是天主}。\BibleRef{格后 1:22} 又说：\Quote{那使耶稣从死者中复活的那一位，也要藉那住在你们内的圣神，使你们有死的身体复活}。\BibleRef{罗 8:11} 还有，在写给\Person{弟茂德}的信中：\Quote{你所受的善托付，要因那住在我们内的圣神保管好}。\par

33. 圣神是自立体、生活、发言并预言，这在前面我已多次说过，保禄也明确写给弟茂德说：\Quote{圣神明明地说：在后期，有些人将背弃信德。}\BibleRef{弟前 4:1}——我们不仅在过去的时代，也在我们自己的时代看到分裂，异端者的错误是如此混杂多样。同一保禄又说：\Quote{这在以前的世代没有告诉过世人，如今在圣神内启示给祂的圣宗徒和先知们。}\BibleRef{弗 3:5} 又说：\Quote{因此，正如圣神所说的}\BibleRef{希 3:7}；又说：\Quote{圣神也向我们作证。} 他又向正义的士兵呼喊：\Quote{并戴上救恩的盔，拿着圣神的剑，即天主的话，以各种祈祷和恳求。}\BibleRef{弗 6:17} 又说：\Quote{不要醉酒，醉酒使人淫乱；却要充满圣神，以圣咏、圣歌及属神的歌曲彼此对谈。} 又说：\Quote{主耶稣基督的恩宠，天主的爱，和圣神的相通，常与你们众人同在。}\BibleRef{格后 13:14}\par

34. 藉所有这些证据，以及更多未提及的，圣神的位格、圣化及有效的大能得以确立，为能领会的人；因为若我想引用保禄十四封书信中关于圣神仍存的内容，我的讲论时间将不够用，他在其中以如此多样、全面而虔诚的方式教导了。并且属于圣神本身的大能，必须赐予我们宽恕我们所遗漏的，因为时日不多，并更完善地将剩余的知识印在你们听者心中；而透过勤读圣经，你们中勤勉者得以领悟这些事，并且此时，从这些演讲和以前告诉你们的内容，更坚定地持守信德：\Quote{我信唯一全能天主圣父；我信我们的主耶稣基督，祂的独生子；我信圣神，护慰者。} 虽然\Quote{神}这个词和名号在圣经中共同用于祂们——因为关于圣父说：\Quote{天主是神}\BibleRef{若 4:24}，正如若望福音所记载；关于圣子：\Quote{在我们面前的神，主基督}，正如耶肋米亚先知所说；关于圣神：\Quote{护慰者，圣神}\BibleRef{若 14:25}，如所说——但信经中信条的排列，若虔诚理解，也驳斥了撒柏流的错误。因此我们在讲论中回到那现在紧迫且对你们有益之点。\par

35. 小心，免得你像\Person{西满}一样，怀着虚伪之心来到施行圣洗者面前，你的心却不寻求真理。我们职责是警告，你们职责是保护自己。如果你\Quote{因信德站立}\BibleRef{罗 11:20}，你是有福的；如果你因不信跌倒，从今天起抛开不信，并接受全备的确信。因为在圣洗之时，当你来到主教、司铎或执事前——（因为其恩宠无处不在，在村庄和城市，对低贱者与高贵者，对为奴者与自由人，因为这恩宠不是出于人，而是藉着人来自天主的恩赐）——走近圣洗施行者时，但走近时，不要想你看见的那人的面容，而要记住我们正在谈论的这圣神。因为祂已准备好要封印你的灵魂，祂将赐给你那使恶神战栗的印记，属天而神圣的印记，正如经上所写：\Quote{在祂内你们也相信了，并因所应许的圣神受了封印。}\par

36. 然而祂考验灵魂。祂不把珍珠丢在猪前；如果你虚伪，虽然现在人们为你付洗，圣神却不会为你付洗。但如果你怀着信心走近，虽然人施行那可见的，圣神却赐予那不可见的。你正来到一个伟大的考验，一场伟大的集合，在这一小时里，你若浪费它，你的灾难将无法挽回；但若你蒙恩，你的灵魂将被光照，你将领受你未曾有的大能，你将领受令恶神畏惧的武器；你若不离弃你的武器，而将印记保持在你灵魂上，没有恶神会接近你；因它们将被吓倒；的确，恶神是因天主之神而被驱逐的。\par

37. 如果你相信，你不只领受罪的赦免，也要行超越人之力的事。愿你也能蒙受先知之恩！因为你将按你能力的尺度领受恩宠，而非按我言词的尺度；因为我能说的或许微不足道，你却能领受更大的；因为信德是大事。护慰者将一生作你的护守者；祂将照顾你，如同照顾自己的士兵；为你出行、进入和敌人的阴谋。祂将赐给你各种恩赐，若你不因罪而令祂忧伤；因为经上记着：\Quote{不要叫天主的圣神忧郁，你们原是在祂内受了印记，以待救赎的日子。}\BibleRef{弗 4:30} 那么，亲爱的，保守恩宠是什么意思？预备好领受恩宠，既已领受，就不要丢弃它。\par

38. 愿那藉圣神通过先知们发言的万有的天主，那在五旬节之日于此处派遣圣神到宗徒们身上的那一位，此时也派遣祂到你们身上；并且藉祂也保守我们，将祂的恩惠共同赐给我们众人，使我们永远结出圣神的果实：\Quote{爱德、喜乐、平安、忍耐、仁慈、良善、信实、柔和、节制}\BibleRef{迦 5:22}，在我们的主耶稣基督内。藉着祂，并与祂，同圣神，愿光荣归于圣父，从今直到永远。阿们。\par

\SourceRecord{Translated by Edwin Hamilton Gifford.；Nicene and Post-Nicene Fathers, Second Series；Vol. 7.；Edited by Philip Schaff and Henry Wace.；Buffalo, NY: Christian Literature Publishing Co.,；1894.；<http://www.newadvent.org/fathers/310117.htm>.}

% structure: p0001 source=h1 target=section reason=page-title

\section{教理讲授十八}

\Emph{\Quote{并信唯一、至圣、至公、从宗徒传下来的教会}，以及\Quote{肉身的复活}和\Quote{永生}}\par

\BibleRef{则 37:1}\par

\BibleQuote{上主的手临于我，上主的神引领我出去，把我放在平原中，那里布满了骨头。}\par

一切善工的根源是复活的希望；因为对报酬的期待激励灵魂行善。因为每个劳动者，若能预见劳苦的奖赏，就愿意忍受辛劳；但若徒劳无功，心身都会很快消沉。期待奖赏的士兵愿意作战，但没有人愿意为一个漠视其属下、不尊崇其辛劳的国王去死。同样，每个相信复活的灵魂自然谨慎自守；但若不信复活，就会放任自流，走向丧亡。相信身体将复活以再起的人，会珍惜自己的衣裳，不以淫行玷污它；但不信复活的人，便放纵于淫行，滥用自己的身体，仿佛不是自己的。因此，信德在死人复活上，是至圣至公教会的一条伟大诫命和教义；这是伟大且极为必要的教义，虽被许多人反对，却确实由真理所证实。希腊人反驳它，撒玛黎雅人不信它，异端者曲解它；反对虽多，真理却是一致的。\par

如今希腊人和撒玛黎雅人一同这样反驳我们：死人倒下，腐烂，全部化为蛆虫；蛆虫也死了；身体遭遇这样的腐朽和毁灭，如何再复活？遭遇海难者被鱼吞噬，而这些鱼又被吞噬。与野兽搏斗的人，连骨头都被碾碎，被熊和狮子吃掉。秃鹫和乌鸦啄食未埋葬死者的肉，然后飞向世界各地；那么身体如何收集？因为吞食了尸体的飞鸟，有的可能死在印度，有的死在波斯，有的死在哥特人之地。另有一些人被火烧毁，骨灰被雨水或风吹散；身体又如何重新聚合？\par

对你，可怜软弱的人啊，印度离哥特人之地很远，西班牙离波斯很远；但对于天主，祂将整个大地握在掌中（\BibleRef{依 40:12}），万有都近在咫尺。不要因此将软弱归给天主，以你的软弱作比较，而应默思祂的大能。那么，太阳——天主的一个微小化工——凭它的一束光芒就能温暖整个世界；天主所造的空气环绕世上万物；而创造太阳和空气的天主，难道远离世界吗？设想各种植物的种子混合在一起（因为你信心软弱，我举的例子也是软弱的），这些不同的种子装在你的一只手中；那么，你作为一个人，分辨你手中的东西，依其本性收集每一粒种子，放回其种类中，是难事还是易事？你能分辨手中的东西，而天主不能分辨祂手中的东西，并将它们放回原处吗？请思考，否认这一点岂非不敬？\par

此外，请留意正义的原则本身，看看你自己的情况。你有不同种类的仆人，有的好，有的坏；你因此尊崇善者，惩罚恶者。如果你是一位法官，你会对善者给予赞扬，对犯法者施以惩罚。那么，作为凡人，你尚且施行正义，而天主，永远不变的万有君王，难道没有报应的正义吗？不，否认它是不敬的。因为请思考我所说的：许多杀人犯躺在床上未受惩罚而死；天主的正义在哪里？是的，一个犯下五十起谋杀罪的凶犯往往只被斩首一次；那么他为那四十九起罪在何处受罚呢？除非在现世之后有审判和报应，你就是在指控天主不义。然而，不要因审判的延迟而惊奇；没有参赛者在竞赛结束之前被加冕或受辱；竞赛的主持者不会在人们还在奋力时给他们加冕，而是等待所有参赛者完赛，然后在他们之间裁决，分发奖品和花冠。同样，天主在现世的争斗持续期间，只是部分地扶助义人，但事后会将赏报完全赐予他们。\par

但如果按照你们的说法，没有死人复活，为何你们定盗墓者的罪？因为若身体灭亡，且没有复活可盼望，为何破坏坟墓者要受惩罚？你看到，尽管你口里否认，但你心中仍保留着不可磨灭的复活本能。\par

6. 再者，一棵树被砍伐后，还能再发芽；人若被砍伐后，就不能再发芽吗？麦子被撒下又收割，尚且为打谷场而存留；人从这世界被收割后，难道就不为打谷场而存留吗？葡萄树或其他树木的枝条，被干净地剪下并移植后，还能复活结果；而为了人的缘故存在的这一切，人却要落入土中不再复活吗？比较一下，是重新塑造一个原先不存在的雕像更难，还是将已经倒下的雕像重新铸成同样的形状更难？天主从无中创造了我们，难道不能使那已存在却倒下的复活吗？但是，你身为外邦人，不相信圣经关于复活的记载；那么请从自然界的类比中来思考这些事，并从今日所见之事来理解。例如，小麦或某种其他谷物被撒下；种子落在地里，死了、腐烂了，从此不再适合食用。但那腐烂的，却发芽长出青绿；虽然撒下时很小，却长得非常美丽。麦子原是为我们而造的；因为麦子和所有种子都不是为自己，而是为我们的使用而被造的；那么，那些为我们而造的事物在死后尚且获得生命，而我们这些它们为之而被造的人，死后却不再复活吗？\par

7. 如你所见，现在正是冬天，树木好像死了一样：无花果树的叶子在哪里？葡萄树的果实又在哪里？它们在冬天是死的，但在春天却变绿；当季节来到时，它们仿佛从死亡状态中恢复了生机。因为天主知道你们的不信，每年都在这些可见的事物中行一个复活；好叫你们看见无生命之物所发生的事，从而相信有生命和理性之物。此外，苍蝇和蜜蜂常常淹死在水里，但过一段时间又活过来；冬眠的睡鼠整个冬天一动不动，夏天又恢复过来（因为为你微小的思想，提供了这样的例子）；那赐予无理性、卑贱之物超自然生命的天主，岂不将此恩赐给我们这些祂为他们而造的人吗？\par

8. 但是外邦人寻求更显明的死人复活；他们说，即使这些生物复活了，它们却并未完全腐朽；他们要求清楚地看到某种生物在完全腐烂后复活。天主知道人的不信，为此预备了一种鸟，称为不死鸟。正如\Person{克莱孟}所记载的，以及其他许多人所述，这种鸟是独一无二的，每五百年来到\Place{埃及}地，显明复活的事，不是在荒僻之处，以免这奥秘的发生无人知晓，而是出现在一座著名的城中，以使人们甚至可以触摸那原本会被否认的事。因为它用乳香、没药和其他香料为自己造一口棺材，当它的年限满足时，就进入其中，显然死去并腐烂了。然后从死鸟腐烂的肉中生出一条虫，这条虫长大后变成一只鸟——不要不信，因为你看到蜜蜂的后代也是从虫形成的，而且从完全是液体的蛋中，你看到鸟的翅膀、骨头和筋腱生出。此后，前述的不死鸟长出羽毛，成为一只完全长大的不死鸟，像之前一样，飞上天空，如同它死时一样，向人显明死人复活的最明显证据。不死鸟诚然是奇妙的鸟，但没有理性，从来没有赞颂过天主；它在天空中飞翔，却不知道天主的独生子是谁。那么，难道复活从死里被赐给了这不认识造物主的无理性生物，而我们这些归光荣于天主并遵守祂诫命的人，却得不到复活吗？\par

9. 但是，既然不死鸟的迹象遥远且不常见，而人们仍然不相信我们的复活，那就从你们每天所见的事中再取一个证明。一百年或两百年前，我们所有说话的人和听道的人，在哪里呢？我们难道不知道我们身体物质的基础吗？难道不知道我们从微弱、无形和简单的元素被孕育，从简单而微弱的东西形成一个活人吗？并且那微弱的东西如何变成肉，转化为强壮的筋、明亮的眼睛、敏锐的鼻子、听力的耳朵、说话的舌头、跳动的心脏、忙碌的手、快速的脚，以及各种肢体？那曾经微弱的东西如何成为造船匠、建筑者、建筑师、各种工艺的工匠、士兵、统治者、立法者、国王？那么，天主从不完全的材料创造了我们，难道不能在我们朽坏之后使我们复活吗？那从卑贱之物如此塑造身体的天主，难道不能使倒下的身体再次复活吗？那塑造不存在之物的天主，难道不能使已经存在却倒下的东西复活吗？\par

10. 再从天空及其光体每月可见的实例中，取一个死人复活的明显证明。月亮的身体完全消失，以致再也看不见它的任何部分，但它又再充满，恢复到原来的状态；为了完全证明这件事，月亮在若干年的周期中发生月食，明显变成血色，却又恢复它光亮的身体：天主这样安排，为的是使你，这由血形成的人，不至于拒绝相信死人复活，反而能从月亮所见的事上，也相信你自己。因此，你要用这些作为反驳外邦人的论据；因为对于不接受圣经的人，你要用未经书写的武器，即仅仅通过推理和论证来与他们争辩；因为他们不知道\Person{梅瑟}是谁，不知道\Person{依撒意亚}，不知道\WorkTitle{福音}，也不知道\Person{保禄}。\par

11. 现在转向撒玛黎雅人，他们只接受\WorkTitle{法律}，不接受\WorkTitle{先知}。对他们而言，刚才从\WorkTitle{厄则克耳}所读的经文似乎无效，因为，如我所说，他们不接受任何先知；那么，我们如何说服撒玛黎雅人呢？让我们到\WorkTitle{法律}的著作中去。天主对梅瑟说：\Quote{我是亚巴郎的天主，依撒格的天主，雅各伯的天主}；这必定指那些具有存有和实体的。因为如果亚巴郎已经终结，依撒格和雅各伯亦然，那么祂就是不存在者的天主。何时一位君王说过：我是我所没有的士兵的君王？何时有人展示过他所没有的财富？因此，亚巴郎、依撒格和雅各伯必须存在，好让天主成为存有者的天主；因为祂没有说：\Quote{我曾是他们的天主}，而是\Quote{我是}。而且有审判，亚巴郎向主说时表明：\Quote{审判全地的，岂能不执行公义？} \BibleRef{创 18:25}\par

12. 但愚昧的撒玛黎雅人又对此提出异议，说亚巴郎、依撒格和雅各伯的灵魂或许继续存在，但他们的身体不可能复活。那么，义人梅瑟的棍杖能变成蛇，义人的身体就不能活着并复活吗？那是违反本性的，而这些就不能按本性恢复吗？再者，亚郎的棍杖，虽已被砍下并枯死，却发芽了，\Quote{不待水气} \BibleRef{约 14:9}，虽在屋顶下，却如田间开花；虽置于干旱之处，一夜之间就生出多年浇灌之植物的花和果实。亚郎的棍杖犹如从死中复活，而亚郎本人就不会复活吗？天主在木头中行奇迹以保障他获得大司祭职，难道不会赐予亚郎本人复活吗？一个女人也违反本性变成了盐；肉变成了盐；肉就不能恢复为肉吗？罗特的妻子成了盐柱，亚巴郎的妻子就不能复活吗？梅瑟的手，一小时内变得像雪又复原，是靠什么能力？当然是靠天主的命令。那时祂的命令有力量，现在就没有力量了吗？\par

13. 人啊，最无知的撒玛黎雅人，起初人究竟从何而来？去\WorkTitle{圣经}的第一卷书，就是你们也接受的：\Quote{天主用地上的尘土形成了人。} \BibleRef{创 2:7} 尘土变成了肉，肉就不能再恢复为肉吗？还必须问你们：诸天、大地、海洋从何而来？太阳、月亮和星辰从何而来？飞鸟和水族从何而来？地上的一切活物从何而来？如此众多的万物从无中被造出来，而我们这些带着天主肖像的人就不能复活吗？这实在是充满不信，不信者应受严厉谴责；当亚巴郎称呼主为\Quote{审判全地者}时，学法律的人却不信；当经上写人出于尘土时，读者却不信。\par

14. 这些问题是为他们，不信者准备的；但先知的话是为我们相信的人。然而，有些使用先知书的人也不信所写的，并引用经文来反对我们：\Quote{不虔敬的人在审判中不得兴起}，以及\Quote{人若下到坟墓，就不得再上来} \BibleRef{约 7:9}，还有\Quote{死人不能赞美你，主}——因为他们对写得好的经文做了恶用——因此，简单扼要地，按目前所可能的，予以回应。因为如果说\Quote{不虔敬的人在审判中不得兴起}，这表明他们兴起不是受审判，而是受惩罚；因为天主不需要冗长的审查，不虔敬者复活后随即就是惩罚。如果说\Quote{死人不能赞美你，主}，这表明，只在今生才有悔改和赦免的时机，那些享有悔改时机的人才会\Quote{赞美主}；死后，对于死在罪中的人，不再有赞美以领受恩惠，而是哀哭；因为赞美属于感恩者，而受鞭打者属于哀号。因此，义人那时献上赞美；但死在罪中的人没有进一步认罪的机会。\par

15. 至于那段经文：\Quote{人若下到坟墓，就不得再上来}，请注意下文，因为经上写着：\Quote{他不得再上来，也不再回到自己的家中}。因为整个世界要过去，每座房屋都要被拆毁，他怎能回到自己的家中呢？从此将有新的、不同的大地。但他们应该听约伯的话：\Quote{因为树还有希望；若被砍下，仍可发芽，嫩枝也不止息。虽然它的根在地里衰老，干在土中死去；然而一有水气，它就会萌芽，长出枝条如新栽的树。但人一死，就去了；凡人一旦倒下，就不再存在了？} \BibleRef{约 14:7} 仿佛在抗议和责备（因为这样读\Quote{不再存在}应带疑问语气）；他说，既然树倒下后又复活，难道为它们而造了所有树的人，不会复活吗？为使你不以为我在曲解经文，请读下文；因为在以疑问句说了\Quote{凡人一旦倒下，就不再存在？}之后，他说：\Quote{因为人若死了，必再活}；紧接着他又说：\Quote{我要等待，直到我重新被造}；又在别处说：\Quote{谁将在地上使承受这些的我的皮肤兴起？}。依撒意亚先知说：\Quote{死人要复活，在坟墓中的要醒起。} \BibleRef{依 26:19} 而厄则克耳先知现在对我们非常清楚地说：\Quote{看，我要打开你们的坟墓，把你们从坟墓中领上来。} \BibleRef{则 37:12} 达尼尔说：\Quote{许多睡在尘土中的人要醒来，有的进入永生，有的进入永远的羞辱。} \BibleRef{达 12:2}\par

16. 有许多圣经为死者的复活作证；因为关于这事还有许多其他的说法。但现在仅为了纪念，我们要略略提及第四天拉匝禄的复活；由于时间短促，也仅仅提及那寡妇的儿子被复活的事，而为了提醒你，让我提及会堂长的女儿，以及岩石的崩裂，和\BibleQuote{许多长眠的圣者的身体复活了}\BibleRef{玛 27:52}，他们的坟墓被打开。但特别要记住：\BibleQuote{基督已从死者中复活了}\BibleRef{格前 15:20}。我只顺便提及\Person{厄里亚}和他所复活的寡妇的儿子；也提及\Person{厄里叟}，他两次复活死者；一次在他生前，一次在他死后。因为当他活着时，他藉自己的灵魂施行了复活\BibleRef{列下 4:34}；但这不仅为使义人的灵魂受到尊荣，也为使人相信在义人的身体内也含有一种能力，那被抛入\Person{厄里叟}坟墓中的尸体，一碰到先知的尸体就复活了，先知的尸体做了灵魂的工作，那已死并被埋葬的给了死者生命，虽然它给了生命，自己却仍留在死者中。为什么？免得若\Person{厄里叟}复活，这工作被归功于他的灵魂；而是为表明，即使灵魂不在场，圣人的身体内仍有一种德能，因为那义人的灵魂曾多年居住其中，并使用它作为工具。让我们不要愚昧地不信，好像这事没有发生过；因为如果手帕和围裙，只是外在的东西，触摸病人的身体就能使病人起来，那么先知的尸体岂不更能复活死者？\par

17. 关于这些事例，我们本可详述每一事的奇妙情节；但你们已因预备日的附加斋戒和守夜感到疲倦，所以暂且满足于这些简述；这些话有如稀疏播下的种子，好使你们如同最肥沃的田地接受种子，在结实时增多。但要记住，宗徒们也复活了死者；\Person{伯多禄}在\Place{约培}复活了\Person{塔彼达}，\Person{保禄}在\Place{特洛阿}复活了\Person{欧提科}；其余众宗徒也是如此，尽管每人所行的奇事并未全部记载。此外，要记住\Person{保禄}在\WorkTitle{格林多前书}中，针对那些说\BibleQuote{死人怎么复活呢？他们将带着什么样的身体来呢？}\BibleRef{格前 15:35}的人所说的一切话。以及他如何说：\BibleQuote{因为死人若不复活，基督也就没有复活}；并将那些不信的人称为\BibleQuote{糊涂人}；并记住他在那里关于死者复活的全部教导，以及他写给\Place{得撒洛尼}人的信：\BibleQuote{弟兄们，论到长眠的人，我们不愿你们不知道，免得你们忧伤，像其他没有望德的人一样}\BibleRef{得前 4:13}及其下文；但主要是：\BibleQuote{那在基督内死了的人要先复活}。\par

18. 但要特别留意\Person{保禄}是多么一针见血地说：\BibleQuote{这必朽坏的必须穿上不朽坏的，这必死的必须穿上不死的}\BibleRef{格前 15:53}。因为这身体复活时，将不再像现在这样软弱；而是复活同一个身体，但藉穿上不朽坏而被改造——如铁与火融合成为火，或者更好说，如那复活我们的主所知道的那样。因此，这个身体将要复活，但它不会保持现在的样子，而是一个永恒的身体；不再需要像现在这样的食物来维持生命，也不需要阶梯来上升，因为它将被造成\Emph{属神的}，一件奇妙的事，我们无法恰当地言说。经上说：\BibleQuote{那时，义人将发光如太阳}\BibleRef{玛 13:43}，和月亮，\BibleQuote{又如穹苍的光明}\BibleRef{达 12:3}。天主预知人的不信，在夏天赐给小小的虫子从身体中放射光芒，好叫人从可见的事相信那所期望的事；因为那赐予部分者也能赐予全部，那使虫子发光者，将更会光照义人。\par

19. 因此我们都要复活，带着永恒的身体，但并非所有人的身体都一样：因为如果一个人是义人，他将领受一个属天的身体，好能堪与天使交游；但如果一个人是罪人，他将领受一个永恒的身体，适于承受罪的刑罚，好使他永远在火中焚烧，却永不消灭。天主公义地将此分派给两类人；因为我们没有身体就不能做什么。我们用口亵渎，也用口祈祷。我们用身体行奸淫，也用身体守贞洁。我们用手偷窃，也用手施舍；其余也是如此。既然身体在一切事上曾是我们的仆役，将来它也必与我们分享过去的果实。\par

20. 因此，弟兄们，让我们谨慎对待自己的身体，不要滥用它，仿佛它不属于我们。不要像异端者那样说，这身体的外衣不属于我们，而要把它当作自己的来谨慎对待；因为我们必须向主交账，关于通过身体所做的一切事。不要说，没有人看见我；不要想，没有行为的见证者。人的见证常常没有，但那位塑造我们的、无误的\Emph{见证者}，\Emph{在天上}信实地住着，并看见你所做的。罪的污点也留在身体里；因为正如当伤口深入身体时，即使愈合了，疤痕仍然存在，同样，罪伤害灵魂和身体，其疤痕的痕迹留在所有人身上；只有那些领受了圣洗圣事洗涤的人，这些痕迹才被除去。因此，天主藉着圣洗医治灵魂和身体过去的创伤；至于将来的创伤，让我们全体共同提防，好使我们能保持这身体的外衣纯洁，不至于为了短暂的奸淫和感官纵欲或任何其他罪，而失去天堂的救恩，反而能继承天主的永恒国度；愿天主因祂的恩宠，使你们众人堪当承受此国。\par

21. 以上是关于死人复活的证明；现在，让我再次向你们背诵信仰的宣认，你们要在我说话时以全部勤勉宣告它，并记住它。\par

\_\_\_\_\_\_\_\_\_\_\_\_\_\_\_\_\_\_\_\_\_\_\par

22. 我们所背诵的信仰，按顺序包含以下内容：\Quote{并且承认一个为赦罪而忏悔的圣洗；一个圣而公教会；肉身的复活；以及永生。}关于圣洗和忏悔，我已在最早的讲道中谈过；现在关于死人复活的话，是针对信条\Quote{肉身的复活}而说的。现在，让我完成对信条\Quote{一个圣而公教会}仍有待说的话；关于这信条，虽然可以说很多，但我们将简略地谈。\par

23. 它之所以被称为\Quote{公教会}，是因为它遍及全世界，从地极到地极；因为它普遍而完整地教导人应当知晓的一切教义，涉及可见与不可见、天上与地下的事物；因为它使全人类——统治者与被统治者、有学问与无学问的——都归顺于虔敬；因为它普遍地处理并医治灵魂或身体所犯的一切罪恶，并且在自身中拥有每一种称为德行的形式，无论是行为还是言语，以及在每一种属灵恩赐中。\par

24. 它正确地被称为\OriginalTerm{教会}{Ecclesia}，因为它召唤并聚集所有人；正如主在\BibleBook{}{肋未纪}中所说：\BibleQuote{在会幕门口，为全体会众召开集会。}应当注意，\BibleQuote{集会}这个词在圣经中首次使用于此，即主立亚郎为大司祭之时。在\BibleBook{}{申命纪}中，主也对梅瑟说：\BibleQuote{你要将人民聚集到我面前，让他们听我的话，好学习敬畏我。} \BibleRef{申 4:10}他又在提到约版时说：\BibleQuote{上面写满了主在山上从火中与你们说话的一切话，就是在集会的那一天。}好像更清楚地说：在你们被天主召唤并聚集在一起的那一天。\BibleBook{}{圣咏}作者也说：\BibleQuote{主啊，我要在大会中称谢你；在众多民中我要赞美你。}\par

25. 从前，\BibleBook{}{圣咏}作者歌唱：\BibleQuote{你们要在会众中祝福天主，也就是主，你们这些从以色列的源头出来的。}但犹太人因他们对救主所设的阴谋而被抛弃，离开了祂的恩宠，救主便从外邦人中建立了第二所圣教会，即我们基督徒的教会。关于这教会，祂对伯多禄说：\Quote{在这磐石上，我要建立我的教会，地狱的门决不能战胜她。} \BibleRef{玛 16:18}达味预言了这两者，明确地说到那被弃绝的第一个：\BibleQuote{我憎恨恶人的集会。}但说到那被建立的第二个，他在同一篇\BibleBook{}{圣咏}中说：\BibleQuote{上主啊，我喜爱你圣殿的华美。}紧接着：\BibleQuote{在集会中，上主啊，我要赞美你。}因为如今在犹大地的那个教会已被抛弃，基督的教会却遍及全世界；关于它们，\BibleBook{}{圣咏}中说：\BibleQuote{你们要向主唱新歌，在圣人的集会中赞美他。}与此相应，先知也对犹太人说：\Quote{万军的上主说：我不喜欢你们。} \BibleRef{拉 1:10}接着又说：\BibleQuote{从日出之地到日落之处，我的名在外邦人中受尊崇。}关于这个圣而公教会，保禄在写给弟茂德的信中说：\BibleQuote{你要知道应当如何在天主的家中生活，这家就是永生天主的教会，真理的柱石和基础。} \BibleRef{弟前 3:15}\par

26. 但既然\OriginalTerm{教会}{Ecclesia}一词应用于不同的事物（如经上论及厄弗所的戏院群众所写：\BibleQuote{他这样说了，便解散了会众}\BibleRef{宗 19:14}），并且有人可以恰当而真实地说，有一个\Emph{作恶者的教会}——我指的是异端者、\OriginalTerm{马尔西翁派}{Marcionists}和\OriginalTerm{摩尼教徒}{Manichees}等等的聚会——为此，这份信仰如今稳固地交付给了你们这篇信文：\Quote{我信唯一、圣而天主教会}；好使你们避免他们可怜的聚会，并永远留在你们在其中得以重生的圣而天主教会里。如果你们有时旅居在城市中，不要仅仅问主的圣殿在哪里（因为其他亵渎教派也试图称自己的巢穴为主殿），也不要仅仅问教会在哪里，而要问\OriginalTerm{天主教会}{Catholic Church}在哪里。因为这是这座圣教会——我们众人的母亲——特有的名称，她是我们的主耶稣基督、天主的独生子的净配（正如经上所写：\BibleQuote{基督也爱了教会，并为她舍弃了自己}\BibleRef{弗 5:25}，以及其余的一切），并且是\BibleQuote{那在上面的耶路撒冷，她是自由的，是我们众人的母亲}\BibleRef{迦 4:26}的预像和副本；这耶路撒冷先前是不生育的，如今却有许多子女。\par

27. 因为当最初的教会被舍弃后，在第二个教会——即天主教会——天主\Emph{设立了}，如保禄所说：\BibleQuote{第一是宗徒，第二是先知，第三是教师，其次是奇迹，然后是治病的奇恩、助人的、治理人的、各种语言之恩}\BibleRef{格前 12:28}，以及各种美德，即智慧和明悟、节制和正义、仁慈和良善，以及在迫害中不可战胜的忍耐。她\BibleQuote{藉着左右的正义武器，藉着荣誉和羞辱}\BibleRef{格后 6:7}，在昔日迫害和磨难中，用多种多样、绽放忍耐的华冠为圣殉道者加冕，如今在和平时期，藉着天主的恩宠，从\BibleQuote{君王和一切在高位的人}\BibleRef{弟前 2:2}，以及各个种类和民族的人，接受她应得的尊敬。尽管个别国家的君王对其权力有界限，但圣而天主教会独自将其权力扩张到全世界且无限制；\Quote{因为天主立定了她的边界为和平}。但是，如果我愿意谈论与她有关的一切，我还需要好几个小时来讲述。\par

\_\_\_\_\_\_\_\_\_\_\_\_\_\_\_\_\_\_\_\_\_\_\par

28. 在这圣而天主教会中，接受教导并行为正直，我们将获得天国，并继承永生；为此我们也忍受一切劳苦，好能从主那里成为其参与者。因为我们的目标并非小事，我们的努力是为着永生。因此，在宣认信仰时，在\Quote{我信肉身的复活}这词语之后，即死人的复活（我们已论述过），我们也受训相信\Quote{永生}，我们作为基督徒正是为此而奋斗。\par

29. 真正的生命是圣父，祂通过圣子在圣神内，如从泉源般将祂天上的恩赐倾注给众人；并因祂对人的爱，永生的恩福也确凿地应许给我们人类。我们不可不信这可能性，而应着眼祂的大能而非我们的软弱，相信；\Emph{因为在天主，一切都是可能的}。至于这是可能的，并且我们可以期望永生，达尼尔宣告：\Quote{而许多正义的人，要发出光芒，如同星辰，直到永远}。保禄也说：\BibleQuote{这样，我们就必永远与主同在}\BibleRef{得前 4:17}；因为\Emph{永远与主同在}意味着永生。但最清楚的是，救主自己在福音中说：\BibleQuote{这些人要进入永罚，而正义的人要进入永生}\BibleRef{玛 25:46}。\par

30. 关于永生的证据很多。当我们渴望获得永生时，圣经指示我们获得它的途径；由于我们讲论的时间已长，现在我们摆在你们面前的经文只是少数，其余留给勤勉者去查考。它们有时指出，它出于信德；因为经上记载：\BibleQuote{信从子的人，有永生}\BibleRef{若 3:36}，以及后面的话；而祂自己又说：\Quote{我实实在在告诉你们：听我的话，相信派遣我来的那一位，就有永生}，以及其余。另有一次，是由于宣讲福音；因为祂说：\Quote{收割的人已得工钱，并为永生收集果实}。另有一次，是由于为基督之名而殉道和宣认；因为祂说：\Quote{那爱惜自己性命的，反而丧失性命；那恨恶自己性命在现世的，必要保存性命于永生}。又因为爱基督胜过财富或亲属；\BibleQuote{凡为了我的名，舍弃了房屋、或兄弟、或姐妹}\BibleRef{玛 19:29}，以及其余，\BibleQuote{必要承受永生}。此外，还由于遵守诫命，\Quote{不可奸淫，不可杀人}，以及随后的话；正如祂回答那来到祂面前，说\BibleQuote{善师，我该做什么才能得永生}\BibleRef{谷 10:17}的人。更进一步，是由于远离恶行，从此服侍天主；因为保禄说：\BibleQuote{但是现在你们脱离了罪恶，成为天主的奴隶，你们得到的效果是成圣，结局是永生}\BibleRef{罗 6:22}。\par

31. 获得永生的途径很多，虽然我因数量众多而省略了它们。因为主以祂的慈爱开放了，不是一两个门，而是许多门，通过它们进入永生，尽祂所能，使众人毫无阻碍地享受它。如此，我们目前简略地谈论了永生，这是信德中所宣认的最后一条道理，也是其终结；愿我们众人，无论是教师还是聆听者，都因天主的恩宠而享受这永生！\par

\_\_\_\_\_\_\_\_\_\_\_\_\_\_\_\_\_\_\_\_\_\_\par

32. 如今，亲爱的弟兄们，训导之言劝勉你们众人，预备你们的灵魂，以接受天上的恩赐。关于交付给你们明认的圣而宗徒的信仰（\OriginalTerm{圣而宗徒的信仰}{the Holy and Apostolic Faith}），我们藉主的恩宠，在过去的四旬期内，尽可能多地宣讲了这些训导；这并不是说我们该说的都已说尽，因为还有许多要点尚未提及；而这些或许由更卓越的导师思考得更为周全。但现在逾越圣日临近，你们在基督内所爱的，将藉再生之洗（\OriginalTerm{再生之洗}{the Laver of regeneration}）受到光照。因此，若天主愿意，你们将再次受教于那必要的事：你们被召时，应怀着何等虔诚与秩序进入，以及圣洗圣事（\OriginalTerm{圣洗}{Baptism}）的每项神圣奥秘（\OriginalTerm{奥秘}{mysteries}）为何目的而施行，又当以何等敬畏与秩序，从圣洗走向天主的神圣祭台，并享用它属灵与天上的奥秘；为使你们的灵魂，先藉教义之言受到光照，从而在每一细节中发现天主厚赐于你们的恩赐之伟大。\par

33. 在逾越节救恩圣日之后，你们要从周一开始，每日前来，在复活圣所（\Place{复活圣所}{the Holy Place of the Resurrection}）聚会，在那里，若天主许可，你们将聆听其余的训导；在这些训导中，你们将再次受教于所行一切事的原因，并从旧约与新约中获得其证明——首先是关于圣洗前所行之事；其次，你们如何被主藉着水洗，藉着言语（\OriginalTerm{藉着水洗，藉着言语}{by the washing of water with the word}）（\BibleRef{弗 5:26}）洗净了罪过；又如何如同司铎一样，成了基督之名的分享者；以及圣神共融的印记（\OriginalTerm{圣神共融的印记}{the Seal of the fellowship of the Holy Ghost}）是如何赐给你们的；还有关于新约祭台的奥秘（\OriginalTerm{新约祭台的奥秘}{the mysteries at the Altar of the New Testament}），这些奥秘从这地方发源，神圣圣经传授给我们什么，这些奥秘的德能是什么，你们应如何接近它们，何时及如何领受它们；最后，此后你们当如何在言语和行为上相称于这恩宠，使你们众人皆能得享永生。这些事，若天主愿意，都将讲述。\par

34. 最后，我的弟兄们，你们要常常在主内喜乐；我再说，喜乐吧！因为你们的救赎近了，天上的天使大军正等候你们的得救。如今有呼喊者在旷野中的声音说：\Quote{你们当预备主的道路}（\BibleRef{依 40:3}）。先知也呼喊：\Quote{凡口渴的，请到水边来}。紧接着又说：\Quote{你们当听从我，就能吃美物，你们的灵魂必以美味为乐}。再过不久，你们将听到那卓越的教训说：\Quote{照耀吧，照耀吧，新耶路撒冷！因为你的光明已来到}（\BibleRef{依 60:1}）。关于这耶路撒冷，先知曾言：\Quote{以后你必称为正义之城，熙雍，忠信的众城之母}；\Quote{因为法律出自熙雍，上主的言语出自耶路撒冷}，这言语从此处倾注于普世。先知也论及你们对她说：\Quote{你举目四望，看你的儿女聚集而来}；她回答说：\Quote{这些像云彩飞来，像鸽子带着幼雏飞到我这里的是谁？}（说像云彩是因他们的属灵性质，像鸽子是因他们的纯洁）。她又说：\Quote{谁曾听说过这样的事？谁曾见过这样的事？一片土地岂能在一天之内生出？一个民族岂能在一时之间诞生？因为熙雍一经历产痛，便生下了她的儿女}。万物将因主所说的话充满说不尽的喜乐：\Quote{看，我要创造耶路撒冷为喜乐，她的百姓为欢乐}。\par

35. 愿这些话如今也再次在你们身上应验：\Quote{诸天，欢唱吧！大地，喜乐吧！}；接着：\Quote{因为上主怜悯了祂的百姓，安慰了祂百姓中的卑微人}。这要因天主的慈爱而成就，祂对你们说：\Quote{看，我要如云彩般涂抹你的过犯，如密云般涂抹你的罪恶}（\BibleRef{依 44:22}）。但你们既蒙受那称为\Quote{忠信者}之名（关于这名字，经上记载：\Quote{在我的仆人身上，要称呼一个新名，这名字必在地上受祝福}），你们要欢然说：\Quote{我们的主耶稣基督的天主和父当受赞美，祂在基督内，以天上各种属灵的祝福祝福了我们}（\BibleRef{弗 1:3}）；\Quote{在祂内，我们藉着祂的血获得了救赎，罪过的赦免，这是按祂恩宠的丰富，这恩宠祂曾丰厚地倾注在我们身上}，以及随后的话；又说：\Quote{但天主富于慈悲，因祂爱我们的大爱，竟在我们因过犯而死的时候，使我们同基督一起活过来}，等等。同样，你们要赞美一切美善的主，说：\Quote{但当我们的救主天主的良善和祂对人的慈爱显现时，祂救了我们，并不是由于我们本着义德所立的功劳，而是出于祂的怜悯，藉着再生之洗和圣神的更新，这圣神祂藉着我们的救主耶稣基督，丰厚地倾注在我们身上，好使我们因祂的恩宠成义，本着希望成为永生的承继者}（\BibleRef{铎 3:4}）。愿我们的主耶稣基督的天主和父，荣耀的父，\Quote{赐给你们智慧和启示的神恩，使你们认识祂，并光照你们心目的眼睛}（\BibleRef{弗 1:17}），愿祂常常保守你们在善行、善言和善念中；愿光荣、尊威和权能归于祂，藉着我们的主耶稣基督，与圣神同在，从现在直到永远，永无穷尽。阿们。\par

\SourceRecord{Translated by Edwin Hamilton Gifford.；Nicene and Post-Nicene Fathers, Second Series；Vol. 7.；Edited by Philip Schaff and Henry Wace.；Buffalo, NY: Christian Literature Publishing Co.,；1894.；<http://www.newadvent.org/fathers/310118.htm>.}

% structure: p0001 source=h1 target=section reason=page-title

\section{教理讲授第十九讲}

\Emph{奥秘讲授第一讲}\par

经文选自\BibleBook{伯多禄}{前书}，自\BibleQuote{你们要节制，要警醒}起，至书信终。\par

我长久以来一直渴望，哦，教会真正的、亲爱的儿女们，与你们谈论这些属神和天上的奥秘；但我深知眼见远比耳听更有说服力，所以我等待这个时机；好使你们因现在的经验而更易接受我话语的影响，我引领你们进入眼前乐园中更明亮、更芬芳的草地；尤其因为你们在堪当领受神圣且赋予生命的圣洗之后，已准备好接受更神圣的奥秘。因此，既然还剩下为你们摆设更完美训导的餐桌，现在让我们精确地教导你们这些事，好使你们知道在你们领受圣洗的那天晚上，在你们身上所成就的事。\par

首先你们进入洗礼堂的前厅，在那里面向西方，你们听从命令伸出手，并在撒殚面前弃绝他。现在你们必须知道，这个象征在古老历史中已有记载。因为当法老，那个最残酷暴虐的暴君，压迫着希伯来人自由高贵的民族时，天主派遣梅瑟将他们从埃及人的邪恶奴役中领出来。于是门框被涂上羔羊的血，使毁灭者远离有血记的房屋；希伯来人奇妙地得救了。然而，敌人在他们获救后追赶他们\BibleRef{出 14:9}，并看见海奇妙地为他们分开；但他仍继续紧追其后，突然被红海吞没淹没。\par

现在从旧事转向新事，从预象转向实体。那里有天主派遣梅瑟到埃及；这里有基督，由祂的圣父派遣到世界：那里，梅瑟带领受苦的人民出离埃及；这里，基督拯救世上受罪恶压迫的人：那里，羔羊的血是抵御毁灭者的咒语；这里，无玷羔羊耶稣基督的血成为驱退邪恶神灵的法力：那里，暴君追赶那古老的民族直到海边；这里，那胆大无耻、为邪恶之源的邪灵，一直追赶你们直到救恩的水流。古代的暴君被淹死在海中；而这现在的暴君在救恩之水中消失。\par

然而，你们仍奉命向着他，好像他在眼前，伸出手臂说：\Quote{我弃绝你，撒殚}。我也愿意说明为何你们面向西方站立；这是因为有必要。既然西方是感官黑暗的区域，而他是黑暗，他的统治也在黑暗之中，因此，以象征的意义面向西方，你们弃绝那黑暗阴郁的掌权者。那么，你们每人站起来说了什么？\Quote{我弃绝你，撒殚}——你邪恶而最残酷的暴君！意思是：\Quote{我不再惧怕你的势力；因为基督已经击败了它，祂与我同取了血肉，为要藉着这些，祂\Emph{以死亡摧毁死亡}\BibleRef{希 2:14}，使我不至于永远\Emph{受束缚}。}\Quote{我弃绝你}——你狡猾而最阴险的毒蛇。\Quote{我弃绝你}——你这阴谋家，在友谊的伪装下策划了所有悖逆，并在我们原祖父母中引发背信。\Quote{我弃绝你，撒殚}——一切邪恶的制造者和怂恿者。\par

然后，在第二句中，你们受教说：\Quote{以及你的一切作为}。撒殚的作为就是一切罪恶，你们也必须弃绝这些——就像一个人逃脱了暴君，也必然逃脱了他的武器一样。因此，所有各类罪恶都包含在魔鬼的作为中。但要知道：你们所说的一切，特别是在那最令人战栗的时刻，都记录在天主的册子上；因此，当你们做出任何违背这些承诺的事时，你们将被判为\Emph{违犯者}。\BibleRef{迦 2:18} 所以你们弃绝撒殚的作为；我是指一切违背理性的行为和思想。\par

然后你们说：\Quote{以及他的一切虚华}。魔鬼的虚华就是剧场的疯狂、赛马、狩猎以及这一切虚荣：那圣人在祈祷中求脱离这些，向天主说：\Quote{求祢转开我的眼睛，免见虚荣}。不要对剧场的疯狂感兴趣，在那里你们会看到演员轻浮的姿态，伴随着嘲弄和一切不雅，以及柔弱男子的狂舞——也不要在意那些在狩猎中将自己暴露于野兽的人，他们为了满足可怜的口腹之欲；他们为了用肉食服侍自己的肚腹，实际上自己却成了未驯服野兽腹中的食物；公平地说，为了他们自己的神——肚腹，他们在决斗中轻率地抛弃了生命。也要避开赛马，那疯狂而倾覆灵魂的景象。因为这一切都是魔鬼的虚华。\par

此外，在偶像节日中悬挂之物，无论是肉还是饼，或其他类似被不洁神灵呼求所污染之物，都被算作魔鬼的虚华。因为正如圣体圣事的饼和酒，在呼求神圣可敬的天主圣三之前只是普通的饼和酒，而在呼求之后，饼成为基督的身体，酒成为基督的血；同样，这些属于撒殚虚华的肉类，虽然本性简单，却因恶神的呼求而成为亵渎之物。\par

八、此后你念\Quote{且弃绝你一切的服务。}现在魔鬼的服务是在偶像庙宇中的祈祷，为敬拜无生命的偶像而做的事，点灯，或在泉边、河边烧香，正如有些人受梦境或恶灵欺骗，确实求助于此，以为甚至能治疗身体的疾病。不要追随这些事。观鸟、占卜、预兆、护符、叶上书写的咒语、巫术或其他邪术，以及诸如此类的事，都是魔鬼的服务；因此，要躲避它们。因为如果你在弃绝了\Person{撒殚}并与\Person{基督}联合之后，又落入它们的影响之下，你会发现那暴君更为残忍；或许，因为他从前待你如己，解除了你沉重的束缚，但现在被你大大激怒；因此，你将失去\Person{基督}，并体验那另一个。你难道没有听过关于\Person{罗特}和他女儿们的古老历史吗？当他到达山上时，他不是和女儿们一起得救了吗？而他的妻子却变成了一根盐柱，永远竖立作为纪念，以记念她邪恶的意志和她的回转。因此，你要谨慎自己，不要再次转向\Emph{背后的事}，既已手扶着犁，却又转回这世俗行为的咸味；但要逃到山上，逃到\Person{耶稣基督}那里，祂是\Emph{非人手凿出的石头}\BibleRef{达 2:35}，那石头已充满了世界。\par

九、因此，当你弃绝\Person{撒殚}，彻底打破你与他的一切盟约，即那古老的\Emph{与地狱的联盟}\BibleRef{依 28:15}，\Person{天主}的乐园便向你敞开，祂在东方栽植了它，我们的原祖因他的过犯从那里被逐出；而这一点的象征，就是你从西方转向东方，光明之地。然后你被教导要念：\Quote{我信\Person{圣父}、\Person{圣子}、\Person{圣神}，并信一个悔改的圣洗。}关于这些事，我们已经在以前的训诲中，靠着\Person{天主}的恩宠，向你们详细讲论了。\par

十、因此，受这些训诲的保护，\Emph{要警醒}。\Emph{因为}我们的\Emph{仇敌魔鬼}，正如刚才所读的，\Emph{如同咆哮的狮子，巡游寻找可吞食的人}\BibleRef{伯前 5:9}。但是，虽然在以往死亡权势强大且吞灭一切，但在圣洗之盆（重生之洗）中，\Person{天主}已经\Emph{擦去各人脸上的一切眼泪}。因为你不再哀哭，既已脱去旧人；却要庆祝节日，\Emph{穿上救恩之衣}\BibleRef{依 61:10}，就是\Person{耶稣基督}。\par

十一、这些事是在外殿进行的。但是，如果\Person{天主}愿意，在后来的关于\OriginalTerm{奥迹}{Mysteries}的训诲中，当我们进入\OriginalTerm{至圣所}{Holy of Holies}时，我们将在那里知道所行之事的象征意义。愿一切光荣、权能和尊威归于\Person{圣父}，与\Person{圣子}及\Person{圣神}，于无穷世。阿们。\par

\SourceRecord{Translated by Edwin Hamilton Gifford.；Nicene and Post-Nicene Fathers, Second Series；Vol. 7.；Edited by Philip Schaff and Henry Wace.；Buffalo, NY: Christian Literature Publishing Co.,；1894.；<http://www.newadvent.org/fathers/310119.htm>.}

% structure: p0001 source=h1 target=section reason=page-title

\section{要理讲授 第二十讲}

\Emph{（论奥迹·二）}\par

\Emph{论圣洗}\par

\BibleRef{罗 6:3-14}\par

\BibleQuote{岂不知我们这许多受过洗归于耶稣基督的人，是受洗归于祂的死吗？……因为你们不在法律之下，而是在恩宠之下。}\par

1. 这些每日引入\OriginalTerm{奥迹}{Mysteries}的教导，以及新的训诲——这些是新真理的宣布——对我们是有益的；尤其对你们，你们已从旧状态更新为新状态。因此，我必要将昨日讲授的后续摆在你们面前，使你们明白那些在内室中所行的事是何象征。\par

2. 当你们一进入时，你们就脱去外衣；这象征\BibleQuote{脱去\OriginalTerm{旧人}{old man}和他的行为}。\BibleRef{哥 3:9} 你们脱衣后，赤身露体；在这件事上也效法基督，祂在十字架上被剥去衣服，并藉祂的赤身\Emph{脱去祂自己的\OriginalTerm{宰制者和掌权者}{principalities and powers}，并公开在木架上凯旋战胜了它们}。因为既然敌对势力在你们的肢体中筑巢，你们不可再穿那件旧衣裳；我完全不是指这件看得见的衣裳，而是\BibleQuote{那因私欲的迷惑渐渐败坏的旧人}。\BibleRef{弗 4:22} 愿那灵魂既已脱去它，就永不再穿上，却要同基督的新娘在\BibleBook{雅}{歌}中说：\BibleQuote{我脱了外衣，怎能再穿上呢？}\BibleRef{歌 5:3} 何等奇妙！你们在众人面前赤身，却不以为耻；因为你们实在带着那首先受造的亚当的肖像，他在乐园中赤身，却不以为耻。\par

3. 那时，你们被脱去衣服后，就用\OriginalTerm{驱魔油}{exorcised oil}从头顶到脚底被傅油，并成为好橄榄树耶稣基督的分享者。因为你们从\OriginalTerm{野橄榄树}{wild olive-tree}上被砍下，被接在\OriginalTerm{好橄榄树}{good olive-tree}上，并得以分享真橄榄树的肥脂。因此驱魔油是分享基督肥脂的象征，是驱除一切敌对势力痕迹的咒语。因为正如圣人们的呼气，以及呼求天主圣名，像最猛烈的火焰，烧灼并驱逐恶神；同样，这驱魔油，藉呼求天主和祈祷，领受如此德能，不仅烧尽并洗净罪的痕迹，也赶走那恶者的一切无形力量。\par

4. 此后，你们被引至\OriginalTerm{神圣圣洗池}{holy pool of Divine Baptism}，正如基督从十字架上被抬至我们眼前的坟墓。每人被问是否信父、子及圣神之名，并作了那\OriginalTerm{得救的宣认}{saving confession}，然后三次下水，又上升；这里也以象征暗示\OriginalTerm{三天埋葬}{three days burial}。因为正如我们的救世主在地心过了三天三夜，同样你们首次从水中上升，表征基督在地中的第一日，而你们的下降表征夜间；因为正如在夜间的人不再看见，而白昼的人留在光中，同样在下降时，如夜间，你们什么也看不见，而在再上升时，你们如在白昼。在同一时刻，你们既死又生；那救恩之水同时是你们的坟墓和母亲。所罗门论及其他人的话也适合你们；因为他曾说：\BibleQuote{诞生有时，死亡有时}\BibleRef{训 3:2}；但你们，顺序相反，死亡有时，诞生有时；同一时间完成这两者，你们的诞生与你们的死亡并行。\par

5. 奥妙而不可思议的事！我们并非真的死了，并非真的埋葬了，并非真的被钉十字架并复活了；但我们的\OriginalTerm{仿效}{imitation}是\OriginalTerm{预像}{figure}，而我们的救恩是真实的。基督确实被钉十字架，确实埋葬了，确实复活了；祂将这些白白赐予了我们，使我们借仿效分享祂的苦难，而得到真实的救恩。超越的慈爱！基督在祂无玷的手足上承受了钉子，并忍受了痛苦；而我，无需痛苦劳碌，藉与祂苦难的共融，祂白白赐予我救恩。\par

6. 那么，没有人应以为圣洗仅仅是\OriginalTerm{罪的赦免}{remission of sins}的恩宠，甚或仅是\OriginalTerm{义子地位}{adoption}的恩宠；因为若翰的圣洗只赋予罪的赦免：但我们深知，圣洗既洗净我们的罪，又赐予我们圣神的恩赐，同样也是基督苦难的对应。为此，保禄刚才大声呼喊说：\BibleQuote{岂不知我们这许多受洗归于基督耶稣的人，是受洗归于祂的死吗？所以我们藉着洗礼已与祂同葬于死。}\BibleRef{罗 6:3-4} 这些话他是对那些倾向于认为圣洗赐予我们罪的赦免和义子地位，却没有进一步通过预像与基督真实苦难共融的人说的。\par

7. 因此，为使我们明白：基督为我们并为我们的救恩所忍受的一切，是真实地而非表面地受的苦，并且我们也成为祂苦难的分享者，保禄以完全精确的真理喊道：\BibleQuote{如果我们与祂\OriginalTerm{死亡的相似}{likeness of His death}\OriginalTerm{同植}{planted together}，我们也要与祂复活的相似同植}\BibleRef{罗 6:5}。他说\Emph{同植}说得很好。因为既然真葡萄树被栽植于此，我们藉着参与死亡之圣洗，也与祂同植。请你们非常专注地思考宗徒的话。他并没有说：\Quote{如果我们与祂的死同植}，而是说：\Emph{与祂死亡的相似}。因为在基督的情况中，有真实的死亡，因为祂的灵魂确实与祂的身体分离，有真实的埋葬，因为祂神圣的身体被包裹在净亚麻布中；一切都真实地发生在祂身上；但在你们的情况中，只有死亡和苦难的相似，而救恩则不是相似，而是真实。\par

8. 在充分受教于这些事后，我恳求你们牢记在心；这样，我虽不配，也能对你们说：\Quote{\Emph{现在我爱你们}，\Emph{因为你们常常记念我，持守我所传给你们的传授}。}天主，他使你们\Quote{如同从死人中复活}（\BibleRef{罗 6:13}），能够赐予你们\Quote{在新生活中行走}；因为光荣和权能归于他，从现在直到永远。阿们。\par

\SourceRecord{Translated by Edwin Hamilton Gifford.；Nicene and Post-Nicene Fathers, Second Series；Vol. 7.；Edited by Philip Schaff and Henry Wace.；Buffalo, NY: Christian Literature Publishing Co.,；1894.；<http://www.newadvent.org/fathers/310120.htm>.}

% structure: p0001 source=h1 target=section reason=page-title

\section{\WorkTitle{要理讲授第21讲}}

\Emph{论奥秘（三）}\par

\Emph{论傅油}\par

\BibleRef{若一 2:20-28}\par

\Quote{你们从圣者领受了傅油，等等……好使当祂显现时，我们能坦然无惧，在祂来临时不至于在祂面前羞愧。}\par

1. 既然你们已\BibleQuote{在基督内受洗}并\BibleQuote{穿上了基督}\BibleRef{迦 3:27}，你们便成了与天主子相似的人；因为天主\BibleQuote{预定了我们得以嗣子的名分}\BibleRef{弗 1:5}，使我们\BibleQuote{与基督光荣的身体相似}\BibleRef{斐 3:21}。因此，你们既成了\BibleQuote{有分于基督的人}\BibleRef{希 3:14}，便恰当地被称为\Quote{受傅者}，天主论及你们时说过：\BibleQuote{不要触犯我的受傅者}。如今，你们因领受了圣神的预象而成了受傅者；万有都在你们身上藉着仿效而成，因为你们是基督的肖像。祂在\Place{约但河}洗了澡，将祂天主性的芬芳赐给了水，然后从水上来；圣神以祂的圆满降临在祂身上，如同停留在同类上。同样，在你们从圣水之池上来之后，也领受了一种傅油，即基督受傅的预象；这傅油就是圣神。蒙福的\Person{依撒意亚}在关于祂的预言中，以主的口吻说：\BibleQuote{上主的神临到我身上，因为祂给我傅了油；祂派遣我向贫穷人传报喜讯}\BibleRef{依 61:1}。\par

2. 因为基督并非由人用油或物质香膏傅油，而是圣父在预定祂为普世救主之后，以圣神傅了祂，正如\Person{伯多禄}所说：\BibleQuote{纳匝肋人耶稣，天主以圣神傅了祂}\BibleRef{宗 10:38}。\Person{达味}先知也呼喊说：\BibleQuote{天主，祢的宝座永远常存；祢王国的权杖是正直的权杖。祢喜爱正义，憎恨邪恶；因此天主，祢的天主，用喜乐之油傅了祢，胜过祢的同伴}\BibleRef{咏 45:7-8}。正如基督确实被钉十字架、被埋葬、复活，而你们在圣洗中蒙恩在相似中与祂同钉、同葬、同复活，同样，傅油也是如此。祂以理想的喜乐之油受傅，即圣神，之所以称为喜乐之油，因为祂是属灵喜乐的元始；同样，你们也用香膏受傅，成了\Quote{有分于基督}和\Quote{祂的同伴}。\par

3. 但你们要小心，不要以为这是普通的香膏。因为正如圣体圣事之饼在呼求圣神之后，不再是寻常之饼，而是基督的圣体；同样，这傅油在呼求之后，不再是单纯的香膏，也不是（姑且说）普通的香膏，而是基督的恩赐，并因圣神的降临而成为能赋予祂神性之合适的媒介。这傅油象征性地傅在你们的前额和其他感官上；当你们的身体用可见的香膏受傅时，你们的灵魂被圣而赋予生命的圣神所圣化。\par

4. 你们首先在前额受傅，好使你们脱离那第一个悖逆者到处携带的羞耻，并使你们\BibleQuote{以揭开的脸面反映主的荣耀}\BibleRef{格后 3:18}。然后傅在耳朵上，好使你们领受那敏锐聆听神圣奥秘的耳朵，\Person{依撒意亚}曾论及这奥秘说：\BibleQuote{主又赐给我一只耳朵来聆听}\BibleRef{依 50:4}；主耶稣在福音中也说：\BibleQuote{有耳听的，听吧}\BibleRef{玛 11:15}。然后傅在鼻孔上，好使你们领受圣傅油后得以说：\BibleQuote{我们是基督的馨香，在得救的人中，献给天主的馨香}\BibleRef{格后 2:15}。最后傅在胸前，好使你们穿上\Quote{正义的胸甲}，\Quote{对抗魔鬼的诡计}。因为正如基督在受洗和圣神降临之后，出去战胜了仇敌；同样，你们在领受圣洗和神秘的傅油之后，穿上了圣神的全副武装，当对抗仇敌的权势并战胜它，说：\BibleQuote{我靠那加强我力量的基督，能应付一切}\BibleRef{斐 4:13}。\par

5. 你们既被恩赐这圣傅油，便被称为\Quote{基督徒}，并因新生而证实了这个名号。因为在你们被恩赐这恩宠之前，你们本无权拥有这名号，而是正走在成为基督徒的路上。\par

6. 此外，你们要知道在旧约中已有这傅油的预象。当\Person{梅瑟}将天主的命令传授给他的兄弟，立他为大司祭时，先用水洗了他，然后傅油；\Person{亚郎}便被称为\Quote{受傅者}，显然是由于这个预象性的傅油。同样，大司祭在立\Person{撒罗满}为王时，先让他在\Place{基红池}洗浴，然后傅油\BibleRef{列上 1:39}。然而，这些事在他们身上是以预象发生的，但在你们身上却不是预象，而是真实，因为你们确实被圣神傅了油。基督是你们救恩的元始；因为\BibleQuote{祂确实是初果，而你们是面团}\BibleRef{罗 11:16}；但若初果是圣的，显然它的圣德也会传给面团。\par

7. 请保持此圣油无玷：因为它存留在你们内，必要教导你们一切，正如你们刚才听真福\Person{若望}所宣示的，他关于此傅油多有论述。因为这圣物是身体的神圣护卫，也是灵魂的救恩。关于此，真福\Person{依撒意亚}古时曾预言说：\BibleQuote{在这山上}——（他也在别处称教会为山，如他说：\BibleQuote{在末日，上主殿的山必要彰显} \BibleRef{依 2:2}；）——\BibleQuote{在这山上，上主要为万民设宴；他们要喝酒，他们要喝喜乐，他们要用油膏自己。}为使你们确知，请听他对这奥秘的油所说的话：\BibleQuote{将这些事通传于万民，因为上主的计划是为了万民。}因此，你们既已受此圣油傅抹，就应保持它在你们内无玷无瑕，藉善行不断前进，使你们救恩的元首基督耶稣喜悦，愿光荣归于祂，至于无穷之世。阿们。\par

\SourceRecord{Translated by Edwin Hamilton Gifford.；Nicene and Post-Nicene Fathers, Second Series；Vol. 7.；Edited by Philip Schaff and Henry Wace.；Buffalo, NY: Christian Literature Publishing Co.,；1894.；<http://www.newadvent.org/fathers/310121.htm>.}

% structure: p0001 source=h1 target=section reason=page-title

\section{教理讲授 第二十二讲}

\Emph{（论奥迹。第四讲）}\par

\Emph{论基督的圣体和圣血}\par

\BibleRef{格前 11:23}\par

\BibleQuote{我从主领受了，我也传授给你们：主耶稣在被出卖的那一夜，拿起饼来，等等。}\par

单纯真福保禄的教导就足以使你们对神圣奥迹完全确信，因为你们蒙恩与基督成为同一身体和血。你们刚才清楚听见他说：\Quote{我们的主耶稣基督在被出卖的那一夜，拿起饼来，祝谢了，擘开，递给祂的门徒说：\InnerQuote{你们拿去吃，这是我的身体。}又拿起杯来，祝谢了，说：\InnerQuote{你们拿去喝，这是我的血。}}既然祂亲自宣告并说出：\Quote{这是我的身体}，谁还敢再怀疑呢？既然祂亲自肯定并说：\Quote{这是我的血}，谁还会犹豫说这不是祂的血呢？\par

祂曾在加里肋亚的加纳把水变成酒（那酒与血相似），难道祂把酒变成血就不可信吗？当被邀请参加婚宴时，祂行了这奇事；对于\Quote{婚筵之子}（\BibleRef{玛 9:15}），祂岂不是更应被承认为赏赐了祂身体和血的享用吗？\par

因此，让我们满怀确信地领受基督的身体和血；因为在饼的形象下赐给你们祂的身体，在酒的形象下赐给你们祂的血；好使你们藉着领受基督的身体和血，与祂成为同一身体和同一血。这样我们就在我们内带着基督，因为祂的身体和血分布到我们的肢体中；因此，按照真福伯多禄的话，\Quote{我们成为有分于天主性体的人}（\BibleRef{伯后 1:4}）。\par

基督曾在某次与犹太人谈话时说：\Quote{你们若不吃我的肉，不喝我的血，在你们内就没有生命}（\BibleRef{若 6:53}）。他们因没有按属灵意义听祂的话而起了反感，并退去了，以为祂是邀请他们吃血肉。\par

在旧约中也有供饼；但那是属于旧约的，已经终止了；但在新约中却有天上的饼和救恩之杯，圣化灵魂和身体；因为正如饼对应于我们的身体，同样圣言适合于我们的灵魂。\par

因此，不要把饼和酒看作是纯元素，因为他们按照主的宣告，是基督的身体和血；因为虽然感官向你暗示这一点，但让信德坚固你。不要从味道判断这事，但要藉着信德毫不犹疑地确信，基督的身体和血已恩赐给了你们。\par

真福达味也会向你解释这含义，说：\Quote{在我敌人面前，你为我摆设了筵席。}他所说的是这个意思：在你来临之前，恶灵为人预备了污秽、玷污、充满魔鬼影响的筵席；但自从你来临，主啊，\Quote{你为我摆设了筵席}。当人对天主说\Quote{你在我面前摆设了筵席}，他指的不就是那神秘的、属灵的筵席吗？那是天主为我们准备的，与恶灵相对相反。非常正确；因为那筵席与魔鬼相通，而这筵席与天主相通。\Quote{你用油膏了我的头。}祂用油在你的头上膏抹了你的前额，作为从天主所得的印记；使你成为\Quote{刻着\InnerQuote{圣洁归主}的印玺}。而\Quote{你的杯使我沉醉，如同极强的酒一般。}你看见这里所说的杯，就是耶稣拿在手中，祝谢了，并说的：\Quote{这是我的血，为多人流出来，以赦免罪过}（\BibleRef{玛 26:28}）。\par

因此，撒罗满也在\BibleBook{训道}{篇}中暗示这恩宠，说：\Quote{你且来，欢欢喜喜地吃你的饼}（即属灵的饼；\Quote{你且来}，他用以召唤得救和祝福的呼声），\Quote{并快乐地喝你的酒}（即属灵的酒）；\Quote{且让油倾倒在你的头上}（你看，他甚至暗示了圣化之油）；\Quote{让你的衣服时常洁白，因为上主喜悦你的工作}（\BibleRef{训 9:7}）；因为在你领受洗礼之前，你的工作都是\Quote{虚而又虚}。但现在，既然脱去了旧衣，穿上了属灵洁白的新衣，你必须时常身穿白衣；当然，我们不是说你要一直穿白色衣服；而是你必须穿上真正洁白、发光、属灵的衣裳，好使你能与真福依撒意亚一同说：\Quote{我的灵魂要因我的天主而喜乐，因为他给我穿上救恩的衣服，给我披上喜乐的外衣}（\BibleRef{依 61:10}）。\par

你们既学习了这些事，并完全确信那看似饼的并不是饼（虽然味觉感到如此），而是基督的身体；那看似酒的并不是酒（虽然味觉会以为是），而是基督的血；达味昔日曾歌咏说：\Quote{而饼强健人的心，用油使他的面容发光}，你们要因领受这属灵的饼而\Quote{强健你们的心}，并\Quote{使你们灵魂的面容发光}。因此，以纯洁的良心无遮蔽地如此行，愿你们\Quote{如同镜子反映上主的光荣}（\BibleRef{格后 3:18}），并从\Quote{光荣到光荣}，在我们的主耶稣基督内——愿荣耀、权能和光荣归于祂，直到永远。阿们。\par

\SourceRecord{Translated by Edwin Hamilton Gifford.；Nicene and Post-Nicene Fathers, Second Series；Vol. 7.；Edited by Philip Schaff and Henry Wace.；Buffalo, NY: Christian Literature Publishing Co.,；1894.；<http://www.newadvent.org/fathers/310122.htm>.}

% structure: p0001 source=h1 target=section reason=page-title

\section{慕道宣讲 23}

\Emph{(论奥迹. 五)}\par

\Emph{论神圣礼仪与共融。}\par

\BibleRef{伯前 2:1}\par

\Emph{所以，你们要除去一切污秽、一切诡诈和恶言，等等。}\par

1. 由于天主的慈爱，你们在以前的聚会中已经充分听过了关于圣洗、坚振，以及领受基督的圣体圣血的道理；现在需要继续论到接下来次序中的内容，意思是今天要为你们灵性建造的工程加上冠冕。\par

2. 你们已经看见执事将洗手的水递给司铎，并递给围绕天主祭台的司铎们。他这样做完全不是由于身体的污秽；并非如此，因为我们起初并不是带着污秽的身体进入教会的。但洗手是一个象征，表明你们应当远离一切罪恶和违法的事；因为双手是行为的象征，通过洗手，显然我们代表行为的纯洁和无辜。难道你们没有听见真福达味启示这一奥迹，说：\Emph{我要洗手表明无辜，才环绕你的祭台，上主。}？因此，洗手是无罪的象征。\par

3. 然后执事大声呼喊：\Quote{你们要彼此接纳，并要彼此亲吻。}不要以为这个亲吻与普通朋友在公共场所的亲吻相同。并非如此：但这个亲吻使灵魂彼此融合，并为他们求得完全的宽恕。因此，亲吻是我们灵魂彼此相融、驱除所有冤仇的标志。为此原因基督说：\Emph{若你正在祭台上献礼物，在那里想起你的兄弟有什么怨你的事，就把礼物留在祭台前，先去与你的兄弟和好，然后再来献礼物。}所以亲吻是和好，因此是神圣的：正如真福保禄在某处呼喊说：\Emph{你们要以圣吻彼此问候} \BibleRef{格前 16:20}；以及伯多禄：\Emph{用爱德的亲吻} \BibleRef{伯前 3:15}。\par

4. 此后，司铎大声呼喊：\Quote{请举心向上。}因为在那最可敬畏的时刻，我们实在应当将心高举与天主同在，而不应低垂，思念下地和地上的事物。所以司铎实际上命令众人在那一刻放下此生的一切牵挂或家庭忧虑，将心置于天上与仁慈的天主同在。然后你们回答说：\Quote{我们向上主举心。}藉你们的宣认表示同意。但谁也不应到这里来，嘴里说\Quote{我们向上主举心，}而思想中却挂虑此生的牵挂。更确切地说，天主应时刻存于我们的记忆中，但若因人性软弱而无法做到，那么在那时刻，这更应是我们热切努力的事。\par

5. 然后司铎说：\Quote{我们感谢主。}因为我们确实应当感谢，祂召叫了我们这些不配的人，获得了如此大的恩宠；祂在我们是仇敌时使我们与祂和好；祂赐予我们义子的圣神。然后你们说：\Quote{这是理所当然的：}因为感恩时我们做了一件合适且正当的事；但祂所做的不是正当，而是超过正当，祂善待我们，并认为我们配得上如此大的恩惠。\par

6. 此后，我们提及天、地、海，日月星辰，及一切受造物，有理性的和无理性的，可见的与不可见的；天使、总领天使、异能、宰制、率领、掌权、宝座；及有许多面孔的革鲁宾：实际上重复达味的呼喊\Emph{请与我一同颂扬上主}。我们也提及色辣芬，就是依撒意亚在圣神中看见站立在天主宝座周围的，用两个翅膀遮脸，两个翅膀遮脚，两个翅膀飞翔，呼喊说\Emph{圣、圣、圣，万军的上主。} \BibleRef{依 6:2}。我们背诵这从色辣芬传下来的对天主的宣认，原因在于，如此我们得以与天上万军一同参与他们的赞美诗歌。\par

7. 然后，藉这些属灵的赞美诗圣化自己后，我们恳求仁慈的天主派遣祂的圣神降临在祂面前的祭品上，使饼成为基督的身体，酒成为基督的血；因为凡圣神所接触的，确实被圣化并转变。\par

8. 然后，在属灵的祭献、无血的礼仪完成后，在这赎罪祭献上，我们恳求天主赐予教会普遍的和平、世界的福祉；为君王、为士兵和盟友；为患病者、为受苦者；总之，为所有需要援助的人，我们都祈祷并奉献这祭献。\par

9. 然后我们也纪念那些在我们之前安息的人，首先是圣祖、先知、宗徒、殉道者，愿因他们的祈祷和转求，天主接纳我们的恳求。然后也为那些在我们之前安息的圣父和主教们，以及总之所有在过去年间在我们中间安息的人，我们相信，当这神圣而最可敬畏的祭献呈上时，这恳求将为那些灵魂带来极大的益处。\par

10. 我愿意用一个比喻来说服你们。因为我知道许多人说：一个灵魂带着罪或不带罪离开此世，若在祈祷中被纪念，有什么益处呢？假如一位君王流放了某些冒犯他的人，然后那些属于他们的人编织一个花冠献给他，为那些受罚的人求情，他难道不会赦免他们的刑罚吗？同样，当我们为那些安息的人向祂献上恳求时，即使他们是罪人，我们并不是编织花冠，而是奉献为我们罪过而被牺牲的基督，为我们也为他们恳求仁慈的天主。\par

11. 此后，我们以纯洁的良心称天主为父，并念诵救主传授给祂自己的门徒的那篇祷文，说：\Quote{我们的圣父，祢在天上。} 啊，天主何等超越的爱！祂对那些背叛祂、处于极度痛苦中的人，赐予如此完全的罪过赦免，以及如此大的恩宠分施，以致他们竟能称祂为父。\Quote{我们的圣父，祢在天上。} 他们也是一重天，因为\BibleQuote{他们带着天上的肖像} \BibleRef{格前 15:49} ，\BibleQuote{天主在他们内居住并行走} \BibleRef{格后 6:16}。\par

12. \Quote{愿祢的名被尊为圣。} 天主的名本性是圣的，无论我们是否这样说；但由于它在罪人中间有时被亵渎，正如经上所言：\Quote{因着你们，我的名在外邦人中不断被亵渎。} 所以我们祈求天主的名在我们内被尊为圣；并非它从不圣洁变为圣洁，而是因着我们在成圣、行相称圣德之事时，它在我们内成为圣洁。\par

13. \Quote{愿祢的国来临。} 纯洁的灵魂能勇敢地说：\Quote{愿祢的国来临。} 因为谁若听到保禄说：\BibleQuote{所以，不要让罪在你们必死的身体上为王} \BibleRef{罗 6:12}，并在行为、思想和言语上洁净了自己，他就会对天主说：\Quote{愿祢的国来临。}\par

14. \Quote{愿祢的旨意奉行在人间，如同在天上。} 天主神圣而受祝福的天使们承行天主的旨意，正如达味在\BibleBook{}{圣咏}中所说：\Quote{主的天使，请赞美主，你们是执行祂命令的勇士。} 因此，你的祈祷实际上是在说：\Quote{如同天使中祢的旨意得以奉行，同样，也在我内、在地上奉行吧，主。}\par

15. \Quote{求祢今日赐给我们我们日用的食粮。} 这普通的饼并非食粮，而这圣饼才是食粮，即专为灵魂之本体所指定的。因为此饼\BibleQuote{不入肚腹，而是被排泄出去} \BibleRef{玛 15:17}，而是被分施到你的全身，为灵魂和身体益处。但\Quote{今日}之意，他解释为\Quote{每一天}，正如保禄说：\Quote{趁着还有今日} \BibleRef{希 3:15}。\par

16. \Quote{求祢宽恕我们的罪债，如同我们也宽恕欠我们债的人。} 因为我们有许多罪过。我们无论在言语、思想上都有过犯，做了许多应受谴责的事；而且\Quote{若我们说没有罪}，我们就在说谎，正如若望所说。我们与天主立盟约，恳求祂宽恕我们的罪，如同我们也宽恕邻人的债。那么，想想我们领受了什么，又为了什么，就不要推迟或延迟彼此宽恕。得罪我们的事是轻微琐碎的，容易解决；但我们得罪天主的事是重大的，需要唯独祂那样的慈悲。所以务要谨慎，免得因为他人对你轻微的罪过，你为自己关上了天主对你极重罪过的宽恕之门。\par

17. \Quote{主，求祢不要让我们陷入诱惑。} 难道主教导我们这样祈祷，是为了使我们根本不受诱惑吗？那为何在其他地方说：\Quote{未受诱惑的人，是未经证明的人。} 又说：\BibleQuote{我的弟兄们，当你们落在各种诱惑中，要认为是大喜乐} \BibleRef{雅 1:2}？或许\Quote{陷入诱惑}意味着被诱惑所压倒？因为诱惑好像冬天难以渡过的激流。因此，那些在诱惑中不被压倒的人，穿过去，显出自己是优秀的泳者，丝毫未被冲走；而那些不然的人，则陷入其中并被淹没。例如，犹达斯陷入了贪爱金钱的诱惑，却没有游过去，反而被淹没，在身体和灵魂上都勒死了。伯多禄陷入了否认的诱惑；但他进入后，并未被压倒，而是勇敢地游了过去，并从诱惑中被解救出来。再听另一处，一群未受伤害的圣人感谢从诱惑中被解救，他们说：\Quote{天主，祢考验了我们，祢以火炼净我们，如炼银子。祢使我们进入网罗，将重担放在我们腰间。祢使人骑在我们头上；我们经过水火，祢却领我们到安居之地。} 你看他们勇敢地谈论自己经过而未被刺透。\Quote{但祢领我们到安居之地。} 他们到达安居之地，就是从诱惑中被解救。\par

18. \Quote{但救我们脱离凶恶。} 如果\Quote{不要让我们陷入诱惑}的意思是根本不要受试探，祂就不会说\Quote{但救我们脱离凶恶}。凶恶就是我们的仇敌魔鬼，我们祈求从他手中被救。然后祈祷结束后，你说\Quote{阿们}；藉这个\Quote{阿们}，意思是\Quote{但愿如此}，你在那天主教导的祈祷的祈求上盖上了印记。\par

19. 此后，司铎说：\Quote{圣物归于圣者。} 所奉献的礼品是圣的，因为已领受了圣神的降临；你们也是圣的，因为被堪当领受圣神；所以圣物与圣者相应。然后你们说：\Quote{唯一圣者，唯一主，耶稣基督。} 因为那一位确是真圣者，本性圣洁；我们也是圣的，但并非本性的，而是藉着分有、纪律和祈祷。\par

20. 此后，你们听见咏祷者以神圣的旋律邀请你们参与圣奥迹的共融，念诵：\Quote{请尝并看，主是何等甘饴。} 不要依赖你肉身味觉的判断，而要以不可动摇的信德；因为那些品尝的人是受邀来品尝的，不是饼和酒，而是基督的预象之体和血。\par

21. 因此，当你前来时，不要伸开手腕或张开手指，却要用左手为右手作宝座，好迎接那位君王。手掌凹下，接过基督的圣体，并在其上呼\Quote{阿们}。然后，以圣体谨慎地触碰你的双眼，使之圣化，而后领受；务必留心，不可失落任何一部分，因为凡失落的，显然如同失去你自身的一肢。试问，若有人给你金粒，你岂不百般谨慎地握着，提防失落任何一颗而受损？那么，对于比黄金和宝石更贵重的，你岂不更要谨慎守护，不让一点碎屑从你手中掉落？\par

22. 领受了基督的圣体之后，随即也趋近祂的血杯，不要伸开双手，却要躬身，以朝拜和敬畏的神情呼\Quote{阿们}，并藉领受基督的宝血使自己圣化。当唇上还湿润时，以手指沾触，并圣化你的双眼、前额及其他感官。然后等候祈祷，并感谢天主，因祂认为你堪当如此伟大的奥迹。\par

23. 紧守这些毫无玷污的传授，并使自己不陷于过犯。不要自绝于共融，也不要因罪恶的污秽而使自己丧失这些神圣而属灵的奥迹。\BibleQuote{愿赐平安的天主亲自完全圣化你们，并愿你们整个的——灵、魂、身体——在我们主耶稣基督来临时，得以完好无缺，无可指摘。} \BibleRef{得前 5:23}——愿荣耀、尊威与权能归于祂，偕同圣父及圣神，自今至永远，及世之世。阿们。\par

\SourceRecord{Translated by Edwin Hamilton Gifford.；Nicene and Post-Nicene Fathers, Second Series；Vol. 7.；Edited by Philip Schaff and Henry Wace.；Buffalo, NY: Christian Literature Publishing Co.,；1894.；<http://www.newadvent.org/fathers/310123.htm>.}
